% =====================================================================
%  Fondamenti di Unreal Engine, lezione 9 — State Machine, Blend Space
%  e Montage.
%  Nona lezione del corso breve di introduzione a Unreal Engine (focus
%  sui Blueprint), online su Zoom, nel tardo pomeriggio (pausa
%  annunciata con ripresa "alle 18", chiusura verso le 19; "ci sentiamo
%  per martedì prossimo"). Docente titolare (come nelle lezioni
%  precedenti, nella trascrizione non viene mai chiamato per nome); la
%  classe è molto ridotta: sei persone all'inizio, poi cinque, quattro,
%  e alla fine due o tre.
%  Contenuti, nell'ordine della lezione: l'attesa e i corsi di indirizzo
%  degli studenti (3D art, game design, narrative design, character
%  animation); la lezione libera del 25 (menu e UI, o argomenti a
%  scelta); il ripasso dell'Animation Blueprint ABP_Knight
%  (Initialize e Update, Try Get Pawn Owner, Character Movement,
%  GroundSpeed, IsFalling); la State Machine Locomotion (Entry,
%  Conduit, State, Alias, lo stato WalkRun, Output Animation Pose, gli
%  stati Jump e Fall, le transizioni); il Blend Space 1D (asse
%  GroundSpeed da 0 a 600, Max Walk Speed a 300 centimetri al secondo,
%  il Blendspace Player); il problema "questo asset non ha schede
%  ancorate" (Load Layout) e la lingua dell'editor; il personaggio come
%  Default Pawn Class; la camera (Camera e Spring Arm, i 3C, Target Arm
%  Length, Target Offset e Socket Offset, Field of View); i grafi di
%  movimento copiati dal ThirdPersonCharacter; Use Pawn Control
%  Rotation, Orient Rotation to Movement e Use Controller Rotation Yaw;
%  il retargeting delle animazioni Jump e Fall Loop dal manichino di
%  Unreal (Retarget Animations, Export Animation); le regole di
%  transizione (IsFalling, Automatic Rule Based on Sequence Player in
%  State, NOT Boolean, la velocità sulla Z); lo sprint con l'Enhanced
%  Input (Input Mapping Context, Input Action IA_Run, Value Type,
%  Triggers, Add Mapping Context, i pin dell'evento, Set Max Walk
%  Speed, il controller e il Face Button Right); il tentativo con
%  FInterp To; le animazioni da scaricare (l'attacco in root motion,
%  gli strafe in place, il mirror, le diagonali); la pausa e i feedback
%  (il progetto personale del docente con la pallina, inventari e UI,
%  le lezioni con Davide: la gym e il level design); il Blend Space 2D
%  (Calculate Direction, la variabile Direction da -180 a 180, BS2D,
%  la griglia, la posizione delle otto animazioni, Smoothing Time);
%  gli Animation Montage (le animazioni modulari, i Notify e il Play
%  Sound, le Montage Section, gli Slot, Save Cached Pose e Slot
%  DefaultSlot nell'AnimGraph, Play Montage sul tasto F, Disable e
%  Enable Input, la camera che continua a ruotare); il compito sullo
%  StackOBot; la chiusura.
%  Fonte: trascrizione automatica della nona lezione (UE9), circa tre
%  ore. Integrazioni dai libri "Blueprints Visual Scripting for Unreal
%  Engine 5" (Romero, Sewell; capitoli 6 e 17) e "Blueprints Visual
%  Scripting for Unreal Engine", seconda edizione (stessi autori,
%  Unreal Engine 4), e nozioni tecniche generali su Unreal Engine 5,
%  sempre segnalate come approfondimenti.
%
%  Pacchetti richiesti nel preambolo del libro (già presenti nel
%  main.tex): \usepackage{booktabs}, \usepackage{amsmath},
%  \usepackage{amssymb}, \usepackage{tikz} con
%    \usetikzlibrary{shapes.geometric,shapes.symbols,shapes.misc,
%      arrows.meta,positioning,calc,fit}
%  (le librerie vanno caricate nel preambolo, non nel capitolo).
%  Consigliati, come per gli altri capitoli: \usepackage[italian]{babel}
%  e \usepackage[T1]{fontenc} (il capitolo compila anche senza).
%
%  Collocazione nel libro: parte "Fondamenti di Unreal Engine", nono
%  capitolo (numerazione propria 1.9). Nel file principale del libro,
%  subito dopo la riga dell'ottava lezione:
%
%    % =====================================================================
%  Fondamenti di Unreal Engine, lezione 9 — State Machine, Blend Space
%  e Montage.
%  Nona lezione del corso breve di introduzione a Unreal Engine (focus
%  sui Blueprint), online su Zoom, nel tardo pomeriggio (pausa
%  annunciata con ripresa "alle 18", chiusura verso le 19; "ci sentiamo
%  per martedì prossimo"). Docente titolare (come nelle lezioni
%  precedenti, nella trascrizione non viene mai chiamato per nome); la
%  classe è molto ridotta: sei persone all'inizio, poi cinque, quattro,
%  e alla fine due o tre.
%  Contenuti, nell'ordine della lezione: l'attesa e i corsi di indirizzo
%  degli studenti (3D art, game design, narrative design, character
%  animation); la lezione libera del 25 (menu e UI, o argomenti a
%  scelta); il ripasso dell'Animation Blueprint ABP_Knight
%  (Initialize e Update, Try Get Pawn Owner, Character Movement,
%  GroundSpeed, IsFalling); la State Machine Locomotion (Entry,
%  Conduit, State, Alias, lo stato WalkRun, Output Animation Pose, gli
%  stati Jump e Fall, le transizioni); il Blend Space 1D (asse
%  GroundSpeed da 0 a 600, Max Walk Speed a 300 centimetri al secondo,
%  il Blendspace Player); il problema "questo asset non ha schede
%  ancorate" (Load Layout) e la lingua dell'editor; il personaggio come
%  Default Pawn Class; la camera (Camera e Spring Arm, i 3C, Target Arm
%  Length, Target Offset e Socket Offset, Field of View); i grafi di
%  movimento copiati dal ThirdPersonCharacter; Use Pawn Control
%  Rotation, Orient Rotation to Movement e Use Controller Rotation Yaw;
%  il retargeting delle animazioni Jump e Fall Loop dal manichino di
%  Unreal (Retarget Animations, Export Animation); le regole di
%  transizione (IsFalling, Automatic Rule Based on Sequence Player in
%  State, NOT Boolean, la velocità sulla Z); lo sprint con l'Enhanced
%  Input (Input Mapping Context, Input Action IA_Run, Value Type,
%  Triggers, Add Mapping Context, i pin dell'evento, Set Max Walk
%  Speed, il controller e il Face Button Right); il tentativo con
%  FInterp To; le animazioni da scaricare (l'attacco in root motion,
%  gli strafe in place, il mirror, le diagonali); la pausa e i feedback
%  (il progetto personale del docente con la pallina, inventari e UI,
%  le lezioni con Davide: la gym e il level design); il Blend Space 2D
%  (Calculate Direction, la variabile Direction da -180 a 180, BS2D,
%  la griglia, la posizione delle otto animazioni, Smoothing Time);
%  gli Animation Montage (le animazioni modulari, i Notify e il Play
%  Sound, le Montage Section, gli Slot, Save Cached Pose e Slot
%  DefaultSlot nell'AnimGraph, Play Montage sul tasto F, Disable e
%  Enable Input, la camera che continua a ruotare); il compito sullo
%  StackOBot; la chiusura.
%  Fonte: trascrizione automatica della nona lezione (UE9), circa tre
%  ore. Integrazioni dai libri "Blueprints Visual Scripting for Unreal
%  Engine 5" (Romero, Sewell; capitoli 6 e 17) e "Blueprints Visual
%  Scripting for Unreal Engine", seconda edizione (stessi autori,
%  Unreal Engine 4), e nozioni tecniche generali su Unreal Engine 5,
%  sempre segnalate come approfondimenti.
%
%  Pacchetti richiesti nel preambolo del libro (già presenti nel
%  main.tex): \usepackage{booktabs}, \usepackage{amsmath},
%  \usepackage{amssymb}, \usepackage{tikz} con
%    \usetikzlibrary{shapes.geometric,shapes.symbols,shapes.misc,
%      arrows.meta,positioning,calc,fit}
%  (le librerie vanno caricate nel preambolo, non nel capitolo).
%  Consigliati, come per gli altri capitoli: \usepackage[italian]{babel}
%  e \usepackage[T1]{fontenc} (il capitolo compila anche senza).
%
%  Collocazione nel libro: parte "Fondamenti di Unreal Engine", nono
%  capitolo (numerazione propria 1.9). Nel file principale del libro,
%  subito dopo la riga dell'ottava lezione:
%
%    % =====================================================================
%  Fondamenti di Unreal Engine, lezione 9 — State Machine, Blend Space
%  e Montage.
%  Nona lezione del corso breve di introduzione a Unreal Engine (focus
%  sui Blueprint), online su Zoom, nel tardo pomeriggio (pausa
%  annunciata con ripresa "alle 18", chiusura verso le 19; "ci sentiamo
%  per martedì prossimo"). Docente titolare (come nelle lezioni
%  precedenti, nella trascrizione non viene mai chiamato per nome); la
%  classe è molto ridotta: sei persone all'inizio, poi cinque, quattro,
%  e alla fine due o tre.
%  Contenuti, nell'ordine della lezione: l'attesa e i corsi di indirizzo
%  degli studenti (3D art, game design, narrative design, character
%  animation); la lezione libera del 25 (menu e UI, o argomenti a
%  scelta); il ripasso dell'Animation Blueprint ABP_Knight
%  (Initialize e Update, Try Get Pawn Owner, Character Movement,
%  GroundSpeed, IsFalling); la State Machine Locomotion (Entry,
%  Conduit, State, Alias, lo stato WalkRun, Output Animation Pose, gli
%  stati Jump e Fall, le transizioni); il Blend Space 1D (asse
%  GroundSpeed da 0 a 600, Max Walk Speed a 300 centimetri al secondo,
%  il Blendspace Player); il problema "questo asset non ha schede
%  ancorate" (Load Layout) e la lingua dell'editor; il personaggio come
%  Default Pawn Class; la camera (Camera e Spring Arm, i 3C, Target Arm
%  Length, Target Offset e Socket Offset, Field of View); i grafi di
%  movimento copiati dal ThirdPersonCharacter; Use Pawn Control
%  Rotation, Orient Rotation to Movement e Use Controller Rotation Yaw;
%  il retargeting delle animazioni Jump e Fall Loop dal manichino di
%  Unreal (Retarget Animations, Export Animation); le regole di
%  transizione (IsFalling, Automatic Rule Based on Sequence Player in
%  State, NOT Boolean, la velocità sulla Z); lo sprint con l'Enhanced
%  Input (Input Mapping Context, Input Action IA_Run, Value Type,
%  Triggers, Add Mapping Context, i pin dell'evento, Set Max Walk
%  Speed, il controller e il Face Button Right); il tentativo con
%  FInterp To; le animazioni da scaricare (l'attacco in root motion,
%  gli strafe in place, il mirror, le diagonali); la pausa e i feedback
%  (il progetto personale del docente con la pallina, inventari e UI,
%  le lezioni con Davide: la gym e il level design); il Blend Space 2D
%  (Calculate Direction, la variabile Direction da -180 a 180, BS2D,
%  la griglia, la posizione delle otto animazioni, Smoothing Time);
%  gli Animation Montage (le animazioni modulari, i Notify e il Play
%  Sound, le Montage Section, gli Slot, Save Cached Pose e Slot
%  DefaultSlot nell'AnimGraph, Play Montage sul tasto F, Disable e
%  Enable Input, la camera che continua a ruotare); il compito sullo
%  StackOBot; la chiusura.
%  Fonte: trascrizione automatica della nona lezione (UE9), circa tre
%  ore. Integrazioni dai libri "Blueprints Visual Scripting for Unreal
%  Engine 5" (Romero, Sewell; capitoli 6 e 17) e "Blueprints Visual
%  Scripting for Unreal Engine", seconda edizione (stessi autori,
%  Unreal Engine 4), e nozioni tecniche generali su Unreal Engine 5,
%  sempre segnalate come approfondimenti.
%
%  Pacchetti richiesti nel preambolo del libro (già presenti nel
%  main.tex): \usepackage{booktabs}, \usepackage{amsmath},
%  \usepackage{amssymb}, \usepackage{tikz} con
%    \usetikzlibrary{shapes.geometric,shapes.symbols,shapes.misc,
%      arrows.meta,positioning,calc,fit}
%  (le librerie vanno caricate nel preambolo, non nel capitolo).
%  Consigliati, come per gli altri capitoli: \usepackage[italian]{babel}
%  e \usepackage[T1]{fontenc} (il capitolo compila anche senza).
%
%  Collocazione nel libro: parte "Fondamenti di Unreal Engine", nono
%  capitolo (numerazione propria 1.9). Nel file principale del libro,
%  subito dopo la riga dell'ottava lezione:
%
%    % =====================================================================
%  Fondamenti di Unreal Engine, lezione 9 — State Machine, Blend Space
%  e Montage.
%  Nona lezione del corso breve di introduzione a Unreal Engine (focus
%  sui Blueprint), online su Zoom, nel tardo pomeriggio (pausa
%  annunciata con ripresa "alle 18", chiusura verso le 19; "ci sentiamo
%  per martedì prossimo"). Docente titolare (come nelle lezioni
%  precedenti, nella trascrizione non viene mai chiamato per nome); la
%  classe è molto ridotta: sei persone all'inizio, poi cinque, quattro,
%  e alla fine due o tre.
%  Contenuti, nell'ordine della lezione: l'attesa e i corsi di indirizzo
%  degli studenti (3D art, game design, narrative design, character
%  animation); la lezione libera del 25 (menu e UI, o argomenti a
%  scelta); il ripasso dell'Animation Blueprint ABP_Knight
%  (Initialize e Update, Try Get Pawn Owner, Character Movement,
%  GroundSpeed, IsFalling); la State Machine Locomotion (Entry,
%  Conduit, State, Alias, lo stato WalkRun, Output Animation Pose, gli
%  stati Jump e Fall, le transizioni); il Blend Space 1D (asse
%  GroundSpeed da 0 a 600, Max Walk Speed a 300 centimetri al secondo,
%  il Blendspace Player); il problema "questo asset non ha schede
%  ancorate" (Load Layout) e la lingua dell'editor; il personaggio come
%  Default Pawn Class; la camera (Camera e Spring Arm, i 3C, Target Arm
%  Length, Target Offset e Socket Offset, Field of View); i grafi di
%  movimento copiati dal ThirdPersonCharacter; Use Pawn Control
%  Rotation, Orient Rotation to Movement e Use Controller Rotation Yaw;
%  il retargeting delle animazioni Jump e Fall Loop dal manichino di
%  Unreal (Retarget Animations, Export Animation); le regole di
%  transizione (IsFalling, Automatic Rule Based on Sequence Player in
%  State, NOT Boolean, la velocità sulla Z); lo sprint con l'Enhanced
%  Input (Input Mapping Context, Input Action IA_Run, Value Type,
%  Triggers, Add Mapping Context, i pin dell'evento, Set Max Walk
%  Speed, il controller e il Face Button Right); il tentativo con
%  FInterp To; le animazioni da scaricare (l'attacco in root motion,
%  gli strafe in place, il mirror, le diagonali); la pausa e i feedback
%  (il progetto personale del docente con la pallina, inventari e UI,
%  le lezioni con Davide: la gym e il level design); il Blend Space 2D
%  (Calculate Direction, la variabile Direction da -180 a 180, BS2D,
%  la griglia, la posizione delle otto animazioni, Smoothing Time);
%  gli Animation Montage (le animazioni modulari, i Notify e il Play
%  Sound, le Montage Section, gli Slot, Save Cached Pose e Slot
%  DefaultSlot nell'AnimGraph, Play Montage sul tasto F, Disable e
%  Enable Input, la camera che continua a ruotare); il compito sullo
%  StackOBot; la chiusura.
%  Fonte: trascrizione automatica della nona lezione (UE9), circa tre
%  ore. Integrazioni dai libri "Blueprints Visual Scripting for Unreal
%  Engine 5" (Romero, Sewell; capitoli 6 e 17) e "Blueprints Visual
%  Scripting for Unreal Engine", seconda edizione (stessi autori,
%  Unreal Engine 4), e nozioni tecniche generali su Unreal Engine 5,
%  sempre segnalate come approfondimenti.
%
%  Pacchetti richiesti nel preambolo del libro (già presenti nel
%  main.tex): \usepackage{booktabs}, \usepackage{amsmath},
%  \usepackage{amssymb}, \usepackage{tikz} con
%    \usetikzlibrary{shapes.geometric,shapes.symbols,shapes.misc,
%      arrows.meta,positioning,calc,fit}
%  (le librerie vanno caricate nel preambolo, non nel capitolo).
%  Consigliati, come per gli altri capitoli: \usepackage[italian]{babel}
%  e \usepackage[T1]{fontenc} (il capitolo compila anche senza).
%
%  Collocazione nel libro: parte "Fondamenti di Unreal Engine", nono
%  capitolo (numerazione propria 1.9). Nel file principale del libro,
%  subito dopo la riga dell'ottava lezione:
%
%    \input{UE/UE9_state_machine_blend_space_e_montage}
%
%  Tutte le etichette di questo capitolo hanno il prefisso "ue9:"
%  (ue9:cap:..., ue9:sec:..., ue9:fig:..., ue9:tab:...), e tutti gli
%  stili TikZ hanno il prefisso "ue9" (ue9node, ue9exec, ue9data, ...)
%  per non collidere con chiavi native di TikZ né con quelli delle
%  lezioni precedenti. Il capitolo è autonomo: non contiene \ref verso
%  altri capitoli (i rimandi alle lezioni precedenti sono a parole).
%  Le figure sono ricostruzioni schematiche di quanto mostrato e detto a
%  lezione, non copie delle schermate del docente. Dove la trascrizione
%  è troppo rovinata per ricostruire un dettaglio, il capitolo lo dice
%  invece di inventarlo.
% =====================================================================

\chapter[State Machine, Blend Space e Montage]{State Machine, Blend Space e Montage:
dare vita al nuovo personaggio}
\label{ue9:cap:animazioni}

Questa è la \textbf{nona lezione} del corso breve di \textbf{introduzione a Unreal
Engine}, e riprende esattamente dove si era fermata l'ottava. Alla fine della lezione
precedente il cavaliere importato da \textbf{Mixamo} aveva uno scheletro, le sue
animazioni e un \textbf{Animation Blueprint}, \textbf{ABP\_Knight}, che a ogni frame
leggeva dal personaggio due dati: la velocità sul piano (\textbf{GroundSpeed}) e il fatto
di essere in aria o meno (\textbf{IsFalling}). Mancava tutto il resto: chi decide
\emph{quale} animazione mostrare, e quando. ``Siccome ci siamo interrotti l'altra volta,
andiamo diretti a fare le animazioni: l'altra volta ci siamo interrotti al Blend
Space.''

La lezione costruisce, un pezzo alla volta, un personaggio giocabile completo. La
\textbf{State Machine} organizza le animazioni in \emph{stati} (camminare e correre,
saltare, cadere) e in \emph{transizioni} con delle condizioni; il \textbf{Blend Space
1D} mescola fermo, camminata e corsa in base alla velocità; una \textbf{camera} con il
suo \textbf{Spring Arm}, i grafi di movimento presi in prestito dal modello Third Person
e un nuovo input per lo \textbf{sprint}, costruito con l'\textbf{Enhanced Input},
rendono il cavaliere controllabile; il \textbf{retargeting} porta sul suo scheletro le
animazioni di salto e di caduta del manichino di Unreal. Dopo la pausa arrivano il
\textbf{Blend Space 2D}, che aggiunge alla velocità la \textbf{direzione} e permette di
camminare di lato e all'indietro, e gli \textbf{Animation Montage}, le animazioni
``una tantum'' (un attacco con la spada) che si inseriscono sopra il ciclo di
locomozione.

La lezione è \textbf{online}, su Zoom, come le precedenti, e la classe è ridottissima:
``vediamo un po' quanti siete: siamo in sei''; dopo la pausa ``siamo rimasti in
cinque'', poi ``uno, due, tre e quattro''; quando uno studente crasha (``Filippo è andato
via''), ``siamo rimasti noi tre''; alla fine, ``mi sembra che siamo rimasti in due''.
``Pochi ma buoni.'' Il ritmo è quello della dimostrazione dal vivo, interrotta più volte
dai problemi degli studenti, risolti in condivisione di schermo, e da un editor che
``oggi lagga'' in modo vistoso (``non capisco perché mi sta laggando oggi''). La pausa
arriva dopo la scelta delle animazioni da scaricare (``ci rivediamo alle 18''); si
chiude con un compito facoltativo e un ``ci sentiamo per martedì prossimo''.

Tre avvertenze di lettura, le stesse delle lezioni precedenti. La prima: la
trascrizione è automatica e storpia moltissimi nomi. Unreal è ``Arreal'', ``a Real'',
``ArrayL'', ``Array Alt'' e ``a Rial''; il \textbf{Pawn Owner} è ``Pwn Honor'' e
``Pwnowner''; il \textbf{Blend Space} è ``Blaine Space'', ``Blendspace'' e ``BlizzBase'';
lo stato \textbf{WalkRun} è ``walk run'', ``Warcran'', ``war crunch'' e ``walker run'';
l'\textbf{Idle} è ``Idol'' e ``idol''; \textbf{BP\_Knight} è ``Vip Night'', ``Pippi
Night 1'' e ``BP Night''; \textbf{BS\_Knight} è ``BS Night'' (``che dovrebbe essere un
po' più British''); il manichino \textbf{Manny/Quinn} è ``Many Quincy''; il
\textbf{root motion} è ``rut'' e ``outmotion''; lo \textbf{Strafe} è ``Strife'' e
``stair right''; la \textbf{Mesh} è ``Mesci''; l'\textbf{FInterp To} è ``finter2'';
\textbf{Calculate Direction} è ``calcolo interaction''; il \textbf{Save Cached Pose} è
``save cache, pause''; lo \textbf{StackOBot} è ``StaccoBot'' e ``stacco bot''; lo
\textbf{Yaw} è ``io'', ``yo'' e ``YAW''; l'\textbf{Orient Rotation to Movement} è
``orientation to movement'', ``anche rotacion'' e ``orient rotace tu mente tu''; lo
\textbf{Use Pawn Control Rotation} è ``us control pawn rotation'' e ``you spawn control
rotation''. Dove la ricostruzione è sicura il capitolo usa la forma corretta; dove non
lo è riporta la forma della trascrizione e lo dice. La seconda: all'inizio della
registrazione e durante la pausa il riconoscimento vocale produce lunghi blocchi di
testo senza senso e ripetuto (``grazie a tutti per l'attenzione che ci ha spiegato'',
``qual è la risposta che ci può essere?'', ``la situazione è questa, la migliore dal
2005'', la ``dottrina sul mercato del lavoro'', la ``Corte Suprema di Internet'', il
governo di Dubai, una lunga fila di ``di Sarajevo'', la ``struttura del giardino'', ``la
blue map'', un Vittorio Prodi chiamato in causa senza motivo): sono artefatti, non
parole di nessuno, e il capitolo li ignora. La terza: gran parte della lezione è una
dimostrazione dal vivo nell'editor, costruita a voce (``tirate fuori'', ``promote to
variable'', ``ce lo attacchiamo qui''), e in più punti il docente improvvisa, prova,
sbaglia e corregge, soprattutto nella parte sulla camera e in quella finale sul
Montage. Il capitolo ricostruisce i grafi da ciò che il docente dice mentre li
costruisce, distingue i tentativi falliti dalla soluzione finale, e completa i nomi
esatti dei nodi e dei pin di Unreal Engine 5; dove la trascrizione lascia un buco, lo
dichiara.

\paragraph{Come leggere questo capitolo.} Il capitolo segue l'ordine della lezione. Le
figure ricostruiscono la State Machine della locomozione, il Blend Space 1D sulla
velocità, la camera con lo Spring Arm, la catena dell'Enhanced Input, il grafo dello
sprint, il cerchio della direzione, il Blend Space 2D con le sue otto animazioni, la
catena di pose dell'AnimGraph con lo Slot del Montage e il grafo del tasto che fa
partire l'attacco. Le tabelle riassumono le transizioni della State Machine, i
parametri della camera, i valori e i trigger di una Input Action, la posizione delle
animazioni nel Blend Space 2D, gli errori incontrati durante la lezione e i nodi della
giornata. Qua e là, dei paragrafi di \textbf{approfondimento} collegano la lezione ai
libri di Marcos Romero e Brenden Sewell (\emph{Blueprints Visual Scripting for Unreal
Engine 5} e la seconda edizione dello stesso testo, scritta per Unreal Engine 4), i
testi di riferimento per questa parte del libro, e a qualche nozione tecnica generale su
Unreal Engine 5; sono segnalati come tali e non fanno parte della lezione.

Il percorso del capitolo:
\begin{enumerate}
  \item l'\textbf{inizio della lezione}: i corsi degli studenti e la lezione libera del
        25 (sezione~\ref{ue9:sec:inizio});
  \item il \textbf{ripasso} dell'Animation Blueprint (sezione~\ref{ue9:sec:ripasso});
  \item la \textbf{State Machine} della locomozione (sezione~\ref{ue9:sec:statemachine});
  \item il \textbf{Blend Space 1D}, con una parentesi sui problemi dell'editor (sezioni
        \ref{ue9:sec:blend1d} e~\ref{ue9:sec:layout});
  \item il personaggio in gioco: \textbf{Default Pawn Class}, \textbf{camera} e
        \textbf{Spring Arm}, grafi di movimento e rotazioni (sezioni
        da~\ref{ue9:sec:pawnclass} a~\ref{ue9:sec:rotazioni});
  \item il \textbf{retargeting} del salto e della caduta, e le \textbf{regole di
        transizione} (sezioni \ref{ue9:sec:retarget} e~\ref{ue9:sec:transizioni});
  \item lo \textbf{sprint} con l'\textbf{Enhanced Input} e il tentativo con FInterp To
        (sezioni da~\ref{ue9:sec:enhanced} a~\ref{ue9:sec:finterp});
  \item le \textbf{animazioni da scaricare}, la pausa e i \textbf{feedback} (sezioni
        \ref{ue9:sec:download} e~\ref{ue9:sec:feedback});
  \item il \textbf{Blend Space 2D} e la \textbf{direzione} (sezioni
        \ref{ue9:sec:direction} e~\ref{ue9:sec:blend2d});
  \item gli \textbf{Animation Montage}: l'asset, i Notify, le sezioni, gli Slot e il
        tasto che li fa partire (sezioni da~\ref{ue9:sec:montage}
        a~\ref{ue9:sec:playmontage});
  \item il \textbf{compito} sullo StackOBot e la chiusura
        (sezione~\ref{ue9:sec:chiusura}).
\end{enumerate}

% ---------------------------------------------------------------------
\section{Prima della lezione: i corsi degli studenti e la lezione del 25}
\label{ue9:sec:inizio}

La registrazione si apre con qualche minuto di attesa, in cui il riconoscimento vocale
produce soprattutto rumore. Poi i saluti (``ragazzi, come state? Buonasera, buonasera''),
un controllo rapido (``nessun problema sulla telecamera'': stavolta la webcam del docente
funziona) e la decisione di aspettare ``altri cinque minuti, così vediamo che tutti siate
in carreggiata''. Uno studente, Tommaso, ha dei problemi di collegamento. Il docente
sistema le finestre (``oggi c'è il secondo schermo'') e parte: ``bando alle ciance''.

\paragraph{I corsi di indirizzo.} Prima di cominciare, il docente chiede ai presenti
quale corso seguiranno ``quando inizierete a settembre'': il corso di Unreal è infatti
un'introduzione comune, a cui seguono gli indirizzi specialistici. Le risposte dei sei
presenti:
\begin{itemize}
  \item \textbf{3D art};
  \item \textbf{game design} (due studenti);
  \item \textbf{narrative design};
  \item \textbf{character animation} (``perfetto, character animation'').
\end{itemize}
La domanda serve a tarare la lezione: chi farà character animation, in particolare, è
il destinatario naturale di quasi tutto quello che si vedrà oggi.

\paragraph{La lezione libera del 25.} ``Stiamo praticamente arrivando a una lezione che
di solito noi lasciamo abbastanza \emph{free}, come argomenti'': la lezione del
\textbf{25} sarà ``fondamentalmente la mia ultima lezione di base''. Se non ci sono
richieste particolari, si faranno ``i menu'', cioè ``la \textbf{UI}'': ``che si intende
la UI: menu, bottoni, scritte eccetera''. In alternativa, ``se ci sono degli argomenti
particolari'', da concordare fra gli studenti, ``possiamo tranquillamente affrontarli'',
oppure ``dividere anche la lezione in due: fare un piccolo menu e poi fare anche la parte
in più''. ``Nel caso, sentitevi tra di voi.'' La richiesta concreta arriverà durante la
pausa (sezione~\ref{ue9:sec:feedback}).

% ---------------------------------------------------------------------
\section{Il ripasso: l'Animation Blueprint del cavaliere}
\label{ue9:sec:ripasso}

``Eravamo rimasti al nostro bellissimo \textbf{ABP\_Knight}.'' Il docente ha il
cavaliere (\emph{Knight}, il paladino di Mixamo); ``altri di voi avevano scelto il
proprio personaggio''. Nella preview dell'Animation Blueprint il personaggio è ancora
``il nostro bel character in T-pose'': l'AnimGraph non decide ancora nessuna
animazione.

\paragraph{Initialize e Update.} L'ultima cosa fatta era l'evento \textbf{Event Blueprint
Initialize Animation}: ``quando l'animazione si inizializza'', una volta sola. L'altro
evento, \textbf{Event Blueprint Update Animation}, scatta invece ogni volta che
``si cambia qualsiasi particolare all'interno di un'animazione, la rotazione di un braccio
o un minimo particolare'': ``quello che è fondamentalmente il nostro Tick sulle
animazioni''. ``Quando abbiamo visto Begin Play e Tick: questi praticamente sono i loro
corrispettivi.''

\paragraph{Il Pawn Owner e le due referenze.} Nell'Initialize si prende il \textbf{Pawn
Owner} con \textbf{Try Get Pawn Owner}. ``Ricordiamoci: il Pawn Owner è l'entità che
controlla il character'': il docente distingue il \emph{pawn}, ``quello che lo controlla'',
e il \emph{character} stesso, ``che ha la mesh, lo scheletro e tutto quanto''. Con il
\textbf{Cast To} BP\_Knight si ottiene una referenza al cavaliere, e si salvano in due
variabili il cavaliere e il suo \textbf{Character Movement Component}: ``queste due
variabili ci serviranno per comunicare in modo più diretto con il nostro cavaliere''.

\paragraph{Il Character Movement.} ``Dentro il Blueprint, qualsiasi character lo troverete
qui sotto, Character Movement'': selezionandolo, il pannello Details a destra ``presenta
tutte le particolarità di movimento'', per esempio la \textbf{gravità} (\emph{Gravity
Scale}), la \textbf{Max Acceleration}, la \textbf{Max Step Height} (l'altezza massima di
un gradino che il personaggio sale senza saltare). ``Vogliamo farlo andare più veloce,
meno veloce? Se volessimo farlo sprintare, dovremmo lavorare su questa particolarità''; e
infatti ``dopo faremo anche uno sprintino molto veloce, per differenziare le
animazioni''.

\paragraph{GroundSpeed e IsFalling.} Nell'Update, dopo una \textbf{Sequence}, si prende
``la velocity attuale del nostro character'', a ogni frame, e se ne salva la lunghezza in
\textbf{GroundSpeed}. Allo stesso modo, con il nodo \textbf{Is Falling} del Character
Movement, ``andiamo a vedere se effettivamente il nostro character sta cadendo o meno'':
``questo blocco ci dice letteralmente se sta cadendo, e noi riportiamo con un booleano,
true o false''. ``Di conseguenza, tramite questi dati, andremo a differenziare le
animazioni nella nostra State Machine.''

% ---------------------------------------------------------------------
\section{La State Machine della locomozione}
\label{ue9:sec:statemachine}

\paragraph{Creare la State Machine.} La State Machine ``la possiamo invocare
nell'AnimGraph'': tasto destro nell'AnimGraph (``ricordatevi che c'è il blocco
\textbf{Output Pose}''), \textbf{State Machine}, e la si chiama \textbf{Locomotion}:
``questo è il nome classico per la locomozione, appunto il movimento''. ``Una State
Machine, appunto, è una macchina a stati.'' Con un doppio clic si entra nel suo grafo.

\paragraph{Entry, Conduit, State, Alias.} Dentro c'è un solo nodo, \textbf{Entry}. ``Se
noi trasciniamo, abbiamo una freccia che esce fuori''; lasciandola nel vuoto, il menu non
è ``tutto l'elenco gigantesco'' dei nodi dei Blueprint, ma tre sole voci:
\textbf{Conduit}, \textbf{State} e \textbf{Alias}. ``A noi adesso interessa vedere
quelli che sono gli stati.'' Il primo stato, invece di \emph{Locomotion}, si chiama
\textbf{WalkRun}: ``differenziamo''.

\paragraph{Un'altra inception.} ``Ecco a voi che, ancora sempre, un po' un
\emph{inception}'': con un doppio clic si entra anche nello stato. Dentro, un nodo
\textbf{Output Animation Pose}: ``in questo stato risiede una tipologia di output'',
l'animazione che lo stato produce. Il docente vi collega \emph{Walking 1}. La preview
però non si aggiorna, ``perché non abbiamo collegato la State Machine all'output'':
bisogna tornare nell'AnimGraph e collegare il nodo Locomotion all'Output Pose. Fatto
questo, il cavaliere cammina, ma una volta sola: ``l'animazione non è in loop'', e basta
spuntare il loop nei Details del nodo. ``Ecco a voi che cosa fa la State Machine'':
``dentro Locomotion lui entra, dice ok, inizio, passo e vado direttamente a WalkRun''; in
WalkRun ``c'è una richiesta di un output, quindi di un'animazione che noi gli diamo''.

\paragraph{Gli stati Jump e Fall.} ``Ma la cosa interessante della State Machine'' è che
si possono creare altri stati: \textbf{Jump} e \textbf{Fall}. ``A livello di logica, la
State Machine vuole essere'' il racconto di ``come entrano e perché entrano'' le
animazioni nei loro stadi. Una freccia da WalkRun a Jump ``vuol dire che dallo stadio
WalkRun passerà allo stadio Jump''; una freccia inversa vuol dire che dal Jump ``potrà
passare di nuovo al WalkRun''. Lo schema di oggi ha tre stati
(figura~\ref{ue9:fig:statemachine}): si entra in WalkRun; ``se starà cadendo'' (la
variabile IsFalling diventa vera) si passa a Jump o a Fall; ``e far sì che poi, quando
lui finisca di cadere, ritorni'' a WalkRun. ``Quindi fondamentalmente creiamo un loop di
animazione secondo delle condizioni.'' Le animazioni di salto e di caduta non ci sono
ancora: ``la faremo, la ritargetteremo'' (sezione~\ref{ue9:sec:retarget}).

\paragraph{Le transizioni.} ``Come facciamo a mettere queste condizioni? Tramite questi
nodini qua'': sulle frecce fra due stati compare un'icona tonda, la \textbf{transition
rule}, che ``ci dice il comportamento effettivo che dovrà avere'' il personaggio. ``Quando
è che passo da WalkRun a Jump, a livello di animazione? Quando, ad esempio, stai cadendo.
Quando passo da Fall a WalkRun? Se non sta più cadendo. Questa è la logica generale.''

\begin{figure}[htbp]
\centering
\begin{tikzpicture}[font=\footnotesize,
  ue9state/.style={draw, thick, rounded corners=6pt, fill=orange!12, align=center,
                   minimum width=2.4cm, minimum height=1.0cm},
  ue9entry/.style={draw, thick, fill=gray!15, rounded corners=2pt, align=center,
                   minimum width=1.4cm, minimum height=0.7cm},
  ue9tr/.style={-{Stealth[length=2.6mm]}, thick},
  ue9rule/.style={draw, circle, fill=white, inner sep=1.5pt, font=\scriptsize}]
  \node[ue9entry] (e) at (0,0) {Entry};
  \node[ue9state] (w) at (3.4,0) {\textbf{WalkRun}\\{\scriptsize Blend Space}};
  \node[ue9state] (j) at (8.2,1.5) {\textbf{Jump}\\{\scriptsize MM\_Jump}};
  \node[ue9state] (f) at (8.2,-1.5) {\textbf{Fall}\\{\scriptsize MM\_Fall\_Loop}};
  \draw[ue9tr] (e) -- (w);
  \draw[ue9tr] (w.north east) to[bend left=12] node[ue9rule, pos=0.5]{1} (j.west);
  \draw[ue9tr] (j.south) -- node[ue9rule]{2} (f.north);
  \draw[ue9tr] (w.south east) to[bend left=12] node[ue9rule, pos=0.5]{3} (f.west);
  \draw[ue9tr] (f.south west) to[bend left=30] node[ue9rule, pos=0.5]{4} (w.south);
  \node[align=left, text width=12.5cm, anchor=north west, font=\scriptsize] at
    (-1,-2.9) {\textbf{1} WalkRun $\rightarrow$ Jump: IsFalling \textbf{AND}
    VelocityZ $>$ 0 (sta salendo)\\
    \textbf{2} Jump $\rightarrow$ Fall: \emph{Automatic Rule Based on Sequence Player in
    State} (finito il salto, la caduta)\\
    \textbf{3} WalkRun $\rightarrow$ Fall: IsFalling \textbf{AND} VelocityZ $<$ 0 (cade
    senza aver saltato)\\
    \textbf{4} Fall $\rightarrow$ WalkRun: \textbf{NOT} IsFalling (è tornato a terra)};
\end{tikzpicture}
\caption{La State Machine Locomotion alla fine della lezione: tre stati e quattro
transizioni. Le condizioni con la velocità sulla Z sono quelle aggiunte nella seconda
versione (sezione~\ref{ue9:sec:transizioni}); nella prima, la 1 e la 3 erano soltanto
IsFalling. Schema ricostruito.}
\label{ue9:fig:statemachine}
\end{figure}

\paragraph{Approfondimento: Conduit e Alias.} I tre tipi di nodo proposti dal menu della
State Machine hanno ruoli diversi. Uno \textbf{State} contiene un grafo che produce una
posa (un'animazione, un Blend Space, una combinazione). Un \textbf{Conduit} non produce
pose: è un incrocio, uno snodo da cui partono più transizioni, utile quando molti stati
condividono la stessa condizione d'ingresso o d'uscita e non si vuole ripeterla su ogni
freccia. Un \textbf{State Alias} è un ``alias'' di uno o più stati: le transizioni che
partono dall'alias valgono come se partissero da ciascuno degli stati che rappresenta
(per esempio, ``da qualunque stato, se il personaggio muore, vai a Death''), e
risparmiano un groviglio di frecce. Romero e Sewell (capitolo 17) costruiscono la loro
prima State Machine, \emph{Char States}, con due soli stati, \emph{Idle} e
\emph{Moving}, e due regole sulla velocità (``Speed maggiore di 5'' e ``minore di 5''); la
scelta della lezione, un solo stato WalkRun che contiene un Blend Space, sposta la
distinzione fra fermo e in movimento dentro il Blend Space, e lascia alla State Machine
le sole differenze ``qualitative'' (a terra, in salto, in caduta).

% ---------------------------------------------------------------------
\section{Il Blend Space 1D}
\label{ue9:sec:blend1d}

``Visto che WalkRun non è solo un'animazione, è un insieme'': sono tre animazioni,
\textbf{Idle}, \textbf{Walk} e quella che il docente ha chiamato \textbf{Slow Run}.
``Dobbiamo mixarle insieme. Come si fa? Perché non possiamo mettere un'animazione solo
così: dobbiamo far sì che sia dipendente dalla nostra velocità.'' Lo strumento è il
\textbf{Blend Space}.

\paragraph{Creare il Blend Space 1D.} ``Fate solamente il Blend Space 1D.'' Invece di
cercarlo nei menu, si può scrivere nella barra di ricerca del menu del tasto destro
``blend space'', e scegliere \textbf{Blend Space 1D}: ``già vi dice il percorso, che sta
in Animation, Legacy''. L'icona ha un colore ``rosa pelle''. Aperto, si troverebbe il
personaggio in T-pose; ma molti l'avevano già creato la volta precedente (``sì, mi sa che
già l'avevamo fatto'').

\paragraph{Che cos'è un Blend Space.} ``Un Blend Space è l'azione di \emph{blendare}
qualche cosa'', di ``interpolare o comunque mixare le animazioni'': un asset che ``ci dà
la possibilità di poter traslare da un'animazione all'altra''. Si vede subito: dal
pannello \textbf{Asset Browser} dell'editor (``se non lo trovate, ripeto: Window, Asset
Browser'') si trascina un'animazione sul grafico; tenendo premuto \textbf{Shift} la si
aggancia ``su dei punti precisi'' della griglia. Con un'animazione sola, il personaggio
fa solo quella; aggiungendone una seconda, il grafico ``crea una linea di passaggio tra
un'animazione e l'altra''.

\paragraph{Il parametro, come l'alfa.} ``Ma importante: qual è questo dato?'' Il docente
lo collega all'\textbf{alfa} di un Lerp visto in una lezione precedente: ``quel puntatore
che punta, dipendente da uno spazio tra due punti, dove sta al momento''. Qui il valore
dirà ``ehi, guarda che sto camminando, sto andando a velocità 300, oppure sto correndo, a
velocità 600, oppure sono fermo''. Il Blend Space \textbf{1D} ``ci dà solamente la
possibilità di poter usare un asse'': nei Details, \textbf{Horizontal Axis} (la sezione
potrebbe essere chiusa: ``basta che la aprite''). L'asse si chiama \textbf{GroundSpeed},
``come quella che abbiamo trovato all'interno dell'Animation Blueprint''; poi chiede un
\textbf{Minimum Axis Value} e un \textbf{Maximum Axis Value}.

\paragraph{Le velocità sono relative.} ``Importante, ragazzi, quando si va a vedere le
animazioni: dovete sempre farlo relativo, se si parla di velocità'', relativo cioè alla
velocità reale del personaggio. Dove si legge? Nel Blueprint del personaggio,
BP\_Knight, selezionando \textbf{Character Movement}: ``potete andare a vedere tutti i
dati che vi servono''. Quello che serve è la \textbf{Max Walk Speed}, in alto, nella
sezione \emph{Character Movement: Walking}: ``è 600 centimetri al secondo''. ``Unreal
ragiona per centimetri'': il personaggio va a 6 metri al secondo.

\paragraph{Dividere la velocità a metà.} ``Visto che noi abbiamo due stadi, uno di walk e
uno di run'', si vuole che ``la walk sia un movimento normale'': di default il personaggio
cammina, e corre solo se il giocatore lo chiede con un tasto (sezione
\ref{ue9:sec:enhanced}). ``Quindi la dividiamo a metà'': \textbf{Max Walk Speed a 300}.
Nel Blend Space l'asse va da \textbf{0} a \textbf{600}, ``in modo tale che a metà la
nostra walk sia a 300'' (figura~\ref{ue9:fig:blend1d}). A uno studente l'asse arriva solo
a 100: ``ground speed, devi metterlo a 600, lo trovi qua''. E un consiglio: dopo aver
cambiato il massimo, ``sposta i pallini, perché poi ti switcha tutto''; il valore di
ciascun campione si può anche scrivere a mano nei Details.

\begin{figure}[htbp]
\centering
\begin{tikzpicture}[font=\footnotesize,
  ue9pt/.style={draw=black, fill=white, thick, diamond, inner sep=2.2pt}]
  \draw[-{Stealth[length=2.5mm]}, thick] (0,0) -- (11.2,0) node[right]{GroundSpeed};
  \foreach \x/\v in {0/0, 2.5/150, 5/300, 7.5/450, 10/600}
    \draw (\x,0.1) -- (\x,-0.1) node[below, font=\scriptsize]{\v};
  \node[ue9pt] (i) at (0,0) {}; \node[ue9pt] (w) at (5,0) {}; \node[ue9pt] (r) at
    (10,0) {};
  \node[above=6pt of i, font=\scriptsize\bfseries] {Idle};
  \node[above=6pt of w, font=\scriptsize\bfseries, align=center] {Walk\\
    {\normalfont Max Walk Speed di default}};
  \node[above=6pt of r, font=\scriptsize\bfseries, align=center] {Slow Run\\
    {\normalfont con lo sprint}};
  \fill[red!60!black] (3.33,0) circle (0.12);
  \draw[red!60!black, thick, ->] (3.33,-1.2) -- (3.33,-0.25);
  \node[red!50!black, font=\scriptsize, align=center] at (3.33,-1.55) {200: un misto
    fra Idle e Walk, ``un obbrobrio''};
\end{tikzpicture}
\caption{Il Blend Space 1D BS\_Knight: un solo asse, la GroundSpeed, da 0 a 600, con
Idle a 0, Walk a 300 (la Max Walk Speed impostata nel Character Movement) e la corsa a
600 (la velocità dello sprint). Un valore intermedio, come 200, mescola due animazioni
nate per velocità diverse. Schema ricostruito.}
\label{ue9:fig:blend1d}
\end{figure}

\paragraph{Il Blendspace Player.} Tornati in ABP\_Knight (``normale, ragazzi, che vi dia
dei warning, perché ancora non abbiamo le condizioni'': ``lui dice, ma io non ho
condizioni, quindi faccio tutto in una botta''), si entra nello stato WalkRun. Nell'Asset
Browser ``è apparso BS\_Knight''; trascinato nel grafo, ``ci presenta un \textbf{Blend
Space Player}'', diverso dal nodo di un'animazione normale: ha un \textbf{valore in
entrata}, la GroundSpeed. A quel pin si collega un \textbf{Get} della variabile
GroundSpeed, calcolata nell'Event Graph (la lunghezza della velocity), e l'uscita del
Blend Space Player si collega al \textbf{Result} dell'Output Animation Pose, ``in modo
tale che vada in output''.

\paragraph{Get con Ctrl, Set con Alt.} Un promemoria: trascinando una variabile nel
grafo, ``tenete premuto \textbf{Ctrl} e tirate fuori direttamente un Get''; ``se invece
volete fare un Set, tenete premuto \textbf{Alt}''. Nel grafo dello stato l'Alt non
funziona: ``non si può settare all'interno'' (gli stati dell'AnimGraph leggono le
variabili, non le scrivono), ``però nei Blueprint, di base, se premete Alt vi setta la
variabile''.

\paragraph{La prova nella preview.} Compilato, ``il nostro character nella preview adesso
sta in idle, perché giustamente ha velocità 0''. Cambiando il valore di GroundSpeed nella
preview: a \textbf{300} cammina; a \textbf{200} ``sta interpolando tra l'animazione di
idle e l'animazione di walk'': ``mettere un valore a 200 non è buona cosa'', perché
``deve essere relativo al vostro movimento, se no vengono degli obbrobri''; a
\textbf{600} ``il nostro carissimo corre'' (``guarda come è eroico''). ``Quindi adesso
andiamo a cambiare le animazioni tramite un valore.''

\paragraph{Allineare i nodi con Q.} ``Se volete fare un po' di ordine'': selezionati due
o più nodi, il tasto \textbf{Q} li allinea (\emph{Straighten Connections}). ``Tutto bello
ordinato.''

\paragraph{Una prova con un solo stato.} Prima di andare avanti, un test in gioco:
nell'AnimGraph si tolgono le frecce verso Jump e Fall (selezionate e cancellate, o
staccate con \textbf{Alt}+clic sul pin), e si lascia ``solamente WalkRun, che vada in
entrata''.

\paragraph{Approfondimento: il Blend Space nel libro.} Romero e Sewell (capitolo 17)
definiscono il Blend Space come ``un tipo di asset che permette di mescolare animazioni
in base ai valori di uno o due parametri''; nel loro esempio, tratto dal modello Third
Person, Idle è a 0, Walk a 93,75 e Run a 375, con una Max Walk Speed di 375. Il principio
è lo stesso della lezione: i campioni vanno posti dove la velocità reale del personaggio
corrisponde alla velocità ``disegnata'' nell'animazione. Se un'animazione di camminata è
stata animata per una velocità di 150 centimetri al secondo e il personaggio si muove a
300, i piedi ``scivolano'' sul pavimento (\emph{foot sliding}); è il motivo pratico
dietro il ``deve essere relativo'' del docente. Gli autori ricordano anche che il valore
del parametro si può cambiare nell'\emph{Anim Preview Editor}, in basso a destra
nell'editor dell'Animation Blueprint, senza andare in Play: è ciò che il docente fa per
mostrare i valori 200, 300 e 600.

% ---------------------------------------------------------------------
\section{Una parentesi: ``questo asset non ha schede ancorate''}
\label{ue9:sec:layout}

Mentre si sistema la Max Walk Speed, uno studente interviene: ``io ho un problema col
Blueprint: quando apro la scheda Blueprint mi dice: questo asset non ha ancora schede
ancorate''. Il docente ha l'editor in inglese (``di solito gli errori me li ricordo in
inglese'') e chiede di vedere: lo studente condivide prima soltanto una scheda del
browser, poi lo schermo intero.

\paragraph{Il messaggio.} Aprendo BP\_Knight, l'editor del Blueprint è vuoto, senza
pannelli. Il testo completo del messaggio, letto ad alta voce, spiega che l'editor
dell'asset è stato aperto in un sistema di finestre fluttuanti e nascoste, ``non mostra
le schede ancorate nell'interfaccia principale'', e che ``questo comportamento può essere
facilmente ripristinato'' ripristinando il layout dell'editor o modificando le preferenze
perché gli asset si aprano sempre ancorati alla finestra principale (in inglese,
\emph{Editor Preferences $\rightarrow$ General $\rightarrow$ Appearance $\rightarrow$
Asset Editor Open Location}). Inoltre, sopra il grafo compare una nota: il Blueprint è
\emph{data-only}, e c'è un link \textbf{Open Full Blueprint Editor}.

\paragraph{La soluzione.} ``Mi sa che è meglio ripristinare il layout di default'':
dalla barra in alto, \textbf{Window} (in italiano \emph{Finestra}) $\rightarrow$
\textbf{Load Layout} (\emph{Carica il layout}) $\rightarrow$ \textbf{Default Editor
Layout} (\emph{Editor predefinito}). Poi, sulla nota, ``apri l'editor Blueprint
completo''. ``Perfetto, eccolo qua.''

\paragraph{La lingua dell'editor.} ``Comunque, per mettere l'editor in inglese'' senza
riavviare: si cerca \emph{Language} nelle preferenze. Il docente prova prima in
\textbf{Editor Preferences} (sotto \emph{Edit}), scrivendo ``lingua'' e poi ``language''
(``non gli piace che mettiamo due cose insieme''), poi nei \textbf{Project Settings};
alla fine la voce \emph{Language} viene trovata con la ricerca (``ah, ecco perché'').
Dalla trascrizione non si capisce con certezza in quale delle due finestre. Lo studente
ringrazia (``mi hai fatto sentire molto ignorante''; ``verrai buttato fuori dalla
classe, tranquillo''), e il docente annuncia di volerlo usare in inglese anche lui ``per
la prossima volta''.

\paragraph{Approfondimento: dov'è la lingua.} In Unreal Engine 5 la lingua
dell'interfaccia è in \textbf{Edit $\rightarrow$ Editor Preferences $\rightarrow$ General
$\rightarrow$ Region \& Language $\rightarrow$ Editor Language}; nei Project Settings c'è
invece la lingua del \emph{gioco} (la localizzazione), che è un'altra cosa. Lavorare
con l'editor in inglese è una scelta pratica: la documentazione, i tutorial, i libri e i
forum usano i nomi inglesi dei nodi e dei pannelli, e la ricerca dei nodi nel menu del
tasto destro accetta i nomi inglesi anche quando l'interfaccia è tradotta.

% ---------------------------------------------------------------------
\section{Il cavaliere in gioco: la Default Pawn Class}
\label{ue9:sec:pawnclass}

``Adesso vi faccio vedere come impostiamo il player nuovo.'' Si aprono i \textbf{World
Settings} del livello (``qualsiasi cosa che non trovate, c'è sempre Window, e troverete
World Settings''); fra le voci della Game Mode c'è \textbf{Default Pawn Class}.
``Ricordiamo che il Pawn è appunto l'entità che controlla il player.'' Nel menu a
tendina c'è BP\_Knight, il personaggio finito che il docente aveva mostrato alla fine
della lezione precedente, e c'è BP\_Knight1, ``quello nuovo'', che si può anche cercare
scrivendone il nome. Si sceglie BP\_Knight1 e si va in Play.

\paragraph{Il primo problema: nessuna camera.} ``Ecco già un problemino.'' ``Fate
attenzione, perché i character non vengono con la camera'': un Blueprint creato dalla
classe Character ha la capsula, la freccia, la mesh e il Character Movement, ma non una
camera. In Play, la visuale finisce in un punto qualsiasi del corpo (``questo
meraviglioso\dots'', la frase si perde). Si esce dal Play e si apre BP\_Knight1: ``come
potete vedere, non c'è una camera''.

% ---------------------------------------------------------------------
\section{La camera: Camera e Spring Arm}
\label{ue9:sec:camera}

``La camera è impostata su due cose particolari, ragazzi, molto importanti'': la
\textbf{camera} stessa e il suo \textbf{braccio}. ``Immaginatelo come lo stand'': il
braccio dice ``quanto è lontana'' la camera, ``quanto offset ha''.

\paragraph{Aggiungere i due componenti.} Si aggiunge un componente \textbf{Camera}
(``camera normale''): ``eccola qua, che si è già attaccata'' al personaggio, ma in una
posizione che non va bene (``al suo lato B''). Poi uno \textbf{Spring Arm} (il docente
lo chiama prima ``arm length'', poi si corregge: ``perdonate, l'arm length è un
parametro; lo Spring Arm''). La gerarchia è importante: \textbf{lo Spring Arm deve avere
la camera dentro di sé}, cioè la Camera va trascinata sopra lo Spring Arm e diventa sua
figlia. Selezionando lo Spring Arm, nel viewport compare ``un piccolo tracciato rosso'':
``la lunghezza del nostro braccio'', il braccio della camera
(figura~\ref{ue9:fig:springarm}).

\paragraph{I 3C.} ``Ricordatevi, ragazzi, molto importante, perché la camera è
importantissima: ci sono delle figure che si chiamano \textbf{3C}, \emph{Camera, Control
e Character}, che lavorano solo su questo.'' Da qui una regola: ``se volete allontanare la
camera o cambiarle l'offset, \textbf{non cambiate la location della camera}. Non lo fate.
Controllatelo sempre da qui, dallo Spring Arm.''

\paragraph{Target Arm Length.} Il parametro \textbf{Target Arm Length} ``ci permette di
aumentare o diminuire la distanza della camera''. Valore scelto: \textbf{300}.

\paragraph{Target Offset e Socket Offset.} Ci sono due offset, e il docente ha una
preferenza netta. Il \textbf{Target Offset} ``usatelo solamente sulla Z, di poco,
veramente di poco, massimo 20''; ``e fate tutto con il \textbf{Socket Offset}''. Il
motivo: ``se voi cambiate l'offset del target, la camera vi può più facilmente
compenetrare con le cose''. Con il Socket Offset, invece, ``vedrete che la camera va sul
muro e torna'': ``va più vicino'', cioè ``riduce il suo arm length, fondamentalmente''.
Dopo la prima prova in Play (``lo testiamo, ragazzi?''), la camera è un po' alta: il
Socket Offset sulla Z viene abbassato a \textbf{40}.

\begin{figure}[htbp]
\centering
\begin{tikzpicture}[font=\footnotesize]
  % capsula
  \draw[thick, fill=blue!8, rounded corners=14pt] (0,0) rectangle (1.1,2.6);
  \node[align=center, font=\scriptsize] at (0.55,1.1) {Cap-\\sule};
  \node[font=\scriptsize, below] at (0.55,-0.2) {BP\_Knight1};
  \fill (0.55,2.0) circle (0.06);
  \node[font=\scriptsize, anchor=north west] at (0.62,1.95) {pivot};
  % braccio
  \draw[very thick, red!70!black] (0.55,2.0) -- (-4.2,2.0);
  \node[red!60!black, font=\scriptsize, above] at (-1.8,2.0) {Target Arm Length
    = 300};
  % socket offset
  \draw[thick, dashed, -{Stealth[length=2mm]}, green!40!black] (-4.2,2.0) --
    (-4.2,2.8);
  \node[green!40!black, font=\scriptsize, anchor=east, align=right] at (-4.9,2.3)
    {Socket Offset\\Z = 40};
  % camera
  \draw[thick, fill=gray!30] (-4.75,2.6) rectangle (-3.95,3.05);
  \draw[thick, fill=gray!30] (-3.95,2.7) -- (-3.6,2.55) -- (-3.6,3.1) -- (-3.95,2.95);
  \node[font=\scriptsize, above] at (-4.3,3.1) {Camera (figlia)};
  % target offset
  \draw[thick, dashed, -{Stealth[length=2mm]}, orange!70!black] (0.55,2.0) --
    (0.55,2.35);
  \node[orange!60!black, font=\scriptsize, anchor=south, align=center] at (0.55,2.4)
    {Target Offset: solo Z,\\al massimo $\approx$ 20};
  % muro
  \fill[gray!40] (-2.6,-0.2) rectangle (-2.35,1.4);
  \draw[thick] (-2.6,-0.2) rectangle (-2.35,1.4);
  \node[font=\scriptsize, align=center] at (-2.47,-0.6) {muro};
  \draw (-6,-0.2) -- (3.5,-0.2);
  % gerarchia
  \node[draw, thick, rounded corners=3pt, fill=gray!8, align=left, font=\scriptsize,
    anchor=north west] at (3.3,2.2) {\textbf{Components}\\
    Capsule Component\\
    \quad Arrow\\
    \quad Mesh\\
    \quad \textbf{Spring Arm}\\
    \qquad \textbf{Camera}\\
    Character Movement};
\end{tikzpicture}
\caption{La camera del cavaliere. La Camera è figlia dello Spring Arm, e la distanza e
l'altezza si regolano sul braccio, mai spostando la camera. Se fra il personaggio e la
camera si mette un ostacolo (il muro), lo Spring Arm accorcia il braccio e la camera
``torna'' verso il personaggio; con un Target Offset grande, invece, è più facile che la
camera entri negli oggetti. Schema ricostruito, valori della lezione.}
\label{ue9:fig:springarm}
\end{figure}

\paragraph{Il Field of View.} Più avanti, mentre aspetta uno studente, il docente mostra
che ``nella camera è presente anche un po' di dati con cui potete giocarci'', per esempio
il \textbf{Field of View}. Lo abbassa a \textbf{80}: ``sappiate che tanti giochi horror
usano il field of view a 80. Questo è un piccolo tip.''

\begin{table}[htbp]
\centering
\small
\begin{tabular}{@{}p{3.6cm}p{2.0cm}p{8.4cm}@{}}
\toprule
\textbf{Parametro} & \textbf{Valore} & \textbf{Effetto e consiglio del docente} \\
\midrule
Spring Arm $\rightarrow$ Target Arm Length & 300 & distanza della camera dal
personaggio \\
Spring Arm $\rightarrow$ Socket Offset & Z = 40 & spostamento della camera alla fine del
braccio; da preferire \\
Spring Arm $\rightarrow$ Target Offset & Z $\le$ 20 & spostamento dell'origine del
braccio; poco e solo sulla Z, perché fa compenetrare la camera \\
Spring Arm $\rightarrow$ Use Pawn Control Rotation & attivo & il braccio ruota con il
mouse (il controller); va sullo Spring Arm, non sulla Camera \\
Camera $\rightarrow$ Field of View & 80 & angolo di visuale; 80 è tipico dei giochi
horror \\
Camera $\rightarrow$ Location & non toccare & la posizione si regola sempre dallo Spring
Arm \\
\bottomrule
\end{tabular}
\caption{I parametri della camera toccati durante la lezione.}
\label{ue9:tab:camera}
\end{table}

\paragraph{Approfondimento: perché lo Spring Arm si accorcia.} Lo \textbf{Spring Arm
Component} fa, a ogni frame, una traccia (un \emph{sweep} di una piccola sfera, di
raggio \emph{Probe Size}) dal suo punto d'origine verso la posizione desiderata della
camera; se incontra un ostacolo nel canale \emph{Probe Channel} (di default
\emph{Camera}), accorcia il braccio fino al punto di impatto, e lo riallunga quando
l'ostacolo sparisce (la proprietà è \emph{Do Collision Test}). Il \emph{Target Offset}
sposta l'\emph{origine} di quella traccia, nello spazio del mondo: spostarla molto può
portare l'origine dentro un muro o dietro un oggetto, e allora il test non ``vede'' più
l'ostacolo; il \emph{Socket Offset} sposta invece la \emph{fine} del braccio, nello
spazio del braccio stesso, e la traccia resta coerente. È la ragione tecnica del
consiglio del docente. Lo Spring Arm ha anche le opzioni \emph{Enable Camera Lag} ed
\emph{Enable Camera Rotation Lag}, che fanno ``inseguire'' la camera con un leggero
ritardo e rendono più morbidi i movimenti: il docente le sfiora più tardi, nella sezione
del Montage. Quanto ai 3C, nell'industria la sigla indica la disciplina che cura come il
giocatore \emph{vede} (camera), \emph{comanda} (controlli) e \emph{abita} (personaggio)
il gioco: tre cose che si progettano insieme, perché ognuna cambia la percezione delle
altre.

% ---------------------------------------------------------------------
\section{I grafi di movimento, presi in prestito}
\label{ue9:sec:movimento}

In Play, la camera c'è, ma ``non ci possiamo muovere'': ``perché non abbiamo né input né
le funzioni di base di movimento''. ``Questo, ragazzi, non andiamo a farlo direttamente a
mano, ma andiamo a fare quella che è una bellissima'' scorciatoia (la trascrizione dice
``mandragata''): si apre il \textbf{BP\_ThirdPersonCharacter} classico, ``quello che
abbiamo usato'' finora, che contiene ``questi tre principali blocchi che Unreal fa di
base'': il controllo della \textbf{camera} (\emph{Look}), il controllo del
\textbf{movimento} (con \emph{Add Movement Input}) e il \textbf{salto} (\emph{Jump}).

\paragraph{Copiare anche l'Input Mapping.} ``Mi raccomando, prendete anche questo qua
sopra, che è l'\textbf{input mapping}'', cioè il grafo sull'Event Begin Play che aggiunge
l'Input Mapping Context (``andremo a spiegare che cosa si intende''). Tutto con
\textbf{Ctrl+C}, e poi \textbf{Ctrl+V} nell'Event Graph di BP\_Knight1: ``perché queste
sono le funzioni basiche di Unreal''.

\paragraph{IA\_Move.} Un'occhiata all'evento \textbf{IA\_Move} con l'Add Movement Input,
``proprio per fare una lezione sull'input'': ``questo evento dice che è legato a un
tasto, o a dei tasti, che ci riporta due valori, X e Y'', che ``di conseguenza diventano
dei vettori che vanno a creare il movimento sui due assi''. Uno è avanti e indietro
(\emph{forward}, \emph{backward}), l'altro è destra e sinistra.

\paragraph{La camera che ruota male.} Compilato, ``adesso la nostra camera si muove'', ma
in modo strano: si muove ``solo così'', su un asse. ``L'avevo visto un'altra volta'':
``devo usare il control rotation''. La spunta è \textbf{Use Pawn Control Rotation}:
``se lo attaccate, la camera vi va in \emph{free mode}'', ``nel senso che usa la
rotazione del controller'', sia in beccheggio sia in imbardata.

% ---------------------------------------------------------------------
\section{Tre spunte sulla rotazione}
\label{ue9:sec:rotazioni}

\paragraph{Use Pawn Control Rotation, sullo Spring Arm.} Uno studente ha la camera che
``quando vado in Play posso abbassarla e guardare solo il terreno, non vedo proprio
neanche più il personaggio''. Il docente pensa prima a una compenetrazione, poi si fa
condividere lo schermo: ``che camera hai messo? Camera normale, giusto?''. Guardando i
componenti, il problema salta fuori: ``non Use Pawn Control Rotation sulla camera: Use
Pawn Control Rotation sullo \textbf{Spring Arm}''. ``Quindi levalo da lì e mettilo sullo
Spring Arm.'' ``Anche l'altro?'' ``No, no, no: solo sullo Spring Arm.'' La spunta va su
\emph{un} componente solo, il braccio: è il braccio a ruotare intorno al personaggio, e
la camera, che ne è figlia, lo segue.

\paragraph{Orient Rotation to Movement.} In Play, il cavaliere cammina (alla velocità di
300), ma andando a destra ``lui cammina dritto, fondamentalmente'': si sposta di lato
continuando a guardare avanti. ``Questo lo possiamo risolvere con'' \textbf{Orient
Rotation to Movement}, sempre nel Character Movement: ``dal nome, orienta la rotazione al
movimento''. Accanto c'è la \textbf{Rotation Rate}, la velocità con cui il personaggio si
gira (``non è proprio precisissima, però ci piace'').

\paragraph{Use Controller Rotation Yaw.} Il personaggio però continua a non girarsi bene.
``Perché mi ricordo che qua c'era una particolarità che non ricordo'': il docente scorre
le categorie del Blueprint (\emph{Simulate}, \emph{Rotation}, \emph{Character},
\emph{Orientation}, \emph{Physics}, \emph{Controller}\dots) e la trova. ``Per fissare il
nostro character che non si girava, basta che eliminiamo il tick'': selezionato il Pawn
(la radice del Blueprint, \emph{BP\_Knight1 (Self)}), si cerca \emph{yaw} e si toglie
\textbf{Use Controller Rotation Yaw}, ``in modo tale che il character non usi il suo
yaw''. Senza questa spunta, il personaggio non ruota più insieme alla camera: ``andrà
dove siamo direzionati'', ``se io vado a destra, lui si gira completamente verso
destra''. A uno studente non basta: aveva tolto lo Yaw ma non attivato Orient Rotation to
Movement. ``Ok, non ce l'avevo.'' Con tutte e due, ``abbiamo il nostro bel character
fatto per benissimo''.

\paragraph{Approfondimento: chi comanda la rotazione.} Le tre spunte della lezione sono
le tre ``mani'' che possono ruotare un personaggio in terza persona, e vanno combinate
con criterio. \textbf{Use Pawn Control Rotation} (sullo Spring Arm) fa seguire alla
\emph{camera} la rotazione del controller, cioè del mouse o della levetta destra.
\textbf{Use Controller Rotation Pitch/Yaw/Roll} (sul Pawn) fa seguire quella stessa
rotazione al \emph{corpo} del personaggio. \textbf{Orient Rotation to Movement} (nel
Character Movement) fa girare il corpo verso la direzione in cui si muove. La
combinazione ``da avventura'' della prima parte della lezione è: camera libera (Use Pawn
Control Rotation sul braccio), corpo indipendente dalla camera (Use Controller Rotation
Yaw spento), corpo orientato al movimento (Orient Rotation to Movement acceso): è quella
del modello Third Person di Unreal. La combinazione ``da combattimento'' della seconda
parte, con il Blend Space 2D (sezione~\ref{ue9:sec:direction}), è opposta: il corpo
guarda sempre dove guarda la camera (Use Controller Rotation Yaw acceso), e le
animazioni di strafe mostrano lo spostamento laterale. Se sono accese insieme Use
Controller Rotation Yaw e Orient Rotation to Movement, le due ``mani'' si contendono il
personaggio; per questo, nel Character Movement, esiste anche \emph{Use Controller
Desired Rotation}, che fa ruotare il corpo verso la direzione della camera in modo
graduale, con la Rotation Rate (forse il ``desired'' che il docente prova a un certo
punto, senza esito chiaro nella trascrizione).

% ---------------------------------------------------------------------
\section{Il retargeting: il salto del manichino sul cavaliere}
\label{ue9:sec:retarget}

``Adesso, però, guardate una cosa: se ad esempio provo a saltare, non ho nessuna
animazione.'' Il salto funziona (lo fa il Character Movement), ma il cavaliere sale e
scende nella posa della camminata. Le animazioni di salto si potrebbero prendere da
Mixamo, ``anche se potremmo, volendo''; invece il docente usa ``l'opzione di Unreal che
si chiama \textbf{retargeting}''. ``Diciamo che il retargeting, sappiatelo, può uscire
bene, ma può uscire anche abbastanza male: è dipendente dall'animazione. Però andiamo a
farlo comunque, per vedere come si fa.''

\paragraph{I passaggi.}
\begin{enumerate}
  \item Nel Content Browser si cerca \emph{Jump}: si trova l'animazione del salto del
        manichino base di Unreal, \textbf{MM\_Jump}.
  \item Tasto destro sull'animazione $\rightarrow$ \textbf{Retarget Animations}.
  \item La finestra mostra lo scheletro \textbf{sorgente} (\emph{source}), quello da cui
        viene l'animazione: il manichino di Unreal, \textbf{Manny} o \textbf{Quinn}
        (``Many Quincy''). Come \textbf{target} si sceglie lo scheletro del cavaliere, il
        paladino di Mixamo: ``eccolo qua, li mette a confronto''.
  \item Si seleziona \textbf{MM\_Jump} nella lista delle animazioni: nella preview il
        paladino salta. ``Fortunatamente Mixamo usa degli scheletri che sono abbastanza
        compatibili con quello di Unreal'', e il risultato ``è venuto abbastanza bene'';
        ``potete girarci tranquillamente intorno con la viewport''.
  \item Due pulsanti di esportazione. \textbf{Export Retargeted Assets} ``non ci serve,
        ragazzi'': ``vi crea un asset retargettato'', una copia completa degli asset sullo
        scheletro di destinazione. Basta \textbf{Export Animations}: ``esportare
        l'animazione così, come nuova''. ``Ricordate che c'è il source, e vi esporta il
        target.''
  \item Come cartella si sceglie quella del personaggio (\emph{Characters $\rightarrow$
        Knight1 $\rightarrow$ Animation}) e si esporta.
\end{enumerate}
Lo stesso procedimento si ripete per la caduta: \textbf{MM\_Fall\_Loop}, un'animazione
``in loop'' di caduta, retargettata sul paladino ed esportata nella stessa cartella.

\paragraph{Il root motion ``rotto''.} Il risultato del salto ``non è fantastico'':
``il problema di Mixamo è che ha il \textbf{root motion} un po' rotto''. Il docente
mostra il \emph{root}, ``un punto dello scheletro'', l'osso alla base della gerarchia, e
osserva che lo spostamento non è quello atteso; ``potremmo vederlo'', ``ma per vedere come
funziona il retarget va benissimo''. Si potrebbe sistemare in seguito.

\paragraph{Approfondimento: il retargeting in Unreal Engine 5.} Il retargeting trasferisce
un'animazione da uno scheletro a un altro con proporzioni e nomi delle ossa diversi. In
Unreal Engine 5 si basa su due tipi di asset: l'\textbf{IK Rig}, che descrive uno
scheletro raggruppando le ossa in \emph{catene} (spina dorsale, braccia, gambe, testa)
con un osso di \emph{retarget root} (di solito il bacino), e l'\textbf{IK Retargeter},
che mette in corrispondenza le catene di due IK Rig e corregge la differenza fra le pose
di riferimento (la A-pose del manichino e la T-pose di Mixamo). Nelle versioni più recenti il
comando \emph{Retarget Animations} del menu contestuale apre una finestra che crea al
volo questi asset quando gli scheletri sono riconosciuti (come quelli di Mixamo e del
manichino), ed è ciò che si vede a lezione; per risultati migliori, o con scheletri
insoliti, si creano e si regolano a mano. Il problema del root di Mixamo è noto: gli
scheletri di Mixamo non hanno un osso di root a terra, e lo spostamento del personaggio
è registrato sull'osso dei fianchi (\emph{Hips}); un'animazione retargettata può quindi
avere il movimento nel posto ``sbagliato'' per il sistema di root motion di Unreal.

% ---------------------------------------------------------------------
\section{Le regole di transizione}
\label{ue9:sec:transizioni}

Con le due animazioni nuove si aggiorna l'Animation Blueprint: ``dentro il nostro Jump
ci mettiamo l'animazione di jump'' (``niente di particolare, ragazzi''), dentro Fall
l'animazione di caduta. Si ricollegano le frecce. ``Se cammina e corre, può saltare? Sì,
certo. Se salta, dovrà per forza almeno loopare in caduta'', e ``dal Fall dovrà ritornare
al nostro WalkRun''.

\paragraph{Prima versione.} Nelle icone delle transizioni c'è un grafo con un nodo
\textbf{Result}, che ``ci dà un bool'': ``ci dice quando è che posso entrare in questa
animazione''.
\begin{itemize}
  \item \textbf{WalkRun $\rightarrow$ Jump}: un \emph{Get IsFalling} collegato al
        Result. ``Potremmo farlo più preciso dopo''; per esempio, ``possiamo anche dire:
        ha iniziato il jump''.
  \item \textbf{Jump $\rightarrow$ Fall}: qui il grafo vuoto dà errore, ``perché noi
        vogliamo far sì che questo processo tra il Jump e il Fall sia automatico''.
        Selezionata l'icona della transizione, nei Details si spunta \textbf{Automatic
        Rule Based on Sequence Player in State}: ``quando inizia quell'animazione, a
        seguito ci sarà quest'altra''; in altre parole, la transizione scatta da sola
        quando l'animazione dello stato Jump finisce. ``Lo mettiamo così e non ci darà
        problemi.''
  \item \textbf{Fall $\rightarrow$ WalkRun}: ``vuol dire che non sta cadendo'': Get
        IsFalling, poi un \textbf{NOT Boolean}. ``Ricordo che il NOT Boolean è qualcosa
        che ci dà il suo contrario'': deve essere false perché la caduta finisca.
\end{itemize}
Uno studente chiede cosa ci sia fra Jump e Fall: ``se tu clicchi su questo pallino, ci
dice, in Details\dots\ vogliamo che questa animazione parta direttamente se effettivamente
il jump è finito'': \emph{Automatic Rule Based on Sequence Player in State}, ``questo
tick, lo metti vero''.

\paragraph{La prova, e un difetto.} Si prova da un punto alto, ``la nostra piattaforma''.
``Non si nota tantissimo, però effettivamente sta facendo giustamente l'azione di
saltare.'' Ma il docente vuole collegare WalkRun anche direttamente a Fall: il
personaggio ``può cadere anche nel caso in cui lui non abbia spinta, semplicemente cade da
qualche parte'', per esempio scendendo da un bordo. Con le sole regole su IsFalling,
però, le due transizioni in uscita da WalkRun avrebbero la stessa condizione, e Unreal
non saprebbe distinguere un salto da una caduta.

\paragraph{Seconda versione: la velocità sulla Z.} ``Dovremmo andare a differenziare un
altro valore: la Z. Se la Z è maggiore di 0, vuol dire che sta saltando; se la Z invece è
minore di 0, vuol dire che sta cadendo.'' Nell'Event Graph dell'ABP, finora, della
velocity si prendeva solo la lunghezza. Il docente prova prima la strada sbagliata (la
\emph{Jump Z Velocity} del Character Movement, che però è la spinta \emph{iniziale} del
salto, un valore fisso), poi prende la \textbf{Velocity} già salvata, ne fa lo
\textbf{Split Struct Pin} (``lo splittiamo'') e dalla \textbf{Z} fa \textbf{Promote to
Variable}: una nuova float, \textbf{VelocityZ} (``Z va benissimo''). Nelle transizioni:
\begin{itemize}
  \item \textbf{WalkRun $\rightarrow$ Jump}: IsFalling \textbf{AND} VelocityZ maggiore
        di zero, con un nodo \textbf{AND Boolean}. Il docente dice anche ``quando la
        velocità sarà maggiore di 1, più preciso'': una soglia piccola invece di zero
        esatto.
  \item \textbf{WalkRun $\rightarrow$ Fall}: IsFalling e VelocityZ minore di zero.
\end{itemize}
Per tenere in ordine i fili, si possono usare i \emph{reroute node} (``volete collegarli
con questo? È un po' bruttino''). Il risultato: ``se io salto, fa l'animazione''; ``se io
adesso invece cado solamente, lui va diretto in Fall''. Il docente chiude salvando anche
il BP\_ThirdPersonCharacter che l'editor segnala come modificato (``va bene, grazie'').

\begin{table}[htbp]
\centering
\small
\begin{tabular}{@{}p{3.2cm}p{5.3cm}p{5.5cm}@{}}
\toprule
\textbf{Transizione} & \textbf{Regola} & \textbf{Significato} \\
\midrule
Entry $\rightarrow$ WalkRun & nessuna & lo stato iniziale \\
WalkRun $\rightarrow$ Jump & IsFalling AND VelocityZ $>$ 0 & è in aria e sta salendo: ha
saltato \\
Jump $\rightarrow$ Fall & Automatic Rule Based on Sequence Player in State & finita
l'animazione di salto, passa alla caduta \\
WalkRun $\rightarrow$ Fall & IsFalling AND VelocityZ $<$ 0 & è in aria e sta scendendo:
è caduto da un bordo \\
Fall $\rightarrow$ WalkRun & NOT IsFalling & è tornato a terra \\
\bottomrule
\end{tabular}
\caption{Le regole di transizione della State Machine Locomotion (versione finale).}
\label{ue9:tab:transizioni}
\end{table}

\paragraph{Approfondimento: altre regole basate sul tempo.} L'opzione automatica della
lezione è il modo più rapido per dire ``quando l'animazione finisce''. Romero e Sewell
(capitolo 17), per gli stati \emph{StandToProne} e \emph{ProneToStand}, usano invece una
regola esplicita: il nodo \textbf{Time Remaining (ratio)} dell'animazione dello stato,
seguito da un \emph{Less} con 0,1: la transizione scatta quando manca meno del 10\%
dell'animazione. Il vantaggio è il controllo: la transizione può partire un po' prima
della fine e lasciare spazio al \emph{blend} fra i due stati, che ogni transizione ha
(nei Details: \emph{Blend Settings}, con una durata di default di 0,2 secondi).
Ricordano anche che le animazioni ``di una volta'', come il salto, vanno impostate senza
\emph{Loop Animation}: è il caso dello stato Jump, mentre Fall, che deve durare quanto
la caduta, resta in loop.

% ---------------------------------------------------------------------
\section{Lo sprint con l'Enhanced Input}
\label{ue9:sec:enhanced}

``Adesso andremo a mettere il nostro sprintino'': ``un semplicissimo sistemino che passa
la Max Walk Speed da un tot a un altro tot''. Si torna in BP\_Knight1, ``a ragionare a
livello di Blueprint''. ``Ma sì, facciamo anche l'input mapping.''

\paragraph{L'Input Mapping Context.} ``Unreal ha introdotto, quasi dalla 5.2'', quello
che è l'\textbf{Input Mapping Context}, ``una cosa molto interessante di Unreal che ci
permette, ad esempio, di avere degli eventi legati a più input''. Finora, in molte
lezioni, si sono usati gli eventi dei tasti (``abbiamo visto il tasto E, eccolo qua''):
``semplicemente un rilevamento, da parte di Unreal, della tastiera: se è premuto o non è
premuto''. Adesso, invece, ``gli eventi sono legati a più controlli possibili''.

Nel Content Browser si cerca \textbf{IMC} (\emph{Input Mapping Context}). ``È un
\emph{contesto}: un contesto di input che vengono dati in un certo momento, e sappiate
che sono sovrascrivibili.'' L'esempio: ``io devo cambiare dal gioco di avventura 3D al
gioco di guida della macchina''; ``i controlli sono completamente diversi, quindi cambio
contesto'': un Input Mapping Context per la macchina e uno ``per il character che cammina
normalmente''. Aperto \textbf{IMC\_Default}, ci sono i \textbf{Mappings}:
\textbf{IA\_Jump}, \textbf{IA\_Move}, \textbf{IA\_Look}.

\paragraph{L'Input Action IA\_Run.} ``Per far sì di avere un'Input Action, dobbiamo creare
un asset chiamato \textbf{Input Action}.'' Cercando \emph{IA} nel Content Browser ``troverete
un sacco di asset, ma a voi vi interessa questo celestino''. Con il tasto destro,
\textbf{Show in Folder View}, ``in modo tale che voi troviate direttamente la cartella''.
Si duplica IA\_Jump (Ctrl+C, Ctrl+V, o \emph{Duplicate}) e lo si rinomina
\textbf{IA\_Run}.

\paragraph{Dentro un'Input Action.} Aperto IA\_Run, ``ci sono un po' di valori che
possiamo vedere'':
\begin{itemize}
  \item \textbf{Description}: ``semplice descrizione, quindi quello che fa'';
  \item \textbf{Trigger when Paused}: se l'azione deve scattare anche quando il gioco è in
        pausa, ``quindi si blocca il tempo'';
  \item \textbf{Value Type}: ``la tipologia di valore che ci ritorna indietro. Un input ci
        dà un valore'': \textbf{Digital (bool)}, \textbf{Axis1D (float)} (``come abbiamo
        visto, avanti e indietro''), \textbf{Axis2D (Vector2D)} (``più assi, due
        dimensioni'') o \textbf{Axis3D (Vector)} (``tre tipologie di input che ci vengono
        dati''). ``A noi interessa il Run, che è un bool, quindi premuto o non premuto'':
        Digital (bool) ``va benissimo'';
  \item \textbf{Triggers}: ``le azioni che sono già preinstallate in Unreal'', che dicono
        ``ehi, guarda, si triggera se è premuto, o se è rilasciato, o se è sempre
        premuto''. L'elenco letto dal docente: \emph{Chorded Action}, \emph{Combo},
        \emph{Down}, \emph{Hold}, \emph{Hold and Release}, \emph{Pressed}, \emph{Pulse},
        \emph{Released}, \emph{Tap} (la trascrizione dice anche ``Short Deduction'' e
        ``repeat'', non identificabili con certezza). ``Fondamentalmente, questo vuol
        dire che quando si compie questa azione, questo input viene triggerato.'' IA\_Run
        eredita da IA\_Jump la sua configurazione, ``e va benissimo così, perché tanto è
        un'azione molto simile al jump''.
\end{itemize}

\begin{table}[htbp]
\centering
\small
\begin{tabular}{@{}p{3.2cm}p{10.8cm}@{}}
\toprule
\textbf{Voce} & \textbf{Significato} \\
\midrule
Digital (bool) & premuto / non premuto: un pulsante (salto, corsa, interazione) \\
Axis1D (float) & un valore su un asse: un grilletto analogico, avanti/indietro \\
Axis2D (Vector2D) & due assi: il movimento (X, Y) o la visuale del mouse \\
Axis3D (Vector) & tre assi: input 3D (per esempio un controller di movimento) \\
\midrule
Down & scatta a ogni frame finché il tasto è giù \\
Pressed & scatta una volta, quando il tasto viene premuto \\
Released & scatta una volta, quando il tasto viene rilasciato \\
Hold & scatta dopo che il tasto è stato tenuto premuto per un certo tempo \\
Hold and Release & scatta al rilascio, se il tasto è stato tenuto abbastanza \\
Tap & scatta se il tasto viene premuto e rilasciato entro un tempo breve \\
Pulse & scatta a intervalli regolari mentre il tasto è premuto \\
Chorded Action & scatta solo se un'altra azione è attiva insieme (una ``combinazione'') \\
Combo & scatta se una sequenza di azioni viene eseguita nell'ordine \\
\bottomrule
\end{tabular}
\caption{I tipi di valore (Value Type) e i trigger di un'Input Action, letti a lezione
dall'elenco dell'editor. Le descrizioni sono un'integrazione: a lezione il docente
legge i nomi senza spiegarli uno per uno.}
\label{ue9:tab:inputaction}
\end{table}

\paragraph{Il contesto va caricato.} Tornati in BP\_Knight1: ``abbiamo già il contesto
caricato'', con il grafo copiato dal ThirdPersonCharacter. ``Sappiate che se voi non
caricate il contesto, non vi prenderà nessuna cosa'': sull'Event Begin Play, dal Player
Controller si prende l'\emph{Enhanced Input Local Player Subsystem} e si chiama
\textbf{Add Mapping Context} con IMC\_Default (figura~\ref{ue9:fig:enhanced}).

\paragraph{L'evento IA\_Run e i suoi pin.} Nel grafo si invoca \textbf{IA\_Run}; ``se lo
estendete'' (la freccina in basso del nodo), compaiono tutti i pin:
\begin{itemize}
  \item i pin di esecuzione \textbf{Triggered} (``quando è triggerato''), \textbf{Started}
        (``quando inizia a essere triggerato''), \textbf{Ongoing}, \textbf{Canceled},
        \textbf{Completed};
  \item i valori: l'\textbf{Action Value} (``se è true o false, in questo caso''),
        \textbf{Elapsed Seconds} e \textbf{Triggered Seconds} (``quanto tempo è stato
        triggerato''), e l'\textbf{Input Action}, ``la referenza stessa'' all'asset.
\end{itemize}

\paragraph{Set Max Walk Speed.} ``Dobbiamo modificare il Character Movement, ragazzi, che
dovrà definire quanta velocità avremo'': si trascina nel grafo il componente
\textbf{Character Movement} e se ne tira fuori \textbf{Set Max Walk Speed}, due volte.
Sul \textbf{Triggered}, \textbf{600}; sul \textbf{Completed}, \textbf{300}
(figura~\ref{ue9:fig:sprint}).

\paragraph{Il mapping dimenticato.} In Play non succede niente. ``Scusate, ragazzi, mi
sono dimenticato di dare l'input'': IA\_Run esiste, ma non è nel contesto. Si apre
IMC\_Default, si aggiunge un mapping (il \emph{+} accanto a \emph{Mappings}), si sceglie
IA\_Run e gli si assegna il tasto: ``un nostro bello \textbf{Shift}''. ``L'Input Mapping
Context ci permette di mettere quali tipologie di pulsanti facciano'' ciascuna azione.

\paragraph{Lo stesso input, anche con il controller.} ``Io ad esempio qua, ragazzi, ho un
controller'': il docente lo collega (``Unreal riesce a prendere il controller''), aggiunge
un secondo tasto alla stessa azione, clicca sull'icona della tastiera accanto al campo e
preme un pulsante del controller (``mi sono fregato la sorpresa''): ``facciamo tipo Dark
Souls? Premo cerchio'', e l'editor scrive \textbf{Gamepad Face Button Right}. ``Questa
azione, cioè l'azione di correre, è collegata a più input.'' ``Quando faccio progetti in
generale, devo fare anche il controller'': si corre con Shift sulla tastiera ``oppure con
il nostro bel controllino, con quello che vogliamo noi'', ``come è stato designato il
gioco''. Tutto il resto del movimento sul controller ``non ci interessa adesso''.

\paragraph{La prova.} ``Se noi proviamo\dots\ rilascio.'' Con Shift il cavaliere corre, e
il Blend Space passa alla corsa; al rilascio torna a camminare. ``Ecco a voi che abbiamo
creato un nostro bel character nuovo, in cui semplicemente possiamo correre, possiamo
modificare la velocità.''

\begin{figure}[htbp]
\centering
\begin{tikzpicture}[font=\footnotesize,
  ue9box/.style={draw, thick, rounded corners=4pt, align=center, fill=gray!8,
                 minimum height=1.0cm, text width=2.5cm},
  ue9arr/.style={-{Stealth[length=2.5mm]}, thick}]
  \node[ue9box, fill=yellow!15] (k1) at (0,0.7) {Left Shift\\{\scriptsize tastiera}};
  \node[ue9box, fill=yellow!15] (k2) at (0,-0.7) {Gamepad Face\\Button Right};
  \node[ue9box, fill=blue!10, text width=2.8cm] (imc) at (3.8,0) {\textbf{IMC\_Default}\\
    {\scriptsize Input Mapping Context}\\{\scriptsize (mapping IA\_Run)}};
  \node[ue9box, fill=cyan!12] (ia) at (7.6,0) {\textbf{IA\_Run}\\{\scriptsize Input
    Action, Digital (bool)}};
  \node[ue9box, fill=red!12, text width=2.8cm] (ev) at (11.4,0) {evento
    \textbf{IA\_Run}\\{\scriptsize nel Blueprint}\\{\scriptsize Triggered / Completed}};
  \draw[ue9arr] (k1) -- (imc); \draw[ue9arr] (k2) -- (imc);
  \draw[ue9arr] (imc) -- (ia); \draw[ue9arr] (ia) -- (ev);
  \node[ue9box, fill=green!10, text width=4.2cm] (bp) at (3.8,-2.4) {\textbf{Event Begin
    Play} $\rightarrow$ Enhanced Input Local Player Subsystem $\rightarrow$ \textbf{Add
    Mapping Context} (IMC\_Default)};
  \draw[ue9arr, dashed] (bp) -- node[right, font=\scriptsize]{attiva} (imc);
\end{tikzpicture}
\caption{La catena dell'Enhanced Input. I tasti fisici non sono legati direttamente
all'evento: il contesto (IMC) associa più tasti alla stessa azione (IA), e il Blueprint
reagisce all'azione. Senza Add Mapping Context, il contesto non è attivo e nessun input
arriva. Schema ricostruito.}
\label{ue9:fig:enhanced}
\end{figure}

\begin{figure}[htbp]
\centering
\begin{tikzpicture}[font=\footnotesize,
  ue9node/.style={draw, thick, rounded corners=4pt, align=center, fill=gray!8,
                  minimum height=0.8cm, text width=2.4cm},
  ue9ev/.style={ue9node, fill=red!12},
  ue9var/.style={ue9node, fill=blue!10, rounded corners=8pt},
  ue9exec/.style={-{Stealth[length=2.5mm]}, very thick, black!80},
  ue9data/.style={-{Stealth[length=2mm]}, thick, blue!60!black}]
  \node[ue9ev, text width=2.2cm, minimum height=1.8cm] (ev) at (0,0) {\textbf{IA\_Run}\\
    {\scriptsize Enhanced Input Action}};
  \node[ue9node, text width=3.2cm] (s1) at (5,0.8) {\textbf{Set Max Walk Speed}\\
    {\scriptsize Max Walk Speed = 600}};
  \node[ue9node, text width=3.2cm] (s2) at (5,-0.8) {\textbf{Set Max Walk Speed}\\
    {\scriptsize Max Walk Speed = 300}};
  \node[ue9var, text width=2.6cm] (cm) at (9.6,0) {\textbf{Character}\\
    \textbf{Movement}};
  \draw[ue9exec] (ev.east |- s1) -- node[above, font=\scriptsize]{Triggered} (s1);
  \draw[ue9exec] (ev.east |- s2) -- node[below, font=\scriptsize]{Completed} (s2);
  \draw[ue9data] (cm.west) -- ++(-0.5,0) |- node[pos=0.25, right, font=\scriptsize]
    {Target} (s1.east);
  \draw[ue9data] (cm.west) -- ++(-0.5,0) |- (s2.east);
\end{tikzpicture}
\caption{Lo sprint di BP\_Knight1. Premendo Shift (o il cerchio del controller) la
velocità massima sale a 600 e il Blend Space mostra la corsa; al rilascio torna a 300,
la camminata. Schema ricostruito.}
\label{ue9:fig:sprint}
\end{figure}

\paragraph{Approfondimento: Action Mapping ed Enhanced Input.} Il libro di Romero e
Sewell, scritto per le prime versioni di Unreal Engine 5, costruisce lo stesso sprint
con il vecchio sistema: un \emph{Action Mapping} chiamato \emph{Sprint} nei Project
Settings (sezione \emph{Input}), legato a \emph{Left Shift}, e un evento
\emph{InputAction Sprint} con i pin \emph{Pressed} e \emph{Released}, su cui due
\emph{Set Max Walk Speed} portano la velocità a 2200 e poi di nuovo ai 600 di default
(capitolo 6). Il principio è identico; cambia il sistema. L'\textbf{Enhanced Input},
introdotto come plugin in Unreal Engine 5.0 e diventato il sistema predefinito dei
modelli dalla 5.1 (il ``quasi dalla 5.2'' del docente è un ricordo approssimativo), separa
tre livelli: l'\emph{azione} (Input Action, che cosa vuole fare il giocatore), il
\emph{contesto} (Input Mapping Context, quali tasti attivano quali azioni in una certa
situazione) e il \emph{comportamento} (Triggers e \emph{Modifiers}, come per esempio la
\emph{zona morta} di una levetta o l'inversione di un asse). I contesti si aggiungono e
si tolgono a runtime con \emph{Add Mapping Context} e \emph{Remove Mapping Context}, e
hanno una \emph{Priority}: è così che si passa ``dal gioco a piedi alla macchina''
dell'esempio del docente, e così che si costruisce una schermata di rimappatura dei
tasti. Gli autori raccomandano anche di raggruppare i nodi dello sprint in un
\emph{commento} (\emph{Create Comment from Selection}), per tenere ordinato il grafo.

% ---------------------------------------------------------------------
\section{Un tentativo: FInterp To}
\label{ue9:sec:finterp}

``Se volessimo farlo un po' più carino, ragazzi'', il passaggio fra 300 e 600 potrebbe
essere graduale invece che istantaneo. Il blocco è \textbf{FInterp To} (``finter2''):
``fa un'interpolazione di un float''. ``Avete presente quando abbiamo parlato di
interpolazione?'' ``Fa il passaggio tra due numeri. Lui dice: devi passare da questo
valore, cioè \textbf{Current}, a questo valore, \textbf{Target}, in questo tempo'',
\textbf{Delta Time}, per cui il docente consiglia \textbf{Get World Delta Seconds}; e
poi c'è l'\textbf{Interp Speed}, la velocità dell'interpolazione.

\paragraph{Il grafo.} Sul Triggered: \textbf{Current} = Character Movement $\rightarrow$
\textbf{Get Max Walk Speed}; \textbf{Target} = 600; \textbf{Delta Time} = World Delta
Seconds; il risultato nel Set Max Walk Speed. Il docente ricorda che il Delta Seconds si
può moltiplicare per un valore (``ricordiamo che possiamo modificarlo moltiplicandolo
per un valore''), e prova con 10, poi con un decimo. Sul Completed lo stesso, ``ma
semplicemente invertiamo'': il Current parte da 600, il Target è 300. ``Ricordiamoci di
collegare anche al Character Movement.''

\paragraph{Non funziona, e perché.} La corsa resta ``ancora molto veloce'', il passaggio
istantaneo. Il docente prova valori diversi (10, 100: ``è alla norma del pianeta stare
continuamente a vedere se i valori sono giusti'') e poi mette dei \textbf{Print String}:
``ci dà seicento''. La diagnosi: ``fondamentalmente questa è un'azione a triggered
unico, quindi non farà l'interpolazione''. FInterp To fa \emph{un passo}
dell'interpolazione ogni volta che viene eseguito: se viene eseguito una volta sola, fa
un solo passo. ``Lo deve fare tramite un \textbf{timer}'': si rimanda a dopo (``volendo,
ragazzi, si può fare con un timer, però facciamolo magari dopo''), e non verrà ripreso
nella lezione.

\paragraph{Approfondimento: dove mettere un'interpolazione.} Un nodo come FInterp To
(o VInterp To, RInterp To per vettori e rotazioni) va eseguito \emph{a ogni frame}
finché il valore non arriva al target: il suo risultato è il valore ``un po' più vicino''
al target, e il Delta Time serve a renderlo indipendente dal frame rate. Ci sono tre
modi tipici di farlo girare: nel \textbf{Tick}, con una variabile \emph{TargetSpeed}
che i due pin dell'input si limitano a impostare (600 o 300); con un \textbf{timer}
(\emph{Set Timer by Event}, con \emph{Looping}, fermato con \emph{Clear and Invalidate
Timer} quando il valore è arrivato), come suggerisce il docente; o con una
\textbf{Timeline}, vista nella sesta lezione, che va da 300 a 600 con \emph{Play} e
torna indietro con \emph{Reverse}. Con l'Enhanced Input c'è anche una quarta via: il pin
\emph{Triggered} scatta a ogni frame se il trigger è \emph{Down} (o se non c'è nessun
trigger), e una sola volta se il trigger è \emph{Pressed}. IA\_Run, copiata da IA\_Jump,
ha probabilmente un trigger che la fa scattare una volta sola: è ``l'azione a triggered
unico'' del docente.

% ---------------------------------------------------------------------
\section{Prima della pausa: le animazioni da scaricare}
\label{ue9:sec:download}

``Adesso andremo a vedere, se non mi esplode prima Unreal, anche come si fanno i
\textbf{montaggi}''; mancano una ventina di minuti alla pausa. ``Dovremmo andare un attimo
su Mixamo'' (``a quanto pare oggi Mixamo ci mette pure un bel po' a caricare''), per
scaricare due gruppi di animazioni.

\paragraph{Un attacco in root motion.} ``Vorrei che sceglieste un'animazione che sia in
\textbf{root motion}'': ``riprendiamo ad esempio una situazione simile a Dark Souls''.
``Ricordatevi sempre di prendere il vostro character'' su Mixamo (il paladino, nel caso
del docente), perché le animazioni si adattano al personaggio selezionato. Il docente
cerca \emph{attack} e sceglie un attacco con la spada che ``di base, come potete vedere,
non può essere messo in place''. Le impostazioni di download sono quelle della volta
precedente: \textbf{Without Skin}, \textbf{FBX Binary}, \textbf{30} fotogrammi al
secondo. ``Ripeto: non deve essere in place, deve avere il root motion, quindi il
movimento ereditato.'' ``Un'azione unica'': ``ad esempio un attacco con una spada, un
qualche cosa'', ``un'animazione che vi piace di più''.

\paragraph{Le animazioni del Blend Space 2D.} Dopo la pausa si farà il \textbf{Blend
Space 2D}, ``che ci servirà per fare anche il calcolo della direzione del nostro
character''. Servono quindi altre animazioni: ``walk destra, sinistra, dritto e
indietro''. ``Consiglio spassionato'': su Mixamo si chiamano \textbf{Strafe} (``provate a
cercare Strafe''; ``Jog Strafe va benissimo anche''), e queste, al contrario
dell'attacco, ``dovete metterle \textbf{in place}''.

\paragraph{Il mirror.} ``Piccolo consiglio: se volete far sì che siano abbastanza simili
tra di loro, potete fare il \textbf{mirror}.'' Si scarica l'animazione verso sinistra
non specchiata; poi, con l'opzione \emph{Mirror} di Mixamo, ``in automatico andrà a
destra''. Le due animazioni saranno così perfettamente simmetriche.

\paragraph{Quante animazioni.} Su richiesta (``le puoi scrivere per messaggio?''), il
docente le elenca anche in chat: ``in totale sono otto animazioni''. Avanti, indietro,
destra, sinistra, e le quattro \textbf{diagonali}: avanti a destra, avanti a sinistra,
indietro a destra, indietro a sinistra. Uno studente chiede se l'avanti serva, visto che
c'è già la camminata: ``nel senso, possiamo usare tranquillamente quello di avanti; ma in
questo caso abbiamo solamente avanti, dobbiamo trovare anche indietro''. E un'altra: le
diagonali ``soltanto in avanti, o anche indietro?'' ``Anche indietro'': ``se io vado tipo
così, indietro verso destra, e così, indietro verso sinistra''. ``Praticamente è la X,
letteralmente'': le due diagonali incrociate, oltre alla croce degli assi. Il docente
nomina anche alcuni pacchetti di Mixamo con animazioni di locomozione (``Locomotion
Pack''; il resto dei nomi è incomprensibile nella trascrizione). ``Ci rivediamo alle
18.''

% ---------------------------------------------------------------------
\section{La pausa e i feedback}
\label{ue9:sec:feedback}

Al rientro, ``siamo rimasti in cinque'', poi in quattro. Il docente apre uno spazio per
i commenti: ``se avete dei feedback riguardo a questo corso, sparate tranquillamente'',
``siamo anche molto coetanei, anzi credo io di essere anche forse più piccolo di alcuni
di voi''. ``Se avete cose da chiedere, magari di andare più piano, andare più veloce,
oppure un po' come sono strutturate le lezioni''. ``Capisco pure che può essere
effettivamente un po' pesante in alcuni momenti'', ``perché purtroppo i Blueprint, o
Unreal in generale, è divertente quando poi le sai, le cose''.

\paragraph{Le funzioni da imparare a memoria.} Uno studente: ``immagino che sono
tantissime funzioni, quindi sono tutte da imparare a memoria''. ``Quello, eh sì.''

\paragraph{Il progetto del docente: la pallina morbida e dura.} Per mostrare dove porta
la pratica, il docente racconta un suo progetto personale: ``una specie di
\emph{platformer} in cui una pallina diventa morbida e dura, nel senso che diventa di
metallo o di un materiale più morbido''. ``La cosa è che ho simulato tutta la fisica'':
``la fisica non è simulata da Unreal, ma sono io che faccio finta, tramite i Blueprint, di
far fare delle funzioni fisiche''. ``Hai messo proprio te le formule fisiche?'' ``Più o
meno.'' Il meccanismo: si rileva lo \textbf{slope}, ``il quanto è piegato il piano'',
perché ``la meccanica è il fatto che se sei più pesante con la pallina vai più veloce su
uno slope, come nella realtà'': ``una pallina di metallo e una pallina di gomma: quella di
metallo va più veloce, perché è più pesante, comunque più influenzata''. ``È tutto
simulato'': ``sono io che, tramite la detection dello slope, cambio i valori''.

\paragraph{Combinare i blocchi.} ``E quello è ovviamente tutto custom'': ``quando hai
tutte le nozioni in testa, ti dici: ok, so questo blocco A, B, C, D, E, F, G; quale mi
serve? Il blocco C e il blocco A. Li metto insieme e faccio una nuova combinazione.''
``È sempre il fatto di capire come posso interlacciare le cose.'' ``Qua stiamo facendo
cose magari un po' generiche''; poi ``dici: ok, voglio fare simile a quello, ma
leggermente diverso. Voglio fare il character, però con la spada: che mi cambia?''

\paragraph{Inventari e UI.} La richiesta di Federico: ``a me interesserebbero la gestione
degli inventari e le interfacce''. ``Sfortunatamente no'': ``sappi che l'inventario e la
UI sono un mondo a parte''. ``Esistono proprio figure riguardanti questo mondo della UI,
sia designer sia programmatori.'' Si farà ``una cosa come un menu'', ``una cosetta
tranquilla'', ``però per il resto non andremo a fare di più''.

\paragraph{Le lezioni con Davide.} ``Sappiate che il 25 è la mia ultima lezione, e avrete
tutte quelle con \textbf{Davide}, che invece vi farà la parte più di \emph{level}.'' Con
lui si useranno ``tutti gli asset che abbiamo creato: le piattaforme, i pulsanti'', per
lavorare su concetti come la \textbf{gym}, ``il concetto della palestra'', che ``si fa
praticamente in via di sviluppo'': ``è il parco giochi'', ``letteralmente il parco giochi
dove ho tutte le meccaniche, le posso provare tutte quante, vedo se funzionano''. Poi si
costruiscono situazioni: ``voglio fare un livello horror, la bestia che ti insegue: tu
entri in una stanza e devi fuggire dalla bestia''; ``situazione X, situazione Y,
situazione Z'', ``le costruisci e metti insieme i componenti'', ``e vedi quei componenti
che hai programmato come funzionano a livello anche di esperienza, di gioco''. Gli
elementi saranno ``quelli di un classico platform: le piattaforme, i bottoni'', e si
vedranno anche i tool di \textbf{modeling} dell'editor e ``un po' di teoria di level
design'': ``chi fa game design, e a chi interessa tantissimo anche la parte di level:
molto, molto interessante''.

\paragraph{Dove lasciare i feedback.} ``Ripeto, ragazzi, anche a chi sta in differita'':
qualsiasi feedback sulle lezioni va lasciato sul canale \textbf{general} della chat,
``non su introduzione al materiale, perché sennò si fa casino''. Non solo sugli
argomenti, ``ma feedback anche proprio per le lezioni in generale: andare più lento,
andare più veloce, vedere cose più pratiche''.

\paragraph{Approfondimento: la gym.} Nel level design la \textbf{gym} (o \emph{gym
level}, \emph{gray box gym}) è un livello di servizio che non finirà nel gioco: una
stanza in cui ogni meccanica è disposta in isolamento e misurata. Si mettono salti di
lunghezze crescenti per sapere qual è il salto massimo del personaggio, rampe di
pendenze diverse per verificare la \emph{Max Walkable Floor Angle}, gradini di varie
altezze per la \emph{Max Step Height}, corridoi di larghezze diverse per la camera. Le
misure raccolte nella gym diventano le ``regole metriche'' del level design (quanto può
essere largo un burrone, quanto alto un muretto da scavalcare), e i livelli veri si
costruiscono rispettandole. È il ``parco giochi'' del docente; e il problema della
camera che compenetra, visto in questa lezione, è uno dei primi che una gym fa emergere.
% ---------------------------------------------------------------------
\section{La direzione: Calculate Direction}
\label{ue9:sec:direction}

``Abbiamo detto prima che oggi dobbiamo fare il Blend Space 2D.'' Il docente ha importato
otto animazioni, ``cercando di prenderle tutte similari'': \emph{Jog Backward},
\emph{Jog Backward Diagonal}, \emph{Jog Forward}, \emph{Jog Forward Diagonal} e così
via, ``a parte lo Strafe Left, che devo eliminare perché non andava bene come
animazione''. Il Blend Space 1D prende solo la velocità: ``vai più veloce, l'animazione,
benissimo''. Il 2D aggiunge un secondo asse, la \textbf{direzione}, ``in modo tale da non
dover magari direzionare sempre il character'', ma vedere ``che si muove in diagonale, si
muove a sinistra''.

\paragraph{Dove si calcola.} In ABP\_Knight, nell'Event Graph dell'Update, si aggiunge un
pin alla \textbf{Sequence} (il \emph{+}, \emph{Add pin}): \emph{Then 2}, o il successivo.
Questa volta non si parte dal Movement Component, ma dal \textbf{character}: ``perché la
direzione è determinata dalla rotazione del character stesso'', ``cioè la capsula, come è
ruotata''. Dalla variabile del cavaliere si tira fuori \textbf{Get Actor Rotation}.

\paragraph{Il nodo Calculate Direction.} ``Potremmo tranquillamente prenderla così, ma noi,
per semplificarci anche la vita, andiamo a calcolare proprio la direzione con un blocco
chiamato proprio \textbf{Calculate Direction}.'' Il nodo chiede due ingressi: la
\textbf{Velocity} attuale (dal Movement Component, \textbf{Get Velocity}) e la
\textbf{Base Rotation} (la Get Actor Rotation). Restituisce ``un float'', che con
\textbf{Promote to Variable} diventa \textbf{Direction}.

\paragraph{Uno studente perso.} ``Io mi sono perso: Character, Movement Component, da dove
vengono questi due?'' Il docente ripercorre: ``hai fatto questo Event Blueprint Initialize
Animation nell'Event Graph, no?''; ``dovresti avere qua sotto, sotto le variabili, la
variabile del character: la tiri fuori, ne tiri fuori Get Actor Rotation''; dal Movement
Component ``tiri fuori la velocità''; ``da Get Actor Rotation fai Calculate Direction, e
questo diventerà il nostro float''. Un altro studente nota che la variabile del docente si
chiama \emph{Character}, mentre la sua si chiama \textbf{As BP Knight}: ``va benissimo.
Io l'ho solo rinominata.'' Quando si fa Promote to Variable sull'uscita di un cast, Unreal
propone il nome \emph{As} più il nome della classe (``ti dice proprio As BP Knight'').

\paragraph{Da $-180$ a $180$.} ``Molto semplicemente, la direzione sarà tra un range,
sempre senza doverla clampare'' (``il clamp l'abbiamo fatto quasi all'inizio''): fra
\textbf{$-180$ e $180$ gradi}. ``Si parte da 0; a destra vai verso il positivo; quando
arrivi a 180 passa a $-179$''; ``diciamo, non lo fa come un angolo da 0 a 360, ma fa 0,
180, $-180$ e di nuovo 0'' (figura~\ref{ue9:fig:direzione}).

\begin{figure}[htbp]
\centering
\begin{tikzpicture}[font=\footnotesize]
  \draw[thick] (0,0) circle (2.2);
  \foreach \a/\l/\pos in {90/0/above, 45/45/above right, 0/90/right, -45/135/below right,
      -90/{180 / $-180$}/below, -135/{$-135$}/below left, 180/{$-90$}/left,
      135/{$-45$}/above left}
    {\draw[-{Stealth[length=2mm]}, thick, blue!60!black] (0,0) -- (\a:2.2);
     \node[\pos, font=\scriptsize\bfseries] at (\a:2.25) {\l};}
  \fill[gray!40] (0,0) circle (0.35);
  \draw[thick] (0,0) circle (0.35);
  \draw[-{Stealth[length=2mm]}, very thick] (0,0) -- (0,0.6);
  \node[font=\scriptsize, align=left, anchor=west] at (3.4,1.2) {\textbf{0}: avanti\\
    \textbf{positivi}: verso destra\\\textbf{negativi}: verso sinistra\\
    \textbf{$\pm$180}: indietro};
  \node[font=\scriptsize, align=left, anchor=west] at (3.4,-0.9) {la direzione è
    l'angolo fra\\la velocità e il ``davanti''\\del personaggio (Base Rotation)};
\end{tikzpicture}
\caption{Il valore di Direction restituito da Calculate Direction. Il personaggio guarda
in alto; le frecce sono le otto direzioni delle animazioni del Blend Space 2D. Andare
indietro vale sia 180 sia $-180$, a seconda del lato da cui ci si arriva. Schema
ricostruito.}
\label{ue9:fig:direzione}
\end{figure}

\paragraph{Rimettere Use Controller Rotation Yaw.} Per vedere i movimenti laterali,
``importante è rimettere lo \textbf{Use Controller Rotation Yaw}'': ``non vogliamo che il
character ruoti sempre nella direzione in cui lo stiamo indirizzando''. Il corpo deve
seguire la camera, e la direzione del movimento rispetto al corpo deve cambiare.

\paragraph{La verifica con Print String.} Per sicurezza si stacca un attimo il Blend Space
1D (``staccare un attimo le animazioni'') e si stampa Direction con un \textbf{Print
String}, ``che lo farà a ogni frame''. ``Ci dà direttamente nella preview questo valore'';
in Play, ``sarà un po' rotto, ma è normale, ragazzi'': ``$-70$, 70, 160, 0''. ``Mamma mia''
(il valore dipende dalla camera, ``giustamente''). ``Come possiamo vedere, questi valori
rappresentano la nostra direzione attuale.'' Il docente nota anche che il cavaliere
sembra ``leggermente ruotato'' rispetto alla capsula: ``buono a sapersi''.

\paragraph{Approfondimento: che cosa calcola Calculate Direction.} Il nodo, della
libreria dell'Animation Blueprint (\emph{KismetAnimationLibrary}; nelle versioni più
vecchie era una funzione dell'AnimInstance), trasforma la velocità nello spazio locale
del personaggio e ne restituisce l'angolo sul piano orizzontale, in gradi, nell'intervallo
$[-180, 180]$: 0 se la velocità punta dove guarda il personaggio, 90 se punta alla sua
destra, $-90$ a sinistra, $\pm 180$ dietro. A velocità nulla restituisce 0. Con
\emph{Orient Rotation to Movement} acceso il personaggio si gira sempre verso la
velocità, e la direzione resta quasi sempre vicina a zero: il Blend Space 2D avrebbe poco
da fare. Per questo il docente riattiva Use Controller Rotation Yaw.

% ---------------------------------------------------------------------
\section{Il Blend Space 2D}
\label{ue9:sec:blend2d}

\paragraph{Creare l'asset.} Tasto destro $\rightarrow$ \emph{Animation} $\rightarrow$
\textbf{Blend Space}: ``Blend Space senza niente, di base, è 2D; se viene specificato
1D, è 1D''. Unreal chiede lo scheletro: ``voi scegliete il vostro scheletro''. Il nome:
\textbf{BS2D\_Knight}.

\paragraph{I due assi.} ``Come potete vedere, a sinistra e giù sono presenti due nomi'',
per ora \emph{None}. Nei Details:
\begin{itemize}
  \item \textbf{Horizontal Axis}: \textbf{Direction}, da \textbf{$-180$} a
        \textbf{$180$};
  \item \textbf{Vertical Axis}: \textbf{GroundSpeed} (``più che la velocità, la ground
        speed''), da \textbf{0} a \textbf{600}.
\end{itemize}
Il parametro \textbf{Grid Divisions} ``ci permette di aumentare il numero di divisioni
sulla nostra griglia'' (le linee ``non si vedono, perché sono veramente molto poco
chiare''): sono i punti su cui si agganciano i campioni delle animazioni. Il docente ne
mette \textbf{8} su entrambi gli assi. Con 8 divisioni, sull'asse orizzontale ogni passo
vale 45 gradi, esattamente le otto direzioni; sull'asse verticale vale 75.

\paragraph{Un quiz.} ``Domanda, ragazzi: se io sto a 600 di velocità e a 90 di direzione,
dove sto andando?'' ``Ricordatevi che positivo è destra, negativo è sinistra.'' ``A
destra.'' ``Pienamente a destra, perfetto.'' ``E invece a $-135$?'' ``Dietro a
sinistra.'' ``Ottimo.''

\paragraph{Dove va ogni animazione.} Si dispongono i campioni (tabella~\ref{ue9:tab:bs2d} e
figura~\ref{ue9:fig:blend2d}). In basso, a velocità zero, l'\textbf{Idle}: ``mettetene
tre, così abbiamo un range per tutto'' (a $-180$, 0 e 180): ``siamo fermi, siamo fermi''.
In alto, a 600, le otto animazioni. L'animazione all'indietro va messa due volte, a
$-180$ e a $180$: ``stiamo andando indietro sia se abbiamo girato da sinistra, sia da
destra''. Il docente si corregge una volta (``Backward Diagonal L sarà\dots\ no, ho
sbagliato'': a $-135$, non a 90) e invita a controllare i valori nei Details del
campione.

\begin{table}[htbp]
\centering
\small
\begin{tabular}{@{}p{4.4cm}p{2.2cm}p{2.2cm}p{4.6cm}@{}}
\toprule
\textbf{Animazione} & \textbf{Direction} & \textbf{GroundSpeed} & \textbf{Movimento} \\
\midrule
Idle (tre copie) & $-180$, 0, 180 & 0 & fermo \\
Jog Forward & 0 & 600 & avanti \\
Jog Forward Diagonal R & 45 & 600 & avanti a destra \\
Strafe R & 90 & 600 & a destra \\
Jog Backward Diagonal R & 135 & 600 & indietro a destra \\
Jog Backward & 180 e $-180$ & 600 & indietro \\
Jog Backward Diagonal L & $-135$ & 600 & indietro a sinistra \\
Strafe L & $-90$ & 600 & a sinistra \\
Jog Forward Diagonal L & $-45$ & 600 & avanti a sinistra \\
\bottomrule
\end{tabular}
\caption{La posizione dei campioni nel Blend Space 2D BS2D\_Knight. I nomi delle
animazioni sono quelli di Mixamo citati a lezione, per quanto ricostruibili dalla
trascrizione; i valori sono quelli dettati dal docente.}
\label{ue9:tab:bs2d}
\end{table}

\begin{figure}[htbp]
\centering
\begin{tikzpicture}[font=\footnotesize, x=0.032cm, y=0.006cm,
  ue9pt/.style={draw=black, fill=white, thick, diamond, inner sep=1.8pt}]
  % griglia
  \foreach \x in {-180,-135,...,180} \draw[gray!30] (\x,0) -- (\x,600);
  \foreach \y in {0,75,...,600} \draw[gray!30] (-180,\y) -- (180,\y);
  \draw[thick] (-180,0) rectangle (180,600);
  \foreach \x in {-180,-90,0,90,180}
    \node[below, font=\scriptsize] at (\x,-10) {\x};
  \foreach \y in {0,300,600} \node[left, font=\scriptsize] at (-182,\y) {\y};
  \node[below, font=\small] at (0,-80) {Direction};
  \node[rotate=90, font=\small] at (-235,300) {GroundSpeed};
  % campioni
  \foreach \x in {-180,0,180} \node[ue9pt] at (\x,0) {};
  \node[above, font=\scriptsize] at (0,10) {Idle};
  \foreach \x/\l in {-180/Back, -135/BackDiag L, -90/Strafe L, -45/FwdDiag L, 0/Fwd,
      45/FwdDiag R, 90/Strafe R, 135/BackDiag R, 180/Back}
    {\node[ue9pt] at (\x,600) {};
     \node[above, font=\tiny, rotate=35, anchor=south west] at (\x-6,612) {\l};}
  % punto attuale
  \fill[green!50!black] (60,450) circle (2.2pt);
  \node[green!40!black, font=\scriptsize, anchor=west] at (66,420) {(60, 450)};
\end{tikzpicture}
\caption{Il Blend Space 2D del cavaliere: Direction in orizzontale, GroundSpeed in
verticale, griglia 8$\times$8. Le tre copie dell'Idle stanno in basso; le otto direzioni
in alto, con la corsa all'indietro ripetuta ai due estremi. Il punto verde è un esempio:
il personaggio sta andando in avanti a destra, un po' sotto la velocità massima, e Unreal
mescola i quattro campioni più vicini. Schema ricostruito.}
\label{ue9:fig:blend2d}
\end{figure}

\paragraph{La preview con Ctrl.} ``Potete fare la preview delle animazioni con
\textbf{Ctrl}'': tenendolo premuto e trascinando il cursore sul grafico, ``potete vedere
come l'animazione si muove tra tutti gli step''. ``Consiglio sempre di mettersi dietro''
il personaggio nel viewport, ``perché tanto poi, se lo vedete davanti, è difficile'' capire
destra e sinistra. Il docente controlla tutte le direzioni (``vado a sinistra, vado
indietro, vado dritto, vado dritto a destra, vado a destra, giù a destra, indietro.
Perfetto''); una delle animazioni ``mi sa che non va bene'', ma ``sì, sto andando a
sinistra, è solo molto strana''.

\paragraph{Smoothing Time.} ``Piccola cosa, se volete introdurla'': lo \textbf{Smoothing
Time}, ``il tempo di passaggio tra un'animazione e l'altra in generale, per tutte le
animazioni''. Lo si mette a \textbf{0,2} su entrambi gli assi. Il \emph{Damping Ratio}
``non ci interessa''.

\paragraph{Nello stato WalkRun.} ``Facciamo tutto il percorso a ritroso'': AnimGraph,
Locomotion, WalkRun. Al posto del Blend Space 1D si trascina BS2D\_Knight, che chiede
``dammi la Direction e la GroundSpeed'': si collegano i due Get. Compile: ``vediamo che il
nostro Blend Space funziona in preview''.

\paragraph{Le impostazioni del personaggio.} In BP\_Knight1: ``mi raccomando, l'abbiamo
detto prima, ma lo diciamo anche adesso'', si riattiva \textbf{Use Controller Rotation
Yaw}; e nel Character Movement la \textbf{Max Walk Speed} va a \textbf{600}, ``che è
quanto abbiamo messo'' come massimo dell'asse. ``Ecco a voi, ragazzi'': ``se io mi muovo a
destra lui va a destra, se mi muovo a sinistra, anche se l'animazione non è fantastica,
va a sinistra; vado indietro; vado in diagonale''. Nel frattempo ``parte la torretta, che
mi spara'': ``il mio bel cavaliere se ne va a zonzo come gli pare''.

\paragraph{Gli scatti.} C'è ``un piccolo \emph{clip}'', uno scatto, in alcuni passaggi.
Il docente prova a prendere la rotazione della \textbf{Mesh} invece di quella
dell'actor (``facciamo che invece di prendere la direzione del character, prendiamo
quella della Mesh''; ``Get Mesh''), poi ad aumentare lo Smoothing Time della Direction
a \textbf{0,5} (``un po' più tranquillo''), e infine torna al collegamento originale.
La conclusione: ``probabilmente sono dei problemi sull'animazione''; ``mi sa che devo
riscaricare l'animazione, perché snappa un pochino''. Ma il risultato c'è: ``abbiamo il
nostro character che si muove in tutte le direzioni''; ``è un po' veloce, devo essere
sincero'', e la velocità si può abbassare.

\paragraph{Il cerchio, per chi si è perso.} Uno studente chiede di rivedere il Blend
Space 2D. ``Immaginatelo come un cerchio'': ``se stai a 0 di direction, stai andando
dritto; se stai a 180 o $-180$, stai andando dietro''; ``tutti i valori intermedi'' sono le
diagonali e i lati. ``Destra è positivo: 0, 90, 180; sinistra è $-90$, $-180$.'' ``Se vai
in diagonale a destra sarà 45, se vai pienamente a destra sarà 90'', ``e 135 a 600 vuol
dire che andrai praticamente indietro, ma verso destra''. ``Una volta capito che una è
l'opposto dell'altra'', i valori negativi seguono da soli.

\paragraph{Approfondimento: Blend Space 2D e foot sliding.} In un Blend Space 2D Unreal
divide lo spazio dei campioni in triangoli e, per un valore dato, mescola i campioni ai
vertici del triangolo che lo contiene. Ne seguono tre conseguenze pratiche. Primo,
l'Idle ripetuto lungo tutta la riga inferiore (le tre copie del docente) evita che a
bassa velocità il personaggio ``scelga'' una direzione a caso. Secondo, la corsa
all'indietro va duplicata ai due estremi, perché $-180$ e 180 sono lo stesso angolo ma
due punti diversi della griglia: con un solo campione, passando da 179 a $-179$ il
personaggio attraverserebbe in un istante tutto il grafico. Terzo, lo scatto che il
docente vede è spesso proprio quel salto da $+180$ a $-180$, che lo \emph{Smoothing
Time} non può interpolare bene (interpola il valore numerico, non l'angolo); le
alternative sono il \emph{Weight Speed} per campione, oppure le soluzioni più avanzate
di Unreal Engine 5 come l'\emph{Orientation Warping}, che ruota la parte inferiore del
corpo verso la direzione del movimento con quattro sole animazioni cardinali. Infine,
una griglia con l'asse verticale fino a 600 e i campioni tutti a 600 significa che anche
la camminata a 300 è un misto fra Idle e corsa: per avere una vera camminata laterale
servirebbe una seconda riga di campioni a 300, con le animazioni di camminata.

% ---------------------------------------------------------------------
\section{Gli Animation Montage}
\label{ue9:sec:montage}

``Intanto ci abbiamo un'ultima mezz'oretta, ragazzi. Vediamo un attimo i Montage.''
L'animazione di attacco (``l'avevo scaricata, mi sa che non l'ho messa in Knight
Animation''; ``scusate, ragazzi'') è una \textbf{Sword and Shield Attack} di Mixamo:
``questa qua sta in root motion'': ``come potete vedere, eredita il movimento''.

\paragraph{Il problema.} ``Quando vogliamo fare una nostra animazione che di conseguenza
sia attivabile durante un ciclo di animazioni'', non basta aggiungere uno stato: ``dobbiamo
creare quello che viene chiamato un \textbf{Montage}''. (Nel frattempo ``Filippo è andato
via, è crashato'': ``siamo rimasti noi tre''.)

\paragraph{Che cos'è un Montage, dalla documentazione.} Il docente apre la pagina della
documentazione di Unreal sugli \textbf{Animation Montage}. ``Il montaggio di animazioni ci
permette sia di montare delle animazioni\dots'' L'esempio della pagina è la ricarica di
un fucile: un'animazione formata da tre animazioni diverse, \emph{MM Shotgun Reload Drop},
\emph{Reload Load} e \emph{Reload Lock}. ``Mettiamo conto che vogliamo fare un bel gioco
shooter. Il fatto di staccare il caricatore su quattro fucili diversi è praticamente
uguale''; ``non dobbiamo fare quattro animazioni diverse'': ``l'unica cosa che cambia è il
lock, perché magari uno vuole farlo un po' più grosso''. Con un Montage si creano
\textbf{animazioni modulari}: ``guarda, il drop è uguale, il load è uguale, il lock invece
è diverso''. E poi il secondo uso, quello di oggi: ``inserire animazioni in contesti di
cicli di animazione''; ``sono una tantum'': ``premo un tasto e compio questa animazione,
nel mentre che sto camminando''. ``Magari la prossima volta andiamo a vedere pure un po'
più approfondito.''

\paragraph{Creare il Montage.} Tasto destro sull'animazione di attacco $\rightarrow$
\textbf{Create} $\rightarrow$ \textbf{Create AnimMontage}. Il nuovo asset si chiama come
l'animazione con il suffisso \emph{Montage} (\emph{Sword And Shield Attack\_Montage}).
``La cosa bella del Montage è che ci permette di fare un sacco di cose all'interno dei
suoi parametri.''

% ---------------------------------------------------------------------
\section{Dentro il Montage: Notify, sezioni e Slot}
\label{ue9:sec:notify}

\paragraph{I Notify.} ``Una cosa interessante da spiegare sono i \textbf{Notify}'': ``dei
punti, dei puntini, sulla linea del tempo dell'animazione'', che permettono di aggiungere
``un'azione che si triggera'': ``suoni, un VFX, qualsiasi cosa''. ``Quelli più generici
sono già fatti da Unreal, fortunatamente.'' Il docente fa clic destro sulla traccia dei
Notify, \textbf{Add Notify} $\rightarrow$ \textbf{Play Sound}, e come suono sceglie
``l'unico sound che abbiamo'', quello della piattaforma (un Sound Cue). Le tracce dei Notify
possono essere più d'una: con \emph{Add Notify Track} se ne crea un'altra, per esempio
una per i VFX e una per i suoni. In preview (tolto il loop, ``fammelo fare una volta''):
``appena il nostro character arriva a questo punto dell'animazione, compie il suono''.
``Mettiamo conto: in combattimento c'è un VFX, c'è un suono, c'è un tirare fuori la spada,
e quindi si sente la spada che esce fuori dal fodero'': con una linea di Notify ``questo
suono avviene in un preciso momento dell'animazione''. ``Ovviamente, non andremo a usare
il sound effect della piattaforma per il nostro Knight.''

\paragraph{Le Montage Section.} Nel \emph{Montage group} c'è la sezione \textbf{Default}.
Con il tasto destro, \textbf{New Montage Section}: il docente ne crea una chiamata
\textbf{Attack}, nel punto in cui ``lui sta attaccando''. Le sezioni servono ``per
ordinare il lavoro'' e per ``dividere l'animazione in più punti diversi, con più
animazioni''.

\paragraph{Più animazioni nello stesso Montage.} La striscia verde sotto le sezioni è
l'animazione stessa. ``Se vogliamo, possiamo anche introdurre delle animazioni in più.''
Una dimostrazione volutamente assurda: ``io voglio, per qualche motivo arcano, far sì
che quando finisce l'animazione lui salti'': si trascina \textbf{MM\_Jump} nella traccia,
dopo l'attacco. ``È un salto orribile, però farà le due animazioni di conseguenza.'' Il
Montage è ``molto importante per azioni consecutive''. ``Quando si parla con l'animatore,
è importante il fatto di dire: questa animazione deve essere consecutiva a questa'', per
esempio ``mi serve che da questo punto in poi continui l'attacco, oppure faccia un'altra
animazione che ritorna in place''.

\paragraph{L'animazione di ritorno.} Finito l'attacco in root motion, il cavaliere resta dove l'attacco lo ha lasciato e non
torna in posa: manca ``un'animazione che ritorna in place'', cioè in idle. Il docente la
cerca su Mixamo, provando voci come \emph{Sword and Shield Idle}, \emph{Idle
Transition} e \emph{Start Block}; ma ``è un po' difficile'', perché ``bisogna trovare
il cavaliere proprio'', con la stessa posa di fine attacco, altrimenti ``snapperà
malissimo, almeno nel mio caso''.

\paragraph{Gli Slot.} Sotto, c'è il \textbf{Default Group} con lo \textbf{Slot}
\textbf{DefaultSlot}: ``questo ci permette di dire: su quale slot vuoi che questa
animazione si compia?''. DefaultSlot ``riprende tutto il corpo''. ``Potremmo, ad esempio,
voler creare uno slot che riprenda solamente alcune zone del nostro character'': ``vogliamo
che faccia un'azione solo nella parte superiore del corpo'' (sparare o salutare mentre si
corre); ``in quel caso creeremo un altro slot, che prenderà solo le ossa'' di quella
parte. ``Ma in questo caso ci serve tutto il corpo.''

% ---------------------------------------------------------------------
\section{Lo Slot nell'AnimGraph: Save Cached Pose}
\label{ue9:sec:slot}

Un Montage suona solo se l'AnimGraph ha un nodo \textbf{Slot} con lo stesso nome. In
ABP\_Knight, nell'AnimGraph:
\begin{enumerate}
  \item dall'uscita della State Machine Locomotion si crea un \textbf{New Save Cached
        Pose}, chiamato \textbf{LocomotionCache}: ``andiamo a conservarci quello che ha
        fatto la State Machine'', ``come se tutto questo procedimento lo mettesse in una
        memoria'';
  \item si aggiunge \textbf{Use cached pose 'LocomotionCache'};
  \item da lì un nodo \textbf{Slot 'DefaultSlot'} (la trascrizione: ``default slot'');
  \item lo Slot va all'\textbf{Output Pose}.
\end{enumerate}
``Semplicemente mixa insieme i due'': lo slot del Montage e la cache della locomozione
(figura~\ref{ue9:fig:slot}). Quando non c'è nessun Montage in corso, lo Slot lascia
passare la locomozione; quando un Montage su DefaultSlot parte, la sostituisce. Alla
fine della lezione, a uno studente che non legge il nome del nodo (``vedo tutto
pixelato''), il docente conferma: ``tiri fuori Locomotion Cache'', poi ``default slot''.

\begin{figure}[htbp]
\centering
\begin{tikzpicture}[font=\footnotesize,
  ue9pose/.style={draw, thick, rounded corners=4pt, align=center, fill=gray!8,
                  minimum height=1.0cm, text width=2.3cm},
  ue9arr/.style={-{Stealth[length=2.5mm]}, very thick, black!70}]
  \node[ue9pose, fill=orange!12] (sm) at (0,0) {\textbf{Locomotion}\\{\scriptsize State
    Machine}};
  \node[ue9pose, fill=blue!8] (sc) at (3.1,0) {\textbf{Save Cached Pose}\\{\scriptsize
    LocomotionCache}};
  \node[ue9pose, fill=blue!8] (uc) at (6.2,0) {\textbf{Use Cached Pose}\\{\scriptsize
    LocomotionCache}};
  \node[ue9pose, fill=green!10] (sl) at (9.3,0) {\textbf{Slot}\\{\scriptsize
    'DefaultSlot'}};
  \node[ue9pose, fill=red!10, text width=1.8cm] (op) at (12.1,0) {\textbf{Output
    Pose}};
  \draw[ue9arr] (sm) -- (sc); \draw[ue9arr] (uc) -- (sl); \draw[ue9arr] (sl) -- (op);
  \draw[dashed, thick, gray] (sc.east) -- (uc.west);
  \node[ue9pose, fill=yellow!15, text width=3.2cm] (mt) at (9.3,-2.1)
    {\textbf{Anim Montage}\\{\scriptsize Sword and Shield Attack, slot
    DefaultGroup.DefaultSlot}};
  \draw[-{Stealth[length=2.5mm]}, thick, dashed] (mt) -- node[right, font=\scriptsize]
    {quando suona, sostituisce la posa} (sl);
\end{tikzpicture}
\caption{La catena di pose finale dell'AnimGraph di ABP\_Knight. La State Machine produce
la locomozione, che viene salvata in una cache e riletta; lo Slot DefaultSlot la lascia
passare, a meno che un Montage assegnato a quello slot non sia in esecuzione. Schema
ricostruito.}
\label{ue9:fig:slot}
\end{figure}

% ---------------------------------------------------------------------
\section{Far partire l'attacco: Play Montage}
\label{ue9:sec:playmontage}

\paragraph{Il tasto F e il nodo.} Per far partire il Montage, il docente non crea
un'altra Input Action: usa direttamente un evento di tastiera, \textbf{F} (``facciamo con
OK\dots\ no, facciamo F, F va benissimo''). Nel Blueprint del personaggio, il nodo
\textbf{Play Montage}:
\begin{itemize}
  \item \textbf{In Skeletal Mesh Component}: ``in quale Skeletal Mesh Component
        farlo'', cioè la \textbf{Mesh} del personaggio;
  \item \textbf{Montage to Play}: ``quello che abbiamo appena fatto'' (``ce n'è tanto solo
        uno'');
  \item \textbf{Play Rate}: ``possiamo modificare quanto è veloce l'animazione'';
  \item fra le uscite: \textbf{On Completed} (``quando sarà finita''), \textbf{On Blend
        Out} (``mentre finisce''), \textbf{On Interrupted} (``se fosse interrotta''),
        \textbf{On Notify Begin} e \textbf{On Notify End} (``quando un Notify inizia'' e
        ``quando finisce'').
\end{itemize}
``Se noi premiamo F, compiamo un Montage.''

\paragraph{Fermare il personaggio.} ``È orribile'': a parte la mancanza dell'animazione di
ritorno, durante l'attacco ``possiamo muovere il character tranquillamente''. Quindi, sul
tasto F, prima del Play Montage:
\begin{enumerate}
  \item \textbf{Set Max Walk Speed} a \textbf{0}: ``non dobbiamo farlo muovere'';
  \item \textbf{Get Player Controller} $\rightarrow$ \textbf{Disable Input}: ``andiamo a
        disabilitare gli input, qualsiasi''.
\end{enumerate}
Sulle uscite del Play Montage:
\begin{itemize}
  \item \textbf{On Completed}: Get Player Controller $\rightarrow$ \textbf{Enable Input}
        (``on completed solamente'');
  \item \textbf{On Blend Out}: \textbf{Set Max Walk Speed} a \textbf{600}, ``al suo
        massimo''.
\end{itemize}

\paragraph{La camera che si muove ancora.} In Play, però, ``la camera si muove
effettivamente'', e con la camera il personaggio, ``quindi dobbiamo far sì che non sia
dipendente dalla camera''. Segue una serie di tentativi, ragionati ad alta voce:
\begin{itemize}
  \item togliere \emph{Use Pawn Control Rotation} sullo Spring Arm (``no, no'');
  \item l'\emph{Enable Camera Lag} dello Spring Arm;
  \item le spunte \emph{Inherit Pitch}, \emph{Inherit Yaw} e \emph{Inherit Roll} dello
        Spring Arm (``inerit yo e roll''; ``no, lo yaw ci deve essere''; ``forse sì,
        solo il roll va bene''), con un \emph{Set} preso dallo Spring Arm (``Set Inherit
        Yaw'', a false, e di nuovo a true alla fine); ``completamente\dots\ non va'';
  \item \emph{Stop Movement Immediately} e i parametri di input del Player Controller
        (\emph{movement input}, \emph{look input}).
\end{itemize}
La soluzione, alla fine, è sulle due spunte già incontrate nella prima parte: ``non è lo
yaw'' dello Spring Arm, ``ma è l'orientamento al movimento''. Sul tasto F si disattivano
entrambe:
\begin{itemize}
  \item dal \textbf{Character Movement}, \textbf{Set Orient Rotation to Movement} a
        \textbf{false};
  \item sul Pawn stesso, senza collegare nessun target, \textbf{Set Use Controller
        Rotation Yaw} a \textbf{false};
\end{itemize}
e sull'\textbf{On Completed} si rimettono entrambe a \textbf{true}
(figura~\ref{ue9:fig:playmontage}). ``Ottimo. Adesso riusciamo'': il cavaliere attacca
fermo, la camera gira senza trascinarlo, e alla fine riprende la locomozione. ``Abbiamo
fatto il nostro character che attacca e fa la locomotion.'' (Il docente nota anche, di
passaggio, che il cavaliere non interagisce con il pulsante del livello, ``giustamente,
perché è la Mesh''; dalla trascrizione non si ricostruisce meglio il punto.)

A Federico, che non riesce a sistemare la camera, il docente riassume: ``devi giocare con
lo yaw e con l'Orient Rotation to Movement: li disattivi tutte e due, e li riattivi
sull'On Completed''. ``Dello Spring Arm, giusto?'' ``Non dello Spring Arm: questo qua è
del Pawn, lo trovi direttamente senza collegare niente. Invece questo è solo del
Character Movement.''

\begin{figure}[htbp]
\centering
\begin{tikzpicture}[font=\footnotesize,
  ue9node/.style={draw, thick, rounded corners=4pt, align=center, fill=gray!8,
                  minimum height=0.9cm, text width=2.5cm},
  ue9ev/.style={ue9node, fill=red!12},
  ue9exec/.style={-{Stealth[length=2.5mm]}, very thick, black!80}]
  \node[ue9ev, text width=1.4cm] (f) at (0,0) {\textbf{F}\\{\scriptsize Pressed}};
  \node[ue9node] (a) at (2.8,0) {\textbf{Set Max Walk Speed}\\{\scriptsize 0}};
  \node[ue9node] (b) at (5.9,0) {Get Player Controller $\rightarrow$ \textbf{Disable
    Input}};
  \node[ue9node] (c) at (9.0,0) {\textbf{Set Orient Rotation to Movement}\\
    {\scriptsize false}};
  \node[ue9node] (d) at (12.1,0) {\textbf{Set Use Controller Rotation Yaw}\\
    {\scriptsize false}};
  \node[ue9node, fill=yellow!15, text width=3.0cm, minimum height=2.4cm] (pm) at
    (2.8,-3.0) {\textbf{Play Montage}\\{\scriptsize Mesh,\\Sword and Shield
    Attack\_Montage}};
  \node[ue9node] (e) at (8.3,-2.2) {\textbf{Set Max Walk Speed}\\{\scriptsize 600}};
  \node[ue9node, text width=3.0cm] (g) at (8.3,-3.9) {Get Player Controller
    $\rightarrow$ \textbf{Enable Input}};
  \node[ue9node, text width=3.0cm] (h) at (12.3,-3.9) {\textbf{Set Orient\dots} true\\
    \textbf{Set Use Controller Rotation Yaw} true};
  \draw[ue9exec] (f) -- (a); \draw[ue9exec] (a) -- (b); \draw[ue9exec] (b) -- (c);
  \draw[ue9exec] (c) -- (d);
  \draw[ue9exec] (d.south) -- ++(0,-0.6) -| (pm.north);
  \draw[ue9exec] (pm.east |- e) -- node[above, font=\scriptsize]{On Blend Out} (e);
  \draw[ue9exec] (pm.east |- g) -- node[below, font=\scriptsize]{On Completed} (g);
  \draw[ue9exec] (g) -- (h);
\end{tikzpicture}
\caption{L'attacco sul tasto F, nella versione finale. Prima del Montage si blocca il
personaggio (velocità zero, input disabilitato, nessuna rotazione automatica del corpo);
al Blend Out torna la velocità, all'On Completed tornano l'input e le rotazioni. L'ordine
esatto dei primi quattro nodi è ricostruito: a lezione vengono aggiunti uno alla volta,
durante i tentativi. Schema ricostruito.}
\label{ue9:fig:playmontage}
\end{figure}

\paragraph{Approfondimento: Play Montage, Montage Play e il root motion.} Il nodo usato a
lezione, \textbf{Play Montage}, è un nodo \emph{asincrono} (con l'orologio nell'angolo):
le sue uscite scattano più tardi, quando succede l'evento corrispondente, e per questo
può gestire il ``dopo'' del Montage direttamente nel grafo. Esiste anche \emph{Montage
Play}, chiamato sull'\emph{Anim Instance} della mesh, che restituisce solo la durata, e
\emph{Play Anim Montage}, del Character. Due dettagli completano la lezione. Il primo:
un attacco in \textbf{root motion} sposta davvero la capsula solo se il Montage ha il
root motion attivo e l'Animation Blueprint ha la \emph{Root Motion Mode} impostata su
\emph{Root Motion from Montages Only} (il default); se il root dell'animazione di Mixamo è
``rotto'', come visto nel retargeting, il personaggio avanza nella mesh e poi ``torna
indietro'' di scatto a fine animazione. Il secondo: durante un Montage che fa avanzare il
personaggio in root motion non serve azzerare la Max Walk Speed, perché il root motion
prende il posto dell'input di movimento; azzerarla, come fa il docente, è però un modo
semplice e sicuro di impedire al giocatore di muoversi durante un attacco ``in place''.
Infine, \emph{Disable Input} e \emph{Enable Input} agiscono sull'actor e richiedono un
Player Controller: per bloccare solo alcune azioni (il movimento ma non la camera),
l'Enhanced Input offre un'alternativa più fine, cioè rimuovere temporaneamente un
Input Mapping Context, o aggiungerne uno di priorità più alta.

% ---------------------------------------------------------------------
\section{Il compito e la chiusura}
\label{ue9:sec:chiusura}

\paragraph{Feedback in chat.} Una studentessa, Giulia, ha scritto in chat che è rimasta un
po' indietro ``causa trasloco'' e chiede di poter mandare dei feedback; nel frattempo è
già uscita dalla chiamata. ``Le risponderò sul Discord, assolutamente.'' ``Controllerò i
lavori che mandate sul canale.''

\paragraph{Il compito: lo StackOBot.} ``Una cosa che vorrei darvi come task, se vi va, per
la prossima volta'': ``fare un nuovo character totalmente da zero''. Nel progetto c'è già
un personaggio, lo \textbf{StackOBot} (il robottino dei progetti di esempio di Epic), che
``ha queste animazioni qua: vedete che ce le ha già dentro tutte quante''. Lo si trova
nella cartella \emph{StackOBot $\rightarrow$ Blueprints $\rightarrow$ Character}, con la
sua Skeletal Mesh (\emph{SKM\_Bot}) e il suo ABP. Il compito:
\begin{itemize}
  \item creare un \textbf{Blueprint Character} per lo StackOBot (``fare la BP'');
  \item creare il suo \textbf{Animation Blueprint} da zero, ``fare tutto l'Animation
        Blueprint e collegarci tutte le sue animazioni'': il ciclo di locomozione, con
        State Machine e Blend Space, come quello del cavaliere.
\end{itemize}
``Se volete farlo, lo potete mandare poi sulla chat'', ``che tanto magari ci lavoreremo
la prossima volta''. ``Per adesso, poi ne riparliamo.''

\paragraph{I saluti.} ``Se tutto a posto\dots'' Il docente interrompe la condivisione:
``vi ringrazio per l'attenzione, e ringrazio da casa chi è in differita''. ``Ci sentiamo
per martedì prossimo. Va bene, ciao, ragazzi. Buona serata.''

\begin{table}[htbp]
\centering
\small
\begin{tabular}{@{}p{5.4cm}p{8.6cm}@{}}
\toprule
\textbf{Sintomo} & \textbf{Causa e rimedio} \\
\midrule
L'animazione dello stato non si vede nella preview & la State Machine non è collegata
all'Output Pose dell'AnimGraph \\
L'animazione si ferma dopo un ciclo & manca la spunta di loop nel nodo dell'animazione \\
Warning nella State Machine & transizioni senza condizioni: normale finché non si
scrivono le regole \\
Jump $\rightarrow$ Fall dà errore & regola vuota: attivare Automatic Rule Based on
Sequence Player in State \\
Il Blend Space 1D arriva solo a 100 & Maximum Axis Value da impostare a 600; poi
spostare i campioni \\
Animazioni ``mischiate'' a metà (a 200) & i campioni non corrispondono alle velocità
reali: legarli alla Max Walk Speed \\
Alt sul nodo non crea il Set & negli stati dell'AnimGraph le variabili si leggono e
basta \\
``Questo asset non ha schede ancorate'' & Window $\rightarrow$ Load Layout $\rightarrow$
Default Editor Layout, poi Open Full Blueprint Editor \\
In Play non si vede niente / visuale nel corpo & il Character non ha camera: aggiungere
Spring Arm e Camera (figlia) \\
Non ci si muove & mancano gli input: copiare dal ThirdPersonCharacter i grafi di Look,
Move, Jump e l'Add Mapping Context \\
La camera ruota su un solo asse, o guarda solo il terreno & Use Pawn Control Rotation
sulla Camera invece che sullo Spring Arm \\
Il personaggio va di lato guardando avanti & attivare Orient Rotation to Movement e
togliere Use Controller Rotation Yaw (servono tutte e due) \\
La camera compenetra gli oggetti & Target Offset troppo grande: usare il Socket Offset \\
Il salto in root motion retargettato non è pulito & il root degli scheletri di Mixamo:
accettabile per la prova, da correggere \\
La State Machine non distingue salto e caduta & entrambe le regole su IsFalling:
aggiungere VelocityZ maggiore o minore di zero \\
Lo Shift non fa correre & IA\_Run non aggiunta ai mapping di IMC\_Default \\
FInterp To non interpola & eseguito una volta sola: va eseguito a ogni frame (Tick,
timer o Timeline) \\
Direction quasi sempre a zero & Orient Rotation to Movement acceso, Use Controller
Rotation Yaw spento: invertire per il Blend Space 2D \\
Scatti nel Blend Space 2D & animazioni non omogenee, salto fra 180 e $-180$: Smoothing
Time, o riscaricare le animazioni \\
Durante l'attacco il personaggio si muove & Set Max Walk Speed a 0 e Disable Input prima
del Play Montage \\
Durante l'attacco la camera trascina il personaggio & Set Orient Rotation to Movement e
Set Use Controller Rotation Yaw a false, a true sull'On Completed \\
Il Montage non suona & manca il nodo Slot DefaultSlot nell'AnimGraph \\
\bottomrule
\end{tabular}
\caption{I problemi incontrati durante la lezione, con la causa e il rimedio.}
\label{ue9:tab:errori}
\end{table}

\begin{table}[htbp]
\centering
\small
\begin{tabular}{@{}p{4.8cm}p{9.2cm}@{}}
\toprule
\textbf{Nodo, asset o impostazione} & \textbf{Che cosa fa} \\
\midrule
State Machine (AnimGraph) & organizza le animazioni in stati e transizioni \\
Entry / State / Conduit / Alias & ingresso; stato con una posa; snodo senza posa; alias
di più stati \\
Output Animation Pose & la posa prodotta da uno stato \\
Transition Rule (Result) & la condizione booleana per passare da uno stato all'altro \\
Automatic Rule Based on Sequence Player in State & la transizione scatta a fine
animazione \\
NOT Boolean / AND Boolean & negazione; congiunzione di due condizioni \\
Split Struct Pin & divide un vettore nei suoi componenti X, Y, Z \\
Blend Space 1D / Blend Space (2D) & mescola animazioni su uno o due assi \\
Horizontal / Vertical Axis, Grid Divisions & nome, minimo e massimo degli assi;
divisioni della griglia \\
Smoothing Time & tempo di passaggio fra i campioni \\
Blend Space Player & il nodo che suona un Blend Space nell'AnimGraph \\
Calculate Direction & l'angolo, da $-180$ a 180, fra la velocità e la rotazione \\
Get Actor Rotation / Get Velocity & la rotazione dell'actor; la sua velocità \\
Default Pawn Class (World Settings) & il personaggio del giocatore nel livello \\
Spring Arm: Target Arm Length, Socket Offset, Target Offset & lunghezza del braccio;
spostamento in fondo al braccio; spostamento dell'origine \\
Camera: Field of View & angolo di visuale \\
Use Pawn Control Rotation (Spring Arm) & la camera segue il controller \\
Use Controller Rotation Yaw (Pawn) & il corpo segue la rotazione del controller \\
Orient Rotation to Movement, Rotation Rate (Character Movement) & il corpo si gira verso
il movimento, con quella velocità \\
Max Walk Speed / Set Max Walk Speed & velocità massima a piedi (cm/s) \\
Retarget Animations, Export Animations & trasferisce un'animazione su un altro
scheletro \\
Input Mapping Context (IMC) & associa i tasti alle azioni in un contesto \\
Input Action (IA): Value Type, Triggers & l'azione, il tipo di valore, quando scatta \\
Add Mapping Context & attiva un IMC per il giocatore \\
Evento IA: Triggered, Started, Ongoing, Canceled, Completed & le fasi di un'azione \\
FInterp To & un passo di interpolazione fra Current e Target \\
Get World Delta Seconds & il tempo dell'ultimo frame \\
Animation Montage & animazione ``una tantum'', modulare, con sezioni e slot \\
Notify / Play Sound / Notify Track & eventi sulla linea del tempo; un suono; una traccia
di eventi \\
Montage Section & una parte con nome del Montage \\
Slot (DefaultGroup.DefaultSlot) & dove il Montage si applica (tutto il corpo, o parti) \\
Save Cached Pose / Use Cached Pose & salva e rilegge una posa nell'AnimGraph \\
Play Montage & suona un Montage; On Completed, On Blend Out, On Interrupted, Notify
Begin/End \\
Disable Input / Enable Input & toglie e restituisce l'input all'actor \\
\bottomrule
\end{tabular}
\caption{I nodi, gli asset e le impostazioni della nona lezione.}
\label{ue9:tab:nodi}
\end{table}

% ---------------------------------------------------------------------
\section*{Punti chiave del capitolo}
\addcontentsline{toc}{section}{Punti chiave del capitolo}

\begin{itemize}
  \item Una \textbf{State Machine} nell'AnimGraph organizza le animazioni in
        \textbf{stati} (WalkRun, Jump, Fall), ciascuno con la sua \textbf{Output
        Animation Pose}, e in \textbf{transizioni} con regole booleane: un ``loop di
        animazione secondo delle condizioni''. Va collegata all'Output Pose.
  \item Un \textbf{Blend Space 1D} mescola animazioni lungo un asse: per la locomozione,
        la \textbf{GroundSpeed}. I campioni vanno posti alle velocità reali del
        personaggio (Max Walk Speed, in centimetri al secondo): Idle a 0, Walk a 300,
        corsa a 600. Valori intermedi ``inventati'' producono obbrobri.
  \item Un \textbf{Character} non ha la camera: si aggiungono uno \textbf{Spring Arm} e,
        come suo figlio, una \textbf{Camera}. Distanza e altezza si regolano sullo
        Spring Arm (\textbf{Target Arm Length}, \textbf{Socket Offset}), mai spostando la
        camera; il \textbf{Target Offset} solo sulla Z e poco. Sono il lavoro dei
        \textbf{3C}: Camera, Control, Character.
  \item Tre spunte governano la rotazione: \textbf{Use Pawn Control Rotation} sullo
        Spring Arm (la camera segue il mouse), \textbf{Use Controller Rotation Yaw} sul
        Pawn (il corpo segue la camera), \textbf{Orient Rotation to Movement} nel
        Character Movement (il corpo si gira verso il movimento).
  \item Il \textbf{retargeting} porta un'animazione su un altro scheletro (Retarget
        Animations, Export Animations): comodo, ma il risultato dipende dall'animazione,
        e il root di Mixamo dà problemi.
  \item Fra salto e caduta: \textbf{Automatic Rule Based on Sequence Player in State} per
        passare a fine animazione, \textbf{NOT IsFalling} per tornare a terra, e la
        \textbf{velocità sulla Z} per distinguere un salto (Z $>$ 0) da una caduta
        (Z $<$ 0).
  \item L'\textbf{Enhanced Input} separa i tasti dalle azioni: l'\textbf{Input Action}
        (IA\_Run, Digital bool, con i suoi Triggers) è associata a uno o più tasti
        (Shift, il cerchio del controller) in un \textbf{Input Mapping Context}, che va
        attivato con \textbf{Add Mapping Context}. Lo sprint: Set Max Walk Speed a 600
        sul Triggered, a 300 sul Completed.
  \item \textbf{FInterp To} fa un solo passo per esecuzione: per un passaggio graduale
        va eseguito a ogni frame (Tick, timer, Timeline).
  \item Il \textbf{Blend Space 2D} aggiunge la \textbf{direzione} (\textbf{Calculate
        Direction}, da $-180$ a 180: positivo a destra, negativo a sinistra, $\pm 180$
        indietro): otto animazioni a 600, l'Idle ripetuto a 0, la corsa all'indietro ai
        due estremi, griglia 8$\times$8, Smoothing Time. Serve Use Controller Rotation Yaw
        acceso.
  \item Un \textbf{Animation Montage} è un'animazione ``una tantum'' e modulare, con
        \textbf{Notify} (suoni, VFX in un punto preciso), \textbf{sezioni} e
        \textbf{slot}. Per suonarlo, nell'AnimGraph: State Machine, \textbf{Save Cached
        Pose}, Use Cached Pose, \textbf{Slot DefaultSlot}, Output Pose. Si fa partire con
        \textbf{Play Montage}, bloccando prima movimento, input e rotazioni e
        restituendoli su On Blend Out e On Completed.
  \item Compito facoltativo: un nuovo personaggio da zero con lo \textbf{StackOBot},
        Blueprint e Animation Blueprint completi.
\end{itemize}

%
%  Tutte le etichette di questo capitolo hanno il prefisso "ue9:"
%  (ue9:cap:..., ue9:sec:..., ue9:fig:..., ue9:tab:...), e tutti gli
%  stili TikZ hanno il prefisso "ue9" (ue9node, ue9exec, ue9data, ...)
%  per non collidere con chiavi native di TikZ né con quelli delle
%  lezioni precedenti. Il capitolo è autonomo: non contiene \ref verso
%  altri capitoli (i rimandi alle lezioni precedenti sono a parole).
%  Le figure sono ricostruzioni schematiche di quanto mostrato e detto a
%  lezione, non copie delle schermate del docente. Dove la trascrizione
%  è troppo rovinata per ricostruire un dettaglio, il capitolo lo dice
%  invece di inventarlo.
% =====================================================================

\chapter[State Machine, Blend Space e Montage]{State Machine, Blend Space e Montage:
dare vita al nuovo personaggio}
\label{ue9:cap:animazioni}

Questa è la \textbf{nona lezione} del corso breve di \textbf{introduzione a Unreal
Engine}, e riprende esattamente dove si era fermata l'ottava. Alla fine della lezione
precedente il cavaliere importato da \textbf{Mixamo} aveva uno scheletro, le sue
animazioni e un \textbf{Animation Blueprint}, \textbf{ABP\_Knight}, che a ogni frame
leggeva dal personaggio due dati: la velocità sul piano (\textbf{GroundSpeed}) e il fatto
di essere in aria o meno (\textbf{IsFalling}). Mancava tutto il resto: chi decide
\emph{quale} animazione mostrare, e quando. ``Siccome ci siamo interrotti l'altra volta,
andiamo diretti a fare le animazioni: l'altra volta ci siamo interrotti al Blend
Space.''

La lezione costruisce, un pezzo alla volta, un personaggio giocabile completo. La
\textbf{State Machine} organizza le animazioni in \emph{stati} (camminare e correre,
saltare, cadere) e in \emph{transizioni} con delle condizioni; il \textbf{Blend Space
1D} mescola fermo, camminata e corsa in base alla velocità; una \textbf{camera} con il
suo \textbf{Spring Arm}, i grafi di movimento presi in prestito dal modello Third Person
e un nuovo input per lo \textbf{sprint}, costruito con l'\textbf{Enhanced Input},
rendono il cavaliere controllabile; il \textbf{retargeting} porta sul suo scheletro le
animazioni di salto e di caduta del manichino di Unreal. Dopo la pausa arrivano il
\textbf{Blend Space 2D}, che aggiunge alla velocità la \textbf{direzione} e permette di
camminare di lato e all'indietro, e gli \textbf{Animation Montage}, le animazioni
``una tantum'' (un attacco con la spada) che si inseriscono sopra il ciclo di
locomozione.

La lezione è \textbf{online}, su Zoom, come le precedenti, e la classe è ridottissima:
``vediamo un po' quanti siete: siamo in sei''; dopo la pausa ``siamo rimasti in
cinque'', poi ``uno, due, tre e quattro''; quando uno studente crasha (``Filippo è andato
via''), ``siamo rimasti noi tre''; alla fine, ``mi sembra che siamo rimasti in due''.
``Pochi ma buoni.'' Il ritmo è quello della dimostrazione dal vivo, interrotta più volte
dai problemi degli studenti, risolti in condivisione di schermo, e da un editor che
``oggi lagga'' in modo vistoso (``non capisco perché mi sta laggando oggi''). La pausa
arriva dopo la scelta delle animazioni da scaricare (``ci rivediamo alle 18''); si
chiude con un compito facoltativo e un ``ci sentiamo per martedì prossimo''.

Tre avvertenze di lettura, le stesse delle lezioni precedenti. La prima: la
trascrizione è automatica e storpia moltissimi nomi. Unreal è ``Arreal'', ``a Real'',
``ArrayL'', ``Array Alt'' e ``a Rial''; il \textbf{Pawn Owner} è ``Pwn Honor'' e
``Pwnowner''; il \textbf{Blend Space} è ``Blaine Space'', ``Blendspace'' e ``BlizzBase'';
lo stato \textbf{WalkRun} è ``walk run'', ``Warcran'', ``war crunch'' e ``walker run'';
l'\textbf{Idle} è ``Idol'' e ``idol''; \textbf{BP\_Knight} è ``Vip Night'', ``Pippi
Night 1'' e ``BP Night''; \textbf{BS\_Knight} è ``BS Night'' (``che dovrebbe essere un
po' più British''); il manichino \textbf{Manny/Quinn} è ``Many Quincy''; il
\textbf{root motion} è ``rut'' e ``outmotion''; lo \textbf{Strafe} è ``Strife'' e
``stair right''; la \textbf{Mesh} è ``Mesci''; l'\textbf{FInterp To} è ``finter2'';
\textbf{Calculate Direction} è ``calcolo interaction''; il \textbf{Save Cached Pose} è
``save cache, pause''; lo \textbf{StackOBot} è ``StaccoBot'' e ``stacco bot''; lo
\textbf{Yaw} è ``io'', ``yo'' e ``YAW''; l'\textbf{Orient Rotation to Movement} è
``orientation to movement'', ``anche rotacion'' e ``orient rotace tu mente tu''; lo
\textbf{Use Pawn Control Rotation} è ``us control pawn rotation'' e ``you spawn control
rotation''. Dove la ricostruzione è sicura il capitolo usa la forma corretta; dove non
lo è riporta la forma della trascrizione e lo dice. La seconda: all'inizio della
registrazione e durante la pausa il riconoscimento vocale produce lunghi blocchi di
testo senza senso e ripetuto (``grazie a tutti per l'attenzione che ci ha spiegato'',
``qual è la risposta che ci può essere?'', ``la situazione è questa, la migliore dal
2005'', la ``dottrina sul mercato del lavoro'', la ``Corte Suprema di Internet'', il
governo di Dubai, una lunga fila di ``di Sarajevo'', la ``struttura del giardino'', ``la
blue map'', un Vittorio Prodi chiamato in causa senza motivo): sono artefatti, non
parole di nessuno, e il capitolo li ignora. La terza: gran parte della lezione è una
dimostrazione dal vivo nell'editor, costruita a voce (``tirate fuori'', ``promote to
variable'', ``ce lo attacchiamo qui''), e in più punti il docente improvvisa, prova,
sbaglia e corregge, soprattutto nella parte sulla camera e in quella finale sul
Montage. Il capitolo ricostruisce i grafi da ciò che il docente dice mentre li
costruisce, distingue i tentativi falliti dalla soluzione finale, e completa i nomi
esatti dei nodi e dei pin di Unreal Engine 5; dove la trascrizione lascia un buco, lo
dichiara.

\paragraph{Come leggere questo capitolo.} Il capitolo segue l'ordine della lezione. Le
figure ricostruiscono la State Machine della locomozione, il Blend Space 1D sulla
velocità, la camera con lo Spring Arm, la catena dell'Enhanced Input, il grafo dello
sprint, il cerchio della direzione, il Blend Space 2D con le sue otto animazioni, la
catena di pose dell'AnimGraph con lo Slot del Montage e il grafo del tasto che fa
partire l'attacco. Le tabelle riassumono le transizioni della State Machine, i
parametri della camera, i valori e i trigger di una Input Action, la posizione delle
animazioni nel Blend Space 2D, gli errori incontrati durante la lezione e i nodi della
giornata. Qua e là, dei paragrafi di \textbf{approfondimento} collegano la lezione ai
libri di Marcos Romero e Brenden Sewell (\emph{Blueprints Visual Scripting for Unreal
Engine 5} e la seconda edizione dello stesso testo, scritta per Unreal Engine 4), i
testi di riferimento per questa parte del libro, e a qualche nozione tecnica generale su
Unreal Engine 5; sono segnalati come tali e non fanno parte della lezione.

Il percorso del capitolo:
\begin{enumerate}
  \item l'\textbf{inizio della lezione}: i corsi degli studenti e la lezione libera del
        25 (sezione~\ref{ue9:sec:inizio});
  \item il \textbf{ripasso} dell'Animation Blueprint (sezione~\ref{ue9:sec:ripasso});
  \item la \textbf{State Machine} della locomozione (sezione~\ref{ue9:sec:statemachine});
  \item il \textbf{Blend Space 1D}, con una parentesi sui problemi dell'editor (sezioni
        \ref{ue9:sec:blend1d} e~\ref{ue9:sec:layout});
  \item il personaggio in gioco: \textbf{Default Pawn Class}, \textbf{camera} e
        \textbf{Spring Arm}, grafi di movimento e rotazioni (sezioni
        da~\ref{ue9:sec:pawnclass} a~\ref{ue9:sec:rotazioni});
  \item il \textbf{retargeting} del salto e della caduta, e le \textbf{regole di
        transizione} (sezioni \ref{ue9:sec:retarget} e~\ref{ue9:sec:transizioni});
  \item lo \textbf{sprint} con l'\textbf{Enhanced Input} e il tentativo con FInterp To
        (sezioni da~\ref{ue9:sec:enhanced} a~\ref{ue9:sec:finterp});
  \item le \textbf{animazioni da scaricare}, la pausa e i \textbf{feedback} (sezioni
        \ref{ue9:sec:download} e~\ref{ue9:sec:feedback});
  \item il \textbf{Blend Space 2D} e la \textbf{direzione} (sezioni
        \ref{ue9:sec:direction} e~\ref{ue9:sec:blend2d});
  \item gli \textbf{Animation Montage}: l'asset, i Notify, le sezioni, gli Slot e il
        tasto che li fa partire (sezioni da~\ref{ue9:sec:montage}
        a~\ref{ue9:sec:playmontage});
  \item il \textbf{compito} sullo StackOBot e la chiusura
        (sezione~\ref{ue9:sec:chiusura}).
\end{enumerate}

% ---------------------------------------------------------------------
\section{Prima della lezione: i corsi degli studenti e la lezione del 25}
\label{ue9:sec:inizio}

La registrazione si apre con qualche minuto di attesa, in cui il riconoscimento vocale
produce soprattutto rumore. Poi i saluti (``ragazzi, come state? Buonasera, buonasera''),
un controllo rapido (``nessun problema sulla telecamera'': stavolta la webcam del docente
funziona) e la decisione di aspettare ``altri cinque minuti, così vediamo che tutti siate
in carreggiata''. Uno studente, Tommaso, ha dei problemi di collegamento. Il docente
sistema le finestre (``oggi c'è il secondo schermo'') e parte: ``bando alle ciance''.

\paragraph{I corsi di indirizzo.} Prima di cominciare, il docente chiede ai presenti
quale corso seguiranno ``quando inizierete a settembre'': il corso di Unreal è infatti
un'introduzione comune, a cui seguono gli indirizzi specialistici. Le risposte dei sei
presenti:
\begin{itemize}
  \item \textbf{3D art};
  \item \textbf{game design} (due studenti);
  \item \textbf{narrative design};
  \item \textbf{character animation} (``perfetto, character animation'').
\end{itemize}
La domanda serve a tarare la lezione: chi farà character animation, in particolare, è
il destinatario naturale di quasi tutto quello che si vedrà oggi.

\paragraph{La lezione libera del 25.} ``Stiamo praticamente arrivando a una lezione che
di solito noi lasciamo abbastanza \emph{free}, come argomenti'': la lezione del
\textbf{25} sarà ``fondamentalmente la mia ultima lezione di base''. Se non ci sono
richieste particolari, si faranno ``i menu'', cioè ``la \textbf{UI}'': ``che si intende
la UI: menu, bottoni, scritte eccetera''. In alternativa, ``se ci sono degli argomenti
particolari'', da concordare fra gli studenti, ``possiamo tranquillamente affrontarli'',
oppure ``dividere anche la lezione in due: fare un piccolo menu e poi fare anche la parte
in più''. ``Nel caso, sentitevi tra di voi.'' La richiesta concreta arriverà durante la
pausa (sezione~\ref{ue9:sec:feedback}).

% ---------------------------------------------------------------------
\section{Il ripasso: l'Animation Blueprint del cavaliere}
\label{ue9:sec:ripasso}

``Eravamo rimasti al nostro bellissimo \textbf{ABP\_Knight}.'' Il docente ha il
cavaliere (\emph{Knight}, il paladino di Mixamo); ``altri di voi avevano scelto il
proprio personaggio''. Nella preview dell'Animation Blueprint il personaggio è ancora
``il nostro bel character in T-pose'': l'AnimGraph non decide ancora nessuna
animazione.

\paragraph{Initialize e Update.} L'ultima cosa fatta era l'evento \textbf{Event Blueprint
Initialize Animation}: ``quando l'animazione si inizializza'', una volta sola. L'altro
evento, \textbf{Event Blueprint Update Animation}, scatta invece ogni volta che
``si cambia qualsiasi particolare all'interno di un'animazione, la rotazione di un braccio
o un minimo particolare'': ``quello che è fondamentalmente il nostro Tick sulle
animazioni''. ``Quando abbiamo visto Begin Play e Tick: questi praticamente sono i loro
corrispettivi.''

\paragraph{Il Pawn Owner e le due referenze.} Nell'Initialize si prende il \textbf{Pawn
Owner} con \textbf{Try Get Pawn Owner}. ``Ricordiamoci: il Pawn Owner è l'entità che
controlla il character'': il docente distingue il \emph{pawn}, ``quello che lo controlla'',
e il \emph{character} stesso, ``che ha la mesh, lo scheletro e tutto quanto''. Con il
\textbf{Cast To} BP\_Knight si ottiene una referenza al cavaliere, e si salvano in due
variabili il cavaliere e il suo \textbf{Character Movement Component}: ``queste due
variabili ci serviranno per comunicare in modo più diretto con il nostro cavaliere''.

\paragraph{Il Character Movement.} ``Dentro il Blueprint, qualsiasi character lo troverete
qui sotto, Character Movement'': selezionandolo, il pannello Details a destra ``presenta
tutte le particolarità di movimento'', per esempio la \textbf{gravità} (\emph{Gravity
Scale}), la \textbf{Max Acceleration}, la \textbf{Max Step Height} (l'altezza massima di
un gradino che il personaggio sale senza saltare). ``Vogliamo farlo andare più veloce,
meno veloce? Se volessimo farlo sprintare, dovremmo lavorare su questa particolarità''; e
infatti ``dopo faremo anche uno sprintino molto veloce, per differenziare le
animazioni''.

\paragraph{GroundSpeed e IsFalling.} Nell'Update, dopo una \textbf{Sequence}, si prende
``la velocity attuale del nostro character'', a ogni frame, e se ne salva la lunghezza in
\textbf{GroundSpeed}. Allo stesso modo, con il nodo \textbf{Is Falling} del Character
Movement, ``andiamo a vedere se effettivamente il nostro character sta cadendo o meno'':
``questo blocco ci dice letteralmente se sta cadendo, e noi riportiamo con un booleano,
true o false''. ``Di conseguenza, tramite questi dati, andremo a differenziare le
animazioni nella nostra State Machine.''

% ---------------------------------------------------------------------
\section{La State Machine della locomozione}
\label{ue9:sec:statemachine}

\paragraph{Creare la State Machine.} La State Machine ``la possiamo invocare
nell'AnimGraph'': tasto destro nell'AnimGraph (``ricordatevi che c'è il blocco
\textbf{Output Pose}''), \textbf{State Machine}, e la si chiama \textbf{Locomotion}:
``questo è il nome classico per la locomozione, appunto il movimento''. ``Una State
Machine, appunto, è una macchina a stati.'' Con un doppio clic si entra nel suo grafo.

\paragraph{Entry, Conduit, State, Alias.} Dentro c'è un solo nodo, \textbf{Entry}. ``Se
noi trasciniamo, abbiamo una freccia che esce fuori''; lasciandola nel vuoto, il menu non
è ``tutto l'elenco gigantesco'' dei nodi dei Blueprint, ma tre sole voci:
\textbf{Conduit}, \textbf{State} e \textbf{Alias}. ``A noi adesso interessa vedere
quelli che sono gli stati.'' Il primo stato, invece di \emph{Locomotion}, si chiama
\textbf{WalkRun}: ``differenziamo''.

\paragraph{Un'altra inception.} ``Ecco a voi che, ancora sempre, un po' un
\emph{inception}'': con un doppio clic si entra anche nello stato. Dentro, un nodo
\textbf{Output Animation Pose}: ``in questo stato risiede una tipologia di output'',
l'animazione che lo stato produce. Il docente vi collega \emph{Walking 1}. La preview
però non si aggiorna, ``perché non abbiamo collegato la State Machine all'output'':
bisogna tornare nell'AnimGraph e collegare il nodo Locomotion all'Output Pose. Fatto
questo, il cavaliere cammina, ma una volta sola: ``l'animazione non è in loop'', e basta
spuntare il loop nei Details del nodo. ``Ecco a voi che cosa fa la State Machine'':
``dentro Locomotion lui entra, dice ok, inizio, passo e vado direttamente a WalkRun''; in
WalkRun ``c'è una richiesta di un output, quindi di un'animazione che noi gli diamo''.

\paragraph{Gli stati Jump e Fall.} ``Ma la cosa interessante della State Machine'' è che
si possono creare altri stati: \textbf{Jump} e \textbf{Fall}. ``A livello di logica, la
State Machine vuole essere'' il racconto di ``come entrano e perché entrano'' le
animazioni nei loro stadi. Una freccia da WalkRun a Jump ``vuol dire che dallo stadio
WalkRun passerà allo stadio Jump''; una freccia inversa vuol dire che dal Jump ``potrà
passare di nuovo al WalkRun''. Lo schema di oggi ha tre stati
(figura~\ref{ue9:fig:statemachine}): si entra in WalkRun; ``se starà cadendo'' (la
variabile IsFalling diventa vera) si passa a Jump o a Fall; ``e far sì che poi, quando
lui finisca di cadere, ritorni'' a WalkRun. ``Quindi fondamentalmente creiamo un loop di
animazione secondo delle condizioni.'' Le animazioni di salto e di caduta non ci sono
ancora: ``la faremo, la ritargetteremo'' (sezione~\ref{ue9:sec:retarget}).

\paragraph{Le transizioni.} ``Come facciamo a mettere queste condizioni? Tramite questi
nodini qua'': sulle frecce fra due stati compare un'icona tonda, la \textbf{transition
rule}, che ``ci dice il comportamento effettivo che dovrà avere'' il personaggio. ``Quando
è che passo da WalkRun a Jump, a livello di animazione? Quando, ad esempio, stai cadendo.
Quando passo da Fall a WalkRun? Se non sta più cadendo. Questa è la logica generale.''

\begin{figure}[htbp]
\centering
\begin{tikzpicture}[font=\footnotesize,
  ue9state/.style={draw, thick, rounded corners=6pt, fill=orange!12, align=center,
                   minimum width=2.4cm, minimum height=1.0cm},
  ue9entry/.style={draw, thick, fill=gray!15, rounded corners=2pt, align=center,
                   minimum width=1.4cm, minimum height=0.7cm},
  ue9tr/.style={-{Stealth[length=2.6mm]}, thick},
  ue9rule/.style={draw, circle, fill=white, inner sep=1.5pt, font=\scriptsize}]
  \node[ue9entry] (e) at (0,0) {Entry};
  \node[ue9state] (w) at (3.4,0) {\textbf{WalkRun}\\{\scriptsize Blend Space}};
  \node[ue9state] (j) at (8.2,1.5) {\textbf{Jump}\\{\scriptsize MM\_Jump}};
  \node[ue9state] (f) at (8.2,-1.5) {\textbf{Fall}\\{\scriptsize MM\_Fall\_Loop}};
  \draw[ue9tr] (e) -- (w);
  \draw[ue9tr] (w.north east) to[bend left=12] node[ue9rule, pos=0.5]{1} (j.west);
  \draw[ue9tr] (j.south) -- node[ue9rule]{2} (f.north);
  \draw[ue9tr] (w.south east) to[bend left=12] node[ue9rule, pos=0.5]{3} (f.west);
  \draw[ue9tr] (f.south west) to[bend left=30] node[ue9rule, pos=0.5]{4} (w.south);
  \node[align=left, text width=12.5cm, anchor=north west, font=\scriptsize] at
    (-1,-2.9) {\textbf{1} WalkRun $\rightarrow$ Jump: IsFalling \textbf{AND}
    VelocityZ $>$ 0 (sta salendo)\\
    \textbf{2} Jump $\rightarrow$ Fall: \emph{Automatic Rule Based on Sequence Player in
    State} (finito il salto, la caduta)\\
    \textbf{3} WalkRun $\rightarrow$ Fall: IsFalling \textbf{AND} VelocityZ $<$ 0 (cade
    senza aver saltato)\\
    \textbf{4} Fall $\rightarrow$ WalkRun: \textbf{NOT} IsFalling (è tornato a terra)};
\end{tikzpicture}
\caption{La State Machine Locomotion alla fine della lezione: tre stati e quattro
transizioni. Le condizioni con la velocità sulla Z sono quelle aggiunte nella seconda
versione (sezione~\ref{ue9:sec:transizioni}); nella prima, la 1 e la 3 erano soltanto
IsFalling. Schema ricostruito.}
\label{ue9:fig:statemachine}
\end{figure}

\paragraph{Approfondimento: Conduit e Alias.} I tre tipi di nodo proposti dal menu della
State Machine hanno ruoli diversi. Uno \textbf{State} contiene un grafo che produce una
posa (un'animazione, un Blend Space, una combinazione). Un \textbf{Conduit} non produce
pose: è un incrocio, uno snodo da cui partono più transizioni, utile quando molti stati
condividono la stessa condizione d'ingresso o d'uscita e non si vuole ripeterla su ogni
freccia. Un \textbf{State Alias} è un ``alias'' di uno o più stati: le transizioni che
partono dall'alias valgono come se partissero da ciascuno degli stati che rappresenta
(per esempio, ``da qualunque stato, se il personaggio muore, vai a Death''), e
risparmiano un groviglio di frecce. Romero e Sewell (capitolo 17) costruiscono la loro
prima State Machine, \emph{Char States}, con due soli stati, \emph{Idle} e
\emph{Moving}, e due regole sulla velocità (``Speed maggiore di 5'' e ``minore di 5''); la
scelta della lezione, un solo stato WalkRun che contiene un Blend Space, sposta la
distinzione fra fermo e in movimento dentro il Blend Space, e lascia alla State Machine
le sole differenze ``qualitative'' (a terra, in salto, in caduta).

% ---------------------------------------------------------------------
\section{Il Blend Space 1D}
\label{ue9:sec:blend1d}

``Visto che WalkRun non è solo un'animazione, è un insieme'': sono tre animazioni,
\textbf{Idle}, \textbf{Walk} e quella che il docente ha chiamato \textbf{Slow Run}.
``Dobbiamo mixarle insieme. Come si fa? Perché non possiamo mettere un'animazione solo
così: dobbiamo far sì che sia dipendente dalla nostra velocità.'' Lo strumento è il
\textbf{Blend Space}.

\paragraph{Creare il Blend Space 1D.} ``Fate solamente il Blend Space 1D.'' Invece di
cercarlo nei menu, si può scrivere nella barra di ricerca del menu del tasto destro
``blend space'', e scegliere \textbf{Blend Space 1D}: ``già vi dice il percorso, che sta
in Animation, Legacy''. L'icona ha un colore ``rosa pelle''. Aperto, si troverebbe il
personaggio in T-pose; ma molti l'avevano già creato la volta precedente (``sì, mi sa che
già l'avevamo fatto'').

\paragraph{Che cos'è un Blend Space.} ``Un Blend Space è l'azione di \emph{blendare}
qualche cosa'', di ``interpolare o comunque mixare le animazioni'': un asset che ``ci dà
la possibilità di poter traslare da un'animazione all'altra''. Si vede subito: dal
pannello \textbf{Asset Browser} dell'editor (``se non lo trovate, ripeto: Window, Asset
Browser'') si trascina un'animazione sul grafico; tenendo premuto \textbf{Shift} la si
aggancia ``su dei punti precisi'' della griglia. Con un'animazione sola, il personaggio
fa solo quella; aggiungendone una seconda, il grafico ``crea una linea di passaggio tra
un'animazione e l'altra''.

\paragraph{Il parametro, come l'alfa.} ``Ma importante: qual è questo dato?'' Il docente
lo collega all'\textbf{alfa} di un Lerp visto in una lezione precedente: ``quel puntatore
che punta, dipendente da uno spazio tra due punti, dove sta al momento''. Qui il valore
dirà ``ehi, guarda che sto camminando, sto andando a velocità 300, oppure sto correndo, a
velocità 600, oppure sono fermo''. Il Blend Space \textbf{1D} ``ci dà solamente la
possibilità di poter usare un asse'': nei Details, \textbf{Horizontal Axis} (la sezione
potrebbe essere chiusa: ``basta che la aprite''). L'asse si chiama \textbf{GroundSpeed},
``come quella che abbiamo trovato all'interno dell'Animation Blueprint''; poi chiede un
\textbf{Minimum Axis Value} e un \textbf{Maximum Axis Value}.

\paragraph{Le velocità sono relative.} ``Importante, ragazzi, quando si va a vedere le
animazioni: dovete sempre farlo relativo, se si parla di velocità'', relativo cioè alla
velocità reale del personaggio. Dove si legge? Nel Blueprint del personaggio,
BP\_Knight, selezionando \textbf{Character Movement}: ``potete andare a vedere tutti i
dati che vi servono''. Quello che serve è la \textbf{Max Walk Speed}, in alto, nella
sezione \emph{Character Movement: Walking}: ``è 600 centimetri al secondo''. ``Unreal
ragiona per centimetri'': il personaggio va a 6 metri al secondo.

\paragraph{Dividere la velocità a metà.} ``Visto che noi abbiamo due stadi, uno di walk e
uno di run'', si vuole che ``la walk sia un movimento normale'': di default il personaggio
cammina, e corre solo se il giocatore lo chiede con un tasto (sezione
\ref{ue9:sec:enhanced}). ``Quindi la dividiamo a metà'': \textbf{Max Walk Speed a 300}.
Nel Blend Space l'asse va da \textbf{0} a \textbf{600}, ``in modo tale che a metà la
nostra walk sia a 300'' (figura~\ref{ue9:fig:blend1d}). A uno studente l'asse arriva solo
a 100: ``ground speed, devi metterlo a 600, lo trovi qua''. E un consiglio: dopo aver
cambiato il massimo, ``sposta i pallini, perché poi ti switcha tutto''; il valore di
ciascun campione si può anche scrivere a mano nei Details.

\begin{figure}[htbp]
\centering
\begin{tikzpicture}[font=\footnotesize,
  ue9pt/.style={draw=black, fill=white, thick, diamond, inner sep=2.2pt}]
  \draw[-{Stealth[length=2.5mm]}, thick] (0,0) -- (11.2,0) node[right]{GroundSpeed};
  \foreach \x/\v in {0/0, 2.5/150, 5/300, 7.5/450, 10/600}
    \draw (\x,0.1) -- (\x,-0.1) node[below, font=\scriptsize]{\v};
  \node[ue9pt] (i) at (0,0) {}; \node[ue9pt] (w) at (5,0) {}; \node[ue9pt] (r) at
    (10,0) {};
  \node[above=6pt of i, font=\scriptsize\bfseries] {Idle};
  \node[above=6pt of w, font=\scriptsize\bfseries, align=center] {Walk\\
    {\normalfont Max Walk Speed di default}};
  \node[above=6pt of r, font=\scriptsize\bfseries, align=center] {Slow Run\\
    {\normalfont con lo sprint}};
  \fill[red!60!black] (3.33,0) circle (0.12);
  \draw[red!60!black, thick, ->] (3.33,-1.2) -- (3.33,-0.25);
  \node[red!50!black, font=\scriptsize, align=center] at (3.33,-1.55) {200: un misto
    fra Idle e Walk, ``un obbrobrio''};
\end{tikzpicture}
\caption{Il Blend Space 1D BS\_Knight: un solo asse, la GroundSpeed, da 0 a 600, con
Idle a 0, Walk a 300 (la Max Walk Speed impostata nel Character Movement) e la corsa a
600 (la velocità dello sprint). Un valore intermedio, come 200, mescola due animazioni
nate per velocità diverse. Schema ricostruito.}
\label{ue9:fig:blend1d}
\end{figure}

\paragraph{Il Blendspace Player.} Tornati in ABP\_Knight (``normale, ragazzi, che vi dia
dei warning, perché ancora non abbiamo le condizioni'': ``lui dice, ma io non ho
condizioni, quindi faccio tutto in una botta''), si entra nello stato WalkRun. Nell'Asset
Browser ``è apparso BS\_Knight''; trascinato nel grafo, ``ci presenta un \textbf{Blend
Space Player}'', diverso dal nodo di un'animazione normale: ha un \textbf{valore in
entrata}, la GroundSpeed. A quel pin si collega un \textbf{Get} della variabile
GroundSpeed, calcolata nell'Event Graph (la lunghezza della velocity), e l'uscita del
Blend Space Player si collega al \textbf{Result} dell'Output Animation Pose, ``in modo
tale che vada in output''.

\paragraph{Get con Ctrl, Set con Alt.} Un promemoria: trascinando una variabile nel
grafo, ``tenete premuto \textbf{Ctrl} e tirate fuori direttamente un Get''; ``se invece
volete fare un Set, tenete premuto \textbf{Alt}''. Nel grafo dello stato l'Alt non
funziona: ``non si può settare all'interno'' (gli stati dell'AnimGraph leggono le
variabili, non le scrivono), ``però nei Blueprint, di base, se premete Alt vi setta la
variabile''.

\paragraph{La prova nella preview.} Compilato, ``il nostro character nella preview adesso
sta in idle, perché giustamente ha velocità 0''. Cambiando il valore di GroundSpeed nella
preview: a \textbf{300} cammina; a \textbf{200} ``sta interpolando tra l'animazione di
idle e l'animazione di walk'': ``mettere un valore a 200 non è buona cosa'', perché
``deve essere relativo al vostro movimento, se no vengono degli obbrobri''; a
\textbf{600} ``il nostro carissimo corre'' (``guarda come è eroico''). ``Quindi adesso
andiamo a cambiare le animazioni tramite un valore.''

\paragraph{Allineare i nodi con Q.} ``Se volete fare un po' di ordine'': selezionati due
o più nodi, il tasto \textbf{Q} li allinea (\emph{Straighten Connections}). ``Tutto bello
ordinato.''

\paragraph{Una prova con un solo stato.} Prima di andare avanti, un test in gioco:
nell'AnimGraph si tolgono le frecce verso Jump e Fall (selezionate e cancellate, o
staccate con \textbf{Alt}+clic sul pin), e si lascia ``solamente WalkRun, che vada in
entrata''.

\paragraph{Approfondimento: il Blend Space nel libro.} Romero e Sewell (capitolo 17)
definiscono il Blend Space come ``un tipo di asset che permette di mescolare animazioni
in base ai valori di uno o due parametri''; nel loro esempio, tratto dal modello Third
Person, Idle è a 0, Walk a 93,75 e Run a 375, con una Max Walk Speed di 375. Il principio
è lo stesso della lezione: i campioni vanno posti dove la velocità reale del personaggio
corrisponde alla velocità ``disegnata'' nell'animazione. Se un'animazione di camminata è
stata animata per una velocità di 150 centimetri al secondo e il personaggio si muove a
300, i piedi ``scivolano'' sul pavimento (\emph{foot sliding}); è il motivo pratico
dietro il ``deve essere relativo'' del docente. Gli autori ricordano anche che il valore
del parametro si può cambiare nell'\emph{Anim Preview Editor}, in basso a destra
nell'editor dell'Animation Blueprint, senza andare in Play: è ciò che il docente fa per
mostrare i valori 200, 300 e 600.

% ---------------------------------------------------------------------
\section{Una parentesi: ``questo asset non ha schede ancorate''}
\label{ue9:sec:layout}

Mentre si sistema la Max Walk Speed, uno studente interviene: ``io ho un problema col
Blueprint: quando apro la scheda Blueprint mi dice: questo asset non ha ancora schede
ancorate''. Il docente ha l'editor in inglese (``di solito gli errori me li ricordo in
inglese'') e chiede di vedere: lo studente condivide prima soltanto una scheda del
browser, poi lo schermo intero.

\paragraph{Il messaggio.} Aprendo BP\_Knight, l'editor del Blueprint è vuoto, senza
pannelli. Il testo completo del messaggio, letto ad alta voce, spiega che l'editor
dell'asset è stato aperto in un sistema di finestre fluttuanti e nascoste, ``non mostra
le schede ancorate nell'interfaccia principale'', e che ``questo comportamento può essere
facilmente ripristinato'' ripristinando il layout dell'editor o modificando le preferenze
perché gli asset si aprano sempre ancorati alla finestra principale (in inglese,
\emph{Editor Preferences $\rightarrow$ General $\rightarrow$ Appearance $\rightarrow$
Asset Editor Open Location}). Inoltre, sopra il grafo compare una nota: il Blueprint è
\emph{data-only}, e c'è un link \textbf{Open Full Blueprint Editor}.

\paragraph{La soluzione.} ``Mi sa che è meglio ripristinare il layout di default'':
dalla barra in alto, \textbf{Window} (in italiano \emph{Finestra}) $\rightarrow$
\textbf{Load Layout} (\emph{Carica il layout}) $\rightarrow$ \textbf{Default Editor
Layout} (\emph{Editor predefinito}). Poi, sulla nota, ``apri l'editor Blueprint
completo''. ``Perfetto, eccolo qua.''

\paragraph{La lingua dell'editor.} ``Comunque, per mettere l'editor in inglese'' senza
riavviare: si cerca \emph{Language} nelle preferenze. Il docente prova prima in
\textbf{Editor Preferences} (sotto \emph{Edit}), scrivendo ``lingua'' e poi ``language''
(``non gli piace che mettiamo due cose insieme''), poi nei \textbf{Project Settings};
alla fine la voce \emph{Language} viene trovata con la ricerca (``ah, ecco perché'').
Dalla trascrizione non si capisce con certezza in quale delle due finestre. Lo studente
ringrazia (``mi hai fatto sentire molto ignorante''; ``verrai buttato fuori dalla
classe, tranquillo''), e il docente annuncia di volerlo usare in inglese anche lui ``per
la prossima volta''.

\paragraph{Approfondimento: dov'è la lingua.} In Unreal Engine 5 la lingua
dell'interfaccia è in \textbf{Edit $\rightarrow$ Editor Preferences $\rightarrow$ General
$\rightarrow$ Region \& Language $\rightarrow$ Editor Language}; nei Project Settings c'è
invece la lingua del \emph{gioco} (la localizzazione), che è un'altra cosa. Lavorare
con l'editor in inglese è una scelta pratica: la documentazione, i tutorial, i libri e i
forum usano i nomi inglesi dei nodi e dei pannelli, e la ricerca dei nodi nel menu del
tasto destro accetta i nomi inglesi anche quando l'interfaccia è tradotta.

% ---------------------------------------------------------------------
\section{Il cavaliere in gioco: la Default Pawn Class}
\label{ue9:sec:pawnclass}

``Adesso vi faccio vedere come impostiamo il player nuovo.'' Si aprono i \textbf{World
Settings} del livello (``qualsiasi cosa che non trovate, c'è sempre Window, e troverete
World Settings''); fra le voci della Game Mode c'è \textbf{Default Pawn Class}.
``Ricordiamo che il Pawn è appunto l'entità che controlla il player.'' Nel menu a
tendina c'è BP\_Knight, il personaggio finito che il docente aveva mostrato alla fine
della lezione precedente, e c'è BP\_Knight1, ``quello nuovo'', che si può anche cercare
scrivendone il nome. Si sceglie BP\_Knight1 e si va in Play.

\paragraph{Il primo problema: nessuna camera.} ``Ecco già un problemino.'' ``Fate
attenzione, perché i character non vengono con la camera'': un Blueprint creato dalla
classe Character ha la capsula, la freccia, la mesh e il Character Movement, ma non una
camera. In Play, la visuale finisce in un punto qualsiasi del corpo (``questo
meraviglioso\dots'', la frase si perde). Si esce dal Play e si apre BP\_Knight1: ``come
potete vedere, non c'è una camera''.

% ---------------------------------------------------------------------
\section{La camera: Camera e Spring Arm}
\label{ue9:sec:camera}

``La camera è impostata su due cose particolari, ragazzi, molto importanti'': la
\textbf{camera} stessa e il suo \textbf{braccio}. ``Immaginatelo come lo stand'': il
braccio dice ``quanto è lontana'' la camera, ``quanto offset ha''.

\paragraph{Aggiungere i due componenti.} Si aggiunge un componente \textbf{Camera}
(``camera normale''): ``eccola qua, che si è già attaccata'' al personaggio, ma in una
posizione che non va bene (``al suo lato B''). Poi uno \textbf{Spring Arm} (il docente
lo chiama prima ``arm length'', poi si corregge: ``perdonate, l'arm length è un
parametro; lo Spring Arm''). La gerarchia è importante: \textbf{lo Spring Arm deve avere
la camera dentro di sé}, cioè la Camera va trascinata sopra lo Spring Arm e diventa sua
figlia. Selezionando lo Spring Arm, nel viewport compare ``un piccolo tracciato rosso'':
``la lunghezza del nostro braccio'', il braccio della camera
(figura~\ref{ue9:fig:springarm}).

\paragraph{I 3C.} ``Ricordatevi, ragazzi, molto importante, perché la camera è
importantissima: ci sono delle figure che si chiamano \textbf{3C}, \emph{Camera, Control
e Character}, che lavorano solo su questo.'' Da qui una regola: ``se volete allontanare la
camera o cambiarle l'offset, \textbf{non cambiate la location della camera}. Non lo fate.
Controllatelo sempre da qui, dallo Spring Arm.''

\paragraph{Target Arm Length.} Il parametro \textbf{Target Arm Length} ``ci permette di
aumentare o diminuire la distanza della camera''. Valore scelto: \textbf{300}.

\paragraph{Target Offset e Socket Offset.} Ci sono due offset, e il docente ha una
preferenza netta. Il \textbf{Target Offset} ``usatelo solamente sulla Z, di poco,
veramente di poco, massimo 20''; ``e fate tutto con il \textbf{Socket Offset}''. Il
motivo: ``se voi cambiate l'offset del target, la camera vi può più facilmente
compenetrare con le cose''. Con il Socket Offset, invece, ``vedrete che la camera va sul
muro e torna'': ``va più vicino'', cioè ``riduce il suo arm length, fondamentalmente''.
Dopo la prima prova in Play (``lo testiamo, ragazzi?''), la camera è un po' alta: il
Socket Offset sulla Z viene abbassato a \textbf{40}.

\begin{figure}[htbp]
\centering
\begin{tikzpicture}[font=\footnotesize]
  % capsula
  \draw[thick, fill=blue!8, rounded corners=14pt] (0,0) rectangle (1.1,2.6);
  \node[align=center, font=\scriptsize] at (0.55,1.1) {Cap-\\sule};
  \node[font=\scriptsize, below] at (0.55,-0.2) {BP\_Knight1};
  \fill (0.55,2.0) circle (0.06);
  \node[font=\scriptsize, anchor=north west] at (0.62,1.95) {pivot};
  % braccio
  \draw[very thick, red!70!black] (0.55,2.0) -- (-4.2,2.0);
  \node[red!60!black, font=\scriptsize, above] at (-1.8,2.0) {Target Arm Length
    = 300};
  % socket offset
  \draw[thick, dashed, -{Stealth[length=2mm]}, green!40!black] (-4.2,2.0) --
    (-4.2,2.8);
  \node[green!40!black, font=\scriptsize, anchor=east, align=right] at (-4.9,2.3)
    {Socket Offset\\Z = 40};
  % camera
  \draw[thick, fill=gray!30] (-4.75,2.6) rectangle (-3.95,3.05);
  \draw[thick, fill=gray!30] (-3.95,2.7) -- (-3.6,2.55) -- (-3.6,3.1) -- (-3.95,2.95);
  \node[font=\scriptsize, above] at (-4.3,3.1) {Camera (figlia)};
  % target offset
  \draw[thick, dashed, -{Stealth[length=2mm]}, orange!70!black] (0.55,2.0) --
    (0.55,2.35);
  \node[orange!60!black, font=\scriptsize, anchor=south, align=center] at (0.55,2.4)
    {Target Offset: solo Z,\\al massimo $\approx$ 20};
  % muro
  \fill[gray!40] (-2.6,-0.2) rectangle (-2.35,1.4);
  \draw[thick] (-2.6,-0.2) rectangle (-2.35,1.4);
  \node[font=\scriptsize, align=center] at (-2.47,-0.6) {muro};
  \draw (-6,-0.2) -- (3.5,-0.2);
  % gerarchia
  \node[draw, thick, rounded corners=3pt, fill=gray!8, align=left, font=\scriptsize,
    anchor=north west] at (3.3,2.2) {\textbf{Components}\\
    Capsule Component\\
    \quad Arrow\\
    \quad Mesh\\
    \quad \textbf{Spring Arm}\\
    \qquad \textbf{Camera}\\
    Character Movement};
\end{tikzpicture}
\caption{La camera del cavaliere. La Camera è figlia dello Spring Arm, e la distanza e
l'altezza si regolano sul braccio, mai spostando la camera. Se fra il personaggio e la
camera si mette un ostacolo (il muro), lo Spring Arm accorcia il braccio e la camera
``torna'' verso il personaggio; con un Target Offset grande, invece, è più facile che la
camera entri negli oggetti. Schema ricostruito, valori della lezione.}
\label{ue9:fig:springarm}
\end{figure}

\paragraph{Il Field of View.} Più avanti, mentre aspetta uno studente, il docente mostra
che ``nella camera è presente anche un po' di dati con cui potete giocarci'', per esempio
il \textbf{Field of View}. Lo abbassa a \textbf{80}: ``sappiate che tanti giochi horror
usano il field of view a 80. Questo è un piccolo tip.''

\begin{table}[htbp]
\centering
\small
\begin{tabular}{@{}p{3.6cm}p{2.0cm}p{8.4cm}@{}}
\toprule
\textbf{Parametro} & \textbf{Valore} & \textbf{Effetto e consiglio del docente} \\
\midrule
Spring Arm $\rightarrow$ Target Arm Length & 300 & distanza della camera dal
personaggio \\
Spring Arm $\rightarrow$ Socket Offset & Z = 40 & spostamento della camera alla fine del
braccio; da preferire \\
Spring Arm $\rightarrow$ Target Offset & Z $\le$ 20 & spostamento dell'origine del
braccio; poco e solo sulla Z, perché fa compenetrare la camera \\
Spring Arm $\rightarrow$ Use Pawn Control Rotation & attivo & il braccio ruota con il
mouse (il controller); va sullo Spring Arm, non sulla Camera \\
Camera $\rightarrow$ Field of View & 80 & angolo di visuale; 80 è tipico dei giochi
horror \\
Camera $\rightarrow$ Location & non toccare & la posizione si regola sempre dallo Spring
Arm \\
\bottomrule
\end{tabular}
\caption{I parametri della camera toccati durante la lezione.}
\label{ue9:tab:camera}
\end{table}

\paragraph{Approfondimento: perché lo Spring Arm si accorcia.} Lo \textbf{Spring Arm
Component} fa, a ogni frame, una traccia (un \emph{sweep} di una piccola sfera, di
raggio \emph{Probe Size}) dal suo punto d'origine verso la posizione desiderata della
camera; se incontra un ostacolo nel canale \emph{Probe Channel} (di default
\emph{Camera}), accorcia il braccio fino al punto di impatto, e lo riallunga quando
l'ostacolo sparisce (la proprietà è \emph{Do Collision Test}). Il \emph{Target Offset}
sposta l'\emph{origine} di quella traccia, nello spazio del mondo: spostarla molto può
portare l'origine dentro un muro o dietro un oggetto, e allora il test non ``vede'' più
l'ostacolo; il \emph{Socket Offset} sposta invece la \emph{fine} del braccio, nello
spazio del braccio stesso, e la traccia resta coerente. È la ragione tecnica del
consiglio del docente. Lo Spring Arm ha anche le opzioni \emph{Enable Camera Lag} ed
\emph{Enable Camera Rotation Lag}, che fanno ``inseguire'' la camera con un leggero
ritardo e rendono più morbidi i movimenti: il docente le sfiora più tardi, nella sezione
del Montage. Quanto ai 3C, nell'industria la sigla indica la disciplina che cura come il
giocatore \emph{vede} (camera), \emph{comanda} (controlli) e \emph{abita} (personaggio)
il gioco: tre cose che si progettano insieme, perché ognuna cambia la percezione delle
altre.

% ---------------------------------------------------------------------
\section{I grafi di movimento, presi in prestito}
\label{ue9:sec:movimento}

In Play, la camera c'è, ma ``non ci possiamo muovere'': ``perché non abbiamo né input né
le funzioni di base di movimento''. ``Questo, ragazzi, non andiamo a farlo direttamente a
mano, ma andiamo a fare quella che è una bellissima'' scorciatoia (la trascrizione dice
``mandragata''): si apre il \textbf{BP\_ThirdPersonCharacter} classico, ``quello che
abbiamo usato'' finora, che contiene ``questi tre principali blocchi che Unreal fa di
base'': il controllo della \textbf{camera} (\emph{Look}), il controllo del
\textbf{movimento} (con \emph{Add Movement Input}) e il \textbf{salto} (\emph{Jump}).

\paragraph{Copiare anche l'Input Mapping.} ``Mi raccomando, prendete anche questo qua
sopra, che è l'\textbf{input mapping}'', cioè il grafo sull'Event Begin Play che aggiunge
l'Input Mapping Context (``andremo a spiegare che cosa si intende''). Tutto con
\textbf{Ctrl+C}, e poi \textbf{Ctrl+V} nell'Event Graph di BP\_Knight1: ``perché queste
sono le funzioni basiche di Unreal''.

\paragraph{IA\_Move.} Un'occhiata all'evento \textbf{IA\_Move} con l'Add Movement Input,
``proprio per fare una lezione sull'input'': ``questo evento dice che è legato a un
tasto, o a dei tasti, che ci riporta due valori, X e Y'', che ``di conseguenza diventano
dei vettori che vanno a creare il movimento sui due assi''. Uno è avanti e indietro
(\emph{forward}, \emph{backward}), l'altro è destra e sinistra.

\paragraph{La camera che ruota male.} Compilato, ``adesso la nostra camera si muove'', ma
in modo strano: si muove ``solo così'', su un asse. ``L'avevo visto un'altra volta'':
``devo usare il control rotation''. La spunta è \textbf{Use Pawn Control Rotation}:
``se lo attaccate, la camera vi va in \emph{free mode}'', ``nel senso che usa la
rotazione del controller'', sia in beccheggio sia in imbardata.

% ---------------------------------------------------------------------
\section{Tre spunte sulla rotazione}
\label{ue9:sec:rotazioni}

\paragraph{Use Pawn Control Rotation, sullo Spring Arm.} Uno studente ha la camera che
``quando vado in Play posso abbassarla e guardare solo il terreno, non vedo proprio
neanche più il personaggio''. Il docente pensa prima a una compenetrazione, poi si fa
condividere lo schermo: ``che camera hai messo? Camera normale, giusto?''. Guardando i
componenti, il problema salta fuori: ``non Use Pawn Control Rotation sulla camera: Use
Pawn Control Rotation sullo \textbf{Spring Arm}''. ``Quindi levalo da lì e mettilo sullo
Spring Arm.'' ``Anche l'altro?'' ``No, no, no: solo sullo Spring Arm.'' La spunta va su
\emph{un} componente solo, il braccio: è il braccio a ruotare intorno al personaggio, e
la camera, che ne è figlia, lo segue.

\paragraph{Orient Rotation to Movement.} In Play, il cavaliere cammina (alla velocità di
300), ma andando a destra ``lui cammina dritto, fondamentalmente'': si sposta di lato
continuando a guardare avanti. ``Questo lo possiamo risolvere con'' \textbf{Orient
Rotation to Movement}, sempre nel Character Movement: ``dal nome, orienta la rotazione al
movimento''. Accanto c'è la \textbf{Rotation Rate}, la velocità con cui il personaggio si
gira (``non è proprio precisissima, però ci piace'').

\paragraph{Use Controller Rotation Yaw.} Il personaggio però continua a non girarsi bene.
``Perché mi ricordo che qua c'era una particolarità che non ricordo'': il docente scorre
le categorie del Blueprint (\emph{Simulate}, \emph{Rotation}, \emph{Character},
\emph{Orientation}, \emph{Physics}, \emph{Controller}\dots) e la trova. ``Per fissare il
nostro character che non si girava, basta che eliminiamo il tick'': selezionato il Pawn
(la radice del Blueprint, \emph{BP\_Knight1 (Self)}), si cerca \emph{yaw} e si toglie
\textbf{Use Controller Rotation Yaw}, ``in modo tale che il character non usi il suo
yaw''. Senza questa spunta, il personaggio non ruota più insieme alla camera: ``andrà
dove siamo direzionati'', ``se io vado a destra, lui si gira completamente verso
destra''. A uno studente non basta: aveva tolto lo Yaw ma non attivato Orient Rotation to
Movement. ``Ok, non ce l'avevo.'' Con tutte e due, ``abbiamo il nostro bel character
fatto per benissimo''.

\paragraph{Approfondimento: chi comanda la rotazione.} Le tre spunte della lezione sono
le tre ``mani'' che possono ruotare un personaggio in terza persona, e vanno combinate
con criterio. \textbf{Use Pawn Control Rotation} (sullo Spring Arm) fa seguire alla
\emph{camera} la rotazione del controller, cioè del mouse o della levetta destra.
\textbf{Use Controller Rotation Pitch/Yaw/Roll} (sul Pawn) fa seguire quella stessa
rotazione al \emph{corpo} del personaggio. \textbf{Orient Rotation to Movement} (nel
Character Movement) fa girare il corpo verso la direzione in cui si muove. La
combinazione ``da avventura'' della prima parte della lezione è: camera libera (Use Pawn
Control Rotation sul braccio), corpo indipendente dalla camera (Use Controller Rotation
Yaw spento), corpo orientato al movimento (Orient Rotation to Movement acceso): è quella
del modello Third Person di Unreal. La combinazione ``da combattimento'' della seconda
parte, con il Blend Space 2D (sezione~\ref{ue9:sec:direction}), è opposta: il corpo
guarda sempre dove guarda la camera (Use Controller Rotation Yaw acceso), e le
animazioni di strafe mostrano lo spostamento laterale. Se sono accese insieme Use
Controller Rotation Yaw e Orient Rotation to Movement, le due ``mani'' si contendono il
personaggio; per questo, nel Character Movement, esiste anche \emph{Use Controller
Desired Rotation}, che fa ruotare il corpo verso la direzione della camera in modo
graduale, con la Rotation Rate (forse il ``desired'' che il docente prova a un certo
punto, senza esito chiaro nella trascrizione).

% ---------------------------------------------------------------------
\section{Il retargeting: il salto del manichino sul cavaliere}
\label{ue9:sec:retarget}

``Adesso, però, guardate una cosa: se ad esempio provo a saltare, non ho nessuna
animazione.'' Il salto funziona (lo fa il Character Movement), ma il cavaliere sale e
scende nella posa della camminata. Le animazioni di salto si potrebbero prendere da
Mixamo, ``anche se potremmo, volendo''; invece il docente usa ``l'opzione di Unreal che
si chiama \textbf{retargeting}''. ``Diciamo che il retargeting, sappiatelo, può uscire
bene, ma può uscire anche abbastanza male: è dipendente dall'animazione. Però andiamo a
farlo comunque, per vedere come si fa.''

\paragraph{I passaggi.}
\begin{enumerate}
  \item Nel Content Browser si cerca \emph{Jump}: si trova l'animazione del salto del
        manichino base di Unreal, \textbf{MM\_Jump}.
  \item Tasto destro sull'animazione $\rightarrow$ \textbf{Retarget Animations}.
  \item La finestra mostra lo scheletro \textbf{sorgente} (\emph{source}), quello da cui
        viene l'animazione: il manichino di Unreal, \textbf{Manny} o \textbf{Quinn}
        (``Many Quincy''). Come \textbf{target} si sceglie lo scheletro del cavaliere, il
        paladino di Mixamo: ``eccolo qua, li mette a confronto''.
  \item Si seleziona \textbf{MM\_Jump} nella lista delle animazioni: nella preview il
        paladino salta. ``Fortunatamente Mixamo usa degli scheletri che sono abbastanza
        compatibili con quello di Unreal'', e il risultato ``è venuto abbastanza bene'';
        ``potete girarci tranquillamente intorno con la viewport''.
  \item Due pulsanti di esportazione. \textbf{Export Retargeted Assets} ``non ci serve,
        ragazzi'': ``vi crea un asset retargettato'', una copia completa degli asset sullo
        scheletro di destinazione. Basta \textbf{Export Animations}: ``esportare
        l'animazione così, come nuova''. ``Ricordate che c'è il source, e vi esporta il
        target.''
  \item Come cartella si sceglie quella del personaggio (\emph{Characters $\rightarrow$
        Knight1 $\rightarrow$ Animation}) e si esporta.
\end{enumerate}
Lo stesso procedimento si ripete per la caduta: \textbf{MM\_Fall\_Loop}, un'animazione
``in loop'' di caduta, retargettata sul paladino ed esportata nella stessa cartella.

\paragraph{Il root motion ``rotto''.} Il risultato del salto ``non è fantastico'':
``il problema di Mixamo è che ha il \textbf{root motion} un po' rotto''. Il docente
mostra il \emph{root}, ``un punto dello scheletro'', l'osso alla base della gerarchia, e
osserva che lo spostamento non è quello atteso; ``potremmo vederlo'', ``ma per vedere come
funziona il retarget va benissimo''. Si potrebbe sistemare in seguito.

\paragraph{Approfondimento: il retargeting in Unreal Engine 5.} Il retargeting trasferisce
un'animazione da uno scheletro a un altro con proporzioni e nomi delle ossa diversi. In
Unreal Engine 5 si basa su due tipi di asset: l'\textbf{IK Rig}, che descrive uno
scheletro raggruppando le ossa in \emph{catene} (spina dorsale, braccia, gambe, testa)
con un osso di \emph{retarget root} (di solito il bacino), e l'\textbf{IK Retargeter},
che mette in corrispondenza le catene di due IK Rig e corregge la differenza fra le pose
di riferimento (la A-pose del manichino e la T-pose di Mixamo). Nelle versioni più recenti il
comando \emph{Retarget Animations} del menu contestuale apre una finestra che crea al
volo questi asset quando gli scheletri sono riconosciuti (come quelli di Mixamo e del
manichino), ed è ciò che si vede a lezione; per risultati migliori, o con scheletri
insoliti, si creano e si regolano a mano. Il problema del root di Mixamo è noto: gli
scheletri di Mixamo non hanno un osso di root a terra, e lo spostamento del personaggio
è registrato sull'osso dei fianchi (\emph{Hips}); un'animazione retargettata può quindi
avere il movimento nel posto ``sbagliato'' per il sistema di root motion di Unreal.

% ---------------------------------------------------------------------
\section{Le regole di transizione}
\label{ue9:sec:transizioni}

Con le due animazioni nuove si aggiorna l'Animation Blueprint: ``dentro il nostro Jump
ci mettiamo l'animazione di jump'' (``niente di particolare, ragazzi''), dentro Fall
l'animazione di caduta. Si ricollegano le frecce. ``Se cammina e corre, può saltare? Sì,
certo. Se salta, dovrà per forza almeno loopare in caduta'', e ``dal Fall dovrà ritornare
al nostro WalkRun''.

\paragraph{Prima versione.} Nelle icone delle transizioni c'è un grafo con un nodo
\textbf{Result}, che ``ci dà un bool'': ``ci dice quando è che posso entrare in questa
animazione''.
\begin{itemize}
  \item \textbf{WalkRun $\rightarrow$ Jump}: un \emph{Get IsFalling} collegato al
        Result. ``Potremmo farlo più preciso dopo''; per esempio, ``possiamo anche dire:
        ha iniziato il jump''.
  \item \textbf{Jump $\rightarrow$ Fall}: qui il grafo vuoto dà errore, ``perché noi
        vogliamo far sì che questo processo tra il Jump e il Fall sia automatico''.
        Selezionata l'icona della transizione, nei Details si spunta \textbf{Automatic
        Rule Based on Sequence Player in State}: ``quando inizia quell'animazione, a
        seguito ci sarà quest'altra''; in altre parole, la transizione scatta da sola
        quando l'animazione dello stato Jump finisce. ``Lo mettiamo così e non ci darà
        problemi.''
  \item \textbf{Fall $\rightarrow$ WalkRun}: ``vuol dire che non sta cadendo'': Get
        IsFalling, poi un \textbf{NOT Boolean}. ``Ricordo che il NOT Boolean è qualcosa
        che ci dà il suo contrario'': deve essere false perché la caduta finisca.
\end{itemize}
Uno studente chiede cosa ci sia fra Jump e Fall: ``se tu clicchi su questo pallino, ci
dice, in Details\dots\ vogliamo che questa animazione parta direttamente se effettivamente
il jump è finito'': \emph{Automatic Rule Based on Sequence Player in State}, ``questo
tick, lo metti vero''.

\paragraph{La prova, e un difetto.} Si prova da un punto alto, ``la nostra piattaforma''.
``Non si nota tantissimo, però effettivamente sta facendo giustamente l'azione di
saltare.'' Ma il docente vuole collegare WalkRun anche direttamente a Fall: il
personaggio ``può cadere anche nel caso in cui lui non abbia spinta, semplicemente cade da
qualche parte'', per esempio scendendo da un bordo. Con le sole regole su IsFalling,
però, le due transizioni in uscita da WalkRun avrebbero la stessa condizione, e Unreal
non saprebbe distinguere un salto da una caduta.

\paragraph{Seconda versione: la velocità sulla Z.} ``Dovremmo andare a differenziare un
altro valore: la Z. Se la Z è maggiore di 0, vuol dire che sta saltando; se la Z invece è
minore di 0, vuol dire che sta cadendo.'' Nell'Event Graph dell'ABP, finora, della
velocity si prendeva solo la lunghezza. Il docente prova prima la strada sbagliata (la
\emph{Jump Z Velocity} del Character Movement, che però è la spinta \emph{iniziale} del
salto, un valore fisso), poi prende la \textbf{Velocity} già salvata, ne fa lo
\textbf{Split Struct Pin} (``lo splittiamo'') e dalla \textbf{Z} fa \textbf{Promote to
Variable}: una nuova float, \textbf{VelocityZ} (``Z va benissimo''). Nelle transizioni:
\begin{itemize}
  \item \textbf{WalkRun $\rightarrow$ Jump}: IsFalling \textbf{AND} VelocityZ maggiore
        di zero, con un nodo \textbf{AND Boolean}. Il docente dice anche ``quando la
        velocità sarà maggiore di 1, più preciso'': una soglia piccola invece di zero
        esatto.
  \item \textbf{WalkRun $\rightarrow$ Fall}: IsFalling e VelocityZ minore di zero.
\end{itemize}
Per tenere in ordine i fili, si possono usare i \emph{reroute node} (``volete collegarli
con questo? È un po' bruttino''). Il risultato: ``se io salto, fa l'animazione''; ``se io
adesso invece cado solamente, lui va diretto in Fall''. Il docente chiude salvando anche
il BP\_ThirdPersonCharacter che l'editor segnala come modificato (``va bene, grazie'').

\begin{table}[htbp]
\centering
\small
\begin{tabular}{@{}p{3.2cm}p{5.3cm}p{5.5cm}@{}}
\toprule
\textbf{Transizione} & \textbf{Regola} & \textbf{Significato} \\
\midrule
Entry $\rightarrow$ WalkRun & nessuna & lo stato iniziale \\
WalkRun $\rightarrow$ Jump & IsFalling AND VelocityZ $>$ 0 & è in aria e sta salendo: ha
saltato \\
Jump $\rightarrow$ Fall & Automatic Rule Based on Sequence Player in State & finita
l'animazione di salto, passa alla caduta \\
WalkRun $\rightarrow$ Fall & IsFalling AND VelocityZ $<$ 0 & è in aria e sta scendendo:
è caduto da un bordo \\
Fall $\rightarrow$ WalkRun & NOT IsFalling & è tornato a terra \\
\bottomrule
\end{tabular}
\caption{Le regole di transizione della State Machine Locomotion (versione finale).}
\label{ue9:tab:transizioni}
\end{table}

\paragraph{Approfondimento: altre regole basate sul tempo.} L'opzione automatica della
lezione è il modo più rapido per dire ``quando l'animazione finisce''. Romero e Sewell
(capitolo 17), per gli stati \emph{StandToProne} e \emph{ProneToStand}, usano invece una
regola esplicita: il nodo \textbf{Time Remaining (ratio)} dell'animazione dello stato,
seguito da un \emph{Less} con 0,1: la transizione scatta quando manca meno del 10\%
dell'animazione. Il vantaggio è il controllo: la transizione può partire un po' prima
della fine e lasciare spazio al \emph{blend} fra i due stati, che ogni transizione ha
(nei Details: \emph{Blend Settings}, con una durata di default di 0,2 secondi).
Ricordano anche che le animazioni ``di una volta'', come il salto, vanno impostate senza
\emph{Loop Animation}: è il caso dello stato Jump, mentre Fall, che deve durare quanto
la caduta, resta in loop.

% ---------------------------------------------------------------------
\section{Lo sprint con l'Enhanced Input}
\label{ue9:sec:enhanced}

``Adesso andremo a mettere il nostro sprintino'': ``un semplicissimo sistemino che passa
la Max Walk Speed da un tot a un altro tot''. Si torna in BP\_Knight1, ``a ragionare a
livello di Blueprint''. ``Ma sì, facciamo anche l'input mapping.''

\paragraph{L'Input Mapping Context.} ``Unreal ha introdotto, quasi dalla 5.2'', quello
che è l'\textbf{Input Mapping Context}, ``una cosa molto interessante di Unreal che ci
permette, ad esempio, di avere degli eventi legati a più input''. Finora, in molte
lezioni, si sono usati gli eventi dei tasti (``abbiamo visto il tasto E, eccolo qua''):
``semplicemente un rilevamento, da parte di Unreal, della tastiera: se è premuto o non è
premuto''. Adesso, invece, ``gli eventi sono legati a più controlli possibili''.

Nel Content Browser si cerca \textbf{IMC} (\emph{Input Mapping Context}). ``È un
\emph{contesto}: un contesto di input che vengono dati in un certo momento, e sappiate
che sono sovrascrivibili.'' L'esempio: ``io devo cambiare dal gioco di avventura 3D al
gioco di guida della macchina''; ``i controlli sono completamente diversi, quindi cambio
contesto'': un Input Mapping Context per la macchina e uno ``per il character che cammina
normalmente''. Aperto \textbf{IMC\_Default}, ci sono i \textbf{Mappings}:
\textbf{IA\_Jump}, \textbf{IA\_Move}, \textbf{IA\_Look}.

\paragraph{L'Input Action IA\_Run.} ``Per far sì di avere un'Input Action, dobbiamo creare
un asset chiamato \textbf{Input Action}.'' Cercando \emph{IA} nel Content Browser ``troverete
un sacco di asset, ma a voi vi interessa questo celestino''. Con il tasto destro,
\textbf{Show in Folder View}, ``in modo tale che voi troviate direttamente la cartella''.
Si duplica IA\_Jump (Ctrl+C, Ctrl+V, o \emph{Duplicate}) e lo si rinomina
\textbf{IA\_Run}.

\paragraph{Dentro un'Input Action.} Aperto IA\_Run, ``ci sono un po' di valori che
possiamo vedere'':
\begin{itemize}
  \item \textbf{Description}: ``semplice descrizione, quindi quello che fa'';
  \item \textbf{Trigger when Paused}: se l'azione deve scattare anche quando il gioco è in
        pausa, ``quindi si blocca il tempo'';
  \item \textbf{Value Type}: ``la tipologia di valore che ci ritorna indietro. Un input ci
        dà un valore'': \textbf{Digital (bool)}, \textbf{Axis1D (float)} (``come abbiamo
        visto, avanti e indietro''), \textbf{Axis2D (Vector2D)} (``più assi, due
        dimensioni'') o \textbf{Axis3D (Vector)} (``tre tipologie di input che ci vengono
        dati''). ``A noi interessa il Run, che è un bool, quindi premuto o non premuto'':
        Digital (bool) ``va benissimo'';
  \item \textbf{Triggers}: ``le azioni che sono già preinstallate in Unreal'', che dicono
        ``ehi, guarda, si triggera se è premuto, o se è rilasciato, o se è sempre
        premuto''. L'elenco letto dal docente: \emph{Chorded Action}, \emph{Combo},
        \emph{Down}, \emph{Hold}, \emph{Hold and Release}, \emph{Pressed}, \emph{Pulse},
        \emph{Released}, \emph{Tap} (la trascrizione dice anche ``Short Deduction'' e
        ``repeat'', non identificabili con certezza). ``Fondamentalmente, questo vuol
        dire che quando si compie questa azione, questo input viene triggerato.'' IA\_Run
        eredita da IA\_Jump la sua configurazione, ``e va benissimo così, perché tanto è
        un'azione molto simile al jump''.
\end{itemize}

\begin{table}[htbp]
\centering
\small
\begin{tabular}{@{}p{3.2cm}p{10.8cm}@{}}
\toprule
\textbf{Voce} & \textbf{Significato} \\
\midrule
Digital (bool) & premuto / non premuto: un pulsante (salto, corsa, interazione) \\
Axis1D (float) & un valore su un asse: un grilletto analogico, avanti/indietro \\
Axis2D (Vector2D) & due assi: il movimento (X, Y) o la visuale del mouse \\
Axis3D (Vector) & tre assi: input 3D (per esempio un controller di movimento) \\
\midrule
Down & scatta a ogni frame finché il tasto è giù \\
Pressed & scatta una volta, quando il tasto viene premuto \\
Released & scatta una volta, quando il tasto viene rilasciato \\
Hold & scatta dopo che il tasto è stato tenuto premuto per un certo tempo \\
Hold and Release & scatta al rilascio, se il tasto è stato tenuto abbastanza \\
Tap & scatta se il tasto viene premuto e rilasciato entro un tempo breve \\
Pulse & scatta a intervalli regolari mentre il tasto è premuto \\
Chorded Action & scatta solo se un'altra azione è attiva insieme (una ``combinazione'') \\
Combo & scatta se una sequenza di azioni viene eseguita nell'ordine \\
\bottomrule
\end{tabular}
\caption{I tipi di valore (Value Type) e i trigger di un'Input Action, letti a lezione
dall'elenco dell'editor. Le descrizioni sono un'integrazione: a lezione il docente
legge i nomi senza spiegarli uno per uno.}
\label{ue9:tab:inputaction}
\end{table}

\paragraph{Il contesto va caricato.} Tornati in BP\_Knight1: ``abbiamo già il contesto
caricato'', con il grafo copiato dal ThirdPersonCharacter. ``Sappiate che se voi non
caricate il contesto, non vi prenderà nessuna cosa'': sull'Event Begin Play, dal Player
Controller si prende l'\emph{Enhanced Input Local Player Subsystem} e si chiama
\textbf{Add Mapping Context} con IMC\_Default (figura~\ref{ue9:fig:enhanced}).

\paragraph{L'evento IA\_Run e i suoi pin.} Nel grafo si invoca \textbf{IA\_Run}; ``se lo
estendete'' (la freccina in basso del nodo), compaiono tutti i pin:
\begin{itemize}
  \item i pin di esecuzione \textbf{Triggered} (``quando è triggerato''), \textbf{Started}
        (``quando inizia a essere triggerato''), \textbf{Ongoing}, \textbf{Canceled},
        \textbf{Completed};
  \item i valori: l'\textbf{Action Value} (``se è true o false, in questo caso''),
        \textbf{Elapsed Seconds} e \textbf{Triggered Seconds} (``quanto tempo è stato
        triggerato''), e l'\textbf{Input Action}, ``la referenza stessa'' all'asset.
\end{itemize}

\paragraph{Set Max Walk Speed.} ``Dobbiamo modificare il Character Movement, ragazzi, che
dovrà definire quanta velocità avremo'': si trascina nel grafo il componente
\textbf{Character Movement} e se ne tira fuori \textbf{Set Max Walk Speed}, due volte.
Sul \textbf{Triggered}, \textbf{600}; sul \textbf{Completed}, \textbf{300}
(figura~\ref{ue9:fig:sprint}).

\paragraph{Il mapping dimenticato.} In Play non succede niente. ``Scusate, ragazzi, mi
sono dimenticato di dare l'input'': IA\_Run esiste, ma non è nel contesto. Si apre
IMC\_Default, si aggiunge un mapping (il \emph{+} accanto a \emph{Mappings}), si sceglie
IA\_Run e gli si assegna il tasto: ``un nostro bello \textbf{Shift}''. ``L'Input Mapping
Context ci permette di mettere quali tipologie di pulsanti facciano'' ciascuna azione.

\paragraph{Lo stesso input, anche con il controller.} ``Io ad esempio qua, ragazzi, ho un
controller'': il docente lo collega (``Unreal riesce a prendere il controller''), aggiunge
un secondo tasto alla stessa azione, clicca sull'icona della tastiera accanto al campo e
preme un pulsante del controller (``mi sono fregato la sorpresa''): ``facciamo tipo Dark
Souls? Premo cerchio'', e l'editor scrive \textbf{Gamepad Face Button Right}. ``Questa
azione, cioè l'azione di correre, è collegata a più input.'' ``Quando faccio progetti in
generale, devo fare anche il controller'': si corre con Shift sulla tastiera ``oppure con
il nostro bel controllino, con quello che vogliamo noi'', ``come è stato designato il
gioco''. Tutto il resto del movimento sul controller ``non ci interessa adesso''.

\paragraph{La prova.} ``Se noi proviamo\dots\ rilascio.'' Con Shift il cavaliere corre, e
il Blend Space passa alla corsa; al rilascio torna a camminare. ``Ecco a voi che abbiamo
creato un nostro bel character nuovo, in cui semplicemente possiamo correre, possiamo
modificare la velocità.''

\begin{figure}[htbp]
\centering
\begin{tikzpicture}[font=\footnotesize,
  ue9box/.style={draw, thick, rounded corners=4pt, align=center, fill=gray!8,
                 minimum height=1.0cm, text width=2.5cm},
  ue9arr/.style={-{Stealth[length=2.5mm]}, thick}]
  \node[ue9box, fill=yellow!15] (k1) at (0,0.7) {Left Shift\\{\scriptsize tastiera}};
  \node[ue9box, fill=yellow!15] (k2) at (0,-0.7) {Gamepad Face\\Button Right};
  \node[ue9box, fill=blue!10, text width=2.8cm] (imc) at (3.8,0) {\textbf{IMC\_Default}\\
    {\scriptsize Input Mapping Context}\\{\scriptsize (mapping IA\_Run)}};
  \node[ue9box, fill=cyan!12] (ia) at (7.6,0) {\textbf{IA\_Run}\\{\scriptsize Input
    Action, Digital (bool)}};
  \node[ue9box, fill=red!12, text width=2.8cm] (ev) at (11.4,0) {evento
    \textbf{IA\_Run}\\{\scriptsize nel Blueprint}\\{\scriptsize Triggered / Completed}};
  \draw[ue9arr] (k1) -- (imc); \draw[ue9arr] (k2) -- (imc);
  \draw[ue9arr] (imc) -- (ia); \draw[ue9arr] (ia) -- (ev);
  \node[ue9box, fill=green!10, text width=4.2cm] (bp) at (3.8,-2.4) {\textbf{Event Begin
    Play} $\rightarrow$ Enhanced Input Local Player Subsystem $\rightarrow$ \textbf{Add
    Mapping Context} (IMC\_Default)};
  \draw[ue9arr, dashed] (bp) -- node[right, font=\scriptsize]{attiva} (imc);
\end{tikzpicture}
\caption{La catena dell'Enhanced Input. I tasti fisici non sono legati direttamente
all'evento: il contesto (IMC) associa più tasti alla stessa azione (IA), e il Blueprint
reagisce all'azione. Senza Add Mapping Context, il contesto non è attivo e nessun input
arriva. Schema ricostruito.}
\label{ue9:fig:enhanced}
\end{figure}

\begin{figure}[htbp]
\centering
\begin{tikzpicture}[font=\footnotesize,
  ue9node/.style={draw, thick, rounded corners=4pt, align=center, fill=gray!8,
                  minimum height=0.8cm, text width=2.4cm},
  ue9ev/.style={ue9node, fill=red!12},
  ue9var/.style={ue9node, fill=blue!10, rounded corners=8pt},
  ue9exec/.style={-{Stealth[length=2.5mm]}, very thick, black!80},
  ue9data/.style={-{Stealth[length=2mm]}, thick, blue!60!black}]
  \node[ue9ev, text width=2.2cm, minimum height=1.8cm] (ev) at (0,0) {\textbf{IA\_Run}\\
    {\scriptsize Enhanced Input Action}};
  \node[ue9node, text width=3.2cm] (s1) at (5,0.8) {\textbf{Set Max Walk Speed}\\
    {\scriptsize Max Walk Speed = 600}};
  \node[ue9node, text width=3.2cm] (s2) at (5,-0.8) {\textbf{Set Max Walk Speed}\\
    {\scriptsize Max Walk Speed = 300}};
  \node[ue9var, text width=2.6cm] (cm) at (9.6,0) {\textbf{Character}\\
    \textbf{Movement}};
  \draw[ue9exec] (ev.east |- s1) -- node[above, font=\scriptsize]{Triggered} (s1);
  \draw[ue9exec] (ev.east |- s2) -- node[below, font=\scriptsize]{Completed} (s2);
  \draw[ue9data] (cm.west) -- ++(-0.5,0) |- node[pos=0.25, right, font=\scriptsize]
    {Target} (s1.east);
  \draw[ue9data] (cm.west) -- ++(-0.5,0) |- (s2.east);
\end{tikzpicture}
\caption{Lo sprint di BP\_Knight1. Premendo Shift (o il cerchio del controller) la
velocità massima sale a 600 e il Blend Space mostra la corsa; al rilascio torna a 300,
la camminata. Schema ricostruito.}
\label{ue9:fig:sprint}
\end{figure}

\paragraph{Approfondimento: Action Mapping ed Enhanced Input.} Il libro di Romero e
Sewell, scritto per le prime versioni di Unreal Engine 5, costruisce lo stesso sprint
con il vecchio sistema: un \emph{Action Mapping} chiamato \emph{Sprint} nei Project
Settings (sezione \emph{Input}), legato a \emph{Left Shift}, e un evento
\emph{InputAction Sprint} con i pin \emph{Pressed} e \emph{Released}, su cui due
\emph{Set Max Walk Speed} portano la velocità a 2200 e poi di nuovo ai 600 di default
(capitolo 6). Il principio è identico; cambia il sistema. L'\textbf{Enhanced Input},
introdotto come plugin in Unreal Engine 5.0 e diventato il sistema predefinito dei
modelli dalla 5.1 (il ``quasi dalla 5.2'' del docente è un ricordo approssimativo), separa
tre livelli: l'\emph{azione} (Input Action, che cosa vuole fare il giocatore), il
\emph{contesto} (Input Mapping Context, quali tasti attivano quali azioni in una certa
situazione) e il \emph{comportamento} (Triggers e \emph{Modifiers}, come per esempio la
\emph{zona morta} di una levetta o l'inversione di un asse). I contesti si aggiungono e
si tolgono a runtime con \emph{Add Mapping Context} e \emph{Remove Mapping Context}, e
hanno una \emph{Priority}: è così che si passa ``dal gioco a piedi alla macchina''
dell'esempio del docente, e così che si costruisce una schermata di rimappatura dei
tasti. Gli autori raccomandano anche di raggruppare i nodi dello sprint in un
\emph{commento} (\emph{Create Comment from Selection}), per tenere ordinato il grafo.

% ---------------------------------------------------------------------
\section{Un tentativo: FInterp To}
\label{ue9:sec:finterp}

``Se volessimo farlo un po' più carino, ragazzi'', il passaggio fra 300 e 600 potrebbe
essere graduale invece che istantaneo. Il blocco è \textbf{FInterp To} (``finter2''):
``fa un'interpolazione di un float''. ``Avete presente quando abbiamo parlato di
interpolazione?'' ``Fa il passaggio tra due numeri. Lui dice: devi passare da questo
valore, cioè \textbf{Current}, a questo valore, \textbf{Target}, in questo tempo'',
\textbf{Delta Time}, per cui il docente consiglia \textbf{Get World Delta Seconds}; e
poi c'è l'\textbf{Interp Speed}, la velocità dell'interpolazione.

\paragraph{Il grafo.} Sul Triggered: \textbf{Current} = Character Movement $\rightarrow$
\textbf{Get Max Walk Speed}; \textbf{Target} = 600; \textbf{Delta Time} = World Delta
Seconds; il risultato nel Set Max Walk Speed. Il docente ricorda che il Delta Seconds si
può moltiplicare per un valore (``ricordiamo che possiamo modificarlo moltiplicandolo
per un valore''), e prova con 10, poi con un decimo. Sul Completed lo stesso, ``ma
semplicemente invertiamo'': il Current parte da 600, il Target è 300. ``Ricordiamoci di
collegare anche al Character Movement.''

\paragraph{Non funziona, e perché.} La corsa resta ``ancora molto veloce'', il passaggio
istantaneo. Il docente prova valori diversi (10, 100: ``è alla norma del pianeta stare
continuamente a vedere se i valori sono giusti'') e poi mette dei \textbf{Print String}:
``ci dà seicento''. La diagnosi: ``fondamentalmente questa è un'azione a triggered
unico, quindi non farà l'interpolazione''. FInterp To fa \emph{un passo}
dell'interpolazione ogni volta che viene eseguito: se viene eseguito una volta sola, fa
un solo passo. ``Lo deve fare tramite un \textbf{timer}'': si rimanda a dopo (``volendo,
ragazzi, si può fare con un timer, però facciamolo magari dopo''), e non verrà ripreso
nella lezione.

\paragraph{Approfondimento: dove mettere un'interpolazione.} Un nodo come FInterp To
(o VInterp To, RInterp To per vettori e rotazioni) va eseguito \emph{a ogni frame}
finché il valore non arriva al target: il suo risultato è il valore ``un po' più vicino''
al target, e il Delta Time serve a renderlo indipendente dal frame rate. Ci sono tre
modi tipici di farlo girare: nel \textbf{Tick}, con una variabile \emph{TargetSpeed}
che i due pin dell'input si limitano a impostare (600 o 300); con un \textbf{timer}
(\emph{Set Timer by Event}, con \emph{Looping}, fermato con \emph{Clear and Invalidate
Timer} quando il valore è arrivato), come suggerisce il docente; o con una
\textbf{Timeline}, vista nella sesta lezione, che va da 300 a 600 con \emph{Play} e
torna indietro con \emph{Reverse}. Con l'Enhanced Input c'è anche una quarta via: il pin
\emph{Triggered} scatta a ogni frame se il trigger è \emph{Down} (o se non c'è nessun
trigger), e una sola volta se il trigger è \emph{Pressed}. IA\_Run, copiata da IA\_Jump,
ha probabilmente un trigger che la fa scattare una volta sola: è ``l'azione a triggered
unico'' del docente.

% ---------------------------------------------------------------------
\section{Prima della pausa: le animazioni da scaricare}
\label{ue9:sec:download}

``Adesso andremo a vedere, se non mi esplode prima Unreal, anche come si fanno i
\textbf{montaggi}''; mancano una ventina di minuti alla pausa. ``Dovremmo andare un attimo
su Mixamo'' (``a quanto pare oggi Mixamo ci mette pure un bel po' a caricare''), per
scaricare due gruppi di animazioni.

\paragraph{Un attacco in root motion.} ``Vorrei che sceglieste un'animazione che sia in
\textbf{root motion}'': ``riprendiamo ad esempio una situazione simile a Dark Souls''.
``Ricordatevi sempre di prendere il vostro character'' su Mixamo (il paladino, nel caso
del docente), perché le animazioni si adattano al personaggio selezionato. Il docente
cerca \emph{attack} e sceglie un attacco con la spada che ``di base, come potete vedere,
non può essere messo in place''. Le impostazioni di download sono quelle della volta
precedente: \textbf{Without Skin}, \textbf{FBX Binary}, \textbf{30} fotogrammi al
secondo. ``Ripeto: non deve essere in place, deve avere il root motion, quindi il
movimento ereditato.'' ``Un'azione unica'': ``ad esempio un attacco con una spada, un
qualche cosa'', ``un'animazione che vi piace di più''.

\paragraph{Le animazioni del Blend Space 2D.} Dopo la pausa si farà il \textbf{Blend
Space 2D}, ``che ci servirà per fare anche il calcolo della direzione del nostro
character''. Servono quindi altre animazioni: ``walk destra, sinistra, dritto e
indietro''. ``Consiglio spassionato'': su Mixamo si chiamano \textbf{Strafe} (``provate a
cercare Strafe''; ``Jog Strafe va benissimo anche''), e queste, al contrario
dell'attacco, ``dovete metterle \textbf{in place}''.

\paragraph{Il mirror.} ``Piccolo consiglio: se volete far sì che siano abbastanza simili
tra di loro, potete fare il \textbf{mirror}.'' Si scarica l'animazione verso sinistra
non specchiata; poi, con l'opzione \emph{Mirror} di Mixamo, ``in automatico andrà a
destra''. Le due animazioni saranno così perfettamente simmetriche.

\paragraph{Quante animazioni.} Su richiesta (``le puoi scrivere per messaggio?''), il
docente le elenca anche in chat: ``in totale sono otto animazioni''. Avanti, indietro,
destra, sinistra, e le quattro \textbf{diagonali}: avanti a destra, avanti a sinistra,
indietro a destra, indietro a sinistra. Uno studente chiede se l'avanti serva, visto che
c'è già la camminata: ``nel senso, possiamo usare tranquillamente quello di avanti; ma in
questo caso abbiamo solamente avanti, dobbiamo trovare anche indietro''. E un'altra: le
diagonali ``soltanto in avanti, o anche indietro?'' ``Anche indietro'': ``se io vado tipo
così, indietro verso destra, e così, indietro verso sinistra''. ``Praticamente è la X,
letteralmente'': le due diagonali incrociate, oltre alla croce degli assi. Il docente
nomina anche alcuni pacchetti di Mixamo con animazioni di locomozione (``Locomotion
Pack''; il resto dei nomi è incomprensibile nella trascrizione). ``Ci rivediamo alle
18.''

% ---------------------------------------------------------------------
\section{La pausa e i feedback}
\label{ue9:sec:feedback}

Al rientro, ``siamo rimasti in cinque'', poi in quattro. Il docente apre uno spazio per
i commenti: ``se avete dei feedback riguardo a questo corso, sparate tranquillamente'',
``siamo anche molto coetanei, anzi credo io di essere anche forse più piccolo di alcuni
di voi''. ``Se avete cose da chiedere, magari di andare più piano, andare più veloce,
oppure un po' come sono strutturate le lezioni''. ``Capisco pure che può essere
effettivamente un po' pesante in alcuni momenti'', ``perché purtroppo i Blueprint, o
Unreal in generale, è divertente quando poi le sai, le cose''.

\paragraph{Le funzioni da imparare a memoria.} Uno studente: ``immagino che sono
tantissime funzioni, quindi sono tutte da imparare a memoria''. ``Quello, eh sì.''

\paragraph{Il progetto del docente: la pallina morbida e dura.} Per mostrare dove porta
la pratica, il docente racconta un suo progetto personale: ``una specie di
\emph{platformer} in cui una pallina diventa morbida e dura, nel senso che diventa di
metallo o di un materiale più morbido''. ``La cosa è che ho simulato tutta la fisica'':
``la fisica non è simulata da Unreal, ma sono io che faccio finta, tramite i Blueprint, di
far fare delle funzioni fisiche''. ``Hai messo proprio te le formule fisiche?'' ``Più o
meno.'' Il meccanismo: si rileva lo \textbf{slope}, ``il quanto è piegato il piano'',
perché ``la meccanica è il fatto che se sei più pesante con la pallina vai più veloce su
uno slope, come nella realtà'': ``una pallina di metallo e una pallina di gomma: quella di
metallo va più veloce, perché è più pesante, comunque più influenzata''. ``È tutto
simulato'': ``sono io che, tramite la detection dello slope, cambio i valori''.

\paragraph{Combinare i blocchi.} ``E quello è ovviamente tutto custom'': ``quando hai
tutte le nozioni in testa, ti dici: ok, so questo blocco A, B, C, D, E, F, G; quale mi
serve? Il blocco C e il blocco A. Li metto insieme e faccio una nuova combinazione.''
``È sempre il fatto di capire come posso interlacciare le cose.'' ``Qua stiamo facendo
cose magari un po' generiche''; poi ``dici: ok, voglio fare simile a quello, ma
leggermente diverso. Voglio fare il character, però con la spada: che mi cambia?''

\paragraph{Inventari e UI.} La richiesta di Federico: ``a me interesserebbero la gestione
degli inventari e le interfacce''. ``Sfortunatamente no'': ``sappi che l'inventario e la
UI sono un mondo a parte''. ``Esistono proprio figure riguardanti questo mondo della UI,
sia designer sia programmatori.'' Si farà ``una cosa come un menu'', ``una cosetta
tranquilla'', ``però per il resto non andremo a fare di più''.

\paragraph{Le lezioni con Davide.} ``Sappiate che il 25 è la mia ultima lezione, e avrete
tutte quelle con \textbf{Davide}, che invece vi farà la parte più di \emph{level}.'' Con
lui si useranno ``tutti gli asset che abbiamo creato: le piattaforme, i pulsanti'', per
lavorare su concetti come la \textbf{gym}, ``il concetto della palestra'', che ``si fa
praticamente in via di sviluppo'': ``è il parco giochi'', ``letteralmente il parco giochi
dove ho tutte le meccaniche, le posso provare tutte quante, vedo se funzionano''. Poi si
costruiscono situazioni: ``voglio fare un livello horror, la bestia che ti insegue: tu
entri in una stanza e devi fuggire dalla bestia''; ``situazione X, situazione Y,
situazione Z'', ``le costruisci e metti insieme i componenti'', ``e vedi quei componenti
che hai programmato come funzionano a livello anche di esperienza, di gioco''. Gli
elementi saranno ``quelli di un classico platform: le piattaforme, i bottoni'', e si
vedranno anche i tool di \textbf{modeling} dell'editor e ``un po' di teoria di level
design'': ``chi fa game design, e a chi interessa tantissimo anche la parte di level:
molto, molto interessante''.

\paragraph{Dove lasciare i feedback.} ``Ripeto, ragazzi, anche a chi sta in differita'':
qualsiasi feedback sulle lezioni va lasciato sul canale \textbf{general} della chat,
``non su introduzione al materiale, perché sennò si fa casino''. Non solo sugli
argomenti, ``ma feedback anche proprio per le lezioni in generale: andare più lento,
andare più veloce, vedere cose più pratiche''.

\paragraph{Approfondimento: la gym.} Nel level design la \textbf{gym} (o \emph{gym
level}, \emph{gray box gym}) è un livello di servizio che non finirà nel gioco: una
stanza in cui ogni meccanica è disposta in isolamento e misurata. Si mettono salti di
lunghezze crescenti per sapere qual è il salto massimo del personaggio, rampe di
pendenze diverse per verificare la \emph{Max Walkable Floor Angle}, gradini di varie
altezze per la \emph{Max Step Height}, corridoi di larghezze diverse per la camera. Le
misure raccolte nella gym diventano le ``regole metriche'' del level design (quanto può
essere largo un burrone, quanto alto un muretto da scavalcare), e i livelli veri si
costruiscono rispettandole. È il ``parco giochi'' del docente; e il problema della
camera che compenetra, visto in questa lezione, è uno dei primi che una gym fa emergere.
% ---------------------------------------------------------------------
\section{La direzione: Calculate Direction}
\label{ue9:sec:direction}

``Abbiamo detto prima che oggi dobbiamo fare il Blend Space 2D.'' Il docente ha importato
otto animazioni, ``cercando di prenderle tutte similari'': \emph{Jog Backward},
\emph{Jog Backward Diagonal}, \emph{Jog Forward}, \emph{Jog Forward Diagonal} e così
via, ``a parte lo Strafe Left, che devo eliminare perché non andava bene come
animazione''. Il Blend Space 1D prende solo la velocità: ``vai più veloce, l'animazione,
benissimo''. Il 2D aggiunge un secondo asse, la \textbf{direzione}, ``in modo tale da non
dover magari direzionare sempre il character'', ma vedere ``che si muove in diagonale, si
muove a sinistra''.

\paragraph{Dove si calcola.} In ABP\_Knight, nell'Event Graph dell'Update, si aggiunge un
pin alla \textbf{Sequence} (il \emph{+}, \emph{Add pin}): \emph{Then 2}, o il successivo.
Questa volta non si parte dal Movement Component, ma dal \textbf{character}: ``perché la
direzione è determinata dalla rotazione del character stesso'', ``cioè la capsula, come è
ruotata''. Dalla variabile del cavaliere si tira fuori \textbf{Get Actor Rotation}.

\paragraph{Il nodo Calculate Direction.} ``Potremmo tranquillamente prenderla così, ma noi,
per semplificarci anche la vita, andiamo a calcolare proprio la direzione con un blocco
chiamato proprio \textbf{Calculate Direction}.'' Il nodo chiede due ingressi: la
\textbf{Velocity} attuale (dal Movement Component, \textbf{Get Velocity}) e la
\textbf{Base Rotation} (la Get Actor Rotation). Restituisce ``un float'', che con
\textbf{Promote to Variable} diventa \textbf{Direction}.

\paragraph{Uno studente perso.} ``Io mi sono perso: Character, Movement Component, da dove
vengono questi due?'' Il docente ripercorre: ``hai fatto questo Event Blueprint Initialize
Animation nell'Event Graph, no?''; ``dovresti avere qua sotto, sotto le variabili, la
variabile del character: la tiri fuori, ne tiri fuori Get Actor Rotation''; dal Movement
Component ``tiri fuori la velocità''; ``da Get Actor Rotation fai Calculate Direction, e
questo diventerà il nostro float''. Un altro studente nota che la variabile del docente si
chiama \emph{Character}, mentre la sua si chiama \textbf{As BP Knight}: ``va benissimo.
Io l'ho solo rinominata.'' Quando si fa Promote to Variable sull'uscita di un cast, Unreal
propone il nome \emph{As} più il nome della classe (``ti dice proprio As BP Knight'').

\paragraph{Da $-180$ a $180$.} ``Molto semplicemente, la direzione sarà tra un range,
sempre senza doverla clampare'' (``il clamp l'abbiamo fatto quasi all'inizio''): fra
\textbf{$-180$ e $180$ gradi}. ``Si parte da 0; a destra vai verso il positivo; quando
arrivi a 180 passa a $-179$''; ``diciamo, non lo fa come un angolo da 0 a 360, ma fa 0,
180, $-180$ e di nuovo 0'' (figura~\ref{ue9:fig:direzione}).

\begin{figure}[htbp]
\centering
\begin{tikzpicture}[font=\footnotesize]
  \draw[thick] (0,0) circle (2.2);
  \foreach \a/\l/\pos in {90/0/above, 45/45/above right, 0/90/right, -45/135/below right,
      -90/{180 / $-180$}/below, -135/{$-135$}/below left, 180/{$-90$}/left,
      135/{$-45$}/above left}
    {\draw[-{Stealth[length=2mm]}, thick, blue!60!black] (0,0) -- (\a:2.2);
     \node[\pos, font=\scriptsize\bfseries] at (\a:2.25) {\l};}
  \fill[gray!40] (0,0) circle (0.35);
  \draw[thick] (0,0) circle (0.35);
  \draw[-{Stealth[length=2mm]}, very thick] (0,0) -- (0,0.6);
  \node[font=\scriptsize, align=left, anchor=west] at (3.4,1.2) {\textbf{0}: avanti\\
    \textbf{positivi}: verso destra\\\textbf{negativi}: verso sinistra\\
    \textbf{$\pm$180}: indietro};
  \node[font=\scriptsize, align=left, anchor=west] at (3.4,-0.9) {la direzione è
    l'angolo fra\\la velocità e il ``davanti''\\del personaggio (Base Rotation)};
\end{tikzpicture}
\caption{Il valore di Direction restituito da Calculate Direction. Il personaggio guarda
in alto; le frecce sono le otto direzioni delle animazioni del Blend Space 2D. Andare
indietro vale sia 180 sia $-180$, a seconda del lato da cui ci si arriva. Schema
ricostruito.}
\label{ue9:fig:direzione}
\end{figure}

\paragraph{Rimettere Use Controller Rotation Yaw.} Per vedere i movimenti laterali,
``importante è rimettere lo \textbf{Use Controller Rotation Yaw}'': ``non vogliamo che il
character ruoti sempre nella direzione in cui lo stiamo indirizzando''. Il corpo deve
seguire la camera, e la direzione del movimento rispetto al corpo deve cambiare.

\paragraph{La verifica con Print String.} Per sicurezza si stacca un attimo il Blend Space
1D (``staccare un attimo le animazioni'') e si stampa Direction con un \textbf{Print
String}, ``che lo farà a ogni frame''. ``Ci dà direttamente nella preview questo valore'';
in Play, ``sarà un po' rotto, ma è normale, ragazzi'': ``$-70$, 70, 160, 0''. ``Mamma mia''
(il valore dipende dalla camera, ``giustamente''). ``Come possiamo vedere, questi valori
rappresentano la nostra direzione attuale.'' Il docente nota anche che il cavaliere
sembra ``leggermente ruotato'' rispetto alla capsula: ``buono a sapersi''.

\paragraph{Approfondimento: che cosa calcola Calculate Direction.} Il nodo, della
libreria dell'Animation Blueprint (\emph{KismetAnimationLibrary}; nelle versioni più
vecchie era una funzione dell'AnimInstance), trasforma la velocità nello spazio locale
del personaggio e ne restituisce l'angolo sul piano orizzontale, in gradi, nell'intervallo
$[-180, 180]$: 0 se la velocità punta dove guarda il personaggio, 90 se punta alla sua
destra, $-90$ a sinistra, $\pm 180$ dietro. A velocità nulla restituisce 0. Con
\emph{Orient Rotation to Movement} acceso il personaggio si gira sempre verso la
velocità, e la direzione resta quasi sempre vicina a zero: il Blend Space 2D avrebbe poco
da fare. Per questo il docente riattiva Use Controller Rotation Yaw.

% ---------------------------------------------------------------------
\section{Il Blend Space 2D}
\label{ue9:sec:blend2d}

\paragraph{Creare l'asset.} Tasto destro $\rightarrow$ \emph{Animation} $\rightarrow$
\textbf{Blend Space}: ``Blend Space senza niente, di base, è 2D; se viene specificato
1D, è 1D''. Unreal chiede lo scheletro: ``voi scegliete il vostro scheletro''. Il nome:
\textbf{BS2D\_Knight}.

\paragraph{I due assi.} ``Come potete vedere, a sinistra e giù sono presenti due nomi'',
per ora \emph{None}. Nei Details:
\begin{itemize}
  \item \textbf{Horizontal Axis}: \textbf{Direction}, da \textbf{$-180$} a
        \textbf{$180$};
  \item \textbf{Vertical Axis}: \textbf{GroundSpeed} (``più che la velocità, la ground
        speed''), da \textbf{0} a \textbf{600}.
\end{itemize}
Il parametro \textbf{Grid Divisions} ``ci permette di aumentare il numero di divisioni
sulla nostra griglia'' (le linee ``non si vedono, perché sono veramente molto poco
chiare''): sono i punti su cui si agganciano i campioni delle animazioni. Il docente ne
mette \textbf{8} su entrambi gli assi. Con 8 divisioni, sull'asse orizzontale ogni passo
vale 45 gradi, esattamente le otto direzioni; sull'asse verticale vale 75.

\paragraph{Un quiz.} ``Domanda, ragazzi: se io sto a 600 di velocità e a 90 di direzione,
dove sto andando?'' ``Ricordatevi che positivo è destra, negativo è sinistra.'' ``A
destra.'' ``Pienamente a destra, perfetto.'' ``E invece a $-135$?'' ``Dietro a
sinistra.'' ``Ottimo.''

\paragraph{Dove va ogni animazione.} Si dispongono i campioni (tabella~\ref{ue9:tab:bs2d} e
figura~\ref{ue9:fig:blend2d}). In basso, a velocità zero, l'\textbf{Idle}: ``mettetene
tre, così abbiamo un range per tutto'' (a $-180$, 0 e 180): ``siamo fermi, siamo fermi''.
In alto, a 600, le otto animazioni. L'animazione all'indietro va messa due volte, a
$-180$ e a $180$: ``stiamo andando indietro sia se abbiamo girato da sinistra, sia da
destra''. Il docente si corregge una volta (``Backward Diagonal L sarà\dots\ no, ho
sbagliato'': a $-135$, non a 90) e invita a controllare i valori nei Details del
campione.

\begin{table}[htbp]
\centering
\small
\begin{tabular}{@{}p{4.4cm}p{2.2cm}p{2.2cm}p{4.6cm}@{}}
\toprule
\textbf{Animazione} & \textbf{Direction} & \textbf{GroundSpeed} & \textbf{Movimento} \\
\midrule
Idle (tre copie) & $-180$, 0, 180 & 0 & fermo \\
Jog Forward & 0 & 600 & avanti \\
Jog Forward Diagonal R & 45 & 600 & avanti a destra \\
Strafe R & 90 & 600 & a destra \\
Jog Backward Diagonal R & 135 & 600 & indietro a destra \\
Jog Backward & 180 e $-180$ & 600 & indietro \\
Jog Backward Diagonal L & $-135$ & 600 & indietro a sinistra \\
Strafe L & $-90$ & 600 & a sinistra \\
Jog Forward Diagonal L & $-45$ & 600 & avanti a sinistra \\
\bottomrule
\end{tabular}
\caption{La posizione dei campioni nel Blend Space 2D BS2D\_Knight. I nomi delle
animazioni sono quelli di Mixamo citati a lezione, per quanto ricostruibili dalla
trascrizione; i valori sono quelli dettati dal docente.}
\label{ue9:tab:bs2d}
\end{table}

\begin{figure}[htbp]
\centering
\begin{tikzpicture}[font=\footnotesize, x=0.032cm, y=0.006cm,
  ue9pt/.style={draw=black, fill=white, thick, diamond, inner sep=1.8pt}]
  % griglia
  \foreach \x in {-180,-135,...,180} \draw[gray!30] (\x,0) -- (\x,600);
  \foreach \y in {0,75,...,600} \draw[gray!30] (-180,\y) -- (180,\y);
  \draw[thick] (-180,0) rectangle (180,600);
  \foreach \x in {-180,-90,0,90,180}
    \node[below, font=\scriptsize] at (\x,-10) {\x};
  \foreach \y in {0,300,600} \node[left, font=\scriptsize] at (-182,\y) {\y};
  \node[below, font=\small] at (0,-80) {Direction};
  \node[rotate=90, font=\small] at (-235,300) {GroundSpeed};
  % campioni
  \foreach \x in {-180,0,180} \node[ue9pt] at (\x,0) {};
  \node[above, font=\scriptsize] at (0,10) {Idle};
  \foreach \x/\l in {-180/Back, -135/BackDiag L, -90/Strafe L, -45/FwdDiag L, 0/Fwd,
      45/FwdDiag R, 90/Strafe R, 135/BackDiag R, 180/Back}
    {\node[ue9pt] at (\x,600) {};
     \node[above, font=\tiny, rotate=35, anchor=south west] at (\x-6,612) {\l};}
  % punto attuale
  \fill[green!50!black] (60,450) circle (2.2pt);
  \node[green!40!black, font=\scriptsize, anchor=west] at (66,420) {(60, 450)};
\end{tikzpicture}
\caption{Il Blend Space 2D del cavaliere: Direction in orizzontale, GroundSpeed in
verticale, griglia 8$\times$8. Le tre copie dell'Idle stanno in basso; le otto direzioni
in alto, con la corsa all'indietro ripetuta ai due estremi. Il punto verde è un esempio:
il personaggio sta andando in avanti a destra, un po' sotto la velocità massima, e Unreal
mescola i quattro campioni più vicini. Schema ricostruito.}
\label{ue9:fig:blend2d}
\end{figure}

\paragraph{La preview con Ctrl.} ``Potete fare la preview delle animazioni con
\textbf{Ctrl}'': tenendolo premuto e trascinando il cursore sul grafico, ``potete vedere
come l'animazione si muove tra tutti gli step''. ``Consiglio sempre di mettersi dietro''
il personaggio nel viewport, ``perché tanto poi, se lo vedete davanti, è difficile'' capire
destra e sinistra. Il docente controlla tutte le direzioni (``vado a sinistra, vado
indietro, vado dritto, vado dritto a destra, vado a destra, giù a destra, indietro.
Perfetto''); una delle animazioni ``mi sa che non va bene'', ma ``sì, sto andando a
sinistra, è solo molto strana''.

\paragraph{Smoothing Time.} ``Piccola cosa, se volete introdurla'': lo \textbf{Smoothing
Time}, ``il tempo di passaggio tra un'animazione e l'altra in generale, per tutte le
animazioni''. Lo si mette a \textbf{0,2} su entrambi gli assi. Il \emph{Damping Ratio}
``non ci interessa''.

\paragraph{Nello stato WalkRun.} ``Facciamo tutto il percorso a ritroso'': AnimGraph,
Locomotion, WalkRun. Al posto del Blend Space 1D si trascina BS2D\_Knight, che chiede
``dammi la Direction e la GroundSpeed'': si collegano i due Get. Compile: ``vediamo che il
nostro Blend Space funziona in preview''.

\paragraph{Le impostazioni del personaggio.} In BP\_Knight1: ``mi raccomando, l'abbiamo
detto prima, ma lo diciamo anche adesso'', si riattiva \textbf{Use Controller Rotation
Yaw}; e nel Character Movement la \textbf{Max Walk Speed} va a \textbf{600}, ``che è
quanto abbiamo messo'' come massimo dell'asse. ``Ecco a voi, ragazzi'': ``se io mi muovo a
destra lui va a destra, se mi muovo a sinistra, anche se l'animazione non è fantastica,
va a sinistra; vado indietro; vado in diagonale''. Nel frattempo ``parte la torretta, che
mi spara'': ``il mio bel cavaliere se ne va a zonzo come gli pare''.

\paragraph{Gli scatti.} C'è ``un piccolo \emph{clip}'', uno scatto, in alcuni passaggi.
Il docente prova a prendere la rotazione della \textbf{Mesh} invece di quella
dell'actor (``facciamo che invece di prendere la direzione del character, prendiamo
quella della Mesh''; ``Get Mesh''), poi ad aumentare lo Smoothing Time della Direction
a \textbf{0,5} (``un po' più tranquillo''), e infine torna al collegamento originale.
La conclusione: ``probabilmente sono dei problemi sull'animazione''; ``mi sa che devo
riscaricare l'animazione, perché snappa un pochino''. Ma il risultato c'è: ``abbiamo il
nostro character che si muove in tutte le direzioni''; ``è un po' veloce, devo essere
sincero'', e la velocità si può abbassare.

\paragraph{Il cerchio, per chi si è perso.} Uno studente chiede di rivedere il Blend
Space 2D. ``Immaginatelo come un cerchio'': ``se stai a 0 di direction, stai andando
dritto; se stai a 180 o $-180$, stai andando dietro''; ``tutti i valori intermedi'' sono le
diagonali e i lati. ``Destra è positivo: 0, 90, 180; sinistra è $-90$, $-180$.'' ``Se vai
in diagonale a destra sarà 45, se vai pienamente a destra sarà 90'', ``e 135 a 600 vuol
dire che andrai praticamente indietro, ma verso destra''. ``Una volta capito che una è
l'opposto dell'altra'', i valori negativi seguono da soli.

\paragraph{Approfondimento: Blend Space 2D e foot sliding.} In un Blend Space 2D Unreal
divide lo spazio dei campioni in triangoli e, per un valore dato, mescola i campioni ai
vertici del triangolo che lo contiene. Ne seguono tre conseguenze pratiche. Primo,
l'Idle ripetuto lungo tutta la riga inferiore (le tre copie del docente) evita che a
bassa velocità il personaggio ``scelga'' una direzione a caso. Secondo, la corsa
all'indietro va duplicata ai due estremi, perché $-180$ e 180 sono lo stesso angolo ma
due punti diversi della griglia: con un solo campione, passando da 179 a $-179$ il
personaggio attraverserebbe in un istante tutto il grafico. Terzo, lo scatto che il
docente vede è spesso proprio quel salto da $+180$ a $-180$, che lo \emph{Smoothing
Time} non può interpolare bene (interpola il valore numerico, non l'angolo); le
alternative sono il \emph{Weight Speed} per campione, oppure le soluzioni più avanzate
di Unreal Engine 5 come l'\emph{Orientation Warping}, che ruota la parte inferiore del
corpo verso la direzione del movimento con quattro sole animazioni cardinali. Infine,
una griglia con l'asse verticale fino a 600 e i campioni tutti a 600 significa che anche
la camminata a 300 è un misto fra Idle e corsa: per avere una vera camminata laterale
servirebbe una seconda riga di campioni a 300, con le animazioni di camminata.

% ---------------------------------------------------------------------
\section{Gli Animation Montage}
\label{ue9:sec:montage}

``Intanto ci abbiamo un'ultima mezz'oretta, ragazzi. Vediamo un attimo i Montage.''
L'animazione di attacco (``l'avevo scaricata, mi sa che non l'ho messa in Knight
Animation''; ``scusate, ragazzi'') è una \textbf{Sword and Shield Attack} di Mixamo:
``questa qua sta in root motion'': ``come potete vedere, eredita il movimento''.

\paragraph{Il problema.} ``Quando vogliamo fare una nostra animazione che di conseguenza
sia attivabile durante un ciclo di animazioni'', non basta aggiungere uno stato: ``dobbiamo
creare quello che viene chiamato un \textbf{Montage}''. (Nel frattempo ``Filippo è andato
via, è crashato'': ``siamo rimasti noi tre''.)

\paragraph{Che cos'è un Montage, dalla documentazione.} Il docente apre la pagina della
documentazione di Unreal sugli \textbf{Animation Montage}. ``Il montaggio di animazioni ci
permette sia di montare delle animazioni\dots'' L'esempio della pagina è la ricarica di
un fucile: un'animazione formata da tre animazioni diverse, \emph{MM Shotgun Reload Drop},
\emph{Reload Load} e \emph{Reload Lock}. ``Mettiamo conto che vogliamo fare un bel gioco
shooter. Il fatto di staccare il caricatore su quattro fucili diversi è praticamente
uguale''; ``non dobbiamo fare quattro animazioni diverse'': ``l'unica cosa che cambia è il
lock, perché magari uno vuole farlo un po' più grosso''. Con un Montage si creano
\textbf{animazioni modulari}: ``guarda, il drop è uguale, il load è uguale, il lock invece
è diverso''. E poi il secondo uso, quello di oggi: ``inserire animazioni in contesti di
cicli di animazione''; ``sono una tantum'': ``premo un tasto e compio questa animazione,
nel mentre che sto camminando''. ``Magari la prossima volta andiamo a vedere pure un po'
più approfondito.''

\paragraph{Creare il Montage.} Tasto destro sull'animazione di attacco $\rightarrow$
\textbf{Create} $\rightarrow$ \textbf{Create AnimMontage}. Il nuovo asset si chiama come
l'animazione con il suffisso \emph{Montage} (\emph{Sword And Shield Attack\_Montage}).
``La cosa bella del Montage è che ci permette di fare un sacco di cose all'interno dei
suoi parametri.''

% ---------------------------------------------------------------------
\section{Dentro il Montage: Notify, sezioni e Slot}
\label{ue9:sec:notify}

\paragraph{I Notify.} ``Una cosa interessante da spiegare sono i \textbf{Notify}'': ``dei
punti, dei puntini, sulla linea del tempo dell'animazione'', che permettono di aggiungere
``un'azione che si triggera'': ``suoni, un VFX, qualsiasi cosa''. ``Quelli più generici
sono già fatti da Unreal, fortunatamente.'' Il docente fa clic destro sulla traccia dei
Notify, \textbf{Add Notify} $\rightarrow$ \textbf{Play Sound}, e come suono sceglie
``l'unico sound che abbiamo'', quello della piattaforma (un Sound Cue). Le tracce dei Notify
possono essere più d'una: con \emph{Add Notify Track} se ne crea un'altra, per esempio
una per i VFX e una per i suoni. In preview (tolto il loop, ``fammelo fare una volta''):
``appena il nostro character arriva a questo punto dell'animazione, compie il suono''.
``Mettiamo conto: in combattimento c'è un VFX, c'è un suono, c'è un tirare fuori la spada,
e quindi si sente la spada che esce fuori dal fodero'': con una linea di Notify ``questo
suono avviene in un preciso momento dell'animazione''. ``Ovviamente, non andremo a usare
il sound effect della piattaforma per il nostro Knight.''

\paragraph{Le Montage Section.} Nel \emph{Montage group} c'è la sezione \textbf{Default}.
Con il tasto destro, \textbf{New Montage Section}: il docente ne crea una chiamata
\textbf{Attack}, nel punto in cui ``lui sta attaccando''. Le sezioni servono ``per
ordinare il lavoro'' e per ``dividere l'animazione in più punti diversi, con più
animazioni''.

\paragraph{Più animazioni nello stesso Montage.} La striscia verde sotto le sezioni è
l'animazione stessa. ``Se vogliamo, possiamo anche introdurre delle animazioni in più.''
Una dimostrazione volutamente assurda: ``io voglio, per qualche motivo arcano, far sì
che quando finisce l'animazione lui salti'': si trascina \textbf{MM\_Jump} nella traccia,
dopo l'attacco. ``È un salto orribile, però farà le due animazioni di conseguenza.'' Il
Montage è ``molto importante per azioni consecutive''. ``Quando si parla con l'animatore,
è importante il fatto di dire: questa animazione deve essere consecutiva a questa'', per
esempio ``mi serve che da questo punto in poi continui l'attacco, oppure faccia un'altra
animazione che ritorna in place''.

\paragraph{L'animazione di ritorno.} Finito l'attacco in root motion, il cavaliere resta dove l'attacco lo ha lasciato e non
torna in posa: manca ``un'animazione che ritorna in place'', cioè in idle. Il docente la
cerca su Mixamo, provando voci come \emph{Sword and Shield Idle}, \emph{Idle
Transition} e \emph{Start Block}; ma ``è un po' difficile'', perché ``bisogna trovare
il cavaliere proprio'', con la stessa posa di fine attacco, altrimenti ``snapperà
malissimo, almeno nel mio caso''.

\paragraph{Gli Slot.} Sotto, c'è il \textbf{Default Group} con lo \textbf{Slot}
\textbf{DefaultSlot}: ``questo ci permette di dire: su quale slot vuoi che questa
animazione si compia?''. DefaultSlot ``riprende tutto il corpo''. ``Potremmo, ad esempio,
voler creare uno slot che riprenda solamente alcune zone del nostro character'': ``vogliamo
che faccia un'azione solo nella parte superiore del corpo'' (sparare o salutare mentre si
corre); ``in quel caso creeremo un altro slot, che prenderà solo le ossa'' di quella
parte. ``Ma in questo caso ci serve tutto il corpo.''

% ---------------------------------------------------------------------
\section{Lo Slot nell'AnimGraph: Save Cached Pose}
\label{ue9:sec:slot}

Un Montage suona solo se l'AnimGraph ha un nodo \textbf{Slot} con lo stesso nome. In
ABP\_Knight, nell'AnimGraph:
\begin{enumerate}
  \item dall'uscita della State Machine Locomotion si crea un \textbf{New Save Cached
        Pose}, chiamato \textbf{LocomotionCache}: ``andiamo a conservarci quello che ha
        fatto la State Machine'', ``come se tutto questo procedimento lo mettesse in una
        memoria'';
  \item si aggiunge \textbf{Use cached pose 'LocomotionCache'};
  \item da lì un nodo \textbf{Slot 'DefaultSlot'} (la trascrizione: ``default slot'');
  \item lo Slot va all'\textbf{Output Pose}.
\end{enumerate}
``Semplicemente mixa insieme i due'': lo slot del Montage e la cache della locomozione
(figura~\ref{ue9:fig:slot}). Quando non c'è nessun Montage in corso, lo Slot lascia
passare la locomozione; quando un Montage su DefaultSlot parte, la sostituisce. Alla
fine della lezione, a uno studente che non legge il nome del nodo (``vedo tutto
pixelato''), il docente conferma: ``tiri fuori Locomotion Cache'', poi ``default slot''.

\begin{figure}[htbp]
\centering
\begin{tikzpicture}[font=\footnotesize,
  ue9pose/.style={draw, thick, rounded corners=4pt, align=center, fill=gray!8,
                  minimum height=1.0cm, text width=2.3cm},
  ue9arr/.style={-{Stealth[length=2.5mm]}, very thick, black!70}]
  \node[ue9pose, fill=orange!12] (sm) at (0,0) {\textbf{Locomotion}\\{\scriptsize State
    Machine}};
  \node[ue9pose, fill=blue!8] (sc) at (3.1,0) {\textbf{Save Cached Pose}\\{\scriptsize
    LocomotionCache}};
  \node[ue9pose, fill=blue!8] (uc) at (6.2,0) {\textbf{Use Cached Pose}\\{\scriptsize
    LocomotionCache}};
  \node[ue9pose, fill=green!10] (sl) at (9.3,0) {\textbf{Slot}\\{\scriptsize
    'DefaultSlot'}};
  \node[ue9pose, fill=red!10, text width=1.8cm] (op) at (12.1,0) {\textbf{Output
    Pose}};
  \draw[ue9arr] (sm) -- (sc); \draw[ue9arr] (uc) -- (sl); \draw[ue9arr] (sl) -- (op);
  \draw[dashed, thick, gray] (sc.east) -- (uc.west);
  \node[ue9pose, fill=yellow!15, text width=3.2cm] (mt) at (9.3,-2.1)
    {\textbf{Anim Montage}\\{\scriptsize Sword and Shield Attack, slot
    DefaultGroup.DefaultSlot}};
  \draw[-{Stealth[length=2.5mm]}, thick, dashed] (mt) -- node[right, font=\scriptsize]
    {quando suona, sostituisce la posa} (sl);
\end{tikzpicture}
\caption{La catena di pose finale dell'AnimGraph di ABP\_Knight. La State Machine produce
la locomozione, che viene salvata in una cache e riletta; lo Slot DefaultSlot la lascia
passare, a meno che un Montage assegnato a quello slot non sia in esecuzione. Schema
ricostruito.}
\label{ue9:fig:slot}
\end{figure}

% ---------------------------------------------------------------------
\section{Far partire l'attacco: Play Montage}
\label{ue9:sec:playmontage}

\paragraph{Il tasto F e il nodo.} Per far partire il Montage, il docente non crea
un'altra Input Action: usa direttamente un evento di tastiera, \textbf{F} (``facciamo con
OK\dots\ no, facciamo F, F va benissimo''). Nel Blueprint del personaggio, il nodo
\textbf{Play Montage}:
\begin{itemize}
  \item \textbf{In Skeletal Mesh Component}: ``in quale Skeletal Mesh Component
        farlo'', cioè la \textbf{Mesh} del personaggio;
  \item \textbf{Montage to Play}: ``quello che abbiamo appena fatto'' (``ce n'è tanto solo
        uno'');
  \item \textbf{Play Rate}: ``possiamo modificare quanto è veloce l'animazione'';
  \item fra le uscite: \textbf{On Completed} (``quando sarà finita''), \textbf{On Blend
        Out} (``mentre finisce''), \textbf{On Interrupted} (``se fosse interrotta''),
        \textbf{On Notify Begin} e \textbf{On Notify End} (``quando un Notify inizia'' e
        ``quando finisce'').
\end{itemize}
``Se noi premiamo F, compiamo un Montage.''

\paragraph{Fermare il personaggio.} ``È orribile'': a parte la mancanza dell'animazione di
ritorno, durante l'attacco ``possiamo muovere il character tranquillamente''. Quindi, sul
tasto F, prima del Play Montage:
\begin{enumerate}
  \item \textbf{Set Max Walk Speed} a \textbf{0}: ``non dobbiamo farlo muovere'';
  \item \textbf{Get Player Controller} $\rightarrow$ \textbf{Disable Input}: ``andiamo a
        disabilitare gli input, qualsiasi''.
\end{enumerate}
Sulle uscite del Play Montage:
\begin{itemize}
  \item \textbf{On Completed}: Get Player Controller $\rightarrow$ \textbf{Enable Input}
        (``on completed solamente'');
  \item \textbf{On Blend Out}: \textbf{Set Max Walk Speed} a \textbf{600}, ``al suo
        massimo''.
\end{itemize}

\paragraph{La camera che si muove ancora.} In Play, però, ``la camera si muove
effettivamente'', e con la camera il personaggio, ``quindi dobbiamo far sì che non sia
dipendente dalla camera''. Segue una serie di tentativi, ragionati ad alta voce:
\begin{itemize}
  \item togliere \emph{Use Pawn Control Rotation} sullo Spring Arm (``no, no'');
  \item l'\emph{Enable Camera Lag} dello Spring Arm;
  \item le spunte \emph{Inherit Pitch}, \emph{Inherit Yaw} e \emph{Inherit Roll} dello
        Spring Arm (``inerit yo e roll''; ``no, lo yaw ci deve essere''; ``forse sì,
        solo il roll va bene''), con un \emph{Set} preso dallo Spring Arm (``Set Inherit
        Yaw'', a false, e di nuovo a true alla fine); ``completamente\dots\ non va'';
  \item \emph{Stop Movement Immediately} e i parametri di input del Player Controller
        (\emph{movement input}, \emph{look input}).
\end{itemize}
La soluzione, alla fine, è sulle due spunte già incontrate nella prima parte: ``non è lo
yaw'' dello Spring Arm, ``ma è l'orientamento al movimento''. Sul tasto F si disattivano
entrambe:
\begin{itemize}
  \item dal \textbf{Character Movement}, \textbf{Set Orient Rotation to Movement} a
        \textbf{false};
  \item sul Pawn stesso, senza collegare nessun target, \textbf{Set Use Controller
        Rotation Yaw} a \textbf{false};
\end{itemize}
e sull'\textbf{On Completed} si rimettono entrambe a \textbf{true}
(figura~\ref{ue9:fig:playmontage}). ``Ottimo. Adesso riusciamo'': il cavaliere attacca
fermo, la camera gira senza trascinarlo, e alla fine riprende la locomozione. ``Abbiamo
fatto il nostro character che attacca e fa la locomotion.'' (Il docente nota anche, di
passaggio, che il cavaliere non interagisce con il pulsante del livello, ``giustamente,
perché è la Mesh''; dalla trascrizione non si ricostruisce meglio il punto.)

A Federico, che non riesce a sistemare la camera, il docente riassume: ``devi giocare con
lo yaw e con l'Orient Rotation to Movement: li disattivi tutte e due, e li riattivi
sull'On Completed''. ``Dello Spring Arm, giusto?'' ``Non dello Spring Arm: questo qua è
del Pawn, lo trovi direttamente senza collegare niente. Invece questo è solo del
Character Movement.''

\begin{figure}[htbp]
\centering
\begin{tikzpicture}[font=\footnotesize,
  ue9node/.style={draw, thick, rounded corners=4pt, align=center, fill=gray!8,
                  minimum height=0.9cm, text width=2.5cm},
  ue9ev/.style={ue9node, fill=red!12},
  ue9exec/.style={-{Stealth[length=2.5mm]}, very thick, black!80}]
  \node[ue9ev, text width=1.4cm] (f) at (0,0) {\textbf{F}\\{\scriptsize Pressed}};
  \node[ue9node] (a) at (2.8,0) {\textbf{Set Max Walk Speed}\\{\scriptsize 0}};
  \node[ue9node] (b) at (5.9,0) {Get Player Controller $\rightarrow$ \textbf{Disable
    Input}};
  \node[ue9node] (c) at (9.0,0) {\textbf{Set Orient Rotation to Movement}\\
    {\scriptsize false}};
  \node[ue9node] (d) at (12.1,0) {\textbf{Set Use Controller Rotation Yaw}\\
    {\scriptsize false}};
  \node[ue9node, fill=yellow!15, text width=3.0cm, minimum height=2.4cm] (pm) at
    (2.8,-3.0) {\textbf{Play Montage}\\{\scriptsize Mesh,\\Sword and Shield
    Attack\_Montage}};
  \node[ue9node] (e) at (8.3,-2.2) {\textbf{Set Max Walk Speed}\\{\scriptsize 600}};
  \node[ue9node, text width=3.0cm] (g) at (8.3,-3.9) {Get Player Controller
    $\rightarrow$ \textbf{Enable Input}};
  \node[ue9node, text width=3.0cm] (h) at (12.3,-3.9) {\textbf{Set Orient\dots} true\\
    \textbf{Set Use Controller Rotation Yaw} true};
  \draw[ue9exec] (f) -- (a); \draw[ue9exec] (a) -- (b); \draw[ue9exec] (b) -- (c);
  \draw[ue9exec] (c) -- (d);
  \draw[ue9exec] (d.south) -- ++(0,-0.6) -| (pm.north);
  \draw[ue9exec] (pm.east |- e) -- node[above, font=\scriptsize]{On Blend Out} (e);
  \draw[ue9exec] (pm.east |- g) -- node[below, font=\scriptsize]{On Completed} (g);
  \draw[ue9exec] (g) -- (h);
\end{tikzpicture}
\caption{L'attacco sul tasto F, nella versione finale. Prima del Montage si blocca il
personaggio (velocità zero, input disabilitato, nessuna rotazione automatica del corpo);
al Blend Out torna la velocità, all'On Completed tornano l'input e le rotazioni. L'ordine
esatto dei primi quattro nodi è ricostruito: a lezione vengono aggiunti uno alla volta,
durante i tentativi. Schema ricostruito.}
\label{ue9:fig:playmontage}
\end{figure}

\paragraph{Approfondimento: Play Montage, Montage Play e il root motion.} Il nodo usato a
lezione, \textbf{Play Montage}, è un nodo \emph{asincrono} (con l'orologio nell'angolo):
le sue uscite scattano più tardi, quando succede l'evento corrispondente, e per questo
può gestire il ``dopo'' del Montage direttamente nel grafo. Esiste anche \emph{Montage
Play}, chiamato sull'\emph{Anim Instance} della mesh, che restituisce solo la durata, e
\emph{Play Anim Montage}, del Character. Due dettagli completano la lezione. Il primo:
un attacco in \textbf{root motion} sposta davvero la capsula solo se il Montage ha il
root motion attivo e l'Animation Blueprint ha la \emph{Root Motion Mode} impostata su
\emph{Root Motion from Montages Only} (il default); se il root dell'animazione di Mixamo è
``rotto'', come visto nel retargeting, il personaggio avanza nella mesh e poi ``torna
indietro'' di scatto a fine animazione. Il secondo: durante un Montage che fa avanzare il
personaggio in root motion non serve azzerare la Max Walk Speed, perché il root motion
prende il posto dell'input di movimento; azzerarla, come fa il docente, è però un modo
semplice e sicuro di impedire al giocatore di muoversi durante un attacco ``in place''.
Infine, \emph{Disable Input} e \emph{Enable Input} agiscono sull'actor e richiedono un
Player Controller: per bloccare solo alcune azioni (il movimento ma non la camera),
l'Enhanced Input offre un'alternativa più fine, cioè rimuovere temporaneamente un
Input Mapping Context, o aggiungerne uno di priorità più alta.

% ---------------------------------------------------------------------
\section{Il compito e la chiusura}
\label{ue9:sec:chiusura}

\paragraph{Feedback in chat.} Una studentessa, Giulia, ha scritto in chat che è rimasta un
po' indietro ``causa trasloco'' e chiede di poter mandare dei feedback; nel frattempo è
già uscita dalla chiamata. ``Le risponderò sul Discord, assolutamente.'' ``Controllerò i
lavori che mandate sul canale.''

\paragraph{Il compito: lo StackOBot.} ``Una cosa che vorrei darvi come task, se vi va, per
la prossima volta'': ``fare un nuovo character totalmente da zero''. Nel progetto c'è già
un personaggio, lo \textbf{StackOBot} (il robottino dei progetti di esempio di Epic), che
``ha queste animazioni qua: vedete che ce le ha già dentro tutte quante''. Lo si trova
nella cartella \emph{StackOBot $\rightarrow$ Blueprints $\rightarrow$ Character}, con la
sua Skeletal Mesh (\emph{SKM\_Bot}) e il suo ABP. Il compito:
\begin{itemize}
  \item creare un \textbf{Blueprint Character} per lo StackOBot (``fare la BP'');
  \item creare il suo \textbf{Animation Blueprint} da zero, ``fare tutto l'Animation
        Blueprint e collegarci tutte le sue animazioni'': il ciclo di locomozione, con
        State Machine e Blend Space, come quello del cavaliere.
\end{itemize}
``Se volete farlo, lo potete mandare poi sulla chat'', ``che tanto magari ci lavoreremo
la prossima volta''. ``Per adesso, poi ne riparliamo.''

\paragraph{I saluti.} ``Se tutto a posto\dots'' Il docente interrompe la condivisione:
``vi ringrazio per l'attenzione, e ringrazio da casa chi è in differita''. ``Ci sentiamo
per martedì prossimo. Va bene, ciao, ragazzi. Buona serata.''

\begin{table}[htbp]
\centering
\small
\begin{tabular}{@{}p{5.4cm}p{8.6cm}@{}}
\toprule
\textbf{Sintomo} & \textbf{Causa e rimedio} \\
\midrule
L'animazione dello stato non si vede nella preview & la State Machine non è collegata
all'Output Pose dell'AnimGraph \\
L'animazione si ferma dopo un ciclo & manca la spunta di loop nel nodo dell'animazione \\
Warning nella State Machine & transizioni senza condizioni: normale finché non si
scrivono le regole \\
Jump $\rightarrow$ Fall dà errore & regola vuota: attivare Automatic Rule Based on
Sequence Player in State \\
Il Blend Space 1D arriva solo a 100 & Maximum Axis Value da impostare a 600; poi
spostare i campioni \\
Animazioni ``mischiate'' a metà (a 200) & i campioni non corrispondono alle velocità
reali: legarli alla Max Walk Speed \\
Alt sul nodo non crea il Set & negli stati dell'AnimGraph le variabili si leggono e
basta \\
``Questo asset non ha schede ancorate'' & Window $\rightarrow$ Load Layout $\rightarrow$
Default Editor Layout, poi Open Full Blueprint Editor \\
In Play non si vede niente / visuale nel corpo & il Character non ha camera: aggiungere
Spring Arm e Camera (figlia) \\
Non ci si muove & mancano gli input: copiare dal ThirdPersonCharacter i grafi di Look,
Move, Jump e l'Add Mapping Context \\
La camera ruota su un solo asse, o guarda solo il terreno & Use Pawn Control Rotation
sulla Camera invece che sullo Spring Arm \\
Il personaggio va di lato guardando avanti & attivare Orient Rotation to Movement e
togliere Use Controller Rotation Yaw (servono tutte e due) \\
La camera compenetra gli oggetti & Target Offset troppo grande: usare il Socket Offset \\
Il salto in root motion retargettato non è pulito & il root degli scheletri di Mixamo:
accettabile per la prova, da correggere \\
La State Machine non distingue salto e caduta & entrambe le regole su IsFalling:
aggiungere VelocityZ maggiore o minore di zero \\
Lo Shift non fa correre & IA\_Run non aggiunta ai mapping di IMC\_Default \\
FInterp To non interpola & eseguito una volta sola: va eseguito a ogni frame (Tick,
timer o Timeline) \\
Direction quasi sempre a zero & Orient Rotation to Movement acceso, Use Controller
Rotation Yaw spento: invertire per il Blend Space 2D \\
Scatti nel Blend Space 2D & animazioni non omogenee, salto fra 180 e $-180$: Smoothing
Time, o riscaricare le animazioni \\
Durante l'attacco il personaggio si muove & Set Max Walk Speed a 0 e Disable Input prima
del Play Montage \\
Durante l'attacco la camera trascina il personaggio & Set Orient Rotation to Movement e
Set Use Controller Rotation Yaw a false, a true sull'On Completed \\
Il Montage non suona & manca il nodo Slot DefaultSlot nell'AnimGraph \\
\bottomrule
\end{tabular}
\caption{I problemi incontrati durante la lezione, con la causa e il rimedio.}
\label{ue9:tab:errori}
\end{table}

\begin{table}[htbp]
\centering
\small
\begin{tabular}{@{}p{4.8cm}p{9.2cm}@{}}
\toprule
\textbf{Nodo, asset o impostazione} & \textbf{Che cosa fa} \\
\midrule
State Machine (AnimGraph) & organizza le animazioni in stati e transizioni \\
Entry / State / Conduit / Alias & ingresso; stato con una posa; snodo senza posa; alias
di più stati \\
Output Animation Pose & la posa prodotta da uno stato \\
Transition Rule (Result) & la condizione booleana per passare da uno stato all'altro \\
Automatic Rule Based on Sequence Player in State & la transizione scatta a fine
animazione \\
NOT Boolean / AND Boolean & negazione; congiunzione di due condizioni \\
Split Struct Pin & divide un vettore nei suoi componenti X, Y, Z \\
Blend Space 1D / Blend Space (2D) & mescola animazioni su uno o due assi \\
Horizontal / Vertical Axis, Grid Divisions & nome, minimo e massimo degli assi;
divisioni della griglia \\
Smoothing Time & tempo di passaggio fra i campioni \\
Blend Space Player & il nodo che suona un Blend Space nell'AnimGraph \\
Calculate Direction & l'angolo, da $-180$ a 180, fra la velocità e la rotazione \\
Get Actor Rotation / Get Velocity & la rotazione dell'actor; la sua velocità \\
Default Pawn Class (World Settings) & il personaggio del giocatore nel livello \\
Spring Arm: Target Arm Length, Socket Offset, Target Offset & lunghezza del braccio;
spostamento in fondo al braccio; spostamento dell'origine \\
Camera: Field of View & angolo di visuale \\
Use Pawn Control Rotation (Spring Arm) & la camera segue il controller \\
Use Controller Rotation Yaw (Pawn) & il corpo segue la rotazione del controller \\
Orient Rotation to Movement, Rotation Rate (Character Movement) & il corpo si gira verso
il movimento, con quella velocità \\
Max Walk Speed / Set Max Walk Speed & velocità massima a piedi (cm/s) \\
Retarget Animations, Export Animations & trasferisce un'animazione su un altro
scheletro \\
Input Mapping Context (IMC) & associa i tasti alle azioni in un contesto \\
Input Action (IA): Value Type, Triggers & l'azione, il tipo di valore, quando scatta \\
Add Mapping Context & attiva un IMC per il giocatore \\
Evento IA: Triggered, Started, Ongoing, Canceled, Completed & le fasi di un'azione \\
FInterp To & un passo di interpolazione fra Current e Target \\
Get World Delta Seconds & il tempo dell'ultimo frame \\
Animation Montage & animazione ``una tantum'', modulare, con sezioni e slot \\
Notify / Play Sound / Notify Track & eventi sulla linea del tempo; un suono; una traccia
di eventi \\
Montage Section & una parte con nome del Montage \\
Slot (DefaultGroup.DefaultSlot) & dove il Montage si applica (tutto il corpo, o parti) \\
Save Cached Pose / Use Cached Pose & salva e rilegge una posa nell'AnimGraph \\
Play Montage & suona un Montage; On Completed, On Blend Out, On Interrupted, Notify
Begin/End \\
Disable Input / Enable Input & toglie e restituisce l'input all'actor \\
\bottomrule
\end{tabular}
\caption{I nodi, gli asset e le impostazioni della nona lezione.}
\label{ue9:tab:nodi}
\end{table}

% ---------------------------------------------------------------------
\section*{Punti chiave del capitolo}
\addcontentsline{toc}{section}{Punti chiave del capitolo}

\begin{itemize}
  \item Una \textbf{State Machine} nell'AnimGraph organizza le animazioni in
        \textbf{stati} (WalkRun, Jump, Fall), ciascuno con la sua \textbf{Output
        Animation Pose}, e in \textbf{transizioni} con regole booleane: un ``loop di
        animazione secondo delle condizioni''. Va collegata all'Output Pose.
  \item Un \textbf{Blend Space 1D} mescola animazioni lungo un asse: per la locomozione,
        la \textbf{GroundSpeed}. I campioni vanno posti alle velocità reali del
        personaggio (Max Walk Speed, in centimetri al secondo): Idle a 0, Walk a 300,
        corsa a 600. Valori intermedi ``inventati'' producono obbrobri.
  \item Un \textbf{Character} non ha la camera: si aggiungono uno \textbf{Spring Arm} e,
        come suo figlio, una \textbf{Camera}. Distanza e altezza si regolano sullo
        Spring Arm (\textbf{Target Arm Length}, \textbf{Socket Offset}), mai spostando la
        camera; il \textbf{Target Offset} solo sulla Z e poco. Sono il lavoro dei
        \textbf{3C}: Camera, Control, Character.
  \item Tre spunte governano la rotazione: \textbf{Use Pawn Control Rotation} sullo
        Spring Arm (la camera segue il mouse), \textbf{Use Controller Rotation Yaw} sul
        Pawn (il corpo segue la camera), \textbf{Orient Rotation to Movement} nel
        Character Movement (il corpo si gira verso il movimento).
  \item Il \textbf{retargeting} porta un'animazione su un altro scheletro (Retarget
        Animations, Export Animations): comodo, ma il risultato dipende dall'animazione,
        e il root di Mixamo dà problemi.
  \item Fra salto e caduta: \textbf{Automatic Rule Based on Sequence Player in State} per
        passare a fine animazione, \textbf{NOT IsFalling} per tornare a terra, e la
        \textbf{velocità sulla Z} per distinguere un salto (Z $>$ 0) da una caduta
        (Z $<$ 0).
  \item L'\textbf{Enhanced Input} separa i tasti dalle azioni: l'\textbf{Input Action}
        (IA\_Run, Digital bool, con i suoi Triggers) è associata a uno o più tasti
        (Shift, il cerchio del controller) in un \textbf{Input Mapping Context}, che va
        attivato con \textbf{Add Mapping Context}. Lo sprint: Set Max Walk Speed a 600
        sul Triggered, a 300 sul Completed.
  \item \textbf{FInterp To} fa un solo passo per esecuzione: per un passaggio graduale
        va eseguito a ogni frame (Tick, timer, Timeline).
  \item Il \textbf{Blend Space 2D} aggiunge la \textbf{direzione} (\textbf{Calculate
        Direction}, da $-180$ a 180: positivo a destra, negativo a sinistra, $\pm 180$
        indietro): otto animazioni a 600, l'Idle ripetuto a 0, la corsa all'indietro ai
        due estremi, griglia 8$\times$8, Smoothing Time. Serve Use Controller Rotation Yaw
        acceso.
  \item Un \textbf{Animation Montage} è un'animazione ``una tantum'' e modulare, con
        \textbf{Notify} (suoni, VFX in un punto preciso), \textbf{sezioni} e
        \textbf{slot}. Per suonarlo, nell'AnimGraph: State Machine, \textbf{Save Cached
        Pose}, Use Cached Pose, \textbf{Slot DefaultSlot}, Output Pose. Si fa partire con
        \textbf{Play Montage}, bloccando prima movimento, input e rotazioni e
        restituendoli su On Blend Out e On Completed.
  \item Compito facoltativo: un nuovo personaggio da zero con lo \textbf{StackOBot},
        Blueprint e Animation Blueprint completi.
\end{itemize}

%
%  Tutte le etichette di questo capitolo hanno il prefisso "ue9:"
%  (ue9:cap:..., ue9:sec:..., ue9:fig:..., ue9:tab:...), e tutti gli
%  stili TikZ hanno il prefisso "ue9" (ue9node, ue9exec, ue9data, ...)
%  per non collidere con chiavi native di TikZ né con quelli delle
%  lezioni precedenti. Il capitolo è autonomo: non contiene \ref verso
%  altri capitoli (i rimandi alle lezioni precedenti sono a parole).
%  Le figure sono ricostruzioni schematiche di quanto mostrato e detto a
%  lezione, non copie delle schermate del docente. Dove la trascrizione
%  è troppo rovinata per ricostruire un dettaglio, il capitolo lo dice
%  invece di inventarlo.
% =====================================================================

\chapter[State Machine, Blend Space e Montage]{State Machine, Blend Space e Montage:
dare vita al nuovo personaggio}
\label{ue9:cap:animazioni}

Questa è la \textbf{nona lezione} del corso breve di \textbf{introduzione a Unreal
Engine}, e riprende esattamente dove si era fermata l'ottava. Alla fine della lezione
precedente il cavaliere importato da \textbf{Mixamo} aveva uno scheletro, le sue
animazioni e un \textbf{Animation Blueprint}, \textbf{ABP\_Knight}, che a ogni frame
leggeva dal personaggio due dati: la velocità sul piano (\textbf{GroundSpeed}) e il fatto
di essere in aria o meno (\textbf{IsFalling}). Mancava tutto il resto: chi decide
\emph{quale} animazione mostrare, e quando. ``Siccome ci siamo interrotti l'altra volta,
andiamo diretti a fare le animazioni: l'altra volta ci siamo interrotti al Blend
Space.''

La lezione costruisce, un pezzo alla volta, un personaggio giocabile completo. La
\textbf{State Machine} organizza le animazioni in \emph{stati} (camminare e correre,
saltare, cadere) e in \emph{transizioni} con delle condizioni; il \textbf{Blend Space
1D} mescola fermo, camminata e corsa in base alla velocità; una \textbf{camera} con il
suo \textbf{Spring Arm}, i grafi di movimento presi in prestito dal modello Third Person
e un nuovo input per lo \textbf{sprint}, costruito con l'\textbf{Enhanced Input},
rendono il cavaliere controllabile; il \textbf{retargeting} porta sul suo scheletro le
animazioni di salto e di caduta del manichino di Unreal. Dopo la pausa arrivano il
\textbf{Blend Space 2D}, che aggiunge alla velocità la \textbf{direzione} e permette di
camminare di lato e all'indietro, e gli \textbf{Animation Montage}, le animazioni
``una tantum'' (un attacco con la spada) che si inseriscono sopra il ciclo di
locomozione.

La lezione è \textbf{online}, su Zoom, come le precedenti, e la classe è ridottissima:
``vediamo un po' quanti siete: siamo in sei''; dopo la pausa ``siamo rimasti in
cinque'', poi ``uno, due, tre e quattro''; quando uno studente crasha (``Filippo è andato
via''), ``siamo rimasti noi tre''; alla fine, ``mi sembra che siamo rimasti in due''.
``Pochi ma buoni.'' Il ritmo è quello della dimostrazione dal vivo, interrotta più volte
dai problemi degli studenti, risolti in condivisione di schermo, e da un editor che
``oggi lagga'' in modo vistoso (``non capisco perché mi sta laggando oggi''). La pausa
arriva dopo la scelta delle animazioni da scaricare (``ci rivediamo alle 18''); si
chiude con un compito facoltativo e un ``ci sentiamo per martedì prossimo''.

Tre avvertenze di lettura, le stesse delle lezioni precedenti. La prima: la
trascrizione è automatica e storpia moltissimi nomi. Unreal è ``Arreal'', ``a Real'',
``ArrayL'', ``Array Alt'' e ``a Rial''; il \textbf{Pawn Owner} è ``Pwn Honor'' e
``Pwnowner''; il \textbf{Blend Space} è ``Blaine Space'', ``Blendspace'' e ``BlizzBase'';
lo stato \textbf{WalkRun} è ``walk run'', ``Warcran'', ``war crunch'' e ``walker run'';
l'\textbf{Idle} è ``Idol'' e ``idol''; \textbf{BP\_Knight} è ``Vip Night'', ``Pippi
Night 1'' e ``BP Night''; \textbf{BS\_Knight} è ``BS Night'' (``che dovrebbe essere un
po' più British''); il manichino \textbf{Manny/Quinn} è ``Many Quincy''; il
\textbf{root motion} è ``rut'' e ``outmotion''; lo \textbf{Strafe} è ``Strife'' e
``stair right''; la \textbf{Mesh} è ``Mesci''; l'\textbf{FInterp To} è ``finter2'';
\textbf{Calculate Direction} è ``calcolo interaction''; il \textbf{Save Cached Pose} è
``save cache, pause''; lo \textbf{StackOBot} è ``StaccoBot'' e ``stacco bot''; lo
\textbf{Yaw} è ``io'', ``yo'' e ``YAW''; l'\textbf{Orient Rotation to Movement} è
``orientation to movement'', ``anche rotacion'' e ``orient rotace tu mente tu''; lo
\textbf{Use Pawn Control Rotation} è ``us control pawn rotation'' e ``you spawn control
rotation''. Dove la ricostruzione è sicura il capitolo usa la forma corretta; dove non
lo è riporta la forma della trascrizione e lo dice. La seconda: all'inizio della
registrazione e durante la pausa il riconoscimento vocale produce lunghi blocchi di
testo senza senso e ripetuto (``grazie a tutti per l'attenzione che ci ha spiegato'',
``qual è la risposta che ci può essere?'', ``la situazione è questa, la migliore dal
2005'', la ``dottrina sul mercato del lavoro'', la ``Corte Suprema di Internet'', il
governo di Dubai, una lunga fila di ``di Sarajevo'', la ``struttura del giardino'', ``la
blue map'', un Vittorio Prodi chiamato in causa senza motivo): sono artefatti, non
parole di nessuno, e il capitolo li ignora. La terza: gran parte della lezione è una
dimostrazione dal vivo nell'editor, costruita a voce (``tirate fuori'', ``promote to
variable'', ``ce lo attacchiamo qui''), e in più punti il docente improvvisa, prova,
sbaglia e corregge, soprattutto nella parte sulla camera e in quella finale sul
Montage. Il capitolo ricostruisce i grafi da ciò che il docente dice mentre li
costruisce, distingue i tentativi falliti dalla soluzione finale, e completa i nomi
esatti dei nodi e dei pin di Unreal Engine 5; dove la trascrizione lascia un buco, lo
dichiara.

\paragraph{Come leggere questo capitolo.} Il capitolo segue l'ordine della lezione. Le
figure ricostruiscono la State Machine della locomozione, il Blend Space 1D sulla
velocità, la camera con lo Spring Arm, la catena dell'Enhanced Input, il grafo dello
sprint, il cerchio della direzione, il Blend Space 2D con le sue otto animazioni, la
catena di pose dell'AnimGraph con lo Slot del Montage e il grafo del tasto che fa
partire l'attacco. Le tabelle riassumono le transizioni della State Machine, i
parametri della camera, i valori e i trigger di una Input Action, la posizione delle
animazioni nel Blend Space 2D, gli errori incontrati durante la lezione e i nodi della
giornata. Qua e là, dei paragrafi di \textbf{approfondimento} collegano la lezione ai
libri di Marcos Romero e Brenden Sewell (\emph{Blueprints Visual Scripting for Unreal
Engine 5} e la seconda edizione dello stesso testo, scritta per Unreal Engine 4), i
testi di riferimento per questa parte del libro, e a qualche nozione tecnica generale su
Unreal Engine 5; sono segnalati come tali e non fanno parte della lezione.

Il percorso del capitolo:
\begin{enumerate}
  \item l'\textbf{inizio della lezione}: i corsi degli studenti e la lezione libera del
        25 (sezione~\ref{ue9:sec:inizio});
  \item il \textbf{ripasso} dell'Animation Blueprint (sezione~\ref{ue9:sec:ripasso});
  \item la \textbf{State Machine} della locomozione (sezione~\ref{ue9:sec:statemachine});
  \item il \textbf{Blend Space 1D}, con una parentesi sui problemi dell'editor (sezioni
        \ref{ue9:sec:blend1d} e~\ref{ue9:sec:layout});
  \item il personaggio in gioco: \textbf{Default Pawn Class}, \textbf{camera} e
        \textbf{Spring Arm}, grafi di movimento e rotazioni (sezioni
        da~\ref{ue9:sec:pawnclass} a~\ref{ue9:sec:rotazioni});
  \item il \textbf{retargeting} del salto e della caduta, e le \textbf{regole di
        transizione} (sezioni \ref{ue9:sec:retarget} e~\ref{ue9:sec:transizioni});
  \item lo \textbf{sprint} con l'\textbf{Enhanced Input} e il tentativo con FInterp To
        (sezioni da~\ref{ue9:sec:enhanced} a~\ref{ue9:sec:finterp});
  \item le \textbf{animazioni da scaricare}, la pausa e i \textbf{feedback} (sezioni
        \ref{ue9:sec:download} e~\ref{ue9:sec:feedback});
  \item il \textbf{Blend Space 2D} e la \textbf{direzione} (sezioni
        \ref{ue9:sec:direction} e~\ref{ue9:sec:blend2d});
  \item gli \textbf{Animation Montage}: l'asset, i Notify, le sezioni, gli Slot e il
        tasto che li fa partire (sezioni da~\ref{ue9:sec:montage}
        a~\ref{ue9:sec:playmontage});
  \item il \textbf{compito} sullo StackOBot e la chiusura
        (sezione~\ref{ue9:sec:chiusura}).
\end{enumerate}

% ---------------------------------------------------------------------
\section{Prima della lezione: i corsi degli studenti e la lezione del 25}
\label{ue9:sec:inizio}

La registrazione si apre con qualche minuto di attesa, in cui il riconoscimento vocale
produce soprattutto rumore. Poi i saluti (``ragazzi, come state? Buonasera, buonasera''),
un controllo rapido (``nessun problema sulla telecamera'': stavolta la webcam del docente
funziona) e la decisione di aspettare ``altri cinque minuti, così vediamo che tutti siate
in carreggiata''. Uno studente, Tommaso, ha dei problemi di collegamento. Il docente
sistema le finestre (``oggi c'è il secondo schermo'') e parte: ``bando alle ciance''.

\paragraph{I corsi di indirizzo.} Prima di cominciare, il docente chiede ai presenti
quale corso seguiranno ``quando inizierete a settembre'': il corso di Unreal è infatti
un'introduzione comune, a cui seguono gli indirizzi specialistici. Le risposte dei sei
presenti:
\begin{itemize}
  \item \textbf{3D art};
  \item \textbf{game design} (due studenti);
  \item \textbf{narrative design};
  \item \textbf{character animation} (``perfetto, character animation'').
\end{itemize}
La domanda serve a tarare la lezione: chi farà character animation, in particolare, è
il destinatario naturale di quasi tutto quello che si vedrà oggi.

\paragraph{La lezione libera del 25.} ``Stiamo praticamente arrivando a una lezione che
di solito noi lasciamo abbastanza \emph{free}, come argomenti'': la lezione del
\textbf{25} sarà ``fondamentalmente la mia ultima lezione di base''. Se non ci sono
richieste particolari, si faranno ``i menu'', cioè ``la \textbf{UI}'': ``che si intende
la UI: menu, bottoni, scritte eccetera''. In alternativa, ``se ci sono degli argomenti
particolari'', da concordare fra gli studenti, ``possiamo tranquillamente affrontarli'',
oppure ``dividere anche la lezione in due: fare un piccolo menu e poi fare anche la parte
in più''. ``Nel caso, sentitevi tra di voi.'' La richiesta concreta arriverà durante la
pausa (sezione~\ref{ue9:sec:feedback}).

% ---------------------------------------------------------------------
\section{Il ripasso: l'Animation Blueprint del cavaliere}
\label{ue9:sec:ripasso}

``Eravamo rimasti al nostro bellissimo \textbf{ABP\_Knight}.'' Il docente ha il
cavaliere (\emph{Knight}, il paladino di Mixamo); ``altri di voi avevano scelto il
proprio personaggio''. Nella preview dell'Animation Blueprint il personaggio è ancora
``il nostro bel character in T-pose'': l'AnimGraph non decide ancora nessuna
animazione.

\paragraph{Initialize e Update.} L'ultima cosa fatta era l'evento \textbf{Event Blueprint
Initialize Animation}: ``quando l'animazione si inizializza'', una volta sola. L'altro
evento, \textbf{Event Blueprint Update Animation}, scatta invece ogni volta che
``si cambia qualsiasi particolare all'interno di un'animazione, la rotazione di un braccio
o un minimo particolare'': ``quello che è fondamentalmente il nostro Tick sulle
animazioni''. ``Quando abbiamo visto Begin Play e Tick: questi praticamente sono i loro
corrispettivi.''

\paragraph{Il Pawn Owner e le due referenze.} Nell'Initialize si prende il \textbf{Pawn
Owner} con \textbf{Try Get Pawn Owner}. ``Ricordiamoci: il Pawn Owner è l'entità che
controlla il character'': il docente distingue il \emph{pawn}, ``quello che lo controlla'',
e il \emph{character} stesso, ``che ha la mesh, lo scheletro e tutto quanto''. Con il
\textbf{Cast To} BP\_Knight si ottiene una referenza al cavaliere, e si salvano in due
variabili il cavaliere e il suo \textbf{Character Movement Component}: ``queste due
variabili ci serviranno per comunicare in modo più diretto con il nostro cavaliere''.

\paragraph{Il Character Movement.} ``Dentro il Blueprint, qualsiasi character lo troverete
qui sotto, Character Movement'': selezionandolo, il pannello Details a destra ``presenta
tutte le particolarità di movimento'', per esempio la \textbf{gravità} (\emph{Gravity
Scale}), la \textbf{Max Acceleration}, la \textbf{Max Step Height} (l'altezza massima di
un gradino che il personaggio sale senza saltare). ``Vogliamo farlo andare più veloce,
meno veloce? Se volessimo farlo sprintare, dovremmo lavorare su questa particolarità''; e
infatti ``dopo faremo anche uno sprintino molto veloce, per differenziare le
animazioni''.

\paragraph{GroundSpeed e IsFalling.} Nell'Update, dopo una \textbf{Sequence}, si prende
``la velocity attuale del nostro character'', a ogni frame, e se ne salva la lunghezza in
\textbf{GroundSpeed}. Allo stesso modo, con il nodo \textbf{Is Falling} del Character
Movement, ``andiamo a vedere se effettivamente il nostro character sta cadendo o meno'':
``questo blocco ci dice letteralmente se sta cadendo, e noi riportiamo con un booleano,
true o false''. ``Di conseguenza, tramite questi dati, andremo a differenziare le
animazioni nella nostra State Machine.''

% ---------------------------------------------------------------------
\section{La State Machine della locomozione}
\label{ue9:sec:statemachine}

\paragraph{Creare la State Machine.} La State Machine ``la possiamo invocare
nell'AnimGraph'': tasto destro nell'AnimGraph (``ricordatevi che c'è il blocco
\textbf{Output Pose}''), \textbf{State Machine}, e la si chiama \textbf{Locomotion}:
``questo è il nome classico per la locomozione, appunto il movimento''. ``Una State
Machine, appunto, è una macchina a stati.'' Con un doppio clic si entra nel suo grafo.

\paragraph{Entry, Conduit, State, Alias.} Dentro c'è un solo nodo, \textbf{Entry}. ``Se
noi trasciniamo, abbiamo una freccia che esce fuori''; lasciandola nel vuoto, il menu non
è ``tutto l'elenco gigantesco'' dei nodi dei Blueprint, ma tre sole voci:
\textbf{Conduit}, \textbf{State} e \textbf{Alias}. ``A noi adesso interessa vedere
quelli che sono gli stati.'' Il primo stato, invece di \emph{Locomotion}, si chiama
\textbf{WalkRun}: ``differenziamo''.

\paragraph{Un'altra inception.} ``Ecco a voi che, ancora sempre, un po' un
\emph{inception}'': con un doppio clic si entra anche nello stato. Dentro, un nodo
\textbf{Output Animation Pose}: ``in questo stato risiede una tipologia di output'',
l'animazione che lo stato produce. Il docente vi collega \emph{Walking 1}. La preview
però non si aggiorna, ``perché non abbiamo collegato la State Machine all'output'':
bisogna tornare nell'AnimGraph e collegare il nodo Locomotion all'Output Pose. Fatto
questo, il cavaliere cammina, ma una volta sola: ``l'animazione non è in loop'', e basta
spuntare il loop nei Details del nodo. ``Ecco a voi che cosa fa la State Machine'':
``dentro Locomotion lui entra, dice ok, inizio, passo e vado direttamente a WalkRun''; in
WalkRun ``c'è una richiesta di un output, quindi di un'animazione che noi gli diamo''.

\paragraph{Gli stati Jump e Fall.} ``Ma la cosa interessante della State Machine'' è che
si possono creare altri stati: \textbf{Jump} e \textbf{Fall}. ``A livello di logica, la
State Machine vuole essere'' il racconto di ``come entrano e perché entrano'' le
animazioni nei loro stadi. Una freccia da WalkRun a Jump ``vuol dire che dallo stadio
WalkRun passerà allo stadio Jump''; una freccia inversa vuol dire che dal Jump ``potrà
passare di nuovo al WalkRun''. Lo schema di oggi ha tre stati
(figura~\ref{ue9:fig:statemachine}): si entra in WalkRun; ``se starà cadendo'' (la
variabile IsFalling diventa vera) si passa a Jump o a Fall; ``e far sì che poi, quando
lui finisca di cadere, ritorni'' a WalkRun. ``Quindi fondamentalmente creiamo un loop di
animazione secondo delle condizioni.'' Le animazioni di salto e di caduta non ci sono
ancora: ``la faremo, la ritargetteremo'' (sezione~\ref{ue9:sec:retarget}).

\paragraph{Le transizioni.} ``Come facciamo a mettere queste condizioni? Tramite questi
nodini qua'': sulle frecce fra due stati compare un'icona tonda, la \textbf{transition
rule}, che ``ci dice il comportamento effettivo che dovrà avere'' il personaggio. ``Quando
è che passo da WalkRun a Jump, a livello di animazione? Quando, ad esempio, stai cadendo.
Quando passo da Fall a WalkRun? Se non sta più cadendo. Questa è la logica generale.''

\begin{figure}[htbp]
\centering
\begin{tikzpicture}[font=\footnotesize,
  ue9state/.style={draw, thick, rounded corners=6pt, fill=orange!12, align=center,
                   minimum width=2.4cm, minimum height=1.0cm},
  ue9entry/.style={draw, thick, fill=gray!15, rounded corners=2pt, align=center,
                   minimum width=1.4cm, minimum height=0.7cm},
  ue9tr/.style={-{Stealth[length=2.6mm]}, thick},
  ue9rule/.style={draw, circle, fill=white, inner sep=1.5pt, font=\scriptsize}]
  \node[ue9entry] (e) at (0,0) {Entry};
  \node[ue9state] (w) at (3.4,0) {\textbf{WalkRun}\\{\scriptsize Blend Space}};
  \node[ue9state] (j) at (8.2,1.5) {\textbf{Jump}\\{\scriptsize MM\_Jump}};
  \node[ue9state] (f) at (8.2,-1.5) {\textbf{Fall}\\{\scriptsize MM\_Fall\_Loop}};
  \draw[ue9tr] (e) -- (w);
  \draw[ue9tr] (w.north east) to[bend left=12] node[ue9rule, pos=0.5]{1} (j.west);
  \draw[ue9tr] (j.south) -- node[ue9rule]{2} (f.north);
  \draw[ue9tr] (w.south east) to[bend left=12] node[ue9rule, pos=0.5]{3} (f.west);
  \draw[ue9tr] (f.south west) to[bend left=30] node[ue9rule, pos=0.5]{4} (w.south);
  \node[align=left, text width=12.5cm, anchor=north west, font=\scriptsize] at
    (-1,-2.9) {\textbf{1} WalkRun $\rightarrow$ Jump: IsFalling \textbf{AND}
    VelocityZ $>$ 0 (sta salendo)\\
    \textbf{2} Jump $\rightarrow$ Fall: \emph{Automatic Rule Based on Sequence Player in
    State} (finito il salto, la caduta)\\
    \textbf{3} WalkRun $\rightarrow$ Fall: IsFalling \textbf{AND} VelocityZ $<$ 0 (cade
    senza aver saltato)\\
    \textbf{4} Fall $\rightarrow$ WalkRun: \textbf{NOT} IsFalling (è tornato a terra)};
\end{tikzpicture}
\caption{La State Machine Locomotion alla fine della lezione: tre stati e quattro
transizioni. Le condizioni con la velocità sulla Z sono quelle aggiunte nella seconda
versione (sezione~\ref{ue9:sec:transizioni}); nella prima, la 1 e la 3 erano soltanto
IsFalling. Schema ricostruito.}
\label{ue9:fig:statemachine}
\end{figure}

\paragraph{Approfondimento: Conduit e Alias.} I tre tipi di nodo proposti dal menu della
State Machine hanno ruoli diversi. Uno \textbf{State} contiene un grafo che produce una
posa (un'animazione, un Blend Space, una combinazione). Un \textbf{Conduit} non produce
pose: è un incrocio, uno snodo da cui partono più transizioni, utile quando molti stati
condividono la stessa condizione d'ingresso o d'uscita e non si vuole ripeterla su ogni
freccia. Un \textbf{State Alias} è un ``alias'' di uno o più stati: le transizioni che
partono dall'alias valgono come se partissero da ciascuno degli stati che rappresenta
(per esempio, ``da qualunque stato, se il personaggio muore, vai a Death''), e
risparmiano un groviglio di frecce. Romero e Sewell (capitolo 17) costruiscono la loro
prima State Machine, \emph{Char States}, con due soli stati, \emph{Idle} e
\emph{Moving}, e due regole sulla velocità (``Speed maggiore di 5'' e ``minore di 5''); la
scelta della lezione, un solo stato WalkRun che contiene un Blend Space, sposta la
distinzione fra fermo e in movimento dentro il Blend Space, e lascia alla State Machine
le sole differenze ``qualitative'' (a terra, in salto, in caduta).

% ---------------------------------------------------------------------
\section{Il Blend Space 1D}
\label{ue9:sec:blend1d}

``Visto che WalkRun non è solo un'animazione, è un insieme'': sono tre animazioni,
\textbf{Idle}, \textbf{Walk} e quella che il docente ha chiamato \textbf{Slow Run}.
``Dobbiamo mixarle insieme. Come si fa? Perché non possiamo mettere un'animazione solo
così: dobbiamo far sì che sia dipendente dalla nostra velocità.'' Lo strumento è il
\textbf{Blend Space}.

\paragraph{Creare il Blend Space 1D.} ``Fate solamente il Blend Space 1D.'' Invece di
cercarlo nei menu, si può scrivere nella barra di ricerca del menu del tasto destro
``blend space'', e scegliere \textbf{Blend Space 1D}: ``già vi dice il percorso, che sta
in Animation, Legacy''. L'icona ha un colore ``rosa pelle''. Aperto, si troverebbe il
personaggio in T-pose; ma molti l'avevano già creato la volta precedente (``sì, mi sa che
già l'avevamo fatto'').

\paragraph{Che cos'è un Blend Space.} ``Un Blend Space è l'azione di \emph{blendare}
qualche cosa'', di ``interpolare o comunque mixare le animazioni'': un asset che ``ci dà
la possibilità di poter traslare da un'animazione all'altra''. Si vede subito: dal
pannello \textbf{Asset Browser} dell'editor (``se non lo trovate, ripeto: Window, Asset
Browser'') si trascina un'animazione sul grafico; tenendo premuto \textbf{Shift} la si
aggancia ``su dei punti precisi'' della griglia. Con un'animazione sola, il personaggio
fa solo quella; aggiungendone una seconda, il grafico ``crea una linea di passaggio tra
un'animazione e l'altra''.

\paragraph{Il parametro, come l'alfa.} ``Ma importante: qual è questo dato?'' Il docente
lo collega all'\textbf{alfa} di un Lerp visto in una lezione precedente: ``quel puntatore
che punta, dipendente da uno spazio tra due punti, dove sta al momento''. Qui il valore
dirà ``ehi, guarda che sto camminando, sto andando a velocità 300, oppure sto correndo, a
velocità 600, oppure sono fermo''. Il Blend Space \textbf{1D} ``ci dà solamente la
possibilità di poter usare un asse'': nei Details, \textbf{Horizontal Axis} (la sezione
potrebbe essere chiusa: ``basta che la aprite''). L'asse si chiama \textbf{GroundSpeed},
``come quella che abbiamo trovato all'interno dell'Animation Blueprint''; poi chiede un
\textbf{Minimum Axis Value} e un \textbf{Maximum Axis Value}.

\paragraph{Le velocità sono relative.} ``Importante, ragazzi, quando si va a vedere le
animazioni: dovete sempre farlo relativo, se si parla di velocità'', relativo cioè alla
velocità reale del personaggio. Dove si legge? Nel Blueprint del personaggio,
BP\_Knight, selezionando \textbf{Character Movement}: ``potete andare a vedere tutti i
dati che vi servono''. Quello che serve è la \textbf{Max Walk Speed}, in alto, nella
sezione \emph{Character Movement: Walking}: ``è 600 centimetri al secondo''. ``Unreal
ragiona per centimetri'': il personaggio va a 6 metri al secondo.

\paragraph{Dividere la velocità a metà.} ``Visto che noi abbiamo due stadi, uno di walk e
uno di run'', si vuole che ``la walk sia un movimento normale'': di default il personaggio
cammina, e corre solo se il giocatore lo chiede con un tasto (sezione
\ref{ue9:sec:enhanced}). ``Quindi la dividiamo a metà'': \textbf{Max Walk Speed a 300}.
Nel Blend Space l'asse va da \textbf{0} a \textbf{600}, ``in modo tale che a metà la
nostra walk sia a 300'' (figura~\ref{ue9:fig:blend1d}). A uno studente l'asse arriva solo
a 100: ``ground speed, devi metterlo a 600, lo trovi qua''. E un consiglio: dopo aver
cambiato il massimo, ``sposta i pallini, perché poi ti switcha tutto''; il valore di
ciascun campione si può anche scrivere a mano nei Details.

\begin{figure}[htbp]
\centering
\begin{tikzpicture}[font=\footnotesize,
  ue9pt/.style={draw=black, fill=white, thick, diamond, inner sep=2.2pt}]
  \draw[-{Stealth[length=2.5mm]}, thick] (0,0) -- (11.2,0) node[right]{GroundSpeed};
  \foreach \x/\v in {0/0, 2.5/150, 5/300, 7.5/450, 10/600}
    \draw (\x,0.1) -- (\x,-0.1) node[below, font=\scriptsize]{\v};
  \node[ue9pt] (i) at (0,0) {}; \node[ue9pt] (w) at (5,0) {}; \node[ue9pt] (r) at
    (10,0) {};
  \node[above=6pt of i, font=\scriptsize\bfseries] {Idle};
  \node[above=6pt of w, font=\scriptsize\bfseries, align=center] {Walk\\
    {\normalfont Max Walk Speed di default}};
  \node[above=6pt of r, font=\scriptsize\bfseries, align=center] {Slow Run\\
    {\normalfont con lo sprint}};
  \fill[red!60!black] (3.33,0) circle (0.12);
  \draw[red!60!black, thick, ->] (3.33,-1.2) -- (3.33,-0.25);
  \node[red!50!black, font=\scriptsize, align=center] at (3.33,-1.55) {200: un misto
    fra Idle e Walk, ``un obbrobrio''};
\end{tikzpicture}
\caption{Il Blend Space 1D BS\_Knight: un solo asse, la GroundSpeed, da 0 a 600, con
Idle a 0, Walk a 300 (la Max Walk Speed impostata nel Character Movement) e la corsa a
600 (la velocità dello sprint). Un valore intermedio, come 200, mescola due animazioni
nate per velocità diverse. Schema ricostruito.}
\label{ue9:fig:blend1d}
\end{figure}

\paragraph{Il Blendspace Player.} Tornati in ABP\_Knight (``normale, ragazzi, che vi dia
dei warning, perché ancora non abbiamo le condizioni'': ``lui dice, ma io non ho
condizioni, quindi faccio tutto in una botta''), si entra nello stato WalkRun. Nell'Asset
Browser ``è apparso BS\_Knight''; trascinato nel grafo, ``ci presenta un \textbf{Blend
Space Player}'', diverso dal nodo di un'animazione normale: ha un \textbf{valore in
entrata}, la GroundSpeed. A quel pin si collega un \textbf{Get} della variabile
GroundSpeed, calcolata nell'Event Graph (la lunghezza della velocity), e l'uscita del
Blend Space Player si collega al \textbf{Result} dell'Output Animation Pose, ``in modo
tale che vada in output''.

\paragraph{Get con Ctrl, Set con Alt.} Un promemoria: trascinando una variabile nel
grafo, ``tenete premuto \textbf{Ctrl} e tirate fuori direttamente un Get''; ``se invece
volete fare un Set, tenete premuto \textbf{Alt}''. Nel grafo dello stato l'Alt non
funziona: ``non si può settare all'interno'' (gli stati dell'AnimGraph leggono le
variabili, non le scrivono), ``però nei Blueprint, di base, se premete Alt vi setta la
variabile''.

\paragraph{La prova nella preview.} Compilato, ``il nostro character nella preview adesso
sta in idle, perché giustamente ha velocità 0''. Cambiando il valore di GroundSpeed nella
preview: a \textbf{300} cammina; a \textbf{200} ``sta interpolando tra l'animazione di
idle e l'animazione di walk'': ``mettere un valore a 200 non è buona cosa'', perché
``deve essere relativo al vostro movimento, se no vengono degli obbrobri''; a
\textbf{600} ``il nostro carissimo corre'' (``guarda come è eroico''). ``Quindi adesso
andiamo a cambiare le animazioni tramite un valore.''

\paragraph{Allineare i nodi con Q.} ``Se volete fare un po' di ordine'': selezionati due
o più nodi, il tasto \textbf{Q} li allinea (\emph{Straighten Connections}). ``Tutto bello
ordinato.''

\paragraph{Una prova con un solo stato.} Prima di andare avanti, un test in gioco:
nell'AnimGraph si tolgono le frecce verso Jump e Fall (selezionate e cancellate, o
staccate con \textbf{Alt}+clic sul pin), e si lascia ``solamente WalkRun, che vada in
entrata''.

\paragraph{Approfondimento: il Blend Space nel libro.} Romero e Sewell (capitolo 17)
definiscono il Blend Space come ``un tipo di asset che permette di mescolare animazioni
in base ai valori di uno o due parametri''; nel loro esempio, tratto dal modello Third
Person, Idle è a 0, Walk a 93,75 e Run a 375, con una Max Walk Speed di 375. Il principio
è lo stesso della lezione: i campioni vanno posti dove la velocità reale del personaggio
corrisponde alla velocità ``disegnata'' nell'animazione. Se un'animazione di camminata è
stata animata per una velocità di 150 centimetri al secondo e il personaggio si muove a
300, i piedi ``scivolano'' sul pavimento (\emph{foot sliding}); è il motivo pratico
dietro il ``deve essere relativo'' del docente. Gli autori ricordano anche che il valore
del parametro si può cambiare nell'\emph{Anim Preview Editor}, in basso a destra
nell'editor dell'Animation Blueprint, senza andare in Play: è ciò che il docente fa per
mostrare i valori 200, 300 e 600.

% ---------------------------------------------------------------------
\section{Una parentesi: ``questo asset non ha schede ancorate''}
\label{ue9:sec:layout}

Mentre si sistema la Max Walk Speed, uno studente interviene: ``io ho un problema col
Blueprint: quando apro la scheda Blueprint mi dice: questo asset non ha ancora schede
ancorate''. Il docente ha l'editor in inglese (``di solito gli errori me li ricordo in
inglese'') e chiede di vedere: lo studente condivide prima soltanto una scheda del
browser, poi lo schermo intero.

\paragraph{Il messaggio.} Aprendo BP\_Knight, l'editor del Blueprint è vuoto, senza
pannelli. Il testo completo del messaggio, letto ad alta voce, spiega che l'editor
dell'asset è stato aperto in un sistema di finestre fluttuanti e nascoste, ``non mostra
le schede ancorate nell'interfaccia principale'', e che ``questo comportamento può essere
facilmente ripristinato'' ripristinando il layout dell'editor o modificando le preferenze
perché gli asset si aprano sempre ancorati alla finestra principale (in inglese,
\emph{Editor Preferences $\rightarrow$ General $\rightarrow$ Appearance $\rightarrow$
Asset Editor Open Location}). Inoltre, sopra il grafo compare una nota: il Blueprint è
\emph{data-only}, e c'è un link \textbf{Open Full Blueprint Editor}.

\paragraph{La soluzione.} ``Mi sa che è meglio ripristinare il layout di default'':
dalla barra in alto, \textbf{Window} (in italiano \emph{Finestra}) $\rightarrow$
\textbf{Load Layout} (\emph{Carica il layout}) $\rightarrow$ \textbf{Default Editor
Layout} (\emph{Editor predefinito}). Poi, sulla nota, ``apri l'editor Blueprint
completo''. ``Perfetto, eccolo qua.''

\paragraph{La lingua dell'editor.} ``Comunque, per mettere l'editor in inglese'' senza
riavviare: si cerca \emph{Language} nelle preferenze. Il docente prova prima in
\textbf{Editor Preferences} (sotto \emph{Edit}), scrivendo ``lingua'' e poi ``language''
(``non gli piace che mettiamo due cose insieme''), poi nei \textbf{Project Settings};
alla fine la voce \emph{Language} viene trovata con la ricerca (``ah, ecco perché'').
Dalla trascrizione non si capisce con certezza in quale delle due finestre. Lo studente
ringrazia (``mi hai fatto sentire molto ignorante''; ``verrai buttato fuori dalla
classe, tranquillo''), e il docente annuncia di volerlo usare in inglese anche lui ``per
la prossima volta''.

\paragraph{Approfondimento: dov'è la lingua.} In Unreal Engine 5 la lingua
dell'interfaccia è in \textbf{Edit $\rightarrow$ Editor Preferences $\rightarrow$ General
$\rightarrow$ Region \& Language $\rightarrow$ Editor Language}; nei Project Settings c'è
invece la lingua del \emph{gioco} (la localizzazione), che è un'altra cosa. Lavorare
con l'editor in inglese è una scelta pratica: la documentazione, i tutorial, i libri e i
forum usano i nomi inglesi dei nodi e dei pannelli, e la ricerca dei nodi nel menu del
tasto destro accetta i nomi inglesi anche quando l'interfaccia è tradotta.

% ---------------------------------------------------------------------
\section{Il cavaliere in gioco: la Default Pawn Class}
\label{ue9:sec:pawnclass}

``Adesso vi faccio vedere come impostiamo il player nuovo.'' Si aprono i \textbf{World
Settings} del livello (``qualsiasi cosa che non trovate, c'è sempre Window, e troverete
World Settings''); fra le voci della Game Mode c'è \textbf{Default Pawn Class}.
``Ricordiamo che il Pawn è appunto l'entità che controlla il player.'' Nel menu a
tendina c'è BP\_Knight, il personaggio finito che il docente aveva mostrato alla fine
della lezione precedente, e c'è BP\_Knight1, ``quello nuovo'', che si può anche cercare
scrivendone il nome. Si sceglie BP\_Knight1 e si va in Play.

\paragraph{Il primo problema: nessuna camera.} ``Ecco già un problemino.'' ``Fate
attenzione, perché i character non vengono con la camera'': un Blueprint creato dalla
classe Character ha la capsula, la freccia, la mesh e il Character Movement, ma non una
camera. In Play, la visuale finisce in un punto qualsiasi del corpo (``questo
meraviglioso\dots'', la frase si perde). Si esce dal Play e si apre BP\_Knight1: ``come
potete vedere, non c'è una camera''.

% ---------------------------------------------------------------------
\section{La camera: Camera e Spring Arm}
\label{ue9:sec:camera}

``La camera è impostata su due cose particolari, ragazzi, molto importanti'': la
\textbf{camera} stessa e il suo \textbf{braccio}. ``Immaginatelo come lo stand'': il
braccio dice ``quanto è lontana'' la camera, ``quanto offset ha''.

\paragraph{Aggiungere i due componenti.} Si aggiunge un componente \textbf{Camera}
(``camera normale''): ``eccola qua, che si è già attaccata'' al personaggio, ma in una
posizione che non va bene (``al suo lato B''). Poi uno \textbf{Spring Arm} (il docente
lo chiama prima ``arm length'', poi si corregge: ``perdonate, l'arm length è un
parametro; lo Spring Arm''). La gerarchia è importante: \textbf{lo Spring Arm deve avere
la camera dentro di sé}, cioè la Camera va trascinata sopra lo Spring Arm e diventa sua
figlia. Selezionando lo Spring Arm, nel viewport compare ``un piccolo tracciato rosso'':
``la lunghezza del nostro braccio'', il braccio della camera
(figura~\ref{ue9:fig:springarm}).

\paragraph{I 3C.} ``Ricordatevi, ragazzi, molto importante, perché la camera è
importantissima: ci sono delle figure che si chiamano \textbf{3C}, \emph{Camera, Control
e Character}, che lavorano solo su questo.'' Da qui una regola: ``se volete allontanare la
camera o cambiarle l'offset, \textbf{non cambiate la location della camera}. Non lo fate.
Controllatelo sempre da qui, dallo Spring Arm.''

\paragraph{Target Arm Length.} Il parametro \textbf{Target Arm Length} ``ci permette di
aumentare o diminuire la distanza della camera''. Valore scelto: \textbf{300}.

\paragraph{Target Offset e Socket Offset.} Ci sono due offset, e il docente ha una
preferenza netta. Il \textbf{Target Offset} ``usatelo solamente sulla Z, di poco,
veramente di poco, massimo 20''; ``e fate tutto con il \textbf{Socket Offset}''. Il
motivo: ``se voi cambiate l'offset del target, la camera vi può più facilmente
compenetrare con le cose''. Con il Socket Offset, invece, ``vedrete che la camera va sul
muro e torna'': ``va più vicino'', cioè ``riduce il suo arm length, fondamentalmente''.
Dopo la prima prova in Play (``lo testiamo, ragazzi?''), la camera è un po' alta: il
Socket Offset sulla Z viene abbassato a \textbf{40}.

\begin{figure}[htbp]
\centering
\begin{tikzpicture}[font=\footnotesize]
  % capsula
  \draw[thick, fill=blue!8, rounded corners=14pt] (0,0) rectangle (1.1,2.6);
  \node[align=center, font=\scriptsize] at (0.55,1.1) {Cap-\\sule};
  \node[font=\scriptsize, below] at (0.55,-0.2) {BP\_Knight1};
  \fill (0.55,2.0) circle (0.06);
  \node[font=\scriptsize, anchor=north west] at (0.62,1.95) {pivot};
  % braccio
  \draw[very thick, red!70!black] (0.55,2.0) -- (-4.2,2.0);
  \node[red!60!black, font=\scriptsize, above] at (-1.8,2.0) {Target Arm Length
    = 300};
  % socket offset
  \draw[thick, dashed, -{Stealth[length=2mm]}, green!40!black] (-4.2,2.0) --
    (-4.2,2.8);
  \node[green!40!black, font=\scriptsize, anchor=east, align=right] at (-4.9,2.3)
    {Socket Offset\\Z = 40};
  % camera
  \draw[thick, fill=gray!30] (-4.75,2.6) rectangle (-3.95,3.05);
  \draw[thick, fill=gray!30] (-3.95,2.7) -- (-3.6,2.55) -- (-3.6,3.1) -- (-3.95,2.95);
  \node[font=\scriptsize, above] at (-4.3,3.1) {Camera (figlia)};
  % target offset
  \draw[thick, dashed, -{Stealth[length=2mm]}, orange!70!black] (0.55,2.0) --
    (0.55,2.35);
  \node[orange!60!black, font=\scriptsize, anchor=south, align=center] at (0.55,2.4)
    {Target Offset: solo Z,\\al massimo $\approx$ 20};
  % muro
  \fill[gray!40] (-2.6,-0.2) rectangle (-2.35,1.4);
  \draw[thick] (-2.6,-0.2) rectangle (-2.35,1.4);
  \node[font=\scriptsize, align=center] at (-2.47,-0.6) {muro};
  \draw (-6,-0.2) -- (3.5,-0.2);
  % gerarchia
  \node[draw, thick, rounded corners=3pt, fill=gray!8, align=left, font=\scriptsize,
    anchor=north west] at (3.3,2.2) {\textbf{Components}\\
    Capsule Component\\
    \quad Arrow\\
    \quad Mesh\\
    \quad \textbf{Spring Arm}\\
    \qquad \textbf{Camera}\\
    Character Movement};
\end{tikzpicture}
\caption{La camera del cavaliere. La Camera è figlia dello Spring Arm, e la distanza e
l'altezza si regolano sul braccio, mai spostando la camera. Se fra il personaggio e la
camera si mette un ostacolo (il muro), lo Spring Arm accorcia il braccio e la camera
``torna'' verso il personaggio; con un Target Offset grande, invece, è più facile che la
camera entri negli oggetti. Schema ricostruito, valori della lezione.}
\label{ue9:fig:springarm}
\end{figure}

\paragraph{Il Field of View.} Più avanti, mentre aspetta uno studente, il docente mostra
che ``nella camera è presente anche un po' di dati con cui potete giocarci'', per esempio
il \textbf{Field of View}. Lo abbassa a \textbf{80}: ``sappiate che tanti giochi horror
usano il field of view a 80. Questo è un piccolo tip.''

\begin{table}[htbp]
\centering
\small
\begin{tabular}{@{}p{3.6cm}p{2.0cm}p{8.4cm}@{}}
\toprule
\textbf{Parametro} & \textbf{Valore} & \textbf{Effetto e consiglio del docente} \\
\midrule
Spring Arm $\rightarrow$ Target Arm Length & 300 & distanza della camera dal
personaggio \\
Spring Arm $\rightarrow$ Socket Offset & Z = 40 & spostamento della camera alla fine del
braccio; da preferire \\
Spring Arm $\rightarrow$ Target Offset & Z $\le$ 20 & spostamento dell'origine del
braccio; poco e solo sulla Z, perché fa compenetrare la camera \\
Spring Arm $\rightarrow$ Use Pawn Control Rotation & attivo & il braccio ruota con il
mouse (il controller); va sullo Spring Arm, non sulla Camera \\
Camera $\rightarrow$ Field of View & 80 & angolo di visuale; 80 è tipico dei giochi
horror \\
Camera $\rightarrow$ Location & non toccare & la posizione si regola sempre dallo Spring
Arm \\
\bottomrule
\end{tabular}
\caption{I parametri della camera toccati durante la lezione.}
\label{ue9:tab:camera}
\end{table}

\paragraph{Approfondimento: perché lo Spring Arm si accorcia.} Lo \textbf{Spring Arm
Component} fa, a ogni frame, una traccia (un \emph{sweep} di una piccola sfera, di
raggio \emph{Probe Size}) dal suo punto d'origine verso la posizione desiderata della
camera; se incontra un ostacolo nel canale \emph{Probe Channel} (di default
\emph{Camera}), accorcia il braccio fino al punto di impatto, e lo riallunga quando
l'ostacolo sparisce (la proprietà è \emph{Do Collision Test}). Il \emph{Target Offset}
sposta l'\emph{origine} di quella traccia, nello spazio del mondo: spostarla molto può
portare l'origine dentro un muro o dietro un oggetto, e allora il test non ``vede'' più
l'ostacolo; il \emph{Socket Offset} sposta invece la \emph{fine} del braccio, nello
spazio del braccio stesso, e la traccia resta coerente. È la ragione tecnica del
consiglio del docente. Lo Spring Arm ha anche le opzioni \emph{Enable Camera Lag} ed
\emph{Enable Camera Rotation Lag}, che fanno ``inseguire'' la camera con un leggero
ritardo e rendono più morbidi i movimenti: il docente le sfiora più tardi, nella sezione
del Montage. Quanto ai 3C, nell'industria la sigla indica la disciplina che cura come il
giocatore \emph{vede} (camera), \emph{comanda} (controlli) e \emph{abita} (personaggio)
il gioco: tre cose che si progettano insieme, perché ognuna cambia la percezione delle
altre.

% ---------------------------------------------------------------------
\section{I grafi di movimento, presi in prestito}
\label{ue9:sec:movimento}

In Play, la camera c'è, ma ``non ci possiamo muovere'': ``perché non abbiamo né input né
le funzioni di base di movimento''. ``Questo, ragazzi, non andiamo a farlo direttamente a
mano, ma andiamo a fare quella che è una bellissima'' scorciatoia (la trascrizione dice
``mandragata''): si apre il \textbf{BP\_ThirdPersonCharacter} classico, ``quello che
abbiamo usato'' finora, che contiene ``questi tre principali blocchi che Unreal fa di
base'': il controllo della \textbf{camera} (\emph{Look}), il controllo del
\textbf{movimento} (con \emph{Add Movement Input}) e il \textbf{salto} (\emph{Jump}).

\paragraph{Copiare anche l'Input Mapping.} ``Mi raccomando, prendete anche questo qua
sopra, che è l'\textbf{input mapping}'', cioè il grafo sull'Event Begin Play che aggiunge
l'Input Mapping Context (``andremo a spiegare che cosa si intende''). Tutto con
\textbf{Ctrl+C}, e poi \textbf{Ctrl+V} nell'Event Graph di BP\_Knight1: ``perché queste
sono le funzioni basiche di Unreal''.

\paragraph{IA\_Move.} Un'occhiata all'evento \textbf{IA\_Move} con l'Add Movement Input,
``proprio per fare una lezione sull'input'': ``questo evento dice che è legato a un
tasto, o a dei tasti, che ci riporta due valori, X e Y'', che ``di conseguenza diventano
dei vettori che vanno a creare il movimento sui due assi''. Uno è avanti e indietro
(\emph{forward}, \emph{backward}), l'altro è destra e sinistra.

\paragraph{La camera che ruota male.} Compilato, ``adesso la nostra camera si muove'', ma
in modo strano: si muove ``solo così'', su un asse. ``L'avevo visto un'altra volta'':
``devo usare il control rotation''. La spunta è \textbf{Use Pawn Control Rotation}:
``se lo attaccate, la camera vi va in \emph{free mode}'', ``nel senso che usa la
rotazione del controller'', sia in beccheggio sia in imbardata.

% ---------------------------------------------------------------------
\section{Tre spunte sulla rotazione}
\label{ue9:sec:rotazioni}

\paragraph{Use Pawn Control Rotation, sullo Spring Arm.} Uno studente ha la camera che
``quando vado in Play posso abbassarla e guardare solo il terreno, non vedo proprio
neanche più il personaggio''. Il docente pensa prima a una compenetrazione, poi si fa
condividere lo schermo: ``che camera hai messo? Camera normale, giusto?''. Guardando i
componenti, il problema salta fuori: ``non Use Pawn Control Rotation sulla camera: Use
Pawn Control Rotation sullo \textbf{Spring Arm}''. ``Quindi levalo da lì e mettilo sullo
Spring Arm.'' ``Anche l'altro?'' ``No, no, no: solo sullo Spring Arm.'' La spunta va su
\emph{un} componente solo, il braccio: è il braccio a ruotare intorno al personaggio, e
la camera, che ne è figlia, lo segue.

\paragraph{Orient Rotation to Movement.} In Play, il cavaliere cammina (alla velocità di
300), ma andando a destra ``lui cammina dritto, fondamentalmente'': si sposta di lato
continuando a guardare avanti. ``Questo lo possiamo risolvere con'' \textbf{Orient
Rotation to Movement}, sempre nel Character Movement: ``dal nome, orienta la rotazione al
movimento''. Accanto c'è la \textbf{Rotation Rate}, la velocità con cui il personaggio si
gira (``non è proprio precisissima, però ci piace'').

\paragraph{Use Controller Rotation Yaw.} Il personaggio però continua a non girarsi bene.
``Perché mi ricordo che qua c'era una particolarità che non ricordo'': il docente scorre
le categorie del Blueprint (\emph{Simulate}, \emph{Rotation}, \emph{Character},
\emph{Orientation}, \emph{Physics}, \emph{Controller}\dots) e la trova. ``Per fissare il
nostro character che non si girava, basta che eliminiamo il tick'': selezionato il Pawn
(la radice del Blueprint, \emph{BP\_Knight1 (Self)}), si cerca \emph{yaw} e si toglie
\textbf{Use Controller Rotation Yaw}, ``in modo tale che il character non usi il suo
yaw''. Senza questa spunta, il personaggio non ruota più insieme alla camera: ``andrà
dove siamo direzionati'', ``se io vado a destra, lui si gira completamente verso
destra''. A uno studente non basta: aveva tolto lo Yaw ma non attivato Orient Rotation to
Movement. ``Ok, non ce l'avevo.'' Con tutte e due, ``abbiamo il nostro bel character
fatto per benissimo''.

\paragraph{Approfondimento: chi comanda la rotazione.} Le tre spunte della lezione sono
le tre ``mani'' che possono ruotare un personaggio in terza persona, e vanno combinate
con criterio. \textbf{Use Pawn Control Rotation} (sullo Spring Arm) fa seguire alla
\emph{camera} la rotazione del controller, cioè del mouse o della levetta destra.
\textbf{Use Controller Rotation Pitch/Yaw/Roll} (sul Pawn) fa seguire quella stessa
rotazione al \emph{corpo} del personaggio. \textbf{Orient Rotation to Movement} (nel
Character Movement) fa girare il corpo verso la direzione in cui si muove. La
combinazione ``da avventura'' della prima parte della lezione è: camera libera (Use Pawn
Control Rotation sul braccio), corpo indipendente dalla camera (Use Controller Rotation
Yaw spento), corpo orientato al movimento (Orient Rotation to Movement acceso): è quella
del modello Third Person di Unreal. La combinazione ``da combattimento'' della seconda
parte, con il Blend Space 2D (sezione~\ref{ue9:sec:direction}), è opposta: il corpo
guarda sempre dove guarda la camera (Use Controller Rotation Yaw acceso), e le
animazioni di strafe mostrano lo spostamento laterale. Se sono accese insieme Use
Controller Rotation Yaw e Orient Rotation to Movement, le due ``mani'' si contendono il
personaggio; per questo, nel Character Movement, esiste anche \emph{Use Controller
Desired Rotation}, che fa ruotare il corpo verso la direzione della camera in modo
graduale, con la Rotation Rate (forse il ``desired'' che il docente prova a un certo
punto, senza esito chiaro nella trascrizione).

% ---------------------------------------------------------------------
\section{Il retargeting: il salto del manichino sul cavaliere}
\label{ue9:sec:retarget}

``Adesso, però, guardate una cosa: se ad esempio provo a saltare, non ho nessuna
animazione.'' Il salto funziona (lo fa il Character Movement), ma il cavaliere sale e
scende nella posa della camminata. Le animazioni di salto si potrebbero prendere da
Mixamo, ``anche se potremmo, volendo''; invece il docente usa ``l'opzione di Unreal che
si chiama \textbf{retargeting}''. ``Diciamo che il retargeting, sappiatelo, può uscire
bene, ma può uscire anche abbastanza male: è dipendente dall'animazione. Però andiamo a
farlo comunque, per vedere come si fa.''

\paragraph{I passaggi.}
\begin{enumerate}
  \item Nel Content Browser si cerca \emph{Jump}: si trova l'animazione del salto del
        manichino base di Unreal, \textbf{MM\_Jump}.
  \item Tasto destro sull'animazione $\rightarrow$ \textbf{Retarget Animations}.
  \item La finestra mostra lo scheletro \textbf{sorgente} (\emph{source}), quello da cui
        viene l'animazione: il manichino di Unreal, \textbf{Manny} o \textbf{Quinn}
        (``Many Quincy''). Come \textbf{target} si sceglie lo scheletro del cavaliere, il
        paladino di Mixamo: ``eccolo qua, li mette a confronto''.
  \item Si seleziona \textbf{MM\_Jump} nella lista delle animazioni: nella preview il
        paladino salta. ``Fortunatamente Mixamo usa degli scheletri che sono abbastanza
        compatibili con quello di Unreal'', e il risultato ``è venuto abbastanza bene'';
        ``potete girarci tranquillamente intorno con la viewport''.
  \item Due pulsanti di esportazione. \textbf{Export Retargeted Assets} ``non ci serve,
        ragazzi'': ``vi crea un asset retargettato'', una copia completa degli asset sullo
        scheletro di destinazione. Basta \textbf{Export Animations}: ``esportare
        l'animazione così, come nuova''. ``Ricordate che c'è il source, e vi esporta il
        target.''
  \item Come cartella si sceglie quella del personaggio (\emph{Characters $\rightarrow$
        Knight1 $\rightarrow$ Animation}) e si esporta.
\end{enumerate}
Lo stesso procedimento si ripete per la caduta: \textbf{MM\_Fall\_Loop}, un'animazione
``in loop'' di caduta, retargettata sul paladino ed esportata nella stessa cartella.

\paragraph{Il root motion ``rotto''.} Il risultato del salto ``non è fantastico'':
``il problema di Mixamo è che ha il \textbf{root motion} un po' rotto''. Il docente
mostra il \emph{root}, ``un punto dello scheletro'', l'osso alla base della gerarchia, e
osserva che lo spostamento non è quello atteso; ``potremmo vederlo'', ``ma per vedere come
funziona il retarget va benissimo''. Si potrebbe sistemare in seguito.

\paragraph{Approfondimento: il retargeting in Unreal Engine 5.} Il retargeting trasferisce
un'animazione da uno scheletro a un altro con proporzioni e nomi delle ossa diversi. In
Unreal Engine 5 si basa su due tipi di asset: l'\textbf{IK Rig}, che descrive uno
scheletro raggruppando le ossa in \emph{catene} (spina dorsale, braccia, gambe, testa)
con un osso di \emph{retarget root} (di solito il bacino), e l'\textbf{IK Retargeter},
che mette in corrispondenza le catene di due IK Rig e corregge la differenza fra le pose
di riferimento (la A-pose del manichino e la T-pose di Mixamo). Nelle versioni più recenti il
comando \emph{Retarget Animations} del menu contestuale apre una finestra che crea al
volo questi asset quando gli scheletri sono riconosciuti (come quelli di Mixamo e del
manichino), ed è ciò che si vede a lezione; per risultati migliori, o con scheletri
insoliti, si creano e si regolano a mano. Il problema del root di Mixamo è noto: gli
scheletri di Mixamo non hanno un osso di root a terra, e lo spostamento del personaggio
è registrato sull'osso dei fianchi (\emph{Hips}); un'animazione retargettata può quindi
avere il movimento nel posto ``sbagliato'' per il sistema di root motion di Unreal.

% ---------------------------------------------------------------------
\section{Le regole di transizione}
\label{ue9:sec:transizioni}

Con le due animazioni nuove si aggiorna l'Animation Blueprint: ``dentro il nostro Jump
ci mettiamo l'animazione di jump'' (``niente di particolare, ragazzi''), dentro Fall
l'animazione di caduta. Si ricollegano le frecce. ``Se cammina e corre, può saltare? Sì,
certo. Se salta, dovrà per forza almeno loopare in caduta'', e ``dal Fall dovrà ritornare
al nostro WalkRun''.

\paragraph{Prima versione.} Nelle icone delle transizioni c'è un grafo con un nodo
\textbf{Result}, che ``ci dà un bool'': ``ci dice quando è che posso entrare in questa
animazione''.
\begin{itemize}
  \item \textbf{WalkRun $\rightarrow$ Jump}: un \emph{Get IsFalling} collegato al
        Result. ``Potremmo farlo più preciso dopo''; per esempio, ``possiamo anche dire:
        ha iniziato il jump''.
  \item \textbf{Jump $\rightarrow$ Fall}: qui il grafo vuoto dà errore, ``perché noi
        vogliamo far sì che questo processo tra il Jump e il Fall sia automatico''.
        Selezionata l'icona della transizione, nei Details si spunta \textbf{Automatic
        Rule Based on Sequence Player in State}: ``quando inizia quell'animazione, a
        seguito ci sarà quest'altra''; in altre parole, la transizione scatta da sola
        quando l'animazione dello stato Jump finisce. ``Lo mettiamo così e non ci darà
        problemi.''
  \item \textbf{Fall $\rightarrow$ WalkRun}: ``vuol dire che non sta cadendo'': Get
        IsFalling, poi un \textbf{NOT Boolean}. ``Ricordo che il NOT Boolean è qualcosa
        che ci dà il suo contrario'': deve essere false perché la caduta finisca.
\end{itemize}
Uno studente chiede cosa ci sia fra Jump e Fall: ``se tu clicchi su questo pallino, ci
dice, in Details\dots\ vogliamo che questa animazione parta direttamente se effettivamente
il jump è finito'': \emph{Automatic Rule Based on Sequence Player in State}, ``questo
tick, lo metti vero''.

\paragraph{La prova, e un difetto.} Si prova da un punto alto, ``la nostra piattaforma''.
``Non si nota tantissimo, però effettivamente sta facendo giustamente l'azione di
saltare.'' Ma il docente vuole collegare WalkRun anche direttamente a Fall: il
personaggio ``può cadere anche nel caso in cui lui non abbia spinta, semplicemente cade da
qualche parte'', per esempio scendendo da un bordo. Con le sole regole su IsFalling,
però, le due transizioni in uscita da WalkRun avrebbero la stessa condizione, e Unreal
non saprebbe distinguere un salto da una caduta.

\paragraph{Seconda versione: la velocità sulla Z.} ``Dovremmo andare a differenziare un
altro valore: la Z. Se la Z è maggiore di 0, vuol dire che sta saltando; se la Z invece è
minore di 0, vuol dire che sta cadendo.'' Nell'Event Graph dell'ABP, finora, della
velocity si prendeva solo la lunghezza. Il docente prova prima la strada sbagliata (la
\emph{Jump Z Velocity} del Character Movement, che però è la spinta \emph{iniziale} del
salto, un valore fisso), poi prende la \textbf{Velocity} già salvata, ne fa lo
\textbf{Split Struct Pin} (``lo splittiamo'') e dalla \textbf{Z} fa \textbf{Promote to
Variable}: una nuova float, \textbf{VelocityZ} (``Z va benissimo''). Nelle transizioni:
\begin{itemize}
  \item \textbf{WalkRun $\rightarrow$ Jump}: IsFalling \textbf{AND} VelocityZ maggiore
        di zero, con un nodo \textbf{AND Boolean}. Il docente dice anche ``quando la
        velocità sarà maggiore di 1, più preciso'': una soglia piccola invece di zero
        esatto.
  \item \textbf{WalkRun $\rightarrow$ Fall}: IsFalling e VelocityZ minore di zero.
\end{itemize}
Per tenere in ordine i fili, si possono usare i \emph{reroute node} (``volete collegarli
con questo? È un po' bruttino''). Il risultato: ``se io salto, fa l'animazione''; ``se io
adesso invece cado solamente, lui va diretto in Fall''. Il docente chiude salvando anche
il BP\_ThirdPersonCharacter che l'editor segnala come modificato (``va bene, grazie'').

\begin{table}[htbp]
\centering
\small
\begin{tabular}{@{}p{3.2cm}p{5.3cm}p{5.5cm}@{}}
\toprule
\textbf{Transizione} & \textbf{Regola} & \textbf{Significato} \\
\midrule
Entry $\rightarrow$ WalkRun & nessuna & lo stato iniziale \\
WalkRun $\rightarrow$ Jump & IsFalling AND VelocityZ $>$ 0 & è in aria e sta salendo: ha
saltato \\
Jump $\rightarrow$ Fall & Automatic Rule Based on Sequence Player in State & finita
l'animazione di salto, passa alla caduta \\
WalkRun $\rightarrow$ Fall & IsFalling AND VelocityZ $<$ 0 & è in aria e sta scendendo:
è caduto da un bordo \\
Fall $\rightarrow$ WalkRun & NOT IsFalling & è tornato a terra \\
\bottomrule
\end{tabular}
\caption{Le regole di transizione della State Machine Locomotion (versione finale).}
\label{ue9:tab:transizioni}
\end{table}

\paragraph{Approfondimento: altre regole basate sul tempo.} L'opzione automatica della
lezione è il modo più rapido per dire ``quando l'animazione finisce''. Romero e Sewell
(capitolo 17), per gli stati \emph{StandToProne} e \emph{ProneToStand}, usano invece una
regola esplicita: il nodo \textbf{Time Remaining (ratio)} dell'animazione dello stato,
seguito da un \emph{Less} con 0,1: la transizione scatta quando manca meno del 10\%
dell'animazione. Il vantaggio è il controllo: la transizione può partire un po' prima
della fine e lasciare spazio al \emph{blend} fra i due stati, che ogni transizione ha
(nei Details: \emph{Blend Settings}, con una durata di default di 0,2 secondi).
Ricordano anche che le animazioni ``di una volta'', come il salto, vanno impostate senza
\emph{Loop Animation}: è il caso dello stato Jump, mentre Fall, che deve durare quanto
la caduta, resta in loop.

% ---------------------------------------------------------------------
\section{Lo sprint con l'Enhanced Input}
\label{ue9:sec:enhanced}

``Adesso andremo a mettere il nostro sprintino'': ``un semplicissimo sistemino che passa
la Max Walk Speed da un tot a un altro tot''. Si torna in BP\_Knight1, ``a ragionare a
livello di Blueprint''. ``Ma sì, facciamo anche l'input mapping.''

\paragraph{L'Input Mapping Context.} ``Unreal ha introdotto, quasi dalla 5.2'', quello
che è l'\textbf{Input Mapping Context}, ``una cosa molto interessante di Unreal che ci
permette, ad esempio, di avere degli eventi legati a più input''. Finora, in molte
lezioni, si sono usati gli eventi dei tasti (``abbiamo visto il tasto E, eccolo qua''):
``semplicemente un rilevamento, da parte di Unreal, della tastiera: se è premuto o non è
premuto''. Adesso, invece, ``gli eventi sono legati a più controlli possibili''.

Nel Content Browser si cerca \textbf{IMC} (\emph{Input Mapping Context}). ``È un
\emph{contesto}: un contesto di input che vengono dati in un certo momento, e sappiate
che sono sovrascrivibili.'' L'esempio: ``io devo cambiare dal gioco di avventura 3D al
gioco di guida della macchina''; ``i controlli sono completamente diversi, quindi cambio
contesto'': un Input Mapping Context per la macchina e uno ``per il character che cammina
normalmente''. Aperto \textbf{IMC\_Default}, ci sono i \textbf{Mappings}:
\textbf{IA\_Jump}, \textbf{IA\_Move}, \textbf{IA\_Look}.

\paragraph{L'Input Action IA\_Run.} ``Per far sì di avere un'Input Action, dobbiamo creare
un asset chiamato \textbf{Input Action}.'' Cercando \emph{IA} nel Content Browser ``troverete
un sacco di asset, ma a voi vi interessa questo celestino''. Con il tasto destro,
\textbf{Show in Folder View}, ``in modo tale che voi troviate direttamente la cartella''.
Si duplica IA\_Jump (Ctrl+C, Ctrl+V, o \emph{Duplicate}) e lo si rinomina
\textbf{IA\_Run}.

\paragraph{Dentro un'Input Action.} Aperto IA\_Run, ``ci sono un po' di valori che
possiamo vedere'':
\begin{itemize}
  \item \textbf{Description}: ``semplice descrizione, quindi quello che fa'';
  \item \textbf{Trigger when Paused}: se l'azione deve scattare anche quando il gioco è in
        pausa, ``quindi si blocca il tempo'';
  \item \textbf{Value Type}: ``la tipologia di valore che ci ritorna indietro. Un input ci
        dà un valore'': \textbf{Digital (bool)}, \textbf{Axis1D (float)} (``come abbiamo
        visto, avanti e indietro''), \textbf{Axis2D (Vector2D)} (``più assi, due
        dimensioni'') o \textbf{Axis3D (Vector)} (``tre tipologie di input che ci vengono
        dati''). ``A noi interessa il Run, che è un bool, quindi premuto o non premuto'':
        Digital (bool) ``va benissimo'';
  \item \textbf{Triggers}: ``le azioni che sono già preinstallate in Unreal'', che dicono
        ``ehi, guarda, si triggera se è premuto, o se è rilasciato, o se è sempre
        premuto''. L'elenco letto dal docente: \emph{Chorded Action}, \emph{Combo},
        \emph{Down}, \emph{Hold}, \emph{Hold and Release}, \emph{Pressed}, \emph{Pulse},
        \emph{Released}, \emph{Tap} (la trascrizione dice anche ``Short Deduction'' e
        ``repeat'', non identificabili con certezza). ``Fondamentalmente, questo vuol
        dire che quando si compie questa azione, questo input viene triggerato.'' IA\_Run
        eredita da IA\_Jump la sua configurazione, ``e va benissimo così, perché tanto è
        un'azione molto simile al jump''.
\end{itemize}

\begin{table}[htbp]
\centering
\small
\begin{tabular}{@{}p{3.2cm}p{10.8cm}@{}}
\toprule
\textbf{Voce} & \textbf{Significato} \\
\midrule
Digital (bool) & premuto / non premuto: un pulsante (salto, corsa, interazione) \\
Axis1D (float) & un valore su un asse: un grilletto analogico, avanti/indietro \\
Axis2D (Vector2D) & due assi: il movimento (X, Y) o la visuale del mouse \\
Axis3D (Vector) & tre assi: input 3D (per esempio un controller di movimento) \\
\midrule
Down & scatta a ogni frame finché il tasto è giù \\
Pressed & scatta una volta, quando il tasto viene premuto \\
Released & scatta una volta, quando il tasto viene rilasciato \\
Hold & scatta dopo che il tasto è stato tenuto premuto per un certo tempo \\
Hold and Release & scatta al rilascio, se il tasto è stato tenuto abbastanza \\
Tap & scatta se il tasto viene premuto e rilasciato entro un tempo breve \\
Pulse & scatta a intervalli regolari mentre il tasto è premuto \\
Chorded Action & scatta solo se un'altra azione è attiva insieme (una ``combinazione'') \\
Combo & scatta se una sequenza di azioni viene eseguita nell'ordine \\
\bottomrule
\end{tabular}
\caption{I tipi di valore (Value Type) e i trigger di un'Input Action, letti a lezione
dall'elenco dell'editor. Le descrizioni sono un'integrazione: a lezione il docente
legge i nomi senza spiegarli uno per uno.}
\label{ue9:tab:inputaction}
\end{table}

\paragraph{Il contesto va caricato.} Tornati in BP\_Knight1: ``abbiamo già il contesto
caricato'', con il grafo copiato dal ThirdPersonCharacter. ``Sappiate che se voi non
caricate il contesto, non vi prenderà nessuna cosa'': sull'Event Begin Play, dal Player
Controller si prende l'\emph{Enhanced Input Local Player Subsystem} e si chiama
\textbf{Add Mapping Context} con IMC\_Default (figura~\ref{ue9:fig:enhanced}).

\paragraph{L'evento IA\_Run e i suoi pin.} Nel grafo si invoca \textbf{IA\_Run}; ``se lo
estendete'' (la freccina in basso del nodo), compaiono tutti i pin:
\begin{itemize}
  \item i pin di esecuzione \textbf{Triggered} (``quando è triggerato''), \textbf{Started}
        (``quando inizia a essere triggerato''), \textbf{Ongoing}, \textbf{Canceled},
        \textbf{Completed};
  \item i valori: l'\textbf{Action Value} (``se è true o false, in questo caso''),
        \textbf{Elapsed Seconds} e \textbf{Triggered Seconds} (``quanto tempo è stato
        triggerato''), e l'\textbf{Input Action}, ``la referenza stessa'' all'asset.
\end{itemize}

\paragraph{Set Max Walk Speed.} ``Dobbiamo modificare il Character Movement, ragazzi, che
dovrà definire quanta velocità avremo'': si trascina nel grafo il componente
\textbf{Character Movement} e se ne tira fuori \textbf{Set Max Walk Speed}, due volte.
Sul \textbf{Triggered}, \textbf{600}; sul \textbf{Completed}, \textbf{300}
(figura~\ref{ue9:fig:sprint}).

\paragraph{Il mapping dimenticato.} In Play non succede niente. ``Scusate, ragazzi, mi
sono dimenticato di dare l'input'': IA\_Run esiste, ma non è nel contesto. Si apre
IMC\_Default, si aggiunge un mapping (il \emph{+} accanto a \emph{Mappings}), si sceglie
IA\_Run e gli si assegna il tasto: ``un nostro bello \textbf{Shift}''. ``L'Input Mapping
Context ci permette di mettere quali tipologie di pulsanti facciano'' ciascuna azione.

\paragraph{Lo stesso input, anche con il controller.} ``Io ad esempio qua, ragazzi, ho un
controller'': il docente lo collega (``Unreal riesce a prendere il controller''), aggiunge
un secondo tasto alla stessa azione, clicca sull'icona della tastiera accanto al campo e
preme un pulsante del controller (``mi sono fregato la sorpresa''): ``facciamo tipo Dark
Souls? Premo cerchio'', e l'editor scrive \textbf{Gamepad Face Button Right}. ``Questa
azione, cioè l'azione di correre, è collegata a più input.'' ``Quando faccio progetti in
generale, devo fare anche il controller'': si corre con Shift sulla tastiera ``oppure con
il nostro bel controllino, con quello che vogliamo noi'', ``come è stato designato il
gioco''. Tutto il resto del movimento sul controller ``non ci interessa adesso''.

\paragraph{La prova.} ``Se noi proviamo\dots\ rilascio.'' Con Shift il cavaliere corre, e
il Blend Space passa alla corsa; al rilascio torna a camminare. ``Ecco a voi che abbiamo
creato un nostro bel character nuovo, in cui semplicemente possiamo correre, possiamo
modificare la velocità.''

\begin{figure}[htbp]
\centering
\begin{tikzpicture}[font=\footnotesize,
  ue9box/.style={draw, thick, rounded corners=4pt, align=center, fill=gray!8,
                 minimum height=1.0cm, text width=2.5cm},
  ue9arr/.style={-{Stealth[length=2.5mm]}, thick}]
  \node[ue9box, fill=yellow!15] (k1) at (0,0.7) {Left Shift\\{\scriptsize tastiera}};
  \node[ue9box, fill=yellow!15] (k2) at (0,-0.7) {Gamepad Face\\Button Right};
  \node[ue9box, fill=blue!10, text width=2.8cm] (imc) at (3.8,0) {\textbf{IMC\_Default}\\
    {\scriptsize Input Mapping Context}\\{\scriptsize (mapping IA\_Run)}};
  \node[ue9box, fill=cyan!12] (ia) at (7.6,0) {\textbf{IA\_Run}\\{\scriptsize Input
    Action, Digital (bool)}};
  \node[ue9box, fill=red!12, text width=2.8cm] (ev) at (11.4,0) {evento
    \textbf{IA\_Run}\\{\scriptsize nel Blueprint}\\{\scriptsize Triggered / Completed}};
  \draw[ue9arr] (k1) -- (imc); \draw[ue9arr] (k2) -- (imc);
  \draw[ue9arr] (imc) -- (ia); \draw[ue9arr] (ia) -- (ev);
  \node[ue9box, fill=green!10, text width=4.2cm] (bp) at (3.8,-2.4) {\textbf{Event Begin
    Play} $\rightarrow$ Enhanced Input Local Player Subsystem $\rightarrow$ \textbf{Add
    Mapping Context} (IMC\_Default)};
  \draw[ue9arr, dashed] (bp) -- node[right, font=\scriptsize]{attiva} (imc);
\end{tikzpicture}
\caption{La catena dell'Enhanced Input. I tasti fisici non sono legati direttamente
all'evento: il contesto (IMC) associa più tasti alla stessa azione (IA), e il Blueprint
reagisce all'azione. Senza Add Mapping Context, il contesto non è attivo e nessun input
arriva. Schema ricostruito.}
\label{ue9:fig:enhanced}
\end{figure}

\begin{figure}[htbp]
\centering
\begin{tikzpicture}[font=\footnotesize,
  ue9node/.style={draw, thick, rounded corners=4pt, align=center, fill=gray!8,
                  minimum height=0.8cm, text width=2.4cm},
  ue9ev/.style={ue9node, fill=red!12},
  ue9var/.style={ue9node, fill=blue!10, rounded corners=8pt},
  ue9exec/.style={-{Stealth[length=2.5mm]}, very thick, black!80},
  ue9data/.style={-{Stealth[length=2mm]}, thick, blue!60!black}]
  \node[ue9ev, text width=2.2cm, minimum height=1.8cm] (ev) at (0,0) {\textbf{IA\_Run}\\
    {\scriptsize Enhanced Input Action}};
  \node[ue9node, text width=3.2cm] (s1) at (5,0.8) {\textbf{Set Max Walk Speed}\\
    {\scriptsize Max Walk Speed = 600}};
  \node[ue9node, text width=3.2cm] (s2) at (5,-0.8) {\textbf{Set Max Walk Speed}\\
    {\scriptsize Max Walk Speed = 300}};
  \node[ue9var, text width=2.6cm] (cm) at (9.6,0) {\textbf{Character}\\
    \textbf{Movement}};
  \draw[ue9exec] (ev.east |- s1) -- node[above, font=\scriptsize]{Triggered} (s1);
  \draw[ue9exec] (ev.east |- s2) -- node[below, font=\scriptsize]{Completed} (s2);
  \draw[ue9data] (cm.west) -- ++(-0.5,0) |- node[pos=0.25, right, font=\scriptsize]
    {Target} (s1.east);
  \draw[ue9data] (cm.west) -- ++(-0.5,0) |- (s2.east);
\end{tikzpicture}
\caption{Lo sprint di BP\_Knight1. Premendo Shift (o il cerchio del controller) la
velocità massima sale a 600 e il Blend Space mostra la corsa; al rilascio torna a 300,
la camminata. Schema ricostruito.}
\label{ue9:fig:sprint}
\end{figure}

\paragraph{Approfondimento: Action Mapping ed Enhanced Input.} Il libro di Romero e
Sewell, scritto per le prime versioni di Unreal Engine 5, costruisce lo stesso sprint
con il vecchio sistema: un \emph{Action Mapping} chiamato \emph{Sprint} nei Project
Settings (sezione \emph{Input}), legato a \emph{Left Shift}, e un evento
\emph{InputAction Sprint} con i pin \emph{Pressed} e \emph{Released}, su cui due
\emph{Set Max Walk Speed} portano la velocità a 2200 e poi di nuovo ai 600 di default
(capitolo 6). Il principio è identico; cambia il sistema. L'\textbf{Enhanced Input},
introdotto come plugin in Unreal Engine 5.0 e diventato il sistema predefinito dei
modelli dalla 5.1 (il ``quasi dalla 5.2'' del docente è un ricordo approssimativo), separa
tre livelli: l'\emph{azione} (Input Action, che cosa vuole fare il giocatore), il
\emph{contesto} (Input Mapping Context, quali tasti attivano quali azioni in una certa
situazione) e il \emph{comportamento} (Triggers e \emph{Modifiers}, come per esempio la
\emph{zona morta} di una levetta o l'inversione di un asse). I contesti si aggiungono e
si tolgono a runtime con \emph{Add Mapping Context} e \emph{Remove Mapping Context}, e
hanno una \emph{Priority}: è così che si passa ``dal gioco a piedi alla macchina''
dell'esempio del docente, e così che si costruisce una schermata di rimappatura dei
tasti. Gli autori raccomandano anche di raggruppare i nodi dello sprint in un
\emph{commento} (\emph{Create Comment from Selection}), per tenere ordinato il grafo.

% ---------------------------------------------------------------------
\section{Un tentativo: FInterp To}
\label{ue9:sec:finterp}

``Se volessimo farlo un po' più carino, ragazzi'', il passaggio fra 300 e 600 potrebbe
essere graduale invece che istantaneo. Il blocco è \textbf{FInterp To} (``finter2''):
``fa un'interpolazione di un float''. ``Avete presente quando abbiamo parlato di
interpolazione?'' ``Fa il passaggio tra due numeri. Lui dice: devi passare da questo
valore, cioè \textbf{Current}, a questo valore, \textbf{Target}, in questo tempo'',
\textbf{Delta Time}, per cui il docente consiglia \textbf{Get World Delta Seconds}; e
poi c'è l'\textbf{Interp Speed}, la velocità dell'interpolazione.

\paragraph{Il grafo.} Sul Triggered: \textbf{Current} = Character Movement $\rightarrow$
\textbf{Get Max Walk Speed}; \textbf{Target} = 600; \textbf{Delta Time} = World Delta
Seconds; il risultato nel Set Max Walk Speed. Il docente ricorda che il Delta Seconds si
può moltiplicare per un valore (``ricordiamo che possiamo modificarlo moltiplicandolo
per un valore''), e prova con 10, poi con un decimo. Sul Completed lo stesso, ``ma
semplicemente invertiamo'': il Current parte da 600, il Target è 300. ``Ricordiamoci di
collegare anche al Character Movement.''

\paragraph{Non funziona, e perché.} La corsa resta ``ancora molto veloce'', il passaggio
istantaneo. Il docente prova valori diversi (10, 100: ``è alla norma del pianeta stare
continuamente a vedere se i valori sono giusti'') e poi mette dei \textbf{Print String}:
``ci dà seicento''. La diagnosi: ``fondamentalmente questa è un'azione a triggered
unico, quindi non farà l'interpolazione''. FInterp To fa \emph{un passo}
dell'interpolazione ogni volta che viene eseguito: se viene eseguito una volta sola, fa
un solo passo. ``Lo deve fare tramite un \textbf{timer}'': si rimanda a dopo (``volendo,
ragazzi, si può fare con un timer, però facciamolo magari dopo''), e non verrà ripreso
nella lezione.

\paragraph{Approfondimento: dove mettere un'interpolazione.} Un nodo come FInterp To
(o VInterp To, RInterp To per vettori e rotazioni) va eseguito \emph{a ogni frame}
finché il valore non arriva al target: il suo risultato è il valore ``un po' più vicino''
al target, e il Delta Time serve a renderlo indipendente dal frame rate. Ci sono tre
modi tipici di farlo girare: nel \textbf{Tick}, con una variabile \emph{TargetSpeed}
che i due pin dell'input si limitano a impostare (600 o 300); con un \textbf{timer}
(\emph{Set Timer by Event}, con \emph{Looping}, fermato con \emph{Clear and Invalidate
Timer} quando il valore è arrivato), come suggerisce il docente; o con una
\textbf{Timeline}, vista nella sesta lezione, che va da 300 a 600 con \emph{Play} e
torna indietro con \emph{Reverse}. Con l'Enhanced Input c'è anche una quarta via: il pin
\emph{Triggered} scatta a ogni frame se il trigger è \emph{Down} (o se non c'è nessun
trigger), e una sola volta se il trigger è \emph{Pressed}. IA\_Run, copiata da IA\_Jump,
ha probabilmente un trigger che la fa scattare una volta sola: è ``l'azione a triggered
unico'' del docente.

% ---------------------------------------------------------------------
\section{Prima della pausa: le animazioni da scaricare}
\label{ue9:sec:download}

``Adesso andremo a vedere, se non mi esplode prima Unreal, anche come si fanno i
\textbf{montaggi}''; mancano una ventina di minuti alla pausa. ``Dovremmo andare un attimo
su Mixamo'' (``a quanto pare oggi Mixamo ci mette pure un bel po' a caricare''), per
scaricare due gruppi di animazioni.

\paragraph{Un attacco in root motion.} ``Vorrei che sceglieste un'animazione che sia in
\textbf{root motion}'': ``riprendiamo ad esempio una situazione simile a Dark Souls''.
``Ricordatevi sempre di prendere il vostro character'' su Mixamo (il paladino, nel caso
del docente), perché le animazioni si adattano al personaggio selezionato. Il docente
cerca \emph{attack} e sceglie un attacco con la spada che ``di base, come potete vedere,
non può essere messo in place''. Le impostazioni di download sono quelle della volta
precedente: \textbf{Without Skin}, \textbf{FBX Binary}, \textbf{30} fotogrammi al
secondo. ``Ripeto: non deve essere in place, deve avere il root motion, quindi il
movimento ereditato.'' ``Un'azione unica'': ``ad esempio un attacco con una spada, un
qualche cosa'', ``un'animazione che vi piace di più''.

\paragraph{Le animazioni del Blend Space 2D.} Dopo la pausa si farà il \textbf{Blend
Space 2D}, ``che ci servirà per fare anche il calcolo della direzione del nostro
character''. Servono quindi altre animazioni: ``walk destra, sinistra, dritto e
indietro''. ``Consiglio spassionato'': su Mixamo si chiamano \textbf{Strafe} (``provate a
cercare Strafe''; ``Jog Strafe va benissimo anche''), e queste, al contrario
dell'attacco, ``dovete metterle \textbf{in place}''.

\paragraph{Il mirror.} ``Piccolo consiglio: se volete far sì che siano abbastanza simili
tra di loro, potete fare il \textbf{mirror}.'' Si scarica l'animazione verso sinistra
non specchiata; poi, con l'opzione \emph{Mirror} di Mixamo, ``in automatico andrà a
destra''. Le due animazioni saranno così perfettamente simmetriche.

\paragraph{Quante animazioni.} Su richiesta (``le puoi scrivere per messaggio?''), il
docente le elenca anche in chat: ``in totale sono otto animazioni''. Avanti, indietro,
destra, sinistra, e le quattro \textbf{diagonali}: avanti a destra, avanti a sinistra,
indietro a destra, indietro a sinistra. Uno studente chiede se l'avanti serva, visto che
c'è già la camminata: ``nel senso, possiamo usare tranquillamente quello di avanti; ma in
questo caso abbiamo solamente avanti, dobbiamo trovare anche indietro''. E un'altra: le
diagonali ``soltanto in avanti, o anche indietro?'' ``Anche indietro'': ``se io vado tipo
così, indietro verso destra, e così, indietro verso sinistra''. ``Praticamente è la X,
letteralmente'': le due diagonali incrociate, oltre alla croce degli assi. Il docente
nomina anche alcuni pacchetti di Mixamo con animazioni di locomozione (``Locomotion
Pack''; il resto dei nomi è incomprensibile nella trascrizione). ``Ci rivediamo alle
18.''

% ---------------------------------------------------------------------
\section{La pausa e i feedback}
\label{ue9:sec:feedback}

Al rientro, ``siamo rimasti in cinque'', poi in quattro. Il docente apre uno spazio per
i commenti: ``se avete dei feedback riguardo a questo corso, sparate tranquillamente'',
``siamo anche molto coetanei, anzi credo io di essere anche forse più piccolo di alcuni
di voi''. ``Se avete cose da chiedere, magari di andare più piano, andare più veloce,
oppure un po' come sono strutturate le lezioni''. ``Capisco pure che può essere
effettivamente un po' pesante in alcuni momenti'', ``perché purtroppo i Blueprint, o
Unreal in generale, è divertente quando poi le sai, le cose''.

\paragraph{Le funzioni da imparare a memoria.} Uno studente: ``immagino che sono
tantissime funzioni, quindi sono tutte da imparare a memoria''. ``Quello, eh sì.''

\paragraph{Il progetto del docente: la pallina morbida e dura.} Per mostrare dove porta
la pratica, il docente racconta un suo progetto personale: ``una specie di
\emph{platformer} in cui una pallina diventa morbida e dura, nel senso che diventa di
metallo o di un materiale più morbido''. ``La cosa è che ho simulato tutta la fisica'':
``la fisica non è simulata da Unreal, ma sono io che faccio finta, tramite i Blueprint, di
far fare delle funzioni fisiche''. ``Hai messo proprio te le formule fisiche?'' ``Più o
meno.'' Il meccanismo: si rileva lo \textbf{slope}, ``il quanto è piegato il piano'',
perché ``la meccanica è il fatto che se sei più pesante con la pallina vai più veloce su
uno slope, come nella realtà'': ``una pallina di metallo e una pallina di gomma: quella di
metallo va più veloce, perché è più pesante, comunque più influenzata''. ``È tutto
simulato'': ``sono io che, tramite la detection dello slope, cambio i valori''.

\paragraph{Combinare i blocchi.} ``E quello è ovviamente tutto custom'': ``quando hai
tutte le nozioni in testa, ti dici: ok, so questo blocco A, B, C, D, E, F, G; quale mi
serve? Il blocco C e il blocco A. Li metto insieme e faccio una nuova combinazione.''
``È sempre il fatto di capire come posso interlacciare le cose.'' ``Qua stiamo facendo
cose magari un po' generiche''; poi ``dici: ok, voglio fare simile a quello, ma
leggermente diverso. Voglio fare il character, però con la spada: che mi cambia?''

\paragraph{Inventari e UI.} La richiesta di Federico: ``a me interesserebbero la gestione
degli inventari e le interfacce''. ``Sfortunatamente no'': ``sappi che l'inventario e la
UI sono un mondo a parte''. ``Esistono proprio figure riguardanti questo mondo della UI,
sia designer sia programmatori.'' Si farà ``una cosa come un menu'', ``una cosetta
tranquilla'', ``però per il resto non andremo a fare di più''.

\paragraph{Le lezioni con Davide.} ``Sappiate che il 25 è la mia ultima lezione, e avrete
tutte quelle con \textbf{Davide}, che invece vi farà la parte più di \emph{level}.'' Con
lui si useranno ``tutti gli asset che abbiamo creato: le piattaforme, i pulsanti'', per
lavorare su concetti come la \textbf{gym}, ``il concetto della palestra'', che ``si fa
praticamente in via di sviluppo'': ``è il parco giochi'', ``letteralmente il parco giochi
dove ho tutte le meccaniche, le posso provare tutte quante, vedo se funzionano''. Poi si
costruiscono situazioni: ``voglio fare un livello horror, la bestia che ti insegue: tu
entri in una stanza e devi fuggire dalla bestia''; ``situazione X, situazione Y,
situazione Z'', ``le costruisci e metti insieme i componenti'', ``e vedi quei componenti
che hai programmato come funzionano a livello anche di esperienza, di gioco''. Gli
elementi saranno ``quelli di un classico platform: le piattaforme, i bottoni'', e si
vedranno anche i tool di \textbf{modeling} dell'editor e ``un po' di teoria di level
design'': ``chi fa game design, e a chi interessa tantissimo anche la parte di level:
molto, molto interessante''.

\paragraph{Dove lasciare i feedback.} ``Ripeto, ragazzi, anche a chi sta in differita'':
qualsiasi feedback sulle lezioni va lasciato sul canale \textbf{general} della chat,
``non su introduzione al materiale, perché sennò si fa casino''. Non solo sugli
argomenti, ``ma feedback anche proprio per le lezioni in generale: andare più lento,
andare più veloce, vedere cose più pratiche''.

\paragraph{Approfondimento: la gym.} Nel level design la \textbf{gym} (o \emph{gym
level}, \emph{gray box gym}) è un livello di servizio che non finirà nel gioco: una
stanza in cui ogni meccanica è disposta in isolamento e misurata. Si mettono salti di
lunghezze crescenti per sapere qual è il salto massimo del personaggio, rampe di
pendenze diverse per verificare la \emph{Max Walkable Floor Angle}, gradini di varie
altezze per la \emph{Max Step Height}, corridoi di larghezze diverse per la camera. Le
misure raccolte nella gym diventano le ``regole metriche'' del level design (quanto può
essere largo un burrone, quanto alto un muretto da scavalcare), e i livelli veri si
costruiscono rispettandole. È il ``parco giochi'' del docente; e il problema della
camera che compenetra, visto in questa lezione, è uno dei primi che una gym fa emergere.
% ---------------------------------------------------------------------
\section{La direzione: Calculate Direction}
\label{ue9:sec:direction}

``Abbiamo detto prima che oggi dobbiamo fare il Blend Space 2D.'' Il docente ha importato
otto animazioni, ``cercando di prenderle tutte similari'': \emph{Jog Backward},
\emph{Jog Backward Diagonal}, \emph{Jog Forward}, \emph{Jog Forward Diagonal} e così
via, ``a parte lo Strafe Left, che devo eliminare perché non andava bene come
animazione''. Il Blend Space 1D prende solo la velocità: ``vai più veloce, l'animazione,
benissimo''. Il 2D aggiunge un secondo asse, la \textbf{direzione}, ``in modo tale da non
dover magari direzionare sempre il character'', ma vedere ``che si muove in diagonale, si
muove a sinistra''.

\paragraph{Dove si calcola.} In ABP\_Knight, nell'Event Graph dell'Update, si aggiunge un
pin alla \textbf{Sequence} (il \emph{+}, \emph{Add pin}): \emph{Then 2}, o il successivo.
Questa volta non si parte dal Movement Component, ma dal \textbf{character}: ``perché la
direzione è determinata dalla rotazione del character stesso'', ``cioè la capsula, come è
ruotata''. Dalla variabile del cavaliere si tira fuori \textbf{Get Actor Rotation}.

\paragraph{Il nodo Calculate Direction.} ``Potremmo tranquillamente prenderla così, ma noi,
per semplificarci anche la vita, andiamo a calcolare proprio la direzione con un blocco
chiamato proprio \textbf{Calculate Direction}.'' Il nodo chiede due ingressi: la
\textbf{Velocity} attuale (dal Movement Component, \textbf{Get Velocity}) e la
\textbf{Base Rotation} (la Get Actor Rotation). Restituisce ``un float'', che con
\textbf{Promote to Variable} diventa \textbf{Direction}.

\paragraph{Uno studente perso.} ``Io mi sono perso: Character, Movement Component, da dove
vengono questi due?'' Il docente ripercorre: ``hai fatto questo Event Blueprint Initialize
Animation nell'Event Graph, no?''; ``dovresti avere qua sotto, sotto le variabili, la
variabile del character: la tiri fuori, ne tiri fuori Get Actor Rotation''; dal Movement
Component ``tiri fuori la velocità''; ``da Get Actor Rotation fai Calculate Direction, e
questo diventerà il nostro float''. Un altro studente nota che la variabile del docente si
chiama \emph{Character}, mentre la sua si chiama \textbf{As BP Knight}: ``va benissimo.
Io l'ho solo rinominata.'' Quando si fa Promote to Variable sull'uscita di un cast, Unreal
propone il nome \emph{As} più il nome della classe (``ti dice proprio As BP Knight'').

\paragraph{Da $-180$ a $180$.} ``Molto semplicemente, la direzione sarà tra un range,
sempre senza doverla clampare'' (``il clamp l'abbiamo fatto quasi all'inizio''): fra
\textbf{$-180$ e $180$ gradi}. ``Si parte da 0; a destra vai verso il positivo; quando
arrivi a 180 passa a $-179$''; ``diciamo, non lo fa come un angolo da 0 a 360, ma fa 0,
180, $-180$ e di nuovo 0'' (figura~\ref{ue9:fig:direzione}).

\begin{figure}[htbp]
\centering
\begin{tikzpicture}[font=\footnotesize]
  \draw[thick] (0,0) circle (2.2);
  \foreach \a/\l/\pos in {90/0/above, 45/45/above right, 0/90/right, -45/135/below right,
      -90/{180 / $-180$}/below, -135/{$-135$}/below left, 180/{$-90$}/left,
      135/{$-45$}/above left}
    {\draw[-{Stealth[length=2mm]}, thick, blue!60!black] (0,0) -- (\a:2.2);
     \node[\pos, font=\scriptsize\bfseries] at (\a:2.25) {\l};}
  \fill[gray!40] (0,0) circle (0.35);
  \draw[thick] (0,0) circle (0.35);
  \draw[-{Stealth[length=2mm]}, very thick] (0,0) -- (0,0.6);
  \node[font=\scriptsize, align=left, anchor=west] at (3.4,1.2) {\textbf{0}: avanti\\
    \textbf{positivi}: verso destra\\\textbf{negativi}: verso sinistra\\
    \textbf{$\pm$180}: indietro};
  \node[font=\scriptsize, align=left, anchor=west] at (3.4,-0.9) {la direzione è
    l'angolo fra\\la velocità e il ``davanti''\\del personaggio (Base Rotation)};
\end{tikzpicture}
\caption{Il valore di Direction restituito da Calculate Direction. Il personaggio guarda
in alto; le frecce sono le otto direzioni delle animazioni del Blend Space 2D. Andare
indietro vale sia 180 sia $-180$, a seconda del lato da cui ci si arriva. Schema
ricostruito.}
\label{ue9:fig:direzione}
\end{figure}

\paragraph{Rimettere Use Controller Rotation Yaw.} Per vedere i movimenti laterali,
``importante è rimettere lo \textbf{Use Controller Rotation Yaw}'': ``non vogliamo che il
character ruoti sempre nella direzione in cui lo stiamo indirizzando''. Il corpo deve
seguire la camera, e la direzione del movimento rispetto al corpo deve cambiare.

\paragraph{La verifica con Print String.} Per sicurezza si stacca un attimo il Blend Space
1D (``staccare un attimo le animazioni'') e si stampa Direction con un \textbf{Print
String}, ``che lo farà a ogni frame''. ``Ci dà direttamente nella preview questo valore'';
in Play, ``sarà un po' rotto, ma è normale, ragazzi'': ``$-70$, 70, 160, 0''. ``Mamma mia''
(il valore dipende dalla camera, ``giustamente''). ``Come possiamo vedere, questi valori
rappresentano la nostra direzione attuale.'' Il docente nota anche che il cavaliere
sembra ``leggermente ruotato'' rispetto alla capsula: ``buono a sapersi''.

\paragraph{Approfondimento: che cosa calcola Calculate Direction.} Il nodo, della
libreria dell'Animation Blueprint (\emph{KismetAnimationLibrary}; nelle versioni più
vecchie era una funzione dell'AnimInstance), trasforma la velocità nello spazio locale
del personaggio e ne restituisce l'angolo sul piano orizzontale, in gradi, nell'intervallo
$[-180, 180]$: 0 se la velocità punta dove guarda il personaggio, 90 se punta alla sua
destra, $-90$ a sinistra, $\pm 180$ dietro. A velocità nulla restituisce 0. Con
\emph{Orient Rotation to Movement} acceso il personaggio si gira sempre verso la
velocità, e la direzione resta quasi sempre vicina a zero: il Blend Space 2D avrebbe poco
da fare. Per questo il docente riattiva Use Controller Rotation Yaw.

% ---------------------------------------------------------------------
\section{Il Blend Space 2D}
\label{ue9:sec:blend2d}

\paragraph{Creare l'asset.} Tasto destro $\rightarrow$ \emph{Animation} $\rightarrow$
\textbf{Blend Space}: ``Blend Space senza niente, di base, è 2D; se viene specificato
1D, è 1D''. Unreal chiede lo scheletro: ``voi scegliete il vostro scheletro''. Il nome:
\textbf{BS2D\_Knight}.

\paragraph{I due assi.} ``Come potete vedere, a sinistra e giù sono presenti due nomi'',
per ora \emph{None}. Nei Details:
\begin{itemize}
  \item \textbf{Horizontal Axis}: \textbf{Direction}, da \textbf{$-180$} a
        \textbf{$180$};
  \item \textbf{Vertical Axis}: \textbf{GroundSpeed} (``più che la velocità, la ground
        speed''), da \textbf{0} a \textbf{600}.
\end{itemize}
Il parametro \textbf{Grid Divisions} ``ci permette di aumentare il numero di divisioni
sulla nostra griglia'' (le linee ``non si vedono, perché sono veramente molto poco
chiare''): sono i punti su cui si agganciano i campioni delle animazioni. Il docente ne
mette \textbf{8} su entrambi gli assi. Con 8 divisioni, sull'asse orizzontale ogni passo
vale 45 gradi, esattamente le otto direzioni; sull'asse verticale vale 75.

\paragraph{Un quiz.} ``Domanda, ragazzi: se io sto a 600 di velocità e a 90 di direzione,
dove sto andando?'' ``Ricordatevi che positivo è destra, negativo è sinistra.'' ``A
destra.'' ``Pienamente a destra, perfetto.'' ``E invece a $-135$?'' ``Dietro a
sinistra.'' ``Ottimo.''

\paragraph{Dove va ogni animazione.} Si dispongono i campioni (tabella~\ref{ue9:tab:bs2d} e
figura~\ref{ue9:fig:blend2d}). In basso, a velocità zero, l'\textbf{Idle}: ``mettetene
tre, così abbiamo un range per tutto'' (a $-180$, 0 e 180): ``siamo fermi, siamo fermi''.
In alto, a 600, le otto animazioni. L'animazione all'indietro va messa due volte, a
$-180$ e a $180$: ``stiamo andando indietro sia se abbiamo girato da sinistra, sia da
destra''. Il docente si corregge una volta (``Backward Diagonal L sarà\dots\ no, ho
sbagliato'': a $-135$, non a 90) e invita a controllare i valori nei Details del
campione.

\begin{table}[htbp]
\centering
\small
\begin{tabular}{@{}p{4.4cm}p{2.2cm}p{2.2cm}p{4.6cm}@{}}
\toprule
\textbf{Animazione} & \textbf{Direction} & \textbf{GroundSpeed} & \textbf{Movimento} \\
\midrule
Idle (tre copie) & $-180$, 0, 180 & 0 & fermo \\
Jog Forward & 0 & 600 & avanti \\
Jog Forward Diagonal R & 45 & 600 & avanti a destra \\
Strafe R & 90 & 600 & a destra \\
Jog Backward Diagonal R & 135 & 600 & indietro a destra \\
Jog Backward & 180 e $-180$ & 600 & indietro \\
Jog Backward Diagonal L & $-135$ & 600 & indietro a sinistra \\
Strafe L & $-90$ & 600 & a sinistra \\
Jog Forward Diagonal L & $-45$ & 600 & avanti a sinistra \\
\bottomrule
\end{tabular}
\caption{La posizione dei campioni nel Blend Space 2D BS2D\_Knight. I nomi delle
animazioni sono quelli di Mixamo citati a lezione, per quanto ricostruibili dalla
trascrizione; i valori sono quelli dettati dal docente.}
\label{ue9:tab:bs2d}
\end{table}

\begin{figure}[htbp]
\centering
\begin{tikzpicture}[font=\footnotesize, x=0.032cm, y=0.006cm,
  ue9pt/.style={draw=black, fill=white, thick, diamond, inner sep=1.8pt}]
  % griglia
  \foreach \x in {-180,-135,...,180} \draw[gray!30] (\x,0) -- (\x,600);
  \foreach \y in {0,75,...,600} \draw[gray!30] (-180,\y) -- (180,\y);
  \draw[thick] (-180,0) rectangle (180,600);
  \foreach \x in {-180,-90,0,90,180}
    \node[below, font=\scriptsize] at (\x,-10) {\x};
  \foreach \y in {0,300,600} \node[left, font=\scriptsize] at (-182,\y) {\y};
  \node[below, font=\small] at (0,-80) {Direction};
  \node[rotate=90, font=\small] at (-235,300) {GroundSpeed};
  % campioni
  \foreach \x in {-180,0,180} \node[ue9pt] at (\x,0) {};
  \node[above, font=\scriptsize] at (0,10) {Idle};
  \foreach \x/\l in {-180/Back, -135/BackDiag L, -90/Strafe L, -45/FwdDiag L, 0/Fwd,
      45/FwdDiag R, 90/Strafe R, 135/BackDiag R, 180/Back}
    {\node[ue9pt] at (\x,600) {};
     \node[above, font=\tiny, rotate=35, anchor=south west] at (\x-6,612) {\l};}
  % punto attuale
  \fill[green!50!black] (60,450) circle (2.2pt);
  \node[green!40!black, font=\scriptsize, anchor=west] at (66,420) {(60, 450)};
\end{tikzpicture}
\caption{Il Blend Space 2D del cavaliere: Direction in orizzontale, GroundSpeed in
verticale, griglia 8$\times$8. Le tre copie dell'Idle stanno in basso; le otto direzioni
in alto, con la corsa all'indietro ripetuta ai due estremi. Il punto verde è un esempio:
il personaggio sta andando in avanti a destra, un po' sotto la velocità massima, e Unreal
mescola i quattro campioni più vicini. Schema ricostruito.}
\label{ue9:fig:blend2d}
\end{figure}

\paragraph{La preview con Ctrl.} ``Potete fare la preview delle animazioni con
\textbf{Ctrl}'': tenendolo premuto e trascinando il cursore sul grafico, ``potete vedere
come l'animazione si muove tra tutti gli step''. ``Consiglio sempre di mettersi dietro''
il personaggio nel viewport, ``perché tanto poi, se lo vedete davanti, è difficile'' capire
destra e sinistra. Il docente controlla tutte le direzioni (``vado a sinistra, vado
indietro, vado dritto, vado dritto a destra, vado a destra, giù a destra, indietro.
Perfetto''); una delle animazioni ``mi sa che non va bene'', ma ``sì, sto andando a
sinistra, è solo molto strana''.

\paragraph{Smoothing Time.} ``Piccola cosa, se volete introdurla'': lo \textbf{Smoothing
Time}, ``il tempo di passaggio tra un'animazione e l'altra in generale, per tutte le
animazioni''. Lo si mette a \textbf{0,2} su entrambi gli assi. Il \emph{Damping Ratio}
``non ci interessa''.

\paragraph{Nello stato WalkRun.} ``Facciamo tutto il percorso a ritroso'': AnimGraph,
Locomotion, WalkRun. Al posto del Blend Space 1D si trascina BS2D\_Knight, che chiede
``dammi la Direction e la GroundSpeed'': si collegano i due Get. Compile: ``vediamo che il
nostro Blend Space funziona in preview''.

\paragraph{Le impostazioni del personaggio.} In BP\_Knight1: ``mi raccomando, l'abbiamo
detto prima, ma lo diciamo anche adesso'', si riattiva \textbf{Use Controller Rotation
Yaw}; e nel Character Movement la \textbf{Max Walk Speed} va a \textbf{600}, ``che è
quanto abbiamo messo'' come massimo dell'asse. ``Ecco a voi, ragazzi'': ``se io mi muovo a
destra lui va a destra, se mi muovo a sinistra, anche se l'animazione non è fantastica,
va a sinistra; vado indietro; vado in diagonale''. Nel frattempo ``parte la torretta, che
mi spara'': ``il mio bel cavaliere se ne va a zonzo come gli pare''.

\paragraph{Gli scatti.} C'è ``un piccolo \emph{clip}'', uno scatto, in alcuni passaggi.
Il docente prova a prendere la rotazione della \textbf{Mesh} invece di quella
dell'actor (``facciamo che invece di prendere la direzione del character, prendiamo
quella della Mesh''; ``Get Mesh''), poi ad aumentare lo Smoothing Time della Direction
a \textbf{0,5} (``un po' più tranquillo''), e infine torna al collegamento originale.
La conclusione: ``probabilmente sono dei problemi sull'animazione''; ``mi sa che devo
riscaricare l'animazione, perché snappa un pochino''. Ma il risultato c'è: ``abbiamo il
nostro character che si muove in tutte le direzioni''; ``è un po' veloce, devo essere
sincero'', e la velocità si può abbassare.

\paragraph{Il cerchio, per chi si è perso.} Uno studente chiede di rivedere il Blend
Space 2D. ``Immaginatelo come un cerchio'': ``se stai a 0 di direction, stai andando
dritto; se stai a 180 o $-180$, stai andando dietro''; ``tutti i valori intermedi'' sono le
diagonali e i lati. ``Destra è positivo: 0, 90, 180; sinistra è $-90$, $-180$.'' ``Se vai
in diagonale a destra sarà 45, se vai pienamente a destra sarà 90'', ``e 135 a 600 vuol
dire che andrai praticamente indietro, ma verso destra''. ``Una volta capito che una è
l'opposto dell'altra'', i valori negativi seguono da soli.

\paragraph{Approfondimento: Blend Space 2D e foot sliding.} In un Blend Space 2D Unreal
divide lo spazio dei campioni in triangoli e, per un valore dato, mescola i campioni ai
vertici del triangolo che lo contiene. Ne seguono tre conseguenze pratiche. Primo,
l'Idle ripetuto lungo tutta la riga inferiore (le tre copie del docente) evita che a
bassa velocità il personaggio ``scelga'' una direzione a caso. Secondo, la corsa
all'indietro va duplicata ai due estremi, perché $-180$ e 180 sono lo stesso angolo ma
due punti diversi della griglia: con un solo campione, passando da 179 a $-179$ il
personaggio attraverserebbe in un istante tutto il grafico. Terzo, lo scatto che il
docente vede è spesso proprio quel salto da $+180$ a $-180$, che lo \emph{Smoothing
Time} non può interpolare bene (interpola il valore numerico, non l'angolo); le
alternative sono il \emph{Weight Speed} per campione, oppure le soluzioni più avanzate
di Unreal Engine 5 come l'\emph{Orientation Warping}, che ruota la parte inferiore del
corpo verso la direzione del movimento con quattro sole animazioni cardinali. Infine,
una griglia con l'asse verticale fino a 600 e i campioni tutti a 600 significa che anche
la camminata a 300 è un misto fra Idle e corsa: per avere una vera camminata laterale
servirebbe una seconda riga di campioni a 300, con le animazioni di camminata.

% ---------------------------------------------------------------------
\section{Gli Animation Montage}
\label{ue9:sec:montage}

``Intanto ci abbiamo un'ultima mezz'oretta, ragazzi. Vediamo un attimo i Montage.''
L'animazione di attacco (``l'avevo scaricata, mi sa che non l'ho messa in Knight
Animation''; ``scusate, ragazzi'') è una \textbf{Sword and Shield Attack} di Mixamo:
``questa qua sta in root motion'': ``come potete vedere, eredita il movimento''.

\paragraph{Il problema.} ``Quando vogliamo fare una nostra animazione che di conseguenza
sia attivabile durante un ciclo di animazioni'', non basta aggiungere uno stato: ``dobbiamo
creare quello che viene chiamato un \textbf{Montage}''. (Nel frattempo ``Filippo è andato
via, è crashato'': ``siamo rimasti noi tre''.)

\paragraph{Che cos'è un Montage, dalla documentazione.} Il docente apre la pagina della
documentazione di Unreal sugli \textbf{Animation Montage}. ``Il montaggio di animazioni ci
permette sia di montare delle animazioni\dots'' L'esempio della pagina è la ricarica di
un fucile: un'animazione formata da tre animazioni diverse, \emph{MM Shotgun Reload Drop},
\emph{Reload Load} e \emph{Reload Lock}. ``Mettiamo conto che vogliamo fare un bel gioco
shooter. Il fatto di staccare il caricatore su quattro fucili diversi è praticamente
uguale''; ``non dobbiamo fare quattro animazioni diverse'': ``l'unica cosa che cambia è il
lock, perché magari uno vuole farlo un po' più grosso''. Con un Montage si creano
\textbf{animazioni modulari}: ``guarda, il drop è uguale, il load è uguale, il lock invece
è diverso''. E poi il secondo uso, quello di oggi: ``inserire animazioni in contesti di
cicli di animazione''; ``sono una tantum'': ``premo un tasto e compio questa animazione,
nel mentre che sto camminando''. ``Magari la prossima volta andiamo a vedere pure un po'
più approfondito.''

\paragraph{Creare il Montage.} Tasto destro sull'animazione di attacco $\rightarrow$
\textbf{Create} $\rightarrow$ \textbf{Create AnimMontage}. Il nuovo asset si chiama come
l'animazione con il suffisso \emph{Montage} (\emph{Sword And Shield Attack\_Montage}).
``La cosa bella del Montage è che ci permette di fare un sacco di cose all'interno dei
suoi parametri.''

% ---------------------------------------------------------------------
\section{Dentro il Montage: Notify, sezioni e Slot}
\label{ue9:sec:notify}

\paragraph{I Notify.} ``Una cosa interessante da spiegare sono i \textbf{Notify}'': ``dei
punti, dei puntini, sulla linea del tempo dell'animazione'', che permettono di aggiungere
``un'azione che si triggera'': ``suoni, un VFX, qualsiasi cosa''. ``Quelli più generici
sono già fatti da Unreal, fortunatamente.'' Il docente fa clic destro sulla traccia dei
Notify, \textbf{Add Notify} $\rightarrow$ \textbf{Play Sound}, e come suono sceglie
``l'unico sound che abbiamo'', quello della piattaforma (un Sound Cue). Le tracce dei Notify
possono essere più d'una: con \emph{Add Notify Track} se ne crea un'altra, per esempio
una per i VFX e una per i suoni. In preview (tolto il loop, ``fammelo fare una volta''):
``appena il nostro character arriva a questo punto dell'animazione, compie il suono''.
``Mettiamo conto: in combattimento c'è un VFX, c'è un suono, c'è un tirare fuori la spada,
e quindi si sente la spada che esce fuori dal fodero'': con una linea di Notify ``questo
suono avviene in un preciso momento dell'animazione''. ``Ovviamente, non andremo a usare
il sound effect della piattaforma per il nostro Knight.''

\paragraph{Le Montage Section.} Nel \emph{Montage group} c'è la sezione \textbf{Default}.
Con il tasto destro, \textbf{New Montage Section}: il docente ne crea una chiamata
\textbf{Attack}, nel punto in cui ``lui sta attaccando''. Le sezioni servono ``per
ordinare il lavoro'' e per ``dividere l'animazione in più punti diversi, con più
animazioni''.

\paragraph{Più animazioni nello stesso Montage.} La striscia verde sotto le sezioni è
l'animazione stessa. ``Se vogliamo, possiamo anche introdurre delle animazioni in più.''
Una dimostrazione volutamente assurda: ``io voglio, per qualche motivo arcano, far sì
che quando finisce l'animazione lui salti'': si trascina \textbf{MM\_Jump} nella traccia,
dopo l'attacco. ``È un salto orribile, però farà le due animazioni di conseguenza.'' Il
Montage è ``molto importante per azioni consecutive''. ``Quando si parla con l'animatore,
è importante il fatto di dire: questa animazione deve essere consecutiva a questa'', per
esempio ``mi serve che da questo punto in poi continui l'attacco, oppure faccia un'altra
animazione che ritorna in place''.

\paragraph{L'animazione di ritorno.} Finito l'attacco in root motion, il cavaliere resta dove l'attacco lo ha lasciato e non
torna in posa: manca ``un'animazione che ritorna in place'', cioè in idle. Il docente la
cerca su Mixamo, provando voci come \emph{Sword and Shield Idle}, \emph{Idle
Transition} e \emph{Start Block}; ma ``è un po' difficile'', perché ``bisogna trovare
il cavaliere proprio'', con la stessa posa di fine attacco, altrimenti ``snapperà
malissimo, almeno nel mio caso''.

\paragraph{Gli Slot.} Sotto, c'è il \textbf{Default Group} con lo \textbf{Slot}
\textbf{DefaultSlot}: ``questo ci permette di dire: su quale slot vuoi che questa
animazione si compia?''. DefaultSlot ``riprende tutto il corpo''. ``Potremmo, ad esempio,
voler creare uno slot che riprenda solamente alcune zone del nostro character'': ``vogliamo
che faccia un'azione solo nella parte superiore del corpo'' (sparare o salutare mentre si
corre); ``in quel caso creeremo un altro slot, che prenderà solo le ossa'' di quella
parte. ``Ma in questo caso ci serve tutto il corpo.''

% ---------------------------------------------------------------------
\section{Lo Slot nell'AnimGraph: Save Cached Pose}
\label{ue9:sec:slot}

Un Montage suona solo se l'AnimGraph ha un nodo \textbf{Slot} con lo stesso nome. In
ABP\_Knight, nell'AnimGraph:
\begin{enumerate}
  \item dall'uscita della State Machine Locomotion si crea un \textbf{New Save Cached
        Pose}, chiamato \textbf{LocomotionCache}: ``andiamo a conservarci quello che ha
        fatto la State Machine'', ``come se tutto questo procedimento lo mettesse in una
        memoria'';
  \item si aggiunge \textbf{Use cached pose 'LocomotionCache'};
  \item da lì un nodo \textbf{Slot 'DefaultSlot'} (la trascrizione: ``default slot'');
  \item lo Slot va all'\textbf{Output Pose}.
\end{enumerate}
``Semplicemente mixa insieme i due'': lo slot del Montage e la cache della locomozione
(figura~\ref{ue9:fig:slot}). Quando non c'è nessun Montage in corso, lo Slot lascia
passare la locomozione; quando un Montage su DefaultSlot parte, la sostituisce. Alla
fine della lezione, a uno studente che non legge il nome del nodo (``vedo tutto
pixelato''), il docente conferma: ``tiri fuori Locomotion Cache'', poi ``default slot''.

\begin{figure}[htbp]
\centering
\begin{tikzpicture}[font=\footnotesize,
  ue9pose/.style={draw, thick, rounded corners=4pt, align=center, fill=gray!8,
                  minimum height=1.0cm, text width=2.3cm},
  ue9arr/.style={-{Stealth[length=2.5mm]}, very thick, black!70}]
  \node[ue9pose, fill=orange!12] (sm) at (0,0) {\textbf{Locomotion}\\{\scriptsize State
    Machine}};
  \node[ue9pose, fill=blue!8] (sc) at (3.1,0) {\textbf{Save Cached Pose}\\{\scriptsize
    LocomotionCache}};
  \node[ue9pose, fill=blue!8] (uc) at (6.2,0) {\textbf{Use Cached Pose}\\{\scriptsize
    LocomotionCache}};
  \node[ue9pose, fill=green!10] (sl) at (9.3,0) {\textbf{Slot}\\{\scriptsize
    'DefaultSlot'}};
  \node[ue9pose, fill=red!10, text width=1.8cm] (op) at (12.1,0) {\textbf{Output
    Pose}};
  \draw[ue9arr] (sm) -- (sc); \draw[ue9arr] (uc) -- (sl); \draw[ue9arr] (sl) -- (op);
  \draw[dashed, thick, gray] (sc.east) -- (uc.west);
  \node[ue9pose, fill=yellow!15, text width=3.2cm] (mt) at (9.3,-2.1)
    {\textbf{Anim Montage}\\{\scriptsize Sword and Shield Attack, slot
    DefaultGroup.DefaultSlot}};
  \draw[-{Stealth[length=2.5mm]}, thick, dashed] (mt) -- node[right, font=\scriptsize]
    {quando suona, sostituisce la posa} (sl);
\end{tikzpicture}
\caption{La catena di pose finale dell'AnimGraph di ABP\_Knight. La State Machine produce
la locomozione, che viene salvata in una cache e riletta; lo Slot DefaultSlot la lascia
passare, a meno che un Montage assegnato a quello slot non sia in esecuzione. Schema
ricostruito.}
\label{ue9:fig:slot}
\end{figure}

% ---------------------------------------------------------------------
\section{Far partire l'attacco: Play Montage}
\label{ue9:sec:playmontage}

\paragraph{Il tasto F e il nodo.} Per far partire il Montage, il docente non crea
un'altra Input Action: usa direttamente un evento di tastiera, \textbf{F} (``facciamo con
OK\dots\ no, facciamo F, F va benissimo''). Nel Blueprint del personaggio, il nodo
\textbf{Play Montage}:
\begin{itemize}
  \item \textbf{In Skeletal Mesh Component}: ``in quale Skeletal Mesh Component
        farlo'', cioè la \textbf{Mesh} del personaggio;
  \item \textbf{Montage to Play}: ``quello che abbiamo appena fatto'' (``ce n'è tanto solo
        uno'');
  \item \textbf{Play Rate}: ``possiamo modificare quanto è veloce l'animazione'';
  \item fra le uscite: \textbf{On Completed} (``quando sarà finita''), \textbf{On Blend
        Out} (``mentre finisce''), \textbf{On Interrupted} (``se fosse interrotta''),
        \textbf{On Notify Begin} e \textbf{On Notify End} (``quando un Notify inizia'' e
        ``quando finisce'').
\end{itemize}
``Se noi premiamo F, compiamo un Montage.''

\paragraph{Fermare il personaggio.} ``È orribile'': a parte la mancanza dell'animazione di
ritorno, durante l'attacco ``possiamo muovere il character tranquillamente''. Quindi, sul
tasto F, prima del Play Montage:
\begin{enumerate}
  \item \textbf{Set Max Walk Speed} a \textbf{0}: ``non dobbiamo farlo muovere'';
  \item \textbf{Get Player Controller} $\rightarrow$ \textbf{Disable Input}: ``andiamo a
        disabilitare gli input, qualsiasi''.
\end{enumerate}
Sulle uscite del Play Montage:
\begin{itemize}
  \item \textbf{On Completed}: Get Player Controller $\rightarrow$ \textbf{Enable Input}
        (``on completed solamente'');
  \item \textbf{On Blend Out}: \textbf{Set Max Walk Speed} a \textbf{600}, ``al suo
        massimo''.
\end{itemize}

\paragraph{La camera che si muove ancora.} In Play, però, ``la camera si muove
effettivamente'', e con la camera il personaggio, ``quindi dobbiamo far sì che non sia
dipendente dalla camera''. Segue una serie di tentativi, ragionati ad alta voce:
\begin{itemize}
  \item togliere \emph{Use Pawn Control Rotation} sullo Spring Arm (``no, no'');
  \item l'\emph{Enable Camera Lag} dello Spring Arm;
  \item le spunte \emph{Inherit Pitch}, \emph{Inherit Yaw} e \emph{Inherit Roll} dello
        Spring Arm (``inerit yo e roll''; ``no, lo yaw ci deve essere''; ``forse sì,
        solo il roll va bene''), con un \emph{Set} preso dallo Spring Arm (``Set Inherit
        Yaw'', a false, e di nuovo a true alla fine); ``completamente\dots\ non va'';
  \item \emph{Stop Movement Immediately} e i parametri di input del Player Controller
        (\emph{movement input}, \emph{look input}).
\end{itemize}
La soluzione, alla fine, è sulle due spunte già incontrate nella prima parte: ``non è lo
yaw'' dello Spring Arm, ``ma è l'orientamento al movimento''. Sul tasto F si disattivano
entrambe:
\begin{itemize}
  \item dal \textbf{Character Movement}, \textbf{Set Orient Rotation to Movement} a
        \textbf{false};
  \item sul Pawn stesso, senza collegare nessun target, \textbf{Set Use Controller
        Rotation Yaw} a \textbf{false};
\end{itemize}
e sull'\textbf{On Completed} si rimettono entrambe a \textbf{true}
(figura~\ref{ue9:fig:playmontage}). ``Ottimo. Adesso riusciamo'': il cavaliere attacca
fermo, la camera gira senza trascinarlo, e alla fine riprende la locomozione. ``Abbiamo
fatto il nostro character che attacca e fa la locomotion.'' (Il docente nota anche, di
passaggio, che il cavaliere non interagisce con il pulsante del livello, ``giustamente,
perché è la Mesh''; dalla trascrizione non si ricostruisce meglio il punto.)

A Federico, che non riesce a sistemare la camera, il docente riassume: ``devi giocare con
lo yaw e con l'Orient Rotation to Movement: li disattivi tutte e due, e li riattivi
sull'On Completed''. ``Dello Spring Arm, giusto?'' ``Non dello Spring Arm: questo qua è
del Pawn, lo trovi direttamente senza collegare niente. Invece questo è solo del
Character Movement.''

\begin{figure}[htbp]
\centering
\begin{tikzpicture}[font=\footnotesize,
  ue9node/.style={draw, thick, rounded corners=4pt, align=center, fill=gray!8,
                  minimum height=0.9cm, text width=2.5cm},
  ue9ev/.style={ue9node, fill=red!12},
  ue9exec/.style={-{Stealth[length=2.5mm]}, very thick, black!80}]
  \node[ue9ev, text width=1.4cm] (f) at (0,0) {\textbf{F}\\{\scriptsize Pressed}};
  \node[ue9node] (a) at (2.8,0) {\textbf{Set Max Walk Speed}\\{\scriptsize 0}};
  \node[ue9node] (b) at (5.9,0) {Get Player Controller $\rightarrow$ \textbf{Disable
    Input}};
  \node[ue9node] (c) at (9.0,0) {\textbf{Set Orient Rotation to Movement}\\
    {\scriptsize false}};
  \node[ue9node] (d) at (12.1,0) {\textbf{Set Use Controller Rotation Yaw}\\
    {\scriptsize false}};
  \node[ue9node, fill=yellow!15, text width=3.0cm, minimum height=2.4cm] (pm) at
    (2.8,-3.0) {\textbf{Play Montage}\\{\scriptsize Mesh,\\Sword and Shield
    Attack\_Montage}};
  \node[ue9node] (e) at (8.3,-2.2) {\textbf{Set Max Walk Speed}\\{\scriptsize 600}};
  \node[ue9node, text width=3.0cm] (g) at (8.3,-3.9) {Get Player Controller
    $\rightarrow$ \textbf{Enable Input}};
  \node[ue9node, text width=3.0cm] (h) at (12.3,-3.9) {\textbf{Set Orient\dots} true\\
    \textbf{Set Use Controller Rotation Yaw} true};
  \draw[ue9exec] (f) -- (a); \draw[ue9exec] (a) -- (b); \draw[ue9exec] (b) -- (c);
  \draw[ue9exec] (c) -- (d);
  \draw[ue9exec] (d.south) -- ++(0,-0.6) -| (pm.north);
  \draw[ue9exec] (pm.east |- e) -- node[above, font=\scriptsize]{On Blend Out} (e);
  \draw[ue9exec] (pm.east |- g) -- node[below, font=\scriptsize]{On Completed} (g);
  \draw[ue9exec] (g) -- (h);
\end{tikzpicture}
\caption{L'attacco sul tasto F, nella versione finale. Prima del Montage si blocca il
personaggio (velocità zero, input disabilitato, nessuna rotazione automatica del corpo);
al Blend Out torna la velocità, all'On Completed tornano l'input e le rotazioni. L'ordine
esatto dei primi quattro nodi è ricostruito: a lezione vengono aggiunti uno alla volta,
durante i tentativi. Schema ricostruito.}
\label{ue9:fig:playmontage}
\end{figure}

\paragraph{Approfondimento: Play Montage, Montage Play e il root motion.} Il nodo usato a
lezione, \textbf{Play Montage}, è un nodo \emph{asincrono} (con l'orologio nell'angolo):
le sue uscite scattano più tardi, quando succede l'evento corrispondente, e per questo
può gestire il ``dopo'' del Montage direttamente nel grafo. Esiste anche \emph{Montage
Play}, chiamato sull'\emph{Anim Instance} della mesh, che restituisce solo la durata, e
\emph{Play Anim Montage}, del Character. Due dettagli completano la lezione. Il primo:
un attacco in \textbf{root motion} sposta davvero la capsula solo se il Montage ha il
root motion attivo e l'Animation Blueprint ha la \emph{Root Motion Mode} impostata su
\emph{Root Motion from Montages Only} (il default); se il root dell'animazione di Mixamo è
``rotto'', come visto nel retargeting, il personaggio avanza nella mesh e poi ``torna
indietro'' di scatto a fine animazione. Il secondo: durante un Montage che fa avanzare il
personaggio in root motion non serve azzerare la Max Walk Speed, perché il root motion
prende il posto dell'input di movimento; azzerarla, come fa il docente, è però un modo
semplice e sicuro di impedire al giocatore di muoversi durante un attacco ``in place''.
Infine, \emph{Disable Input} e \emph{Enable Input} agiscono sull'actor e richiedono un
Player Controller: per bloccare solo alcune azioni (il movimento ma non la camera),
l'Enhanced Input offre un'alternativa più fine, cioè rimuovere temporaneamente un
Input Mapping Context, o aggiungerne uno di priorità più alta.

% ---------------------------------------------------------------------
\section{Il compito e la chiusura}
\label{ue9:sec:chiusura}

\paragraph{Feedback in chat.} Una studentessa, Giulia, ha scritto in chat che è rimasta un
po' indietro ``causa trasloco'' e chiede di poter mandare dei feedback; nel frattempo è
già uscita dalla chiamata. ``Le risponderò sul Discord, assolutamente.'' ``Controllerò i
lavori che mandate sul canale.''

\paragraph{Il compito: lo StackOBot.} ``Una cosa che vorrei darvi come task, se vi va, per
la prossima volta'': ``fare un nuovo character totalmente da zero''. Nel progetto c'è già
un personaggio, lo \textbf{StackOBot} (il robottino dei progetti di esempio di Epic), che
``ha queste animazioni qua: vedete che ce le ha già dentro tutte quante''. Lo si trova
nella cartella \emph{StackOBot $\rightarrow$ Blueprints $\rightarrow$ Character}, con la
sua Skeletal Mesh (\emph{SKM\_Bot}) e il suo ABP. Il compito:
\begin{itemize}
  \item creare un \textbf{Blueprint Character} per lo StackOBot (``fare la BP'');
  \item creare il suo \textbf{Animation Blueprint} da zero, ``fare tutto l'Animation
        Blueprint e collegarci tutte le sue animazioni'': il ciclo di locomozione, con
        State Machine e Blend Space, come quello del cavaliere.
\end{itemize}
``Se volete farlo, lo potete mandare poi sulla chat'', ``che tanto magari ci lavoreremo
la prossima volta''. ``Per adesso, poi ne riparliamo.''

\paragraph{I saluti.} ``Se tutto a posto\dots'' Il docente interrompe la condivisione:
``vi ringrazio per l'attenzione, e ringrazio da casa chi è in differita''. ``Ci sentiamo
per martedì prossimo. Va bene, ciao, ragazzi. Buona serata.''

\begin{table}[htbp]
\centering
\small
\begin{tabular}{@{}p{5.4cm}p{8.6cm}@{}}
\toprule
\textbf{Sintomo} & \textbf{Causa e rimedio} \\
\midrule
L'animazione dello stato non si vede nella preview & la State Machine non è collegata
all'Output Pose dell'AnimGraph \\
L'animazione si ferma dopo un ciclo & manca la spunta di loop nel nodo dell'animazione \\
Warning nella State Machine & transizioni senza condizioni: normale finché non si
scrivono le regole \\
Jump $\rightarrow$ Fall dà errore & regola vuota: attivare Automatic Rule Based on
Sequence Player in State \\
Il Blend Space 1D arriva solo a 100 & Maximum Axis Value da impostare a 600; poi
spostare i campioni \\
Animazioni ``mischiate'' a metà (a 200) & i campioni non corrispondono alle velocità
reali: legarli alla Max Walk Speed \\
Alt sul nodo non crea il Set & negli stati dell'AnimGraph le variabili si leggono e
basta \\
``Questo asset non ha schede ancorate'' & Window $\rightarrow$ Load Layout $\rightarrow$
Default Editor Layout, poi Open Full Blueprint Editor \\
In Play non si vede niente / visuale nel corpo & il Character non ha camera: aggiungere
Spring Arm e Camera (figlia) \\
Non ci si muove & mancano gli input: copiare dal ThirdPersonCharacter i grafi di Look,
Move, Jump e l'Add Mapping Context \\
La camera ruota su un solo asse, o guarda solo il terreno & Use Pawn Control Rotation
sulla Camera invece che sullo Spring Arm \\
Il personaggio va di lato guardando avanti & attivare Orient Rotation to Movement e
togliere Use Controller Rotation Yaw (servono tutte e due) \\
La camera compenetra gli oggetti & Target Offset troppo grande: usare il Socket Offset \\
Il salto in root motion retargettato non è pulito & il root degli scheletri di Mixamo:
accettabile per la prova, da correggere \\
La State Machine non distingue salto e caduta & entrambe le regole su IsFalling:
aggiungere VelocityZ maggiore o minore di zero \\
Lo Shift non fa correre & IA\_Run non aggiunta ai mapping di IMC\_Default \\
FInterp To non interpola & eseguito una volta sola: va eseguito a ogni frame (Tick,
timer o Timeline) \\
Direction quasi sempre a zero & Orient Rotation to Movement acceso, Use Controller
Rotation Yaw spento: invertire per il Blend Space 2D \\
Scatti nel Blend Space 2D & animazioni non omogenee, salto fra 180 e $-180$: Smoothing
Time, o riscaricare le animazioni \\
Durante l'attacco il personaggio si muove & Set Max Walk Speed a 0 e Disable Input prima
del Play Montage \\
Durante l'attacco la camera trascina il personaggio & Set Orient Rotation to Movement e
Set Use Controller Rotation Yaw a false, a true sull'On Completed \\
Il Montage non suona & manca il nodo Slot DefaultSlot nell'AnimGraph \\
\bottomrule
\end{tabular}
\caption{I problemi incontrati durante la lezione, con la causa e il rimedio.}
\label{ue9:tab:errori}
\end{table}

\begin{table}[htbp]
\centering
\small
\begin{tabular}{@{}p{4.8cm}p{9.2cm}@{}}
\toprule
\textbf{Nodo, asset o impostazione} & \textbf{Che cosa fa} \\
\midrule
State Machine (AnimGraph) & organizza le animazioni in stati e transizioni \\
Entry / State / Conduit / Alias & ingresso; stato con una posa; snodo senza posa; alias
di più stati \\
Output Animation Pose & la posa prodotta da uno stato \\
Transition Rule (Result) & la condizione booleana per passare da uno stato all'altro \\
Automatic Rule Based on Sequence Player in State & la transizione scatta a fine
animazione \\
NOT Boolean / AND Boolean & negazione; congiunzione di due condizioni \\
Split Struct Pin & divide un vettore nei suoi componenti X, Y, Z \\
Blend Space 1D / Blend Space (2D) & mescola animazioni su uno o due assi \\
Horizontal / Vertical Axis, Grid Divisions & nome, minimo e massimo degli assi;
divisioni della griglia \\
Smoothing Time & tempo di passaggio fra i campioni \\
Blend Space Player & il nodo che suona un Blend Space nell'AnimGraph \\
Calculate Direction & l'angolo, da $-180$ a 180, fra la velocità e la rotazione \\
Get Actor Rotation / Get Velocity & la rotazione dell'actor; la sua velocità \\
Default Pawn Class (World Settings) & il personaggio del giocatore nel livello \\
Spring Arm: Target Arm Length, Socket Offset, Target Offset & lunghezza del braccio;
spostamento in fondo al braccio; spostamento dell'origine \\
Camera: Field of View & angolo di visuale \\
Use Pawn Control Rotation (Spring Arm) & la camera segue il controller \\
Use Controller Rotation Yaw (Pawn) & il corpo segue la rotazione del controller \\
Orient Rotation to Movement, Rotation Rate (Character Movement) & il corpo si gira verso
il movimento, con quella velocità \\
Max Walk Speed / Set Max Walk Speed & velocità massima a piedi (cm/s) \\
Retarget Animations, Export Animations & trasferisce un'animazione su un altro
scheletro \\
Input Mapping Context (IMC) & associa i tasti alle azioni in un contesto \\
Input Action (IA): Value Type, Triggers & l'azione, il tipo di valore, quando scatta \\
Add Mapping Context & attiva un IMC per il giocatore \\
Evento IA: Triggered, Started, Ongoing, Canceled, Completed & le fasi di un'azione \\
FInterp To & un passo di interpolazione fra Current e Target \\
Get World Delta Seconds & il tempo dell'ultimo frame \\
Animation Montage & animazione ``una tantum'', modulare, con sezioni e slot \\
Notify / Play Sound / Notify Track & eventi sulla linea del tempo; un suono; una traccia
di eventi \\
Montage Section & una parte con nome del Montage \\
Slot (DefaultGroup.DefaultSlot) & dove il Montage si applica (tutto il corpo, o parti) \\
Save Cached Pose / Use Cached Pose & salva e rilegge una posa nell'AnimGraph \\
Play Montage & suona un Montage; On Completed, On Blend Out, On Interrupted, Notify
Begin/End \\
Disable Input / Enable Input & toglie e restituisce l'input all'actor \\
\bottomrule
\end{tabular}
\caption{I nodi, gli asset e le impostazioni della nona lezione.}
\label{ue9:tab:nodi}
\end{table}

% ---------------------------------------------------------------------
\section*{Punti chiave del capitolo}
\addcontentsline{toc}{section}{Punti chiave del capitolo}

\begin{itemize}
  \item Una \textbf{State Machine} nell'AnimGraph organizza le animazioni in
        \textbf{stati} (WalkRun, Jump, Fall), ciascuno con la sua \textbf{Output
        Animation Pose}, e in \textbf{transizioni} con regole booleane: un ``loop di
        animazione secondo delle condizioni''. Va collegata all'Output Pose.
  \item Un \textbf{Blend Space 1D} mescola animazioni lungo un asse: per la locomozione,
        la \textbf{GroundSpeed}. I campioni vanno posti alle velocità reali del
        personaggio (Max Walk Speed, in centimetri al secondo): Idle a 0, Walk a 300,
        corsa a 600. Valori intermedi ``inventati'' producono obbrobri.
  \item Un \textbf{Character} non ha la camera: si aggiungono uno \textbf{Spring Arm} e,
        come suo figlio, una \textbf{Camera}. Distanza e altezza si regolano sullo
        Spring Arm (\textbf{Target Arm Length}, \textbf{Socket Offset}), mai spostando la
        camera; il \textbf{Target Offset} solo sulla Z e poco. Sono il lavoro dei
        \textbf{3C}: Camera, Control, Character.
  \item Tre spunte governano la rotazione: \textbf{Use Pawn Control Rotation} sullo
        Spring Arm (la camera segue il mouse), \textbf{Use Controller Rotation Yaw} sul
        Pawn (il corpo segue la camera), \textbf{Orient Rotation to Movement} nel
        Character Movement (il corpo si gira verso il movimento).
  \item Il \textbf{retargeting} porta un'animazione su un altro scheletro (Retarget
        Animations, Export Animations): comodo, ma il risultato dipende dall'animazione,
        e il root di Mixamo dà problemi.
  \item Fra salto e caduta: \textbf{Automatic Rule Based on Sequence Player in State} per
        passare a fine animazione, \textbf{NOT IsFalling} per tornare a terra, e la
        \textbf{velocità sulla Z} per distinguere un salto (Z $>$ 0) da una caduta
        (Z $<$ 0).
  \item L'\textbf{Enhanced Input} separa i tasti dalle azioni: l'\textbf{Input Action}
        (IA\_Run, Digital bool, con i suoi Triggers) è associata a uno o più tasti
        (Shift, il cerchio del controller) in un \textbf{Input Mapping Context}, che va
        attivato con \textbf{Add Mapping Context}. Lo sprint: Set Max Walk Speed a 600
        sul Triggered, a 300 sul Completed.
  \item \textbf{FInterp To} fa un solo passo per esecuzione: per un passaggio graduale
        va eseguito a ogni frame (Tick, timer, Timeline).
  \item Il \textbf{Blend Space 2D} aggiunge la \textbf{direzione} (\textbf{Calculate
        Direction}, da $-180$ a 180: positivo a destra, negativo a sinistra, $\pm 180$
        indietro): otto animazioni a 600, l'Idle ripetuto a 0, la corsa all'indietro ai
        due estremi, griglia 8$\times$8, Smoothing Time. Serve Use Controller Rotation Yaw
        acceso.
  \item Un \textbf{Animation Montage} è un'animazione ``una tantum'' e modulare, con
        \textbf{Notify} (suoni, VFX in un punto preciso), \textbf{sezioni} e
        \textbf{slot}. Per suonarlo, nell'AnimGraph: State Machine, \textbf{Save Cached
        Pose}, Use Cached Pose, \textbf{Slot DefaultSlot}, Output Pose. Si fa partire con
        \textbf{Play Montage}, bloccando prima movimento, input e rotazioni e
        restituendoli su On Blend Out e On Completed.
  \item Compito facoltativo: un nuovo personaggio da zero con lo \textbf{StackOBot},
        Blueprint e Animation Blueprint completi.
\end{itemize}

%
%  Tutte le etichette di questo capitolo hanno il prefisso "ue9:"
%  (ue9:cap:..., ue9:sec:..., ue9:fig:..., ue9:tab:...), e tutti gli
%  stili TikZ hanno il prefisso "ue9" (ue9node, ue9exec, ue9data, ...)
%  per non collidere con chiavi native di TikZ né con quelli delle
%  lezioni precedenti. Il capitolo è autonomo: non contiene \ref verso
%  altri capitoli (i rimandi alle lezioni precedenti sono a parole).
%  Le figure sono ricostruzioni schematiche di quanto mostrato e detto a
%  lezione, non copie delle schermate del docente. Dove la trascrizione
%  è troppo rovinata per ricostruire un dettaglio, il capitolo lo dice
%  invece di inventarlo.
% =====================================================================

\chapter[State Machine, Blend Space e Montage]{State Machine, Blend Space e Montage:
dare vita al nuovo personaggio}
\label{ue9:cap:animazioni}

Questa è la \textbf{nona lezione} del corso breve di \textbf{introduzione a Unreal
Engine}, e riprende esattamente dove si era fermata l'ottava. Alla fine della lezione
precedente il cavaliere importato da \textbf{Mixamo} aveva uno scheletro, le sue
animazioni e un \textbf{Animation Blueprint}, \textbf{ABP\_Knight}, che a ogni frame
leggeva dal personaggio due dati: la velocità sul piano (\textbf{GroundSpeed}) e il fatto
di essere in aria o meno (\textbf{IsFalling}). Mancava tutto il resto: chi decide
\emph{quale} animazione mostrare, e quando. ``Siccome ci siamo interrotti l'altra volta,
andiamo diretti a fare le animazioni: l'altra volta ci siamo interrotti al Blend
Space.''

La lezione costruisce, un pezzo alla volta, un personaggio giocabile completo. La
\textbf{State Machine} organizza le animazioni in \emph{stati} (camminare e correre,
saltare, cadere) e in \emph{transizioni} con delle condizioni; il \textbf{Blend Space
1D} mescola fermo, camminata e corsa in base alla velocità; una \textbf{camera} con il
suo \textbf{Spring Arm}, i grafi di movimento presi in prestito dal modello Third Person
e un nuovo input per lo \textbf{sprint}, costruito con l'\textbf{Enhanced Input},
rendono il cavaliere controllabile; il \textbf{retargeting} porta sul suo scheletro le
animazioni di salto e di caduta del manichino di Unreal. Dopo la pausa arrivano il
\textbf{Blend Space 2D}, che aggiunge alla velocità la \textbf{direzione} e permette di
camminare di lato e all'indietro, e gli \textbf{Animation Montage}, le animazioni
``una tantum'' (un attacco con la spada) che si inseriscono sopra il ciclo di
locomozione.

La lezione è \textbf{online}, su Zoom, come le precedenti, e la classe è ridottissima:
``vediamo un po' quanti siete: siamo in sei''; dopo la pausa ``siamo rimasti in
cinque'', poi ``uno, due, tre e quattro''; quando uno studente crasha (``Filippo è andato
via''), ``siamo rimasti noi tre''; alla fine, ``mi sembra che siamo rimasti in due''.
``Pochi ma buoni.'' Il ritmo è quello della dimostrazione dal vivo, interrotta più volte
dai problemi degli studenti, risolti in condivisione di schermo, e da un editor che
``oggi lagga'' in modo vistoso (``non capisco perché mi sta laggando oggi''). La pausa
arriva dopo la scelta delle animazioni da scaricare (``ci rivediamo alle 18''); si
chiude con un compito facoltativo e un ``ci sentiamo per martedì prossimo''.

Tre avvertenze di lettura, le stesse delle lezioni precedenti. La prima: la
trascrizione è automatica e storpia moltissimi nomi. Unreal è ``Arreal'', ``a Real'',
``ArrayL'', ``Array Alt'' e ``a Rial''; il \textbf{Pawn Owner} è ``Pwn Honor'' e
``Pwnowner''; il \textbf{Blend Space} è ``Blaine Space'', ``Blendspace'' e ``BlizzBase'';
lo stato \textbf{WalkRun} è ``walk run'', ``Warcran'', ``war crunch'' e ``walker run'';
l'\textbf{Idle} è ``Idol'' e ``idol''; \textbf{BP\_Knight} è ``Vip Night'', ``Pippi
Night 1'' e ``BP Night''; \textbf{BS\_Knight} è ``BS Night'' (``che dovrebbe essere un
po' più British''); il manichino \textbf{Manny/Quinn} è ``Many Quincy''; il
\textbf{root motion} è ``rut'' e ``outmotion''; lo \textbf{Strafe} è ``Strife'' e
``stair right''; la \textbf{Mesh} è ``Mesci''; l'\textbf{FInterp To} è ``finter2'';
\textbf{Calculate Direction} è ``calcolo interaction''; il \textbf{Save Cached Pose} è
``save cache, pause''; lo \textbf{StackOBot} è ``StaccoBot'' e ``stacco bot''; lo
\textbf{Yaw} è ``io'', ``yo'' e ``YAW''; l'\textbf{Orient Rotation to Movement} è
``orientation to movement'', ``anche rotacion'' e ``orient rotace tu mente tu''; lo
\textbf{Use Pawn Control Rotation} è ``us control pawn rotation'' e ``you spawn control
rotation''. Dove la ricostruzione è sicura il capitolo usa la forma corretta; dove non
lo è riporta la forma della trascrizione e lo dice. La seconda: all'inizio della
registrazione e durante la pausa il riconoscimento vocale produce lunghi blocchi di
testo senza senso e ripetuto (``grazie a tutti per l'attenzione che ci ha spiegato'',
``qual è la risposta che ci può essere?'', ``la situazione è questa, la migliore dal
2005'', la ``dottrina sul mercato del lavoro'', la ``Corte Suprema di Internet'', il
governo di Dubai, una lunga fila di ``di Sarajevo'', la ``struttura del giardino'', ``la
blue map'', un Vittorio Prodi chiamato in causa senza motivo): sono artefatti, non
parole di nessuno, e il capitolo li ignora. La terza: gran parte della lezione è una
dimostrazione dal vivo nell'editor, costruita a voce (``tirate fuori'', ``promote to
variable'', ``ce lo attacchiamo qui''), e in più punti il docente improvvisa, prova,
sbaglia e corregge, soprattutto nella parte sulla camera e in quella finale sul
Montage. Il capitolo ricostruisce i grafi da ciò che il docente dice mentre li
costruisce, distingue i tentativi falliti dalla soluzione finale, e completa i nomi
esatti dei nodi e dei pin di Unreal Engine 5; dove la trascrizione lascia un buco, lo
dichiara.

\paragraph{Come leggere questo capitolo.} Il capitolo segue l'ordine della lezione. Le
figure ricostruiscono la State Machine della locomozione, il Blend Space 1D sulla
velocità, la camera con lo Spring Arm, la catena dell'Enhanced Input, il grafo dello
sprint, il cerchio della direzione, il Blend Space 2D con le sue otto animazioni, la
catena di pose dell'AnimGraph con lo Slot del Montage e il grafo del tasto che fa
partire l'attacco. Le tabelle riassumono le transizioni della State Machine, i
parametri della camera, i valori e i trigger di una Input Action, la posizione delle
animazioni nel Blend Space 2D, gli errori incontrati durante la lezione e i nodi della
giornata. Qua e là, dei paragrafi di \textbf{approfondimento} collegano la lezione ai
libri di Marcos Romero e Brenden Sewell (\emph{Blueprints Visual Scripting for Unreal
Engine 5} e la seconda edizione dello stesso testo, scritta per Unreal Engine 4), i
testi di riferimento per questa parte del libro, e a qualche nozione tecnica generale su
Unreal Engine 5; sono segnalati come tali e non fanno parte della lezione.

Il percorso del capitolo:
\begin{enumerate}
  \item l'\textbf{inizio della lezione}: i corsi degli studenti e la lezione libera del
        25 (sezione~\ref{ue9:sec:inizio});
  \item il \textbf{ripasso} dell'Animation Blueprint (sezione~\ref{ue9:sec:ripasso});
  \item la \textbf{State Machine} della locomozione (sezione~\ref{ue9:sec:statemachine});
  \item il \textbf{Blend Space 1D}, con una parentesi sui problemi dell'editor (sezioni
        \ref{ue9:sec:blend1d} e~\ref{ue9:sec:layout});
  \item il personaggio in gioco: \textbf{Default Pawn Class}, \textbf{camera} e
        \textbf{Spring Arm}, grafi di movimento e rotazioni (sezioni
        da~\ref{ue9:sec:pawnclass} a~\ref{ue9:sec:rotazioni});
  \item il \textbf{retargeting} del salto e della caduta, e le \textbf{regole di
        transizione} (sezioni \ref{ue9:sec:retarget} e~\ref{ue9:sec:transizioni});
  \item lo \textbf{sprint} con l'\textbf{Enhanced Input} e il tentativo con FInterp To
        (sezioni da~\ref{ue9:sec:enhanced} a~\ref{ue9:sec:finterp});
  \item le \textbf{animazioni da scaricare}, la pausa e i \textbf{feedback} (sezioni
        \ref{ue9:sec:download} e~\ref{ue9:sec:feedback});
  \item il \textbf{Blend Space 2D} e la \textbf{direzione} (sezioni
        \ref{ue9:sec:direction} e~\ref{ue9:sec:blend2d});
  \item gli \textbf{Animation Montage}: l'asset, i Notify, le sezioni, gli Slot e il
        tasto che li fa partire (sezioni da~\ref{ue9:sec:montage}
        a~\ref{ue9:sec:playmontage});
  \item il \textbf{compito} sullo StackOBot e la chiusura
        (sezione~\ref{ue9:sec:chiusura}).
\end{enumerate}

% ---------------------------------------------------------------------
\section{Prima della lezione: i corsi degli studenti e la lezione del 25}
\label{ue9:sec:inizio}

La registrazione si apre con qualche minuto di attesa, in cui il riconoscimento vocale
produce soprattutto rumore. Poi i saluti (``ragazzi, come state? Buonasera, buonasera''),
un controllo rapido (``nessun problema sulla telecamera'': stavolta la webcam del docente
funziona) e la decisione di aspettare ``altri cinque minuti, così vediamo che tutti siate
in carreggiata''. Uno studente, Tommaso, ha dei problemi di collegamento. Il docente
sistema le finestre (``oggi c'è il secondo schermo'') e parte: ``bando alle ciance''.

\paragraph{I corsi di indirizzo.} Prima di cominciare, il docente chiede ai presenti
quale corso seguiranno ``quando inizierete a settembre'': il corso di Unreal è infatti
un'introduzione comune, a cui seguono gli indirizzi specialistici. Le risposte dei sei
presenti:
\begin{itemize}
  \item \textbf{3D art};
  \item \textbf{game design} (due studenti);
  \item \textbf{narrative design};
  \item \textbf{character animation} (``perfetto, character animation'').
\end{itemize}
La domanda serve a tarare la lezione: chi farà character animation, in particolare, è
il destinatario naturale di quasi tutto quello che si vedrà oggi.

\paragraph{La lezione libera del 25.} ``Stiamo praticamente arrivando a una lezione che
di solito noi lasciamo abbastanza \emph{free}, come argomenti'': la lezione del
\textbf{25} sarà ``fondamentalmente la mia ultima lezione di base''. Se non ci sono
richieste particolari, si faranno ``i menu'', cioè ``la \textbf{UI}'': ``che si intende
la UI: menu, bottoni, scritte eccetera''. In alternativa, ``se ci sono degli argomenti
particolari'', da concordare fra gli studenti, ``possiamo tranquillamente affrontarli'',
oppure ``dividere anche la lezione in due: fare un piccolo menu e poi fare anche la parte
in più''. ``Nel caso, sentitevi tra di voi.'' La richiesta concreta arriverà durante la
pausa (sezione~\ref{ue9:sec:feedback}).

% ---------------------------------------------------------------------
\section{Il ripasso: l'Animation Blueprint del cavaliere}
\label{ue9:sec:ripasso}

``Eravamo rimasti al nostro bellissimo \textbf{ABP\_Knight}.'' Il docente ha il
cavaliere (\emph{Knight}, il paladino di Mixamo); ``altri di voi avevano scelto il
proprio personaggio''. Nella preview dell'Animation Blueprint il personaggio è ancora
``il nostro bel character in T-pose'': l'AnimGraph non decide ancora nessuna
animazione.

\paragraph{Initialize e Update.} L'ultima cosa fatta era l'evento \textbf{Event Blueprint
Initialize Animation}: ``quando l'animazione si inizializza'', una volta sola. L'altro
evento, \textbf{Event Blueprint Update Animation}, scatta invece ogni volta che
``si cambia qualsiasi particolare all'interno di un'animazione, la rotazione di un braccio
o un minimo particolare'': ``quello che è fondamentalmente il nostro Tick sulle
animazioni''. ``Quando abbiamo visto Begin Play e Tick: questi praticamente sono i loro
corrispettivi.''

\paragraph{Il Pawn Owner e le due referenze.} Nell'Initialize si prende il \textbf{Pawn
Owner} con \textbf{Try Get Pawn Owner}. ``Ricordiamoci: il Pawn Owner è l'entità che
controlla il character'': il docente distingue il \emph{pawn}, ``quello che lo controlla'',
e il \emph{character} stesso, ``che ha la mesh, lo scheletro e tutto quanto''. Con il
\textbf{Cast To} BP\_Knight si ottiene una referenza al cavaliere, e si salvano in due
variabili il cavaliere e il suo \textbf{Character Movement Component}: ``queste due
variabili ci serviranno per comunicare in modo più diretto con il nostro cavaliere''.

\paragraph{Il Character Movement.} ``Dentro il Blueprint, qualsiasi character lo troverete
qui sotto, Character Movement'': selezionandolo, il pannello Details a destra ``presenta
tutte le particolarità di movimento'', per esempio la \textbf{gravità} (\emph{Gravity
Scale}), la \textbf{Max Acceleration}, la \textbf{Max Step Height} (l'altezza massima di
un gradino che il personaggio sale senza saltare). ``Vogliamo farlo andare più veloce,
meno veloce? Se volessimo farlo sprintare, dovremmo lavorare su questa particolarità''; e
infatti ``dopo faremo anche uno sprintino molto veloce, per differenziare le
animazioni''.

\paragraph{GroundSpeed e IsFalling.} Nell'Update, dopo una \textbf{Sequence}, si prende
``la velocity attuale del nostro character'', a ogni frame, e se ne salva la lunghezza in
\textbf{GroundSpeed}. Allo stesso modo, con il nodo \textbf{Is Falling} del Character
Movement, ``andiamo a vedere se effettivamente il nostro character sta cadendo o meno'':
``questo blocco ci dice letteralmente se sta cadendo, e noi riportiamo con un booleano,
true o false''. ``Di conseguenza, tramite questi dati, andremo a differenziare le
animazioni nella nostra State Machine.''

% ---------------------------------------------------------------------
\section{La State Machine della locomozione}
\label{ue9:sec:statemachine}

\paragraph{Creare la State Machine.} La State Machine ``la possiamo invocare
nell'AnimGraph'': tasto destro nell'AnimGraph (``ricordatevi che c'è il blocco
\textbf{Output Pose}''), \textbf{State Machine}, e la si chiama \textbf{Locomotion}:
``questo è il nome classico per la locomozione, appunto il movimento''. ``Una State
Machine, appunto, è una macchina a stati.'' Con un doppio clic si entra nel suo grafo.

\paragraph{Entry, Conduit, State, Alias.} Dentro c'è un solo nodo, \textbf{Entry}. ``Se
noi trasciniamo, abbiamo una freccia che esce fuori''; lasciandola nel vuoto, il menu non
è ``tutto l'elenco gigantesco'' dei nodi dei Blueprint, ma tre sole voci:
\textbf{Conduit}, \textbf{State} e \textbf{Alias}. ``A noi adesso interessa vedere
quelli che sono gli stati.'' Il primo stato, invece di \emph{Locomotion}, si chiama
\textbf{WalkRun}: ``differenziamo''.

\paragraph{Un'altra inception.} ``Ecco a voi che, ancora sempre, un po' un
\emph{inception}'': con un doppio clic si entra anche nello stato. Dentro, un nodo
\textbf{Output Animation Pose}: ``in questo stato risiede una tipologia di output'',
l'animazione che lo stato produce. Il docente vi collega \emph{Walking 1}. La preview
però non si aggiorna, ``perché non abbiamo collegato la State Machine all'output'':
bisogna tornare nell'AnimGraph e collegare il nodo Locomotion all'Output Pose. Fatto
questo, il cavaliere cammina, ma una volta sola: ``l'animazione non è in loop'', e basta
spuntare il loop nei Details del nodo. ``Ecco a voi che cosa fa la State Machine'':
``dentro Locomotion lui entra, dice ok, inizio, passo e vado direttamente a WalkRun''; in
WalkRun ``c'è una richiesta di un output, quindi di un'animazione che noi gli diamo''.

\paragraph{Gli stati Jump e Fall.} ``Ma la cosa interessante della State Machine'' è che
si possono creare altri stati: \textbf{Jump} e \textbf{Fall}. ``A livello di logica, la
State Machine vuole essere'' il racconto di ``come entrano e perché entrano'' le
animazioni nei loro stadi. Una freccia da WalkRun a Jump ``vuol dire che dallo stadio
WalkRun passerà allo stadio Jump''; una freccia inversa vuol dire che dal Jump ``potrà
passare di nuovo al WalkRun''. Lo schema di oggi ha tre stati
(figura~\ref{ue9:fig:statemachine}): si entra in WalkRun; ``se starà cadendo'' (la
variabile IsFalling diventa vera) si passa a Jump o a Fall; ``e far sì che poi, quando
lui finisca di cadere, ritorni'' a WalkRun. ``Quindi fondamentalmente creiamo un loop di
animazione secondo delle condizioni.'' Le animazioni di salto e di caduta non ci sono
ancora: ``la faremo, la ritargetteremo'' (sezione~\ref{ue9:sec:retarget}).

\paragraph{Le transizioni.} ``Come facciamo a mettere queste condizioni? Tramite questi
nodini qua'': sulle frecce fra due stati compare un'icona tonda, la \textbf{transition
rule}, che ``ci dice il comportamento effettivo che dovrà avere'' il personaggio. ``Quando
è che passo da WalkRun a Jump, a livello di animazione? Quando, ad esempio, stai cadendo.
Quando passo da Fall a WalkRun? Se non sta più cadendo. Questa è la logica generale.''

\begin{figure}[htbp]
\centering
\begin{tikzpicture}[font=\footnotesize,
  ue9state/.style={draw, thick, rounded corners=6pt, fill=orange!12, align=center,
                   minimum width=2.4cm, minimum height=1.0cm},
  ue9entry/.style={draw, thick, fill=gray!15, rounded corners=2pt, align=center,
                   minimum width=1.4cm, minimum height=0.7cm},
  ue9tr/.style={-{Stealth[length=2.6mm]}, thick},
  ue9rule/.style={draw, circle, fill=white, inner sep=1.5pt, font=\scriptsize}]
  \node[ue9entry] (e) at (0,0) {Entry};
  \node[ue9state] (w) at (3.4,0) {\textbf{WalkRun}\\{\scriptsize Blend Space}};
  \node[ue9state] (j) at (8.2,1.5) {\textbf{Jump}\\{\scriptsize MM\_Jump}};
  \node[ue9state] (f) at (8.2,-1.5) {\textbf{Fall}\\{\scriptsize MM\_Fall\_Loop}};
  \draw[ue9tr] (e) -- (w);
  \draw[ue9tr] (w.north east) to[bend left=12] node[ue9rule, pos=0.5]{1} (j.west);
  \draw[ue9tr] (j.south) -- node[ue9rule]{2} (f.north);
  \draw[ue9tr] (w.south east) to[bend left=12] node[ue9rule, pos=0.5]{3} (f.west);
  \draw[ue9tr] (f.south west) to[bend left=30] node[ue9rule, pos=0.5]{4} (w.south);
  \node[align=left, text width=12.5cm, anchor=north west, font=\scriptsize] at
    (-1,-2.9) {\textbf{1} WalkRun $\rightarrow$ Jump: IsFalling \textbf{AND}
    VelocityZ $>$ 0 (sta salendo)\\
    \textbf{2} Jump $\rightarrow$ Fall: \emph{Automatic Rule Based on Sequence Player in
    State} (finito il salto, la caduta)\\
    \textbf{3} WalkRun $\rightarrow$ Fall: IsFalling \textbf{AND} VelocityZ $<$ 0 (cade
    senza aver saltato)\\
    \textbf{4} Fall $\rightarrow$ WalkRun: \textbf{NOT} IsFalling (è tornato a terra)};
\end{tikzpicture}
\caption{La State Machine Locomotion alla fine della lezione: tre stati e quattro
transizioni. Le condizioni con la velocità sulla Z sono quelle aggiunte nella seconda
versione (sezione~\ref{ue9:sec:transizioni}); nella prima, la 1 e la 3 erano soltanto
IsFalling. Schema ricostruito.}
\label{ue9:fig:statemachine}
\end{figure}

\paragraph{Approfondimento: Conduit e Alias.} I tre tipi di nodo proposti dal menu della
State Machine hanno ruoli diversi. Uno \textbf{State} contiene un grafo che produce una
posa (un'animazione, un Blend Space, una combinazione). Un \textbf{Conduit} non produce
pose: è un incrocio, uno snodo da cui partono più transizioni, utile quando molti stati
condividono la stessa condizione d'ingresso o d'uscita e non si vuole ripeterla su ogni
freccia. Un \textbf{State Alias} è un ``alias'' di uno o più stati: le transizioni che
partono dall'alias valgono come se partissero da ciascuno degli stati che rappresenta
(per esempio, ``da qualunque stato, se il personaggio muore, vai a Death''), e
risparmiano un groviglio di frecce. Romero e Sewell (capitolo 17) costruiscono la loro
prima State Machine, \emph{Char States}, con due soli stati, \emph{Idle} e
\emph{Moving}, e due regole sulla velocità (``Speed maggiore di 5'' e ``minore di 5''); la
scelta della lezione, un solo stato WalkRun che contiene un Blend Space, sposta la
distinzione fra fermo e in movimento dentro il Blend Space, e lascia alla State Machine
le sole differenze ``qualitative'' (a terra, in salto, in caduta).

% ---------------------------------------------------------------------
\section{Il Blend Space 1D}
\label{ue9:sec:blend1d}

``Visto che WalkRun non è solo un'animazione, è un insieme'': sono tre animazioni,
\textbf{Idle}, \textbf{Walk} e quella che il docente ha chiamato \textbf{Slow Run}.
``Dobbiamo mixarle insieme. Come si fa? Perché non possiamo mettere un'animazione solo
così: dobbiamo far sì che sia dipendente dalla nostra velocità.'' Lo strumento è il
\textbf{Blend Space}.

\paragraph{Creare il Blend Space 1D.} ``Fate solamente il Blend Space 1D.'' Invece di
cercarlo nei menu, si può scrivere nella barra di ricerca del menu del tasto destro
``blend space'', e scegliere \textbf{Blend Space 1D}: ``già vi dice il percorso, che sta
in Animation, Legacy''. L'icona ha un colore ``rosa pelle''. Aperto, si troverebbe il
personaggio in T-pose; ma molti l'avevano già creato la volta precedente (``sì, mi sa che
già l'avevamo fatto'').

\paragraph{Che cos'è un Blend Space.} ``Un Blend Space è l'azione di \emph{blendare}
qualche cosa'', di ``interpolare o comunque mixare le animazioni'': un asset che ``ci dà
la possibilità di poter traslare da un'animazione all'altra''. Si vede subito: dal
pannello \textbf{Asset Browser} dell'editor (``se non lo trovate, ripeto: Window, Asset
Browser'') si trascina un'animazione sul grafico; tenendo premuto \textbf{Shift} la si
aggancia ``su dei punti precisi'' della griglia. Con un'animazione sola, il personaggio
fa solo quella; aggiungendone una seconda, il grafico ``crea una linea di passaggio tra
un'animazione e l'altra''.

\paragraph{Il parametro, come l'alfa.} ``Ma importante: qual è questo dato?'' Il docente
lo collega all'\textbf{alfa} di un Lerp visto in una lezione precedente: ``quel puntatore
che punta, dipendente da uno spazio tra due punti, dove sta al momento''. Qui il valore
dirà ``ehi, guarda che sto camminando, sto andando a velocità 300, oppure sto correndo, a
velocità 600, oppure sono fermo''. Il Blend Space \textbf{1D} ``ci dà solamente la
possibilità di poter usare un asse'': nei Details, \textbf{Horizontal Axis} (la sezione
potrebbe essere chiusa: ``basta che la aprite''). L'asse si chiama \textbf{GroundSpeed},
``come quella che abbiamo trovato all'interno dell'Animation Blueprint''; poi chiede un
\textbf{Minimum Axis Value} e un \textbf{Maximum Axis Value}.

\paragraph{Le velocità sono relative.} ``Importante, ragazzi, quando si va a vedere le
animazioni: dovete sempre farlo relativo, se si parla di velocità'', relativo cioè alla
velocità reale del personaggio. Dove si legge? Nel Blueprint del personaggio,
BP\_Knight, selezionando \textbf{Character Movement}: ``potete andare a vedere tutti i
dati che vi servono''. Quello che serve è la \textbf{Max Walk Speed}, in alto, nella
sezione \emph{Character Movement: Walking}: ``è 600 centimetri al secondo''. ``Unreal
ragiona per centimetri'': il personaggio va a 6 metri al secondo.

\paragraph{Dividere la velocità a metà.} ``Visto che noi abbiamo due stadi, uno di walk e
uno di run'', si vuole che ``la walk sia un movimento normale'': di default il personaggio
cammina, e corre solo se il giocatore lo chiede con un tasto (sezione
\ref{ue9:sec:enhanced}). ``Quindi la dividiamo a metà'': \textbf{Max Walk Speed a 300}.
Nel Blend Space l'asse va da \textbf{0} a \textbf{600}, ``in modo tale che a metà la
nostra walk sia a 300'' (figura~\ref{ue9:fig:blend1d}). A uno studente l'asse arriva solo
a 100: ``ground speed, devi metterlo a 600, lo trovi qua''. E un consiglio: dopo aver
cambiato il massimo, ``sposta i pallini, perché poi ti switcha tutto''; il valore di
ciascun campione si può anche scrivere a mano nei Details.

\begin{figure}[htbp]
\centering
\begin{tikzpicture}[font=\footnotesize,
  ue9pt/.style={draw=black, fill=white, thick, diamond, inner sep=2.2pt}]
  \draw[-{Stealth[length=2.5mm]}, thick] (0,0) -- (11.2,0) node[right]{GroundSpeed};
  \foreach \x/\v in {0/0, 2.5/150, 5/300, 7.5/450, 10/600}
    \draw (\x,0.1) -- (\x,-0.1) node[below, font=\scriptsize]{\v};
  \node[ue9pt] (i) at (0,0) {}; \node[ue9pt] (w) at (5,0) {}; \node[ue9pt] (r) at
    (10,0) {};
  \node[above=6pt of i, font=\scriptsize\bfseries] {Idle};
  \node[above=6pt of w, font=\scriptsize\bfseries, align=center] {Walk\\
    {\normalfont Max Walk Speed di default}};
  \node[above=6pt of r, font=\scriptsize\bfseries, align=center] {Slow Run\\
    {\normalfont con lo sprint}};
  \fill[red!60!black] (3.33,0) circle (0.12);
  \draw[red!60!black, thick, ->] (3.33,-1.2) -- (3.33,-0.25);
  \node[red!50!black, font=\scriptsize, align=center] at (3.33,-1.55) {200: un misto
    fra Idle e Walk, ``un obbrobrio''};
\end{tikzpicture}
\caption{Il Blend Space 1D BS\_Knight: un solo asse, la GroundSpeed, da 0 a 600, con
Idle a 0, Walk a 300 (la Max Walk Speed impostata nel Character Movement) e la corsa a
600 (la velocità dello sprint). Un valore intermedio, come 200, mescola due animazioni
nate per velocità diverse. Schema ricostruito.}
\label{ue9:fig:blend1d}
\end{figure}

\paragraph{Il Blendspace Player.} Tornati in ABP\_Knight (``normale, ragazzi, che vi dia
dei warning, perché ancora non abbiamo le condizioni'': ``lui dice, ma io non ho
condizioni, quindi faccio tutto in una botta''), si entra nello stato WalkRun. Nell'Asset
Browser ``è apparso BS\_Knight''; trascinato nel grafo, ``ci presenta un \textbf{Blend
Space Player}'', diverso dal nodo di un'animazione normale: ha un \textbf{valore in
entrata}, la GroundSpeed. A quel pin si collega un \textbf{Get} della variabile
GroundSpeed, calcolata nell'Event Graph (la lunghezza della velocity), e l'uscita del
Blend Space Player si collega al \textbf{Result} dell'Output Animation Pose, ``in modo
tale che vada in output''.

\paragraph{Get con Ctrl, Set con Alt.} Un promemoria: trascinando una variabile nel
grafo, ``tenete premuto \textbf{Ctrl} e tirate fuori direttamente un Get''; ``se invece
volete fare un Set, tenete premuto \textbf{Alt}''. Nel grafo dello stato l'Alt non
funziona: ``non si può settare all'interno'' (gli stati dell'AnimGraph leggono le
variabili, non le scrivono), ``però nei Blueprint, di base, se premete Alt vi setta la
variabile''.

\paragraph{La prova nella preview.} Compilato, ``il nostro character nella preview adesso
sta in idle, perché giustamente ha velocità 0''. Cambiando il valore di GroundSpeed nella
preview: a \textbf{300} cammina; a \textbf{200} ``sta interpolando tra l'animazione di
idle e l'animazione di walk'': ``mettere un valore a 200 non è buona cosa'', perché
``deve essere relativo al vostro movimento, se no vengono degli obbrobri''; a
\textbf{600} ``il nostro carissimo corre'' (``guarda come è eroico''). ``Quindi adesso
andiamo a cambiare le animazioni tramite un valore.''

\paragraph{Allineare i nodi con Q.} ``Se volete fare un po' di ordine'': selezionati due
o più nodi, il tasto \textbf{Q} li allinea (\emph{Straighten Connections}). ``Tutto bello
ordinato.''

\paragraph{Una prova con un solo stato.} Prima di andare avanti, un test in gioco:
nell'AnimGraph si tolgono le frecce verso Jump e Fall (selezionate e cancellate, o
staccate con \textbf{Alt}+clic sul pin), e si lascia ``solamente WalkRun, che vada in
entrata''.

\paragraph{Approfondimento: il Blend Space nel libro.} Romero e Sewell (capitolo 17)
definiscono il Blend Space come ``un tipo di asset che permette di mescolare animazioni
in base ai valori di uno o due parametri''; nel loro esempio, tratto dal modello Third
Person, Idle è a 0, Walk a 93,75 e Run a 375, con una Max Walk Speed di 375. Il principio
è lo stesso della lezione: i campioni vanno posti dove la velocità reale del personaggio
corrisponde alla velocità ``disegnata'' nell'animazione. Se un'animazione di camminata è
stata animata per una velocità di 150 centimetri al secondo e il personaggio si muove a
300, i piedi ``scivolano'' sul pavimento (\emph{foot sliding}); è il motivo pratico
dietro il ``deve essere relativo'' del docente. Gli autori ricordano anche che il valore
del parametro si può cambiare nell'\emph{Anim Preview Editor}, in basso a destra
nell'editor dell'Animation Blueprint, senza andare in Play: è ciò che il docente fa per
mostrare i valori 200, 300 e 600.

% ---------------------------------------------------------------------
\section{Una parentesi: ``questo asset non ha schede ancorate''}
\label{ue9:sec:layout}

Mentre si sistema la Max Walk Speed, uno studente interviene: ``io ho un problema col
Blueprint: quando apro la scheda Blueprint mi dice: questo asset non ha ancora schede
ancorate''. Il docente ha l'editor in inglese (``di solito gli errori me li ricordo in
inglese'') e chiede di vedere: lo studente condivide prima soltanto una scheda del
browser, poi lo schermo intero.

\paragraph{Il messaggio.} Aprendo BP\_Knight, l'editor del Blueprint è vuoto, senza
pannelli. Il testo completo del messaggio, letto ad alta voce, spiega che l'editor
dell'asset è stato aperto in un sistema di finestre fluttuanti e nascoste, ``non mostra
le schede ancorate nell'interfaccia principale'', e che ``questo comportamento può essere
facilmente ripristinato'' ripristinando il layout dell'editor o modificando le preferenze
perché gli asset si aprano sempre ancorati alla finestra principale (in inglese,
\emph{Editor Preferences $\rightarrow$ General $\rightarrow$ Appearance $\rightarrow$
Asset Editor Open Location}). Inoltre, sopra il grafo compare una nota: il Blueprint è
\emph{data-only}, e c'è un link \textbf{Open Full Blueprint Editor}.

\paragraph{La soluzione.} ``Mi sa che è meglio ripristinare il layout di default'':
dalla barra in alto, \textbf{Window} (in italiano \emph{Finestra}) $\rightarrow$
\textbf{Load Layout} (\emph{Carica il layout}) $\rightarrow$ \textbf{Default Editor
Layout} (\emph{Editor predefinito}). Poi, sulla nota, ``apri l'editor Blueprint
completo''. ``Perfetto, eccolo qua.''

\paragraph{La lingua dell'editor.} ``Comunque, per mettere l'editor in inglese'' senza
riavviare: si cerca \emph{Language} nelle preferenze. Il docente prova prima in
\textbf{Editor Preferences} (sotto \emph{Edit}), scrivendo ``lingua'' e poi ``language''
(``non gli piace che mettiamo due cose insieme''), poi nei \textbf{Project Settings};
alla fine la voce \emph{Language} viene trovata con la ricerca (``ah, ecco perché'').
Dalla trascrizione non si capisce con certezza in quale delle due finestre. Lo studente
ringrazia (``mi hai fatto sentire molto ignorante''; ``verrai buttato fuori dalla
classe, tranquillo''), e il docente annuncia di volerlo usare in inglese anche lui ``per
la prossima volta''.

\paragraph{Approfondimento: dov'è la lingua.} In Unreal Engine 5 la lingua
dell'interfaccia è in \textbf{Edit $\rightarrow$ Editor Preferences $\rightarrow$ General
$\rightarrow$ Region \& Language $\rightarrow$ Editor Language}; nei Project Settings c'è
invece la lingua del \emph{gioco} (la localizzazione), che è un'altra cosa. Lavorare
con l'editor in inglese è una scelta pratica: la documentazione, i tutorial, i libri e i
forum usano i nomi inglesi dei nodi e dei pannelli, e la ricerca dei nodi nel menu del
tasto destro accetta i nomi inglesi anche quando l'interfaccia è tradotta.

% ---------------------------------------------------------------------
\section{Il cavaliere in gioco: la Default Pawn Class}
\label{ue9:sec:pawnclass}

``Adesso vi faccio vedere come impostiamo il player nuovo.'' Si aprono i \textbf{World
Settings} del livello (``qualsiasi cosa che non trovate, c'è sempre Window, e troverete
World Settings''); fra le voci della Game Mode c'è \textbf{Default Pawn Class}.
``Ricordiamo che il Pawn è appunto l'entità che controlla il player.'' Nel menu a
tendina c'è BP\_Knight, il personaggio finito che il docente aveva mostrato alla fine
della lezione precedente, e c'è BP\_Knight1, ``quello nuovo'', che si può anche cercare
scrivendone il nome. Si sceglie BP\_Knight1 e si va in Play.

\paragraph{Il primo problema: nessuna camera.} ``Ecco già un problemino.'' ``Fate
attenzione, perché i character non vengono con la camera'': un Blueprint creato dalla
classe Character ha la capsula, la freccia, la mesh e il Character Movement, ma non una
camera. In Play, la visuale finisce in un punto qualsiasi del corpo (``questo
meraviglioso\dots'', la frase si perde). Si esce dal Play e si apre BP\_Knight1: ``come
potete vedere, non c'è una camera''.

% ---------------------------------------------------------------------
\section{La camera: Camera e Spring Arm}
\label{ue9:sec:camera}

``La camera è impostata su due cose particolari, ragazzi, molto importanti'': la
\textbf{camera} stessa e il suo \textbf{braccio}. ``Immaginatelo come lo stand'': il
braccio dice ``quanto è lontana'' la camera, ``quanto offset ha''.

\paragraph{Aggiungere i due componenti.} Si aggiunge un componente \textbf{Camera}
(``camera normale''): ``eccola qua, che si è già attaccata'' al personaggio, ma in una
posizione che non va bene (``al suo lato B''). Poi uno \textbf{Spring Arm} (il docente
lo chiama prima ``arm length'', poi si corregge: ``perdonate, l'arm length è un
parametro; lo Spring Arm''). La gerarchia è importante: \textbf{lo Spring Arm deve avere
la camera dentro di sé}, cioè la Camera va trascinata sopra lo Spring Arm e diventa sua
figlia. Selezionando lo Spring Arm, nel viewport compare ``un piccolo tracciato rosso'':
``la lunghezza del nostro braccio'', il braccio della camera
(figura~\ref{ue9:fig:springarm}).

\paragraph{I 3C.} ``Ricordatevi, ragazzi, molto importante, perché la camera è
importantissima: ci sono delle figure che si chiamano \textbf{3C}, \emph{Camera, Control
e Character}, che lavorano solo su questo.'' Da qui una regola: ``se volete allontanare la
camera o cambiarle l'offset, \textbf{non cambiate la location della camera}. Non lo fate.
Controllatelo sempre da qui, dallo Spring Arm.''

\paragraph{Target Arm Length.} Il parametro \textbf{Target Arm Length} ``ci permette di
aumentare o diminuire la distanza della camera''. Valore scelto: \textbf{300}.

\paragraph{Target Offset e Socket Offset.} Ci sono due offset, e il docente ha una
preferenza netta. Il \textbf{Target Offset} ``usatelo solamente sulla Z, di poco,
veramente di poco, massimo 20''; ``e fate tutto con il \textbf{Socket Offset}''. Il
motivo: ``se voi cambiate l'offset del target, la camera vi può più facilmente
compenetrare con le cose''. Con il Socket Offset, invece, ``vedrete che la camera va sul
muro e torna'': ``va più vicino'', cioè ``riduce il suo arm length, fondamentalmente''.
Dopo la prima prova in Play (``lo testiamo, ragazzi?''), la camera è un po' alta: il
Socket Offset sulla Z viene abbassato a \textbf{40}.

\begin{figure}[htbp]
\centering
\begin{tikzpicture}[font=\footnotesize]
  % capsula
  \draw[thick, fill=blue!8, rounded corners=14pt] (0,0) rectangle (1.1,2.6);
  \node[align=center, font=\scriptsize] at (0.55,1.1) {Cap-\\sule};
  \node[font=\scriptsize, below] at (0.55,-0.2) {BP\_Knight1};
  \fill (0.55,2.0) circle (0.06);
  \node[font=\scriptsize, anchor=north west] at (0.62,1.95) {pivot};
  % braccio
  \draw[very thick, red!70!black] (0.55,2.0) -- (-4.2,2.0);
  \node[red!60!black, font=\scriptsize, above] at (-1.8,2.0) {Target Arm Length
    = 300};
  % socket offset
  \draw[thick, dashed, -{Stealth[length=2mm]}, green!40!black] (-4.2,2.0) --
    (-4.2,2.8);
  \node[green!40!black, font=\scriptsize, anchor=east, align=right] at (-4.9,2.3)
    {Socket Offset\\Z = 40};
  % camera
  \draw[thick, fill=gray!30] (-4.75,2.6) rectangle (-3.95,3.05);
  \draw[thick, fill=gray!30] (-3.95,2.7) -- (-3.6,2.55) -- (-3.6,3.1) -- (-3.95,2.95);
  \node[font=\scriptsize, above] at (-4.3,3.1) {Camera (figlia)};
  % target offset
  \draw[thick, dashed, -{Stealth[length=2mm]}, orange!70!black] (0.55,2.0) --
    (0.55,2.35);
  \node[orange!60!black, font=\scriptsize, anchor=south, align=center] at (0.55,2.4)
    {Target Offset: solo Z,\\al massimo $\approx$ 20};
  % muro
  \fill[gray!40] (-2.6,-0.2) rectangle (-2.35,1.4);
  \draw[thick] (-2.6,-0.2) rectangle (-2.35,1.4);
  \node[font=\scriptsize, align=center] at (-2.47,-0.6) {muro};
  \draw (-6,-0.2) -- (3.5,-0.2);
  % gerarchia
  \node[draw, thick, rounded corners=3pt, fill=gray!8, align=left, font=\scriptsize,
    anchor=north west] at (3.3,2.2) {\textbf{Components}\\
    Capsule Component\\
    \quad Arrow\\
    \quad Mesh\\
    \quad \textbf{Spring Arm}\\
    \qquad \textbf{Camera}\\
    Character Movement};
\end{tikzpicture}
\caption{La camera del cavaliere. La Camera è figlia dello Spring Arm, e la distanza e
l'altezza si regolano sul braccio, mai spostando la camera. Se fra il personaggio e la
camera si mette un ostacolo (il muro), lo Spring Arm accorcia il braccio e la camera
``torna'' verso il personaggio; con un Target Offset grande, invece, è più facile che la
camera entri negli oggetti. Schema ricostruito, valori della lezione.}
\label{ue9:fig:springarm}
\end{figure}

\paragraph{Il Field of View.} Più avanti, mentre aspetta uno studente, il docente mostra
che ``nella camera è presente anche un po' di dati con cui potete giocarci'', per esempio
il \textbf{Field of View}. Lo abbassa a \textbf{80}: ``sappiate che tanti giochi horror
usano il field of view a 80. Questo è un piccolo tip.''

\begin{table}[htbp]
\centering
\small
\begin{tabular}{@{}p{3.6cm}p{2.0cm}p{8.4cm}@{}}
\toprule
\textbf{Parametro} & \textbf{Valore} & \textbf{Effetto e consiglio del docente} \\
\midrule
Spring Arm $\rightarrow$ Target Arm Length & 300 & distanza della camera dal
personaggio \\
Spring Arm $\rightarrow$ Socket Offset & Z = 40 & spostamento della camera alla fine del
braccio; da preferire \\
Spring Arm $\rightarrow$ Target Offset & Z $\le$ 20 & spostamento dell'origine del
braccio; poco e solo sulla Z, perché fa compenetrare la camera \\
Spring Arm $\rightarrow$ Use Pawn Control Rotation & attivo & il braccio ruota con il
mouse (il controller); va sullo Spring Arm, non sulla Camera \\
Camera $\rightarrow$ Field of View & 80 & angolo di visuale; 80 è tipico dei giochi
horror \\
Camera $\rightarrow$ Location & non toccare & la posizione si regola sempre dallo Spring
Arm \\
\bottomrule
\end{tabular}
\caption{I parametri della camera toccati durante la lezione.}
\label{ue9:tab:camera}
\end{table}

\paragraph{Approfondimento: perché lo Spring Arm si accorcia.} Lo \textbf{Spring Arm
Component} fa, a ogni frame, una traccia (un \emph{sweep} di una piccola sfera, di
raggio \emph{Probe Size}) dal suo punto d'origine verso la posizione desiderata della
camera; se incontra un ostacolo nel canale \emph{Probe Channel} (di default
\emph{Camera}), accorcia il braccio fino al punto di impatto, e lo riallunga quando
l'ostacolo sparisce (la proprietà è \emph{Do Collision Test}). Il \emph{Target Offset}
sposta l'\emph{origine} di quella traccia, nello spazio del mondo: spostarla molto può
portare l'origine dentro un muro o dietro un oggetto, e allora il test non ``vede'' più
l'ostacolo; il \emph{Socket Offset} sposta invece la \emph{fine} del braccio, nello
spazio del braccio stesso, e la traccia resta coerente. È la ragione tecnica del
consiglio del docente. Lo Spring Arm ha anche le opzioni \emph{Enable Camera Lag} ed
\emph{Enable Camera Rotation Lag}, che fanno ``inseguire'' la camera con un leggero
ritardo e rendono più morbidi i movimenti: il docente le sfiora più tardi, nella sezione
del Montage. Quanto ai 3C, nell'industria la sigla indica la disciplina che cura come il
giocatore \emph{vede} (camera), \emph{comanda} (controlli) e \emph{abita} (personaggio)
il gioco: tre cose che si progettano insieme, perché ognuna cambia la percezione delle
altre.

% ---------------------------------------------------------------------
\section{I grafi di movimento, presi in prestito}
\label{ue9:sec:movimento}

In Play, la camera c'è, ma ``non ci possiamo muovere'': ``perché non abbiamo né input né
le funzioni di base di movimento''. ``Questo, ragazzi, non andiamo a farlo direttamente a
mano, ma andiamo a fare quella che è una bellissima'' scorciatoia (la trascrizione dice
``mandragata''): si apre il \textbf{BP\_ThirdPersonCharacter} classico, ``quello che
abbiamo usato'' finora, che contiene ``questi tre principali blocchi che Unreal fa di
base'': il controllo della \textbf{camera} (\emph{Look}), il controllo del
\textbf{movimento} (con \emph{Add Movement Input}) e il \textbf{salto} (\emph{Jump}).

\paragraph{Copiare anche l'Input Mapping.} ``Mi raccomando, prendete anche questo qua
sopra, che è l'\textbf{input mapping}'', cioè il grafo sull'Event Begin Play che aggiunge
l'Input Mapping Context (``andremo a spiegare che cosa si intende''). Tutto con
\textbf{Ctrl+C}, e poi \textbf{Ctrl+V} nell'Event Graph di BP\_Knight1: ``perché queste
sono le funzioni basiche di Unreal''.

\paragraph{IA\_Move.} Un'occhiata all'evento \textbf{IA\_Move} con l'Add Movement Input,
``proprio per fare una lezione sull'input'': ``questo evento dice che è legato a un
tasto, o a dei tasti, che ci riporta due valori, X e Y'', che ``di conseguenza diventano
dei vettori che vanno a creare il movimento sui due assi''. Uno è avanti e indietro
(\emph{forward}, \emph{backward}), l'altro è destra e sinistra.

\paragraph{La camera che ruota male.} Compilato, ``adesso la nostra camera si muove'', ma
in modo strano: si muove ``solo così'', su un asse. ``L'avevo visto un'altra volta'':
``devo usare il control rotation''. La spunta è \textbf{Use Pawn Control Rotation}:
``se lo attaccate, la camera vi va in \emph{free mode}'', ``nel senso che usa la
rotazione del controller'', sia in beccheggio sia in imbardata.

% ---------------------------------------------------------------------
\section{Tre spunte sulla rotazione}
\label{ue9:sec:rotazioni}

\paragraph{Use Pawn Control Rotation, sullo Spring Arm.} Uno studente ha la camera che
``quando vado in Play posso abbassarla e guardare solo il terreno, non vedo proprio
neanche più il personaggio''. Il docente pensa prima a una compenetrazione, poi si fa
condividere lo schermo: ``che camera hai messo? Camera normale, giusto?''. Guardando i
componenti, il problema salta fuori: ``non Use Pawn Control Rotation sulla camera: Use
Pawn Control Rotation sullo \textbf{Spring Arm}''. ``Quindi levalo da lì e mettilo sullo
Spring Arm.'' ``Anche l'altro?'' ``No, no, no: solo sullo Spring Arm.'' La spunta va su
\emph{un} componente solo, il braccio: è il braccio a ruotare intorno al personaggio, e
la camera, che ne è figlia, lo segue.

\paragraph{Orient Rotation to Movement.} In Play, il cavaliere cammina (alla velocità di
300), ma andando a destra ``lui cammina dritto, fondamentalmente'': si sposta di lato
continuando a guardare avanti. ``Questo lo possiamo risolvere con'' \textbf{Orient
Rotation to Movement}, sempre nel Character Movement: ``dal nome, orienta la rotazione al
movimento''. Accanto c'è la \textbf{Rotation Rate}, la velocità con cui il personaggio si
gira (``non è proprio precisissima, però ci piace'').

\paragraph{Use Controller Rotation Yaw.} Il personaggio però continua a non girarsi bene.
``Perché mi ricordo che qua c'era una particolarità che non ricordo'': il docente scorre
le categorie del Blueprint (\emph{Simulate}, \emph{Rotation}, \emph{Character},
\emph{Orientation}, \emph{Physics}, \emph{Controller}\dots) e la trova. ``Per fissare il
nostro character che non si girava, basta che eliminiamo il tick'': selezionato il Pawn
(la radice del Blueprint, \emph{BP\_Knight1 (Self)}), si cerca \emph{yaw} e si toglie
\textbf{Use Controller Rotation Yaw}, ``in modo tale che il character non usi il suo
yaw''. Senza questa spunta, il personaggio non ruota più insieme alla camera: ``andrà
dove siamo direzionati'', ``se io vado a destra, lui si gira completamente verso
destra''. A uno studente non basta: aveva tolto lo Yaw ma non attivato Orient Rotation to
Movement. ``Ok, non ce l'avevo.'' Con tutte e due, ``abbiamo il nostro bel character
fatto per benissimo''.

\paragraph{Approfondimento: chi comanda la rotazione.} Le tre spunte della lezione sono
le tre ``mani'' che possono ruotare un personaggio in terza persona, e vanno combinate
con criterio. \textbf{Use Pawn Control Rotation} (sullo Spring Arm) fa seguire alla
\emph{camera} la rotazione del controller, cioè del mouse o della levetta destra.
\textbf{Use Controller Rotation Pitch/Yaw/Roll} (sul Pawn) fa seguire quella stessa
rotazione al \emph{corpo} del personaggio. \textbf{Orient Rotation to Movement} (nel
Character Movement) fa girare il corpo verso la direzione in cui si muove. La
combinazione ``da avventura'' della prima parte della lezione è: camera libera (Use Pawn
Control Rotation sul braccio), corpo indipendente dalla camera (Use Controller Rotation
Yaw spento), corpo orientato al movimento (Orient Rotation to Movement acceso): è quella
del modello Third Person di Unreal. La combinazione ``da combattimento'' della seconda
parte, con il Blend Space 2D (sezione~\ref{ue9:sec:direction}), è opposta: il corpo
guarda sempre dove guarda la camera (Use Controller Rotation Yaw acceso), e le
animazioni di strafe mostrano lo spostamento laterale. Se sono accese insieme Use
Controller Rotation Yaw e Orient Rotation to Movement, le due ``mani'' si contendono il
personaggio; per questo, nel Character Movement, esiste anche \emph{Use Controller
Desired Rotation}, che fa ruotare il corpo verso la direzione della camera in modo
graduale, con la Rotation Rate (forse il ``desired'' che il docente prova a un certo
punto, senza esito chiaro nella trascrizione).

% ---------------------------------------------------------------------
\section{Il retargeting: il salto del manichino sul cavaliere}
\label{ue9:sec:retarget}

``Adesso, però, guardate una cosa: se ad esempio provo a saltare, non ho nessuna
animazione.'' Il salto funziona (lo fa il Character Movement), ma il cavaliere sale e
scende nella posa della camminata. Le animazioni di salto si potrebbero prendere da
Mixamo, ``anche se potremmo, volendo''; invece il docente usa ``l'opzione di Unreal che
si chiama \textbf{retargeting}''. ``Diciamo che il retargeting, sappiatelo, può uscire
bene, ma può uscire anche abbastanza male: è dipendente dall'animazione. Però andiamo a
farlo comunque, per vedere come si fa.''

\paragraph{I passaggi.}
\begin{enumerate}
  \item Nel Content Browser si cerca \emph{Jump}: si trova l'animazione del salto del
        manichino base di Unreal, \textbf{MM\_Jump}.
  \item Tasto destro sull'animazione $\rightarrow$ \textbf{Retarget Animations}.
  \item La finestra mostra lo scheletro \textbf{sorgente} (\emph{source}), quello da cui
        viene l'animazione: il manichino di Unreal, \textbf{Manny} o \textbf{Quinn}
        (``Many Quincy''). Come \textbf{target} si sceglie lo scheletro del cavaliere, il
        paladino di Mixamo: ``eccolo qua, li mette a confronto''.
  \item Si seleziona \textbf{MM\_Jump} nella lista delle animazioni: nella preview il
        paladino salta. ``Fortunatamente Mixamo usa degli scheletri che sono abbastanza
        compatibili con quello di Unreal'', e il risultato ``è venuto abbastanza bene'';
        ``potete girarci tranquillamente intorno con la viewport''.
  \item Due pulsanti di esportazione. \textbf{Export Retargeted Assets} ``non ci serve,
        ragazzi'': ``vi crea un asset retargettato'', una copia completa degli asset sullo
        scheletro di destinazione. Basta \textbf{Export Animations}: ``esportare
        l'animazione così, come nuova''. ``Ricordate che c'è il source, e vi esporta il
        target.''
  \item Come cartella si sceglie quella del personaggio (\emph{Characters $\rightarrow$
        Knight1 $\rightarrow$ Animation}) e si esporta.
\end{enumerate}
Lo stesso procedimento si ripete per la caduta: \textbf{MM\_Fall\_Loop}, un'animazione
``in loop'' di caduta, retargettata sul paladino ed esportata nella stessa cartella.

\paragraph{Il root motion ``rotto''.} Il risultato del salto ``non è fantastico'':
``il problema di Mixamo è che ha il \textbf{root motion} un po' rotto''. Il docente
mostra il \emph{root}, ``un punto dello scheletro'', l'osso alla base della gerarchia, e
osserva che lo spostamento non è quello atteso; ``potremmo vederlo'', ``ma per vedere come
funziona il retarget va benissimo''. Si potrebbe sistemare in seguito.

\paragraph{Approfondimento: il retargeting in Unreal Engine 5.} Il retargeting trasferisce
un'animazione da uno scheletro a un altro con proporzioni e nomi delle ossa diversi. In
Unreal Engine 5 si basa su due tipi di asset: l'\textbf{IK Rig}, che descrive uno
scheletro raggruppando le ossa in \emph{catene} (spina dorsale, braccia, gambe, testa)
con un osso di \emph{retarget root} (di solito il bacino), e l'\textbf{IK Retargeter},
che mette in corrispondenza le catene di due IK Rig e corregge la differenza fra le pose
di riferimento (la A-pose del manichino e la T-pose di Mixamo). Nelle versioni più recenti il
comando \emph{Retarget Animations} del menu contestuale apre una finestra che crea al
volo questi asset quando gli scheletri sono riconosciuti (come quelli di Mixamo e del
manichino), ed è ciò che si vede a lezione; per risultati migliori, o con scheletri
insoliti, si creano e si regolano a mano. Il problema del root di Mixamo è noto: gli
scheletri di Mixamo non hanno un osso di root a terra, e lo spostamento del personaggio
è registrato sull'osso dei fianchi (\emph{Hips}); un'animazione retargettata può quindi
avere il movimento nel posto ``sbagliato'' per il sistema di root motion di Unreal.

% ---------------------------------------------------------------------
\section{Le regole di transizione}
\label{ue9:sec:transizioni}

Con le due animazioni nuove si aggiorna l'Animation Blueprint: ``dentro il nostro Jump
ci mettiamo l'animazione di jump'' (``niente di particolare, ragazzi''), dentro Fall
l'animazione di caduta. Si ricollegano le frecce. ``Se cammina e corre, può saltare? Sì,
certo. Se salta, dovrà per forza almeno loopare in caduta'', e ``dal Fall dovrà ritornare
al nostro WalkRun''.

\paragraph{Prima versione.} Nelle icone delle transizioni c'è un grafo con un nodo
\textbf{Result}, che ``ci dà un bool'': ``ci dice quando è che posso entrare in questa
animazione''.
\begin{itemize}
  \item \textbf{WalkRun $\rightarrow$ Jump}: un \emph{Get IsFalling} collegato al
        Result. ``Potremmo farlo più preciso dopo''; per esempio, ``possiamo anche dire:
        ha iniziato il jump''.
  \item \textbf{Jump $\rightarrow$ Fall}: qui il grafo vuoto dà errore, ``perché noi
        vogliamo far sì che questo processo tra il Jump e il Fall sia automatico''.
        Selezionata l'icona della transizione, nei Details si spunta \textbf{Automatic
        Rule Based on Sequence Player in State}: ``quando inizia quell'animazione, a
        seguito ci sarà quest'altra''; in altre parole, la transizione scatta da sola
        quando l'animazione dello stato Jump finisce. ``Lo mettiamo così e non ci darà
        problemi.''
  \item \textbf{Fall $\rightarrow$ WalkRun}: ``vuol dire che non sta cadendo'': Get
        IsFalling, poi un \textbf{NOT Boolean}. ``Ricordo che il NOT Boolean è qualcosa
        che ci dà il suo contrario'': deve essere false perché la caduta finisca.
\end{itemize}
Uno studente chiede cosa ci sia fra Jump e Fall: ``se tu clicchi su questo pallino, ci
dice, in Details\dots\ vogliamo che questa animazione parta direttamente se effettivamente
il jump è finito'': \emph{Automatic Rule Based on Sequence Player in State}, ``questo
tick, lo metti vero''.

\paragraph{La prova, e un difetto.} Si prova da un punto alto, ``la nostra piattaforma''.
``Non si nota tantissimo, però effettivamente sta facendo giustamente l'azione di
saltare.'' Ma il docente vuole collegare WalkRun anche direttamente a Fall: il
personaggio ``può cadere anche nel caso in cui lui non abbia spinta, semplicemente cade da
qualche parte'', per esempio scendendo da un bordo. Con le sole regole su IsFalling,
però, le due transizioni in uscita da WalkRun avrebbero la stessa condizione, e Unreal
non saprebbe distinguere un salto da una caduta.

\paragraph{Seconda versione: la velocità sulla Z.} ``Dovremmo andare a differenziare un
altro valore: la Z. Se la Z è maggiore di 0, vuol dire che sta saltando; se la Z invece è
minore di 0, vuol dire che sta cadendo.'' Nell'Event Graph dell'ABP, finora, della
velocity si prendeva solo la lunghezza. Il docente prova prima la strada sbagliata (la
\emph{Jump Z Velocity} del Character Movement, che però è la spinta \emph{iniziale} del
salto, un valore fisso), poi prende la \textbf{Velocity} già salvata, ne fa lo
\textbf{Split Struct Pin} (``lo splittiamo'') e dalla \textbf{Z} fa \textbf{Promote to
Variable}: una nuova float, \textbf{VelocityZ} (``Z va benissimo''). Nelle transizioni:
\begin{itemize}
  \item \textbf{WalkRun $\rightarrow$ Jump}: IsFalling \textbf{AND} VelocityZ maggiore
        di zero, con un nodo \textbf{AND Boolean}. Il docente dice anche ``quando la
        velocità sarà maggiore di 1, più preciso'': una soglia piccola invece di zero
        esatto.
  \item \textbf{WalkRun $\rightarrow$ Fall}: IsFalling e VelocityZ minore di zero.
\end{itemize}
Per tenere in ordine i fili, si possono usare i \emph{reroute node} (``volete collegarli
con questo? È un po' bruttino''). Il risultato: ``se io salto, fa l'animazione''; ``se io
adesso invece cado solamente, lui va diretto in Fall''. Il docente chiude salvando anche
il BP\_ThirdPersonCharacter che l'editor segnala come modificato (``va bene, grazie'').

\begin{table}[htbp]
\centering
\small
\begin{tabular}{@{}p{3.2cm}p{5.3cm}p{5.5cm}@{}}
\toprule
\textbf{Transizione} & \textbf{Regola} & \textbf{Significato} \\
\midrule
Entry $\rightarrow$ WalkRun & nessuna & lo stato iniziale \\
WalkRun $\rightarrow$ Jump & IsFalling AND VelocityZ $>$ 0 & è in aria e sta salendo: ha
saltato \\
Jump $\rightarrow$ Fall & Automatic Rule Based on Sequence Player in State & finita
l'animazione di salto, passa alla caduta \\
WalkRun $\rightarrow$ Fall & IsFalling AND VelocityZ $<$ 0 & è in aria e sta scendendo:
è caduto da un bordo \\
Fall $\rightarrow$ WalkRun & NOT IsFalling & è tornato a terra \\
\bottomrule
\end{tabular}
\caption{Le regole di transizione della State Machine Locomotion (versione finale).}
\label{ue9:tab:transizioni}
\end{table}

\paragraph{Approfondimento: altre regole basate sul tempo.} L'opzione automatica della
lezione è il modo più rapido per dire ``quando l'animazione finisce''. Romero e Sewell
(capitolo 17), per gli stati \emph{StandToProne} e \emph{ProneToStand}, usano invece una
regola esplicita: il nodo \textbf{Time Remaining (ratio)} dell'animazione dello stato,
seguito da un \emph{Less} con 0,1: la transizione scatta quando manca meno del 10\%
dell'animazione. Il vantaggio è il controllo: la transizione può partire un po' prima
della fine e lasciare spazio al \emph{blend} fra i due stati, che ogni transizione ha
(nei Details: \emph{Blend Settings}, con una durata di default di 0,2 secondi).
Ricordano anche che le animazioni ``di una volta'', come il salto, vanno impostate senza
\emph{Loop Animation}: è il caso dello stato Jump, mentre Fall, che deve durare quanto
la caduta, resta in loop.

% ---------------------------------------------------------------------
\section{Lo sprint con l'Enhanced Input}
\label{ue9:sec:enhanced}

``Adesso andremo a mettere il nostro sprintino'': ``un semplicissimo sistemino che passa
la Max Walk Speed da un tot a un altro tot''. Si torna in BP\_Knight1, ``a ragionare a
livello di Blueprint''. ``Ma sì, facciamo anche l'input mapping.''

\paragraph{L'Input Mapping Context.} ``Unreal ha introdotto, quasi dalla 5.2'', quello
che è l'\textbf{Input Mapping Context}, ``una cosa molto interessante di Unreal che ci
permette, ad esempio, di avere degli eventi legati a più input''. Finora, in molte
lezioni, si sono usati gli eventi dei tasti (``abbiamo visto il tasto E, eccolo qua''):
``semplicemente un rilevamento, da parte di Unreal, della tastiera: se è premuto o non è
premuto''. Adesso, invece, ``gli eventi sono legati a più controlli possibili''.

Nel Content Browser si cerca \textbf{IMC} (\emph{Input Mapping Context}). ``È un
\emph{contesto}: un contesto di input che vengono dati in un certo momento, e sappiate
che sono sovrascrivibili.'' L'esempio: ``io devo cambiare dal gioco di avventura 3D al
gioco di guida della macchina''; ``i controlli sono completamente diversi, quindi cambio
contesto'': un Input Mapping Context per la macchina e uno ``per il character che cammina
normalmente''. Aperto \textbf{IMC\_Default}, ci sono i \textbf{Mappings}:
\textbf{IA\_Jump}, \textbf{IA\_Move}, \textbf{IA\_Look}.

\paragraph{L'Input Action IA\_Run.} ``Per far sì di avere un'Input Action, dobbiamo creare
un asset chiamato \textbf{Input Action}.'' Cercando \emph{IA} nel Content Browser ``troverete
un sacco di asset, ma a voi vi interessa questo celestino''. Con il tasto destro,
\textbf{Show in Folder View}, ``in modo tale che voi troviate direttamente la cartella''.
Si duplica IA\_Jump (Ctrl+C, Ctrl+V, o \emph{Duplicate}) e lo si rinomina
\textbf{IA\_Run}.

\paragraph{Dentro un'Input Action.} Aperto IA\_Run, ``ci sono un po' di valori che
possiamo vedere'':
\begin{itemize}
  \item \textbf{Description}: ``semplice descrizione, quindi quello che fa'';
  \item \textbf{Trigger when Paused}: se l'azione deve scattare anche quando il gioco è in
        pausa, ``quindi si blocca il tempo'';
  \item \textbf{Value Type}: ``la tipologia di valore che ci ritorna indietro. Un input ci
        dà un valore'': \textbf{Digital (bool)}, \textbf{Axis1D (float)} (``come abbiamo
        visto, avanti e indietro''), \textbf{Axis2D (Vector2D)} (``più assi, due
        dimensioni'') o \textbf{Axis3D (Vector)} (``tre tipologie di input che ci vengono
        dati''). ``A noi interessa il Run, che è un bool, quindi premuto o non premuto'':
        Digital (bool) ``va benissimo'';
  \item \textbf{Triggers}: ``le azioni che sono già preinstallate in Unreal'', che dicono
        ``ehi, guarda, si triggera se è premuto, o se è rilasciato, o se è sempre
        premuto''. L'elenco letto dal docente: \emph{Chorded Action}, \emph{Combo},
        \emph{Down}, \emph{Hold}, \emph{Hold and Release}, \emph{Pressed}, \emph{Pulse},
        \emph{Released}, \emph{Tap} (la trascrizione dice anche ``Short Deduction'' e
        ``repeat'', non identificabili con certezza). ``Fondamentalmente, questo vuol
        dire che quando si compie questa azione, questo input viene triggerato.'' IA\_Run
        eredita da IA\_Jump la sua configurazione, ``e va benissimo così, perché tanto è
        un'azione molto simile al jump''.
\end{itemize}

\begin{table}[htbp]
\centering
\small
\begin{tabular}{@{}p{3.2cm}p{10.8cm}@{}}
\toprule
\textbf{Voce} & \textbf{Significato} \\
\midrule
Digital (bool) & premuto / non premuto: un pulsante (salto, corsa, interazione) \\
Axis1D (float) & un valore su un asse: un grilletto analogico, avanti/indietro \\
Axis2D (Vector2D) & due assi: il movimento (X, Y) o la visuale del mouse \\
Axis3D (Vector) & tre assi: input 3D (per esempio un controller di movimento) \\
\midrule
Down & scatta a ogni frame finché il tasto è giù \\
Pressed & scatta una volta, quando il tasto viene premuto \\
Released & scatta una volta, quando il tasto viene rilasciato \\
Hold & scatta dopo che il tasto è stato tenuto premuto per un certo tempo \\
Hold and Release & scatta al rilascio, se il tasto è stato tenuto abbastanza \\
Tap & scatta se il tasto viene premuto e rilasciato entro un tempo breve \\
Pulse & scatta a intervalli regolari mentre il tasto è premuto \\
Chorded Action & scatta solo se un'altra azione è attiva insieme (una ``combinazione'') \\
Combo & scatta se una sequenza di azioni viene eseguita nell'ordine \\
\bottomrule
\end{tabular}
\caption{I tipi di valore (Value Type) e i trigger di un'Input Action, letti a lezione
dall'elenco dell'editor. Le descrizioni sono un'integrazione: a lezione il docente
legge i nomi senza spiegarli uno per uno.}
\label{ue9:tab:inputaction}
\end{table}

\paragraph{Il contesto va caricato.} Tornati in BP\_Knight1: ``abbiamo già il contesto
caricato'', con il grafo copiato dal ThirdPersonCharacter. ``Sappiate che se voi non
caricate il contesto, non vi prenderà nessuna cosa'': sull'Event Begin Play, dal Player
Controller si prende l'\emph{Enhanced Input Local Player Subsystem} e si chiama
\textbf{Add Mapping Context} con IMC\_Default (figura~\ref{ue9:fig:enhanced}).

\paragraph{L'evento IA\_Run e i suoi pin.} Nel grafo si invoca \textbf{IA\_Run}; ``se lo
estendete'' (la freccina in basso del nodo), compaiono tutti i pin:
\begin{itemize}
  \item i pin di esecuzione \textbf{Triggered} (``quando è triggerato''), \textbf{Started}
        (``quando inizia a essere triggerato''), \textbf{Ongoing}, \textbf{Canceled},
        \textbf{Completed};
  \item i valori: l'\textbf{Action Value} (``se è true o false, in questo caso''),
        \textbf{Elapsed Seconds} e \textbf{Triggered Seconds} (``quanto tempo è stato
        triggerato''), e l'\textbf{Input Action}, ``la referenza stessa'' all'asset.
\end{itemize}

\paragraph{Set Max Walk Speed.} ``Dobbiamo modificare il Character Movement, ragazzi, che
dovrà definire quanta velocità avremo'': si trascina nel grafo il componente
\textbf{Character Movement} e se ne tira fuori \textbf{Set Max Walk Speed}, due volte.
Sul \textbf{Triggered}, \textbf{600}; sul \textbf{Completed}, \textbf{300}
(figura~\ref{ue9:fig:sprint}).

\paragraph{Il mapping dimenticato.} In Play non succede niente. ``Scusate, ragazzi, mi
sono dimenticato di dare l'input'': IA\_Run esiste, ma non è nel contesto. Si apre
IMC\_Default, si aggiunge un mapping (il \emph{+} accanto a \emph{Mappings}), si sceglie
IA\_Run e gli si assegna il tasto: ``un nostro bello \textbf{Shift}''. ``L'Input Mapping
Context ci permette di mettere quali tipologie di pulsanti facciano'' ciascuna azione.

\paragraph{Lo stesso input, anche con il controller.} ``Io ad esempio qua, ragazzi, ho un
controller'': il docente lo collega (``Unreal riesce a prendere il controller''), aggiunge
un secondo tasto alla stessa azione, clicca sull'icona della tastiera accanto al campo e
preme un pulsante del controller (``mi sono fregato la sorpresa''): ``facciamo tipo Dark
Souls? Premo cerchio'', e l'editor scrive \textbf{Gamepad Face Button Right}. ``Questa
azione, cioè l'azione di correre, è collegata a più input.'' ``Quando faccio progetti in
generale, devo fare anche il controller'': si corre con Shift sulla tastiera ``oppure con
il nostro bel controllino, con quello che vogliamo noi'', ``come è stato designato il
gioco''. Tutto il resto del movimento sul controller ``non ci interessa adesso''.

\paragraph{La prova.} ``Se noi proviamo\dots\ rilascio.'' Con Shift il cavaliere corre, e
il Blend Space passa alla corsa; al rilascio torna a camminare. ``Ecco a voi che abbiamo
creato un nostro bel character nuovo, in cui semplicemente possiamo correre, possiamo
modificare la velocità.''

\begin{figure}[htbp]
\centering
\begin{tikzpicture}[font=\footnotesize,
  ue9box/.style={draw, thick, rounded corners=4pt, align=center, fill=gray!8,
                 minimum height=1.0cm, text width=2.5cm},
  ue9arr/.style={-{Stealth[length=2.5mm]}, thick}]
  \node[ue9box, fill=yellow!15] (k1) at (0,0.7) {Left Shift\\{\scriptsize tastiera}};
  \node[ue9box, fill=yellow!15] (k2) at (0,-0.7) {Gamepad Face\\Button Right};
  \node[ue9box, fill=blue!10, text width=2.8cm] (imc) at (3.8,0) {\textbf{IMC\_Default}\\
    {\scriptsize Input Mapping Context}\\{\scriptsize (mapping IA\_Run)}};
  \node[ue9box, fill=cyan!12] (ia) at (7.6,0) {\textbf{IA\_Run}\\{\scriptsize Input
    Action, Digital (bool)}};
  \node[ue9box, fill=red!12, text width=2.8cm] (ev) at (11.4,0) {evento
    \textbf{IA\_Run}\\{\scriptsize nel Blueprint}\\{\scriptsize Triggered / Completed}};
  \draw[ue9arr] (k1) -- (imc); \draw[ue9arr] (k2) -- (imc);
  \draw[ue9arr] (imc) -- (ia); \draw[ue9arr] (ia) -- (ev);
  \node[ue9box, fill=green!10, text width=4.2cm] (bp) at (3.8,-2.4) {\textbf{Event Begin
    Play} $\rightarrow$ Enhanced Input Local Player Subsystem $\rightarrow$ \textbf{Add
    Mapping Context} (IMC\_Default)};
  \draw[ue9arr, dashed] (bp) -- node[right, font=\scriptsize]{attiva} (imc);
\end{tikzpicture}
\caption{La catena dell'Enhanced Input. I tasti fisici non sono legati direttamente
all'evento: il contesto (IMC) associa più tasti alla stessa azione (IA), e il Blueprint
reagisce all'azione. Senza Add Mapping Context, il contesto non è attivo e nessun input
arriva. Schema ricostruito.}
\label{ue9:fig:enhanced}
\end{figure}

\begin{figure}[htbp]
\centering
\begin{tikzpicture}[font=\footnotesize,
  ue9node/.style={draw, thick, rounded corners=4pt, align=center, fill=gray!8,
                  minimum height=0.8cm, text width=2.4cm},
  ue9ev/.style={ue9node, fill=red!12},
  ue9var/.style={ue9node, fill=blue!10, rounded corners=8pt},
  ue9exec/.style={-{Stealth[length=2.5mm]}, very thick, black!80},
  ue9data/.style={-{Stealth[length=2mm]}, thick, blue!60!black}]
  \node[ue9ev, text width=2.2cm, minimum height=1.8cm] (ev) at (0,0) {\textbf{IA\_Run}\\
    {\scriptsize Enhanced Input Action}};
  \node[ue9node, text width=3.2cm] (s1) at (5,0.8) {\textbf{Set Max Walk Speed}\\
    {\scriptsize Max Walk Speed = 600}};
  \node[ue9node, text width=3.2cm] (s2) at (5,-0.8) {\textbf{Set Max Walk Speed}\\
    {\scriptsize Max Walk Speed = 300}};
  \node[ue9var, text width=2.6cm] (cm) at (9.6,0) {\textbf{Character}\\
    \textbf{Movement}};
  \draw[ue9exec] (ev.east |- s1) -- node[above, font=\scriptsize]{Triggered} (s1);
  \draw[ue9exec] (ev.east |- s2) -- node[below, font=\scriptsize]{Completed} (s2);
  \draw[ue9data] (cm.west) -- ++(-0.5,0) |- node[pos=0.25, right, font=\scriptsize]
    {Target} (s1.east);
  \draw[ue9data] (cm.west) -- ++(-0.5,0) |- (s2.east);
\end{tikzpicture}
\caption{Lo sprint di BP\_Knight1. Premendo Shift (o il cerchio del controller) la
velocità massima sale a 600 e il Blend Space mostra la corsa; al rilascio torna a 300,
la camminata. Schema ricostruito.}
\label{ue9:fig:sprint}
\end{figure}

\paragraph{Approfondimento: Action Mapping ed Enhanced Input.} Il libro di Romero e
Sewell, scritto per le prime versioni di Unreal Engine 5, costruisce lo stesso sprint
con il vecchio sistema: un \emph{Action Mapping} chiamato \emph{Sprint} nei Project
Settings (sezione \emph{Input}), legato a \emph{Left Shift}, e un evento
\emph{InputAction Sprint} con i pin \emph{Pressed} e \emph{Released}, su cui due
\emph{Set Max Walk Speed} portano la velocità a 2200 e poi di nuovo ai 600 di default
(capitolo 6). Il principio è identico; cambia il sistema. L'\textbf{Enhanced Input},
introdotto come plugin in Unreal Engine 5.0 e diventato il sistema predefinito dei
modelli dalla 5.1 (il ``quasi dalla 5.2'' del docente è un ricordo approssimativo), separa
tre livelli: l'\emph{azione} (Input Action, che cosa vuole fare il giocatore), il
\emph{contesto} (Input Mapping Context, quali tasti attivano quali azioni in una certa
situazione) e il \emph{comportamento} (Triggers e \emph{Modifiers}, come per esempio la
\emph{zona morta} di una levetta o l'inversione di un asse). I contesti si aggiungono e
si tolgono a runtime con \emph{Add Mapping Context} e \emph{Remove Mapping Context}, e
hanno una \emph{Priority}: è così che si passa ``dal gioco a piedi alla macchina''
dell'esempio del docente, e così che si costruisce una schermata di rimappatura dei
tasti. Gli autori raccomandano anche di raggruppare i nodi dello sprint in un
\emph{commento} (\emph{Create Comment from Selection}), per tenere ordinato il grafo.

% ---------------------------------------------------------------------
\section{Un tentativo: FInterp To}
\label{ue9:sec:finterp}

``Se volessimo farlo un po' più carino, ragazzi'', il passaggio fra 300 e 600 potrebbe
essere graduale invece che istantaneo. Il blocco è \textbf{FInterp To} (``finter2''):
``fa un'interpolazione di un float''. ``Avete presente quando abbiamo parlato di
interpolazione?'' ``Fa il passaggio tra due numeri. Lui dice: devi passare da questo
valore, cioè \textbf{Current}, a questo valore, \textbf{Target}, in questo tempo'',
\textbf{Delta Time}, per cui il docente consiglia \textbf{Get World Delta Seconds}; e
poi c'è l'\textbf{Interp Speed}, la velocità dell'interpolazione.

\paragraph{Il grafo.} Sul Triggered: \textbf{Current} = Character Movement $\rightarrow$
\textbf{Get Max Walk Speed}; \textbf{Target} = 600; \textbf{Delta Time} = World Delta
Seconds; il risultato nel Set Max Walk Speed. Il docente ricorda che il Delta Seconds si
può moltiplicare per un valore (``ricordiamo che possiamo modificarlo moltiplicandolo
per un valore''), e prova con 10, poi con un decimo. Sul Completed lo stesso, ``ma
semplicemente invertiamo'': il Current parte da 600, il Target è 300. ``Ricordiamoci di
collegare anche al Character Movement.''

\paragraph{Non funziona, e perché.} La corsa resta ``ancora molto veloce'', il passaggio
istantaneo. Il docente prova valori diversi (10, 100: ``è alla norma del pianeta stare
continuamente a vedere se i valori sono giusti'') e poi mette dei \textbf{Print String}:
``ci dà seicento''. La diagnosi: ``fondamentalmente questa è un'azione a triggered
unico, quindi non farà l'interpolazione''. FInterp To fa \emph{un passo}
dell'interpolazione ogni volta che viene eseguito: se viene eseguito una volta sola, fa
un solo passo. ``Lo deve fare tramite un \textbf{timer}'': si rimanda a dopo (``volendo,
ragazzi, si può fare con un timer, però facciamolo magari dopo''), e non verrà ripreso
nella lezione.

\paragraph{Approfondimento: dove mettere un'interpolazione.} Un nodo come FInterp To
(o VInterp To, RInterp To per vettori e rotazioni) va eseguito \emph{a ogni frame}
finché il valore non arriva al target: il suo risultato è il valore ``un po' più vicino''
al target, e il Delta Time serve a renderlo indipendente dal frame rate. Ci sono tre
modi tipici di farlo girare: nel \textbf{Tick}, con una variabile \emph{TargetSpeed}
che i due pin dell'input si limitano a impostare (600 o 300); con un \textbf{timer}
(\emph{Set Timer by Event}, con \emph{Looping}, fermato con \emph{Clear and Invalidate
Timer} quando il valore è arrivato), come suggerisce il docente; o con una
\textbf{Timeline}, vista nella sesta lezione, che va da 300 a 600 con \emph{Play} e
torna indietro con \emph{Reverse}. Con l'Enhanced Input c'è anche una quarta via: il pin
\emph{Triggered} scatta a ogni frame se il trigger è \emph{Down} (o se non c'è nessun
trigger), e una sola volta se il trigger è \emph{Pressed}. IA\_Run, copiata da IA\_Jump,
ha probabilmente un trigger che la fa scattare una volta sola: è ``l'azione a triggered
unico'' del docente.

% ---------------------------------------------------------------------
\section{Prima della pausa: le animazioni da scaricare}
\label{ue9:sec:download}

``Adesso andremo a vedere, se non mi esplode prima Unreal, anche come si fanno i
\textbf{montaggi}''; mancano una ventina di minuti alla pausa. ``Dovremmo andare un attimo
su Mixamo'' (``a quanto pare oggi Mixamo ci mette pure un bel po' a caricare''), per
scaricare due gruppi di animazioni.

\paragraph{Un attacco in root motion.} ``Vorrei che sceglieste un'animazione che sia in
\textbf{root motion}'': ``riprendiamo ad esempio una situazione simile a Dark Souls''.
``Ricordatevi sempre di prendere il vostro character'' su Mixamo (il paladino, nel caso
del docente), perché le animazioni si adattano al personaggio selezionato. Il docente
cerca \emph{attack} e sceglie un attacco con la spada che ``di base, come potete vedere,
non può essere messo in place''. Le impostazioni di download sono quelle della volta
precedente: \textbf{Without Skin}, \textbf{FBX Binary}, \textbf{30} fotogrammi al
secondo. ``Ripeto: non deve essere in place, deve avere il root motion, quindi il
movimento ereditato.'' ``Un'azione unica'': ``ad esempio un attacco con una spada, un
qualche cosa'', ``un'animazione che vi piace di più''.

\paragraph{Le animazioni del Blend Space 2D.} Dopo la pausa si farà il \textbf{Blend
Space 2D}, ``che ci servirà per fare anche il calcolo della direzione del nostro
character''. Servono quindi altre animazioni: ``walk destra, sinistra, dritto e
indietro''. ``Consiglio spassionato'': su Mixamo si chiamano \textbf{Strafe} (``provate a
cercare Strafe''; ``Jog Strafe va benissimo anche''), e queste, al contrario
dell'attacco, ``dovete metterle \textbf{in place}''.

\paragraph{Il mirror.} ``Piccolo consiglio: se volete far sì che siano abbastanza simili
tra di loro, potete fare il \textbf{mirror}.'' Si scarica l'animazione verso sinistra
non specchiata; poi, con l'opzione \emph{Mirror} di Mixamo, ``in automatico andrà a
destra''. Le due animazioni saranno così perfettamente simmetriche.

\paragraph{Quante animazioni.} Su richiesta (``le puoi scrivere per messaggio?''), il
docente le elenca anche in chat: ``in totale sono otto animazioni''. Avanti, indietro,
destra, sinistra, e le quattro \textbf{diagonali}: avanti a destra, avanti a sinistra,
indietro a destra, indietro a sinistra. Uno studente chiede se l'avanti serva, visto che
c'è già la camminata: ``nel senso, possiamo usare tranquillamente quello di avanti; ma in
questo caso abbiamo solamente avanti, dobbiamo trovare anche indietro''. E un'altra: le
diagonali ``soltanto in avanti, o anche indietro?'' ``Anche indietro'': ``se io vado tipo
così, indietro verso destra, e così, indietro verso sinistra''. ``Praticamente è la X,
letteralmente'': le due diagonali incrociate, oltre alla croce degli assi. Il docente
nomina anche alcuni pacchetti di Mixamo con animazioni di locomozione (``Locomotion
Pack''; il resto dei nomi è incomprensibile nella trascrizione). ``Ci rivediamo alle
18.''

% ---------------------------------------------------------------------
\section{La pausa e i feedback}
\label{ue9:sec:feedback}

Al rientro, ``siamo rimasti in cinque'', poi in quattro. Il docente apre uno spazio per
i commenti: ``se avete dei feedback riguardo a questo corso, sparate tranquillamente'',
``siamo anche molto coetanei, anzi credo io di essere anche forse più piccolo di alcuni
di voi''. ``Se avete cose da chiedere, magari di andare più piano, andare più veloce,
oppure un po' come sono strutturate le lezioni''. ``Capisco pure che può essere
effettivamente un po' pesante in alcuni momenti'', ``perché purtroppo i Blueprint, o
Unreal in generale, è divertente quando poi le sai, le cose''.

\paragraph{Le funzioni da imparare a memoria.} Uno studente: ``immagino che sono
tantissime funzioni, quindi sono tutte da imparare a memoria''. ``Quello, eh sì.''

\paragraph{Il progetto del docente: la pallina morbida e dura.} Per mostrare dove porta
la pratica, il docente racconta un suo progetto personale: ``una specie di
\emph{platformer} in cui una pallina diventa morbida e dura, nel senso che diventa di
metallo o di un materiale più morbido''. ``La cosa è che ho simulato tutta la fisica'':
``la fisica non è simulata da Unreal, ma sono io che faccio finta, tramite i Blueprint, di
far fare delle funzioni fisiche''. ``Hai messo proprio te le formule fisiche?'' ``Più o
meno.'' Il meccanismo: si rileva lo \textbf{slope}, ``il quanto è piegato il piano'',
perché ``la meccanica è il fatto che se sei più pesante con la pallina vai più veloce su
uno slope, come nella realtà'': ``una pallina di metallo e una pallina di gomma: quella di
metallo va più veloce, perché è più pesante, comunque più influenzata''. ``È tutto
simulato'': ``sono io che, tramite la detection dello slope, cambio i valori''.

\paragraph{Combinare i blocchi.} ``E quello è ovviamente tutto custom'': ``quando hai
tutte le nozioni in testa, ti dici: ok, so questo blocco A, B, C, D, E, F, G; quale mi
serve? Il blocco C e il blocco A. Li metto insieme e faccio una nuova combinazione.''
``È sempre il fatto di capire come posso interlacciare le cose.'' ``Qua stiamo facendo
cose magari un po' generiche''; poi ``dici: ok, voglio fare simile a quello, ma
leggermente diverso. Voglio fare il character, però con la spada: che mi cambia?''

\paragraph{Inventari e UI.} La richiesta di Federico: ``a me interesserebbero la gestione
degli inventari e le interfacce''. ``Sfortunatamente no'': ``sappi che l'inventario e la
UI sono un mondo a parte''. ``Esistono proprio figure riguardanti questo mondo della UI,
sia designer sia programmatori.'' Si farà ``una cosa come un menu'', ``una cosetta
tranquilla'', ``però per il resto non andremo a fare di più''.

\paragraph{Le lezioni con Davide.} ``Sappiate che il 25 è la mia ultima lezione, e avrete
tutte quelle con \textbf{Davide}, che invece vi farà la parte più di \emph{level}.'' Con
lui si useranno ``tutti gli asset che abbiamo creato: le piattaforme, i pulsanti'', per
lavorare su concetti come la \textbf{gym}, ``il concetto della palestra'', che ``si fa
praticamente in via di sviluppo'': ``è il parco giochi'', ``letteralmente il parco giochi
dove ho tutte le meccaniche, le posso provare tutte quante, vedo se funzionano''. Poi si
costruiscono situazioni: ``voglio fare un livello horror, la bestia che ti insegue: tu
entri in una stanza e devi fuggire dalla bestia''; ``situazione X, situazione Y,
situazione Z'', ``le costruisci e metti insieme i componenti'', ``e vedi quei componenti
che hai programmato come funzionano a livello anche di esperienza, di gioco''. Gli
elementi saranno ``quelli di un classico platform: le piattaforme, i bottoni'', e si
vedranno anche i tool di \textbf{modeling} dell'editor e ``un po' di teoria di level
design'': ``chi fa game design, e a chi interessa tantissimo anche la parte di level:
molto, molto interessante''.

\paragraph{Dove lasciare i feedback.} ``Ripeto, ragazzi, anche a chi sta in differita'':
qualsiasi feedback sulle lezioni va lasciato sul canale \textbf{general} della chat,
``non su introduzione al materiale, perché sennò si fa casino''. Non solo sugli
argomenti, ``ma feedback anche proprio per le lezioni in generale: andare più lento,
andare più veloce, vedere cose più pratiche''.

\paragraph{Approfondimento: la gym.} Nel level design la \textbf{gym} (o \emph{gym
level}, \emph{gray box gym}) è un livello di servizio che non finirà nel gioco: una
stanza in cui ogni meccanica è disposta in isolamento e misurata. Si mettono salti di
lunghezze crescenti per sapere qual è il salto massimo del personaggio, rampe di
pendenze diverse per verificare la \emph{Max Walkable Floor Angle}, gradini di varie
altezze per la \emph{Max Step Height}, corridoi di larghezze diverse per la camera. Le
misure raccolte nella gym diventano le ``regole metriche'' del level design (quanto può
essere largo un burrone, quanto alto un muretto da scavalcare), e i livelli veri si
costruiscono rispettandole. È il ``parco giochi'' del docente; e il problema della
camera che compenetra, visto in questa lezione, è uno dei primi che una gym fa emergere.
% ---------------------------------------------------------------------
\section{La direzione: Calculate Direction}
\label{ue9:sec:direction}

``Abbiamo detto prima che oggi dobbiamo fare il Blend Space 2D.'' Il docente ha importato
otto animazioni, ``cercando di prenderle tutte similari'': \emph{Jog Backward},
\emph{Jog Backward Diagonal}, \emph{Jog Forward}, \emph{Jog Forward Diagonal} e così
via, ``a parte lo Strafe Left, che devo eliminare perché non andava bene come
animazione''. Il Blend Space 1D prende solo la velocità: ``vai più veloce, l'animazione,
benissimo''. Il 2D aggiunge un secondo asse, la \textbf{direzione}, ``in modo tale da non
dover magari direzionare sempre il character'', ma vedere ``che si muove in diagonale, si
muove a sinistra''.

\paragraph{Dove si calcola.} In ABP\_Knight, nell'Event Graph dell'Update, si aggiunge un
pin alla \textbf{Sequence} (il \emph{+}, \emph{Add pin}): \emph{Then 2}, o il successivo.
Questa volta non si parte dal Movement Component, ma dal \textbf{character}: ``perché la
direzione è determinata dalla rotazione del character stesso'', ``cioè la capsula, come è
ruotata''. Dalla variabile del cavaliere si tira fuori \textbf{Get Actor Rotation}.

\paragraph{Il nodo Calculate Direction.} ``Potremmo tranquillamente prenderla così, ma noi,
per semplificarci anche la vita, andiamo a calcolare proprio la direzione con un blocco
chiamato proprio \textbf{Calculate Direction}.'' Il nodo chiede due ingressi: la
\textbf{Velocity} attuale (dal Movement Component, \textbf{Get Velocity}) e la
\textbf{Base Rotation} (la Get Actor Rotation). Restituisce ``un float'', che con
\textbf{Promote to Variable} diventa \textbf{Direction}.

\paragraph{Uno studente perso.} ``Io mi sono perso: Character, Movement Component, da dove
vengono questi due?'' Il docente ripercorre: ``hai fatto questo Event Blueprint Initialize
Animation nell'Event Graph, no?''; ``dovresti avere qua sotto, sotto le variabili, la
variabile del character: la tiri fuori, ne tiri fuori Get Actor Rotation''; dal Movement
Component ``tiri fuori la velocità''; ``da Get Actor Rotation fai Calculate Direction, e
questo diventerà il nostro float''. Un altro studente nota che la variabile del docente si
chiama \emph{Character}, mentre la sua si chiama \textbf{As BP Knight}: ``va benissimo.
Io l'ho solo rinominata.'' Quando si fa Promote to Variable sull'uscita di un cast, Unreal
propone il nome \emph{As} più il nome della classe (``ti dice proprio As BP Knight'').

\paragraph{Da $-180$ a $180$.} ``Molto semplicemente, la direzione sarà tra un range,
sempre senza doverla clampare'' (``il clamp l'abbiamo fatto quasi all'inizio''): fra
\textbf{$-180$ e $180$ gradi}. ``Si parte da 0; a destra vai verso il positivo; quando
arrivi a 180 passa a $-179$''; ``diciamo, non lo fa come un angolo da 0 a 360, ma fa 0,
180, $-180$ e di nuovo 0'' (figura~\ref{ue9:fig:direzione}).

\begin{figure}[htbp]
\centering
\begin{tikzpicture}[font=\footnotesize]
  \draw[thick] (0,0) circle (2.2);
  \foreach \a/\l/\pos in {90/0/above, 45/45/above right, 0/90/right, -45/135/below right,
      -90/{180 / $-180$}/below, -135/{$-135$}/below left, 180/{$-90$}/left,
      135/{$-45$}/above left}
    {\draw[-{Stealth[length=2mm]}, thick, blue!60!black] (0,0) -- (\a:2.2);
     \node[\pos, font=\scriptsize\bfseries] at (\a:2.25) {\l};}
  \fill[gray!40] (0,0) circle (0.35);
  \draw[thick] (0,0) circle (0.35);
  \draw[-{Stealth[length=2mm]}, very thick] (0,0) -- (0,0.6);
  \node[font=\scriptsize, align=left, anchor=west] at (3.4,1.2) {\textbf{0}: avanti\\
    \textbf{positivi}: verso destra\\\textbf{negativi}: verso sinistra\\
    \textbf{$\pm$180}: indietro};
  \node[font=\scriptsize, align=left, anchor=west] at (3.4,-0.9) {la direzione è
    l'angolo fra\\la velocità e il ``davanti''\\del personaggio (Base Rotation)};
\end{tikzpicture}
\caption{Il valore di Direction restituito da Calculate Direction. Il personaggio guarda
in alto; le frecce sono le otto direzioni delle animazioni del Blend Space 2D. Andare
indietro vale sia 180 sia $-180$, a seconda del lato da cui ci si arriva. Schema
ricostruito.}
\label{ue9:fig:direzione}
\end{figure}

\paragraph{Rimettere Use Controller Rotation Yaw.} Per vedere i movimenti laterali,
``importante è rimettere lo \textbf{Use Controller Rotation Yaw}'': ``non vogliamo che il
character ruoti sempre nella direzione in cui lo stiamo indirizzando''. Il corpo deve
seguire la camera, e la direzione del movimento rispetto al corpo deve cambiare.

\paragraph{La verifica con Print String.} Per sicurezza si stacca un attimo il Blend Space
1D (``staccare un attimo le animazioni'') e si stampa Direction con un \textbf{Print
String}, ``che lo farà a ogni frame''. ``Ci dà direttamente nella preview questo valore'';
in Play, ``sarà un po' rotto, ma è normale, ragazzi'': ``$-70$, 70, 160, 0''. ``Mamma mia''
(il valore dipende dalla camera, ``giustamente''). ``Come possiamo vedere, questi valori
rappresentano la nostra direzione attuale.'' Il docente nota anche che il cavaliere
sembra ``leggermente ruotato'' rispetto alla capsula: ``buono a sapersi''.

\paragraph{Approfondimento: che cosa calcola Calculate Direction.} Il nodo, della
libreria dell'Animation Blueprint (\emph{KismetAnimationLibrary}; nelle versioni più
vecchie era una funzione dell'AnimInstance), trasforma la velocità nello spazio locale
del personaggio e ne restituisce l'angolo sul piano orizzontale, in gradi, nell'intervallo
$[-180, 180]$: 0 se la velocità punta dove guarda il personaggio, 90 se punta alla sua
destra, $-90$ a sinistra, $\pm 180$ dietro. A velocità nulla restituisce 0. Con
\emph{Orient Rotation to Movement} acceso il personaggio si gira sempre verso la
velocità, e la direzione resta quasi sempre vicina a zero: il Blend Space 2D avrebbe poco
da fare. Per questo il docente riattiva Use Controller Rotation Yaw.

% ---------------------------------------------------------------------
\section{Il Blend Space 2D}
\label{ue9:sec:blend2d}

\paragraph{Creare l'asset.} Tasto destro $\rightarrow$ \emph{Animation} $\rightarrow$
\textbf{Blend Space}: ``Blend Space senza niente, di base, è 2D; se viene specificato
1D, è 1D''. Unreal chiede lo scheletro: ``voi scegliete il vostro scheletro''. Il nome:
\textbf{BS2D\_Knight}.

\paragraph{I due assi.} ``Come potete vedere, a sinistra e giù sono presenti due nomi'',
per ora \emph{None}. Nei Details:
\begin{itemize}
  \item \textbf{Horizontal Axis}: \textbf{Direction}, da \textbf{$-180$} a
        \textbf{$180$};
  \item \textbf{Vertical Axis}: \textbf{GroundSpeed} (``più che la velocità, la ground
        speed''), da \textbf{0} a \textbf{600}.
\end{itemize}
Il parametro \textbf{Grid Divisions} ``ci permette di aumentare il numero di divisioni
sulla nostra griglia'' (le linee ``non si vedono, perché sono veramente molto poco
chiare''): sono i punti su cui si agganciano i campioni delle animazioni. Il docente ne
mette \textbf{8} su entrambi gli assi. Con 8 divisioni, sull'asse orizzontale ogni passo
vale 45 gradi, esattamente le otto direzioni; sull'asse verticale vale 75.

\paragraph{Un quiz.} ``Domanda, ragazzi: se io sto a 600 di velocità e a 90 di direzione,
dove sto andando?'' ``Ricordatevi che positivo è destra, negativo è sinistra.'' ``A
destra.'' ``Pienamente a destra, perfetto.'' ``E invece a $-135$?'' ``Dietro a
sinistra.'' ``Ottimo.''

\paragraph{Dove va ogni animazione.} Si dispongono i campioni (tabella~\ref{ue9:tab:bs2d} e
figura~\ref{ue9:fig:blend2d}). In basso, a velocità zero, l'\textbf{Idle}: ``mettetene
tre, così abbiamo un range per tutto'' (a $-180$, 0 e 180): ``siamo fermi, siamo fermi''.
In alto, a 600, le otto animazioni. L'animazione all'indietro va messa due volte, a
$-180$ e a $180$: ``stiamo andando indietro sia se abbiamo girato da sinistra, sia da
destra''. Il docente si corregge una volta (``Backward Diagonal L sarà\dots\ no, ho
sbagliato'': a $-135$, non a 90) e invita a controllare i valori nei Details del
campione.

\begin{table}[htbp]
\centering
\small
\begin{tabular}{@{}p{4.4cm}p{2.2cm}p{2.2cm}p{4.6cm}@{}}
\toprule
\textbf{Animazione} & \textbf{Direction} & \textbf{GroundSpeed} & \textbf{Movimento} \\
\midrule
Idle (tre copie) & $-180$, 0, 180 & 0 & fermo \\
Jog Forward & 0 & 600 & avanti \\
Jog Forward Diagonal R & 45 & 600 & avanti a destra \\
Strafe R & 90 & 600 & a destra \\
Jog Backward Diagonal R & 135 & 600 & indietro a destra \\
Jog Backward & 180 e $-180$ & 600 & indietro \\
Jog Backward Diagonal L & $-135$ & 600 & indietro a sinistra \\
Strafe L & $-90$ & 600 & a sinistra \\
Jog Forward Diagonal L & $-45$ & 600 & avanti a sinistra \\
\bottomrule
\end{tabular}
\caption{La posizione dei campioni nel Blend Space 2D BS2D\_Knight. I nomi delle
animazioni sono quelli di Mixamo citati a lezione, per quanto ricostruibili dalla
trascrizione; i valori sono quelli dettati dal docente.}
\label{ue9:tab:bs2d}
\end{table}

\begin{figure}[htbp]
\centering
\begin{tikzpicture}[font=\footnotesize, x=0.032cm, y=0.006cm,
  ue9pt/.style={draw=black, fill=white, thick, diamond, inner sep=1.8pt}]
  % griglia
  \foreach \x in {-180,-135,...,180} \draw[gray!30] (\x,0) -- (\x,600);
  \foreach \y in {0,75,...,600} \draw[gray!30] (-180,\y) -- (180,\y);
  \draw[thick] (-180,0) rectangle (180,600);
  \foreach \x in {-180,-90,0,90,180}
    \node[below, font=\scriptsize] at (\x,-10) {\x};
  \foreach \y in {0,300,600} \node[left, font=\scriptsize] at (-182,\y) {\y};
  \node[below, font=\small] at (0,-80) {Direction};
  \node[rotate=90, font=\small] at (-235,300) {GroundSpeed};
  % campioni
  \foreach \x in {-180,0,180} \node[ue9pt] at (\x,0) {};
  \node[above, font=\scriptsize] at (0,10) {Idle};
  \foreach \x/\l in {-180/Back, -135/BackDiag L, -90/Strafe L, -45/FwdDiag L, 0/Fwd,
      45/FwdDiag R, 90/Strafe R, 135/BackDiag R, 180/Back}
    {\node[ue9pt] at (\x,600) {};
     \node[above, font=\tiny, rotate=35, anchor=south west] at (\x-6,612) {\l};}
  % punto attuale
  \fill[green!50!black] (60,450) circle (2.2pt);
  \node[green!40!black, font=\scriptsize, anchor=west] at (66,420) {(60, 450)};
\end{tikzpicture}
\caption{Il Blend Space 2D del cavaliere: Direction in orizzontale, GroundSpeed in
verticale, griglia 8$\times$8. Le tre copie dell'Idle stanno in basso; le otto direzioni
in alto, con la corsa all'indietro ripetuta ai due estremi. Il punto verde è un esempio:
il personaggio sta andando in avanti a destra, un po' sotto la velocità massima, e Unreal
mescola i quattro campioni più vicini. Schema ricostruito.}
\label{ue9:fig:blend2d}
\end{figure}

\paragraph{La preview con Ctrl.} ``Potete fare la preview delle animazioni con
\textbf{Ctrl}'': tenendolo premuto e trascinando il cursore sul grafico, ``potete vedere
come l'animazione si muove tra tutti gli step''. ``Consiglio sempre di mettersi dietro''
il personaggio nel viewport, ``perché tanto poi, se lo vedete davanti, è difficile'' capire
destra e sinistra. Il docente controlla tutte le direzioni (``vado a sinistra, vado
indietro, vado dritto, vado dritto a destra, vado a destra, giù a destra, indietro.
Perfetto''); una delle animazioni ``mi sa che non va bene'', ma ``sì, sto andando a
sinistra, è solo molto strana''.

\paragraph{Smoothing Time.} ``Piccola cosa, se volete introdurla'': lo \textbf{Smoothing
Time}, ``il tempo di passaggio tra un'animazione e l'altra in generale, per tutte le
animazioni''. Lo si mette a \textbf{0,2} su entrambi gli assi. Il \emph{Damping Ratio}
``non ci interessa''.

\paragraph{Nello stato WalkRun.} ``Facciamo tutto il percorso a ritroso'': AnimGraph,
Locomotion, WalkRun. Al posto del Blend Space 1D si trascina BS2D\_Knight, che chiede
``dammi la Direction e la GroundSpeed'': si collegano i due Get. Compile: ``vediamo che il
nostro Blend Space funziona in preview''.

\paragraph{Le impostazioni del personaggio.} In BP\_Knight1: ``mi raccomando, l'abbiamo
detto prima, ma lo diciamo anche adesso'', si riattiva \textbf{Use Controller Rotation
Yaw}; e nel Character Movement la \textbf{Max Walk Speed} va a \textbf{600}, ``che è
quanto abbiamo messo'' come massimo dell'asse. ``Ecco a voi, ragazzi'': ``se io mi muovo a
destra lui va a destra, se mi muovo a sinistra, anche se l'animazione non è fantastica,
va a sinistra; vado indietro; vado in diagonale''. Nel frattempo ``parte la torretta, che
mi spara'': ``il mio bel cavaliere se ne va a zonzo come gli pare''.

\paragraph{Gli scatti.} C'è ``un piccolo \emph{clip}'', uno scatto, in alcuni passaggi.
Il docente prova a prendere la rotazione della \textbf{Mesh} invece di quella
dell'actor (``facciamo che invece di prendere la direzione del character, prendiamo
quella della Mesh''; ``Get Mesh''), poi ad aumentare lo Smoothing Time della Direction
a \textbf{0,5} (``un po' più tranquillo''), e infine torna al collegamento originale.
La conclusione: ``probabilmente sono dei problemi sull'animazione''; ``mi sa che devo
riscaricare l'animazione, perché snappa un pochino''. Ma il risultato c'è: ``abbiamo il
nostro character che si muove in tutte le direzioni''; ``è un po' veloce, devo essere
sincero'', e la velocità si può abbassare.

\paragraph{Il cerchio, per chi si è perso.} Uno studente chiede di rivedere il Blend
Space 2D. ``Immaginatelo come un cerchio'': ``se stai a 0 di direction, stai andando
dritto; se stai a 180 o $-180$, stai andando dietro''; ``tutti i valori intermedi'' sono le
diagonali e i lati. ``Destra è positivo: 0, 90, 180; sinistra è $-90$, $-180$.'' ``Se vai
in diagonale a destra sarà 45, se vai pienamente a destra sarà 90'', ``e 135 a 600 vuol
dire che andrai praticamente indietro, ma verso destra''. ``Una volta capito che una è
l'opposto dell'altra'', i valori negativi seguono da soli.

\paragraph{Approfondimento: Blend Space 2D e foot sliding.} In un Blend Space 2D Unreal
divide lo spazio dei campioni in triangoli e, per un valore dato, mescola i campioni ai
vertici del triangolo che lo contiene. Ne seguono tre conseguenze pratiche. Primo,
l'Idle ripetuto lungo tutta la riga inferiore (le tre copie del docente) evita che a
bassa velocità il personaggio ``scelga'' una direzione a caso. Secondo, la corsa
all'indietro va duplicata ai due estremi, perché $-180$ e 180 sono lo stesso angolo ma
due punti diversi della griglia: con un solo campione, passando da 179 a $-179$ il
personaggio attraverserebbe in un istante tutto il grafico. Terzo, lo scatto che il
docente vede è spesso proprio quel salto da $+180$ a $-180$, che lo \emph{Smoothing
Time} non può interpolare bene (interpola il valore numerico, non l'angolo); le
alternative sono il \emph{Weight Speed} per campione, oppure le soluzioni più avanzate
di Unreal Engine 5 come l'\emph{Orientation Warping}, che ruota la parte inferiore del
corpo verso la direzione del movimento con quattro sole animazioni cardinali. Infine,
una griglia con l'asse verticale fino a 600 e i campioni tutti a 600 significa che anche
la camminata a 300 è un misto fra Idle e corsa: per avere una vera camminata laterale
servirebbe una seconda riga di campioni a 300, con le animazioni di camminata.

% ---------------------------------------------------------------------
\section{Gli Animation Montage}
\label{ue9:sec:montage}

``Intanto ci abbiamo un'ultima mezz'oretta, ragazzi. Vediamo un attimo i Montage.''
L'animazione di attacco (``l'avevo scaricata, mi sa che non l'ho messa in Knight
Animation''; ``scusate, ragazzi'') è una \textbf{Sword and Shield Attack} di Mixamo:
``questa qua sta in root motion'': ``come potete vedere, eredita il movimento''.

\paragraph{Il problema.} ``Quando vogliamo fare una nostra animazione che di conseguenza
sia attivabile durante un ciclo di animazioni'', non basta aggiungere uno stato: ``dobbiamo
creare quello che viene chiamato un \textbf{Montage}''. (Nel frattempo ``Filippo è andato
via, è crashato'': ``siamo rimasti noi tre''.)

\paragraph{Che cos'è un Montage, dalla documentazione.} Il docente apre la pagina della
documentazione di Unreal sugli \textbf{Animation Montage}. ``Il montaggio di animazioni ci
permette sia di montare delle animazioni\dots'' L'esempio della pagina è la ricarica di
un fucile: un'animazione formata da tre animazioni diverse, \emph{MM Shotgun Reload Drop},
\emph{Reload Load} e \emph{Reload Lock}. ``Mettiamo conto che vogliamo fare un bel gioco
shooter. Il fatto di staccare il caricatore su quattro fucili diversi è praticamente
uguale''; ``non dobbiamo fare quattro animazioni diverse'': ``l'unica cosa che cambia è il
lock, perché magari uno vuole farlo un po' più grosso''. Con un Montage si creano
\textbf{animazioni modulari}: ``guarda, il drop è uguale, il load è uguale, il lock invece
è diverso''. E poi il secondo uso, quello di oggi: ``inserire animazioni in contesti di
cicli di animazione''; ``sono una tantum'': ``premo un tasto e compio questa animazione,
nel mentre che sto camminando''. ``Magari la prossima volta andiamo a vedere pure un po'
più approfondito.''

\paragraph{Creare il Montage.} Tasto destro sull'animazione di attacco $\rightarrow$
\textbf{Create} $\rightarrow$ \textbf{Create AnimMontage}. Il nuovo asset si chiama come
l'animazione con il suffisso \emph{Montage} (\emph{Sword And Shield Attack\_Montage}).
``La cosa bella del Montage è che ci permette di fare un sacco di cose all'interno dei
suoi parametri.''

% ---------------------------------------------------------------------
\section{Dentro il Montage: Notify, sezioni e Slot}
\label{ue9:sec:notify}

\paragraph{I Notify.} ``Una cosa interessante da spiegare sono i \textbf{Notify}'': ``dei
punti, dei puntini, sulla linea del tempo dell'animazione'', che permettono di aggiungere
``un'azione che si triggera'': ``suoni, un VFX, qualsiasi cosa''. ``Quelli più generici
sono già fatti da Unreal, fortunatamente.'' Il docente fa clic destro sulla traccia dei
Notify, \textbf{Add Notify} $\rightarrow$ \textbf{Play Sound}, e come suono sceglie
``l'unico sound che abbiamo'', quello della piattaforma (un Sound Cue). Le tracce dei Notify
possono essere più d'una: con \emph{Add Notify Track} se ne crea un'altra, per esempio
una per i VFX e una per i suoni. In preview (tolto il loop, ``fammelo fare una volta''):
``appena il nostro character arriva a questo punto dell'animazione, compie il suono''.
``Mettiamo conto: in combattimento c'è un VFX, c'è un suono, c'è un tirare fuori la spada,
e quindi si sente la spada che esce fuori dal fodero'': con una linea di Notify ``questo
suono avviene in un preciso momento dell'animazione''. ``Ovviamente, non andremo a usare
il sound effect della piattaforma per il nostro Knight.''

\paragraph{Le Montage Section.} Nel \emph{Montage group} c'è la sezione \textbf{Default}.
Con il tasto destro, \textbf{New Montage Section}: il docente ne crea una chiamata
\textbf{Attack}, nel punto in cui ``lui sta attaccando''. Le sezioni servono ``per
ordinare il lavoro'' e per ``dividere l'animazione in più punti diversi, con più
animazioni''.

\paragraph{Più animazioni nello stesso Montage.} La striscia verde sotto le sezioni è
l'animazione stessa. ``Se vogliamo, possiamo anche introdurre delle animazioni in più.''
Una dimostrazione volutamente assurda: ``io voglio, per qualche motivo arcano, far sì
che quando finisce l'animazione lui salti'': si trascina \textbf{MM\_Jump} nella traccia,
dopo l'attacco. ``È un salto orribile, però farà le due animazioni di conseguenza.'' Il
Montage è ``molto importante per azioni consecutive''. ``Quando si parla con l'animatore,
è importante il fatto di dire: questa animazione deve essere consecutiva a questa'', per
esempio ``mi serve che da questo punto in poi continui l'attacco, oppure faccia un'altra
animazione che ritorna in place''.

\paragraph{L'animazione di ritorno.} Finito l'attacco in root motion, il cavaliere resta dove l'attacco lo ha lasciato e non
torna in posa: manca ``un'animazione che ritorna in place'', cioè in idle. Il docente la
cerca su Mixamo, provando voci come \emph{Sword and Shield Idle}, \emph{Idle
Transition} e \emph{Start Block}; ma ``è un po' difficile'', perché ``bisogna trovare
il cavaliere proprio'', con la stessa posa di fine attacco, altrimenti ``snapperà
malissimo, almeno nel mio caso''.

\paragraph{Gli Slot.} Sotto, c'è il \textbf{Default Group} con lo \textbf{Slot}
\textbf{DefaultSlot}: ``questo ci permette di dire: su quale slot vuoi che questa
animazione si compia?''. DefaultSlot ``riprende tutto il corpo''. ``Potremmo, ad esempio,
voler creare uno slot che riprenda solamente alcune zone del nostro character'': ``vogliamo
che faccia un'azione solo nella parte superiore del corpo'' (sparare o salutare mentre si
corre); ``in quel caso creeremo un altro slot, che prenderà solo le ossa'' di quella
parte. ``Ma in questo caso ci serve tutto il corpo.''

% ---------------------------------------------------------------------
\section{Lo Slot nell'AnimGraph: Save Cached Pose}
\label{ue9:sec:slot}

Un Montage suona solo se l'AnimGraph ha un nodo \textbf{Slot} con lo stesso nome. In
ABP\_Knight, nell'AnimGraph:
\begin{enumerate}
  \item dall'uscita della State Machine Locomotion si crea un \textbf{New Save Cached
        Pose}, chiamato \textbf{LocomotionCache}: ``andiamo a conservarci quello che ha
        fatto la State Machine'', ``come se tutto questo procedimento lo mettesse in una
        memoria'';
  \item si aggiunge \textbf{Use cached pose 'LocomotionCache'};
  \item da lì un nodo \textbf{Slot 'DefaultSlot'} (la trascrizione: ``default slot'');
  \item lo Slot va all'\textbf{Output Pose}.
\end{enumerate}
``Semplicemente mixa insieme i due'': lo slot del Montage e la cache della locomozione
(figura~\ref{ue9:fig:slot}). Quando non c'è nessun Montage in corso, lo Slot lascia
passare la locomozione; quando un Montage su DefaultSlot parte, la sostituisce. Alla
fine della lezione, a uno studente che non legge il nome del nodo (``vedo tutto
pixelato''), il docente conferma: ``tiri fuori Locomotion Cache'', poi ``default slot''.

\begin{figure}[htbp]
\centering
\begin{tikzpicture}[font=\footnotesize,
  ue9pose/.style={draw, thick, rounded corners=4pt, align=center, fill=gray!8,
                  minimum height=1.0cm, text width=2.3cm},
  ue9arr/.style={-{Stealth[length=2.5mm]}, very thick, black!70}]
  \node[ue9pose, fill=orange!12] (sm) at (0,0) {\textbf{Locomotion}\\{\scriptsize State
    Machine}};
  \node[ue9pose, fill=blue!8] (sc) at (3.1,0) {\textbf{Save Cached Pose}\\{\scriptsize
    LocomotionCache}};
  \node[ue9pose, fill=blue!8] (uc) at (6.2,0) {\textbf{Use Cached Pose}\\{\scriptsize
    LocomotionCache}};
  \node[ue9pose, fill=green!10] (sl) at (9.3,0) {\textbf{Slot}\\{\scriptsize
    'DefaultSlot'}};
  \node[ue9pose, fill=red!10, text width=1.8cm] (op) at (12.1,0) {\textbf{Output
    Pose}};
  \draw[ue9arr] (sm) -- (sc); \draw[ue9arr] (uc) -- (sl); \draw[ue9arr] (sl) -- (op);
  \draw[dashed, thick, gray] (sc.east) -- (uc.west);
  \node[ue9pose, fill=yellow!15, text width=3.2cm] (mt) at (9.3,-2.1)
    {\textbf{Anim Montage}\\{\scriptsize Sword and Shield Attack, slot
    DefaultGroup.DefaultSlot}};
  \draw[-{Stealth[length=2.5mm]}, thick, dashed] (mt) -- node[right, font=\scriptsize]
    {quando suona, sostituisce la posa} (sl);
\end{tikzpicture}
\caption{La catena di pose finale dell'AnimGraph di ABP\_Knight. La State Machine produce
la locomozione, che viene salvata in una cache e riletta; lo Slot DefaultSlot la lascia
passare, a meno che un Montage assegnato a quello slot non sia in esecuzione. Schema
ricostruito.}
\label{ue9:fig:slot}
\end{figure}

% ---------------------------------------------------------------------
\section{Far partire l'attacco: Play Montage}
\label{ue9:sec:playmontage}

\paragraph{Il tasto F e il nodo.} Per far partire il Montage, il docente non crea
un'altra Input Action: usa direttamente un evento di tastiera, \textbf{F} (``facciamo con
OK\dots\ no, facciamo F, F va benissimo''). Nel Blueprint del personaggio, il nodo
\textbf{Play Montage}:
\begin{itemize}
  \item \textbf{In Skeletal Mesh Component}: ``in quale Skeletal Mesh Component
        farlo'', cioè la \textbf{Mesh} del personaggio;
  \item \textbf{Montage to Play}: ``quello che abbiamo appena fatto'' (``ce n'è tanto solo
        uno'');
  \item \textbf{Play Rate}: ``possiamo modificare quanto è veloce l'animazione'';
  \item fra le uscite: \textbf{On Completed} (``quando sarà finita''), \textbf{On Blend
        Out} (``mentre finisce''), \textbf{On Interrupted} (``se fosse interrotta''),
        \textbf{On Notify Begin} e \textbf{On Notify End} (``quando un Notify inizia'' e
        ``quando finisce'').
\end{itemize}
``Se noi premiamo F, compiamo un Montage.''

\paragraph{Fermare il personaggio.} ``È orribile'': a parte la mancanza dell'animazione di
ritorno, durante l'attacco ``possiamo muovere il character tranquillamente''. Quindi, sul
tasto F, prima del Play Montage:
\begin{enumerate}
  \item \textbf{Set Max Walk Speed} a \textbf{0}: ``non dobbiamo farlo muovere'';
  \item \textbf{Get Player Controller} $\rightarrow$ \textbf{Disable Input}: ``andiamo a
        disabilitare gli input, qualsiasi''.
\end{enumerate}
Sulle uscite del Play Montage:
\begin{itemize}
  \item \textbf{On Completed}: Get Player Controller $\rightarrow$ \textbf{Enable Input}
        (``on completed solamente'');
  \item \textbf{On Blend Out}: \textbf{Set Max Walk Speed} a \textbf{600}, ``al suo
        massimo''.
\end{itemize}

\paragraph{La camera che si muove ancora.} In Play, però, ``la camera si muove
effettivamente'', e con la camera il personaggio, ``quindi dobbiamo far sì che non sia
dipendente dalla camera''. Segue una serie di tentativi, ragionati ad alta voce:
\begin{itemize}
  \item togliere \emph{Use Pawn Control Rotation} sullo Spring Arm (``no, no'');
  \item l'\emph{Enable Camera Lag} dello Spring Arm;
  \item le spunte \emph{Inherit Pitch}, \emph{Inherit Yaw} e \emph{Inherit Roll} dello
        Spring Arm (``inerit yo e roll''; ``no, lo yaw ci deve essere''; ``forse sì,
        solo il roll va bene''), con un \emph{Set} preso dallo Spring Arm (``Set Inherit
        Yaw'', a false, e di nuovo a true alla fine); ``completamente\dots\ non va'';
  \item \emph{Stop Movement Immediately} e i parametri di input del Player Controller
        (\emph{movement input}, \emph{look input}).
\end{itemize}
La soluzione, alla fine, è sulle due spunte già incontrate nella prima parte: ``non è lo
yaw'' dello Spring Arm, ``ma è l'orientamento al movimento''. Sul tasto F si disattivano
entrambe:
\begin{itemize}
  \item dal \textbf{Character Movement}, \textbf{Set Orient Rotation to Movement} a
        \textbf{false};
  \item sul Pawn stesso, senza collegare nessun target, \textbf{Set Use Controller
        Rotation Yaw} a \textbf{false};
\end{itemize}
e sull'\textbf{On Completed} si rimettono entrambe a \textbf{true}
(figura~\ref{ue9:fig:playmontage}). ``Ottimo. Adesso riusciamo'': il cavaliere attacca
fermo, la camera gira senza trascinarlo, e alla fine riprende la locomozione. ``Abbiamo
fatto il nostro character che attacca e fa la locomotion.'' (Il docente nota anche, di
passaggio, che il cavaliere non interagisce con il pulsante del livello, ``giustamente,
perché è la Mesh''; dalla trascrizione non si ricostruisce meglio il punto.)

A Federico, che non riesce a sistemare la camera, il docente riassume: ``devi giocare con
lo yaw e con l'Orient Rotation to Movement: li disattivi tutte e due, e li riattivi
sull'On Completed''. ``Dello Spring Arm, giusto?'' ``Non dello Spring Arm: questo qua è
del Pawn, lo trovi direttamente senza collegare niente. Invece questo è solo del
Character Movement.''

\begin{figure}[htbp]
\centering
\begin{tikzpicture}[font=\footnotesize,
  ue9node/.style={draw, thick, rounded corners=4pt, align=center, fill=gray!8,
                  minimum height=0.9cm, text width=2.5cm},
  ue9ev/.style={ue9node, fill=red!12},
  ue9exec/.style={-{Stealth[length=2.5mm]}, very thick, black!80}]
  \node[ue9ev, text width=1.4cm] (f) at (0,0) {\textbf{F}\\{\scriptsize Pressed}};
  \node[ue9node] (a) at (2.8,0) {\textbf{Set Max Walk Speed}\\{\scriptsize 0}};
  \node[ue9node] (b) at (5.9,0) {Get Player Controller $\rightarrow$ \textbf{Disable
    Input}};
  \node[ue9node] (c) at (9.0,0) {\textbf{Set Orient Rotation to Movement}\\
    {\scriptsize false}};
  \node[ue9node] (d) at (12.1,0) {\textbf{Set Use Controller Rotation Yaw}\\
    {\scriptsize false}};
  \node[ue9node, fill=yellow!15, text width=3.0cm, minimum height=2.4cm] (pm) at
    (2.8,-3.0) {\textbf{Play Montage}\\{\scriptsize Mesh,\\Sword and Shield
    Attack\_Montage}};
  \node[ue9node] (e) at (8.3,-2.2) {\textbf{Set Max Walk Speed}\\{\scriptsize 600}};
  \node[ue9node, text width=3.0cm] (g) at (8.3,-3.9) {Get Player Controller
    $\rightarrow$ \textbf{Enable Input}};
  \node[ue9node, text width=3.0cm] (h) at (12.3,-3.9) {\textbf{Set Orient\dots} true\\
    \textbf{Set Use Controller Rotation Yaw} true};
  \draw[ue9exec] (f) -- (a); \draw[ue9exec] (a) -- (b); \draw[ue9exec] (b) -- (c);
  \draw[ue9exec] (c) -- (d);
  \draw[ue9exec] (d.south) -- ++(0,-0.6) -| (pm.north);
  \draw[ue9exec] (pm.east |- e) -- node[above, font=\scriptsize]{On Blend Out} (e);
  \draw[ue9exec] (pm.east |- g) -- node[below, font=\scriptsize]{On Completed} (g);
  \draw[ue9exec] (g) -- (h);
\end{tikzpicture}
\caption{L'attacco sul tasto F, nella versione finale. Prima del Montage si blocca il
personaggio (velocità zero, input disabilitato, nessuna rotazione automatica del corpo);
al Blend Out torna la velocità, all'On Completed tornano l'input e le rotazioni. L'ordine
esatto dei primi quattro nodi è ricostruito: a lezione vengono aggiunti uno alla volta,
durante i tentativi. Schema ricostruito.}
\label{ue9:fig:playmontage}
\end{figure}

\paragraph{Approfondimento: Play Montage, Montage Play e il root motion.} Il nodo usato a
lezione, \textbf{Play Montage}, è un nodo \emph{asincrono} (con l'orologio nell'angolo):
le sue uscite scattano più tardi, quando succede l'evento corrispondente, e per questo
può gestire il ``dopo'' del Montage direttamente nel grafo. Esiste anche \emph{Montage
Play}, chiamato sull'\emph{Anim Instance} della mesh, che restituisce solo la durata, e
\emph{Play Anim Montage}, del Character. Due dettagli completano la lezione. Il primo:
un attacco in \textbf{root motion} sposta davvero la capsula solo se il Montage ha il
root motion attivo e l'Animation Blueprint ha la \emph{Root Motion Mode} impostata su
\emph{Root Motion from Montages Only} (il default); se il root dell'animazione di Mixamo è
``rotto'', come visto nel retargeting, il personaggio avanza nella mesh e poi ``torna
indietro'' di scatto a fine animazione. Il secondo: durante un Montage che fa avanzare il
personaggio in root motion non serve azzerare la Max Walk Speed, perché il root motion
prende il posto dell'input di movimento; azzerarla, come fa il docente, è però un modo
semplice e sicuro di impedire al giocatore di muoversi durante un attacco ``in place''.
Infine, \emph{Disable Input} e \emph{Enable Input} agiscono sull'actor e richiedono un
Player Controller: per bloccare solo alcune azioni (il movimento ma non la camera),
l'Enhanced Input offre un'alternativa più fine, cioè rimuovere temporaneamente un
Input Mapping Context, o aggiungerne uno di priorità più alta.

% ---------------------------------------------------------------------
\section{Il compito e la chiusura}
\label{ue9:sec:chiusura}

\paragraph{Feedback in chat.} Una studentessa, Giulia, ha scritto in chat che è rimasta un
po' indietro ``causa trasloco'' e chiede di poter mandare dei feedback; nel frattempo è
già uscita dalla chiamata. ``Le risponderò sul Discord, assolutamente.'' ``Controllerò i
lavori che mandate sul canale.''

\paragraph{Il compito: lo StackOBot.} ``Una cosa che vorrei darvi come task, se vi va, per
la prossima volta'': ``fare un nuovo character totalmente da zero''. Nel progetto c'è già
un personaggio, lo \textbf{StackOBot} (il robottino dei progetti di esempio di Epic), che
``ha queste animazioni qua: vedete che ce le ha già dentro tutte quante''. Lo si trova
nella cartella \emph{StackOBot $\rightarrow$ Blueprints $\rightarrow$ Character}, con la
sua Skeletal Mesh (\emph{SKM\_Bot}) e il suo ABP. Il compito:
\begin{itemize}
  \item creare un \textbf{Blueprint Character} per lo StackOBot (``fare la BP'');
  \item creare il suo \textbf{Animation Blueprint} da zero, ``fare tutto l'Animation
        Blueprint e collegarci tutte le sue animazioni'': il ciclo di locomozione, con
        State Machine e Blend Space, come quello del cavaliere.
\end{itemize}
``Se volete farlo, lo potete mandare poi sulla chat'', ``che tanto magari ci lavoreremo
la prossima volta''. ``Per adesso, poi ne riparliamo.''

\paragraph{I saluti.} ``Se tutto a posto\dots'' Il docente interrompe la condivisione:
``vi ringrazio per l'attenzione, e ringrazio da casa chi è in differita''. ``Ci sentiamo
per martedì prossimo. Va bene, ciao, ragazzi. Buona serata.''

\begin{table}[htbp]
\centering
\small
\begin{tabular}{@{}p{5.4cm}p{8.6cm}@{}}
\toprule
\textbf{Sintomo} & \textbf{Causa e rimedio} \\
\midrule
L'animazione dello stato non si vede nella preview & la State Machine non è collegata
all'Output Pose dell'AnimGraph \\
L'animazione si ferma dopo un ciclo & manca la spunta di loop nel nodo dell'animazione \\
Warning nella State Machine & transizioni senza condizioni: normale finché non si
scrivono le regole \\
Jump $\rightarrow$ Fall dà errore & regola vuota: attivare Automatic Rule Based on
Sequence Player in State \\
Il Blend Space 1D arriva solo a 100 & Maximum Axis Value da impostare a 600; poi
spostare i campioni \\
Animazioni ``mischiate'' a metà (a 200) & i campioni non corrispondono alle velocità
reali: legarli alla Max Walk Speed \\
Alt sul nodo non crea il Set & negli stati dell'AnimGraph le variabili si leggono e
basta \\
``Questo asset non ha schede ancorate'' & Window $\rightarrow$ Load Layout $\rightarrow$
Default Editor Layout, poi Open Full Blueprint Editor \\
In Play non si vede niente / visuale nel corpo & il Character non ha camera: aggiungere
Spring Arm e Camera (figlia) \\
Non ci si muove & mancano gli input: copiare dal ThirdPersonCharacter i grafi di Look,
Move, Jump e l'Add Mapping Context \\
La camera ruota su un solo asse, o guarda solo il terreno & Use Pawn Control Rotation
sulla Camera invece che sullo Spring Arm \\
Il personaggio va di lato guardando avanti & attivare Orient Rotation to Movement e
togliere Use Controller Rotation Yaw (servono tutte e due) \\
La camera compenetra gli oggetti & Target Offset troppo grande: usare il Socket Offset \\
Il salto in root motion retargettato non è pulito & il root degli scheletri di Mixamo:
accettabile per la prova, da correggere \\
La State Machine non distingue salto e caduta & entrambe le regole su IsFalling:
aggiungere VelocityZ maggiore o minore di zero \\
Lo Shift non fa correre & IA\_Run non aggiunta ai mapping di IMC\_Default \\
FInterp To non interpola & eseguito una volta sola: va eseguito a ogni frame (Tick,
timer o Timeline) \\
Direction quasi sempre a zero & Orient Rotation to Movement acceso, Use Controller
Rotation Yaw spento: invertire per il Blend Space 2D \\
Scatti nel Blend Space 2D & animazioni non omogenee, salto fra 180 e $-180$: Smoothing
Time, o riscaricare le animazioni \\
Durante l'attacco il personaggio si muove & Set Max Walk Speed a 0 e Disable Input prima
del Play Montage \\
Durante l'attacco la camera trascina il personaggio & Set Orient Rotation to Movement e
Set Use Controller Rotation Yaw a false, a true sull'On Completed \\
Il Montage non suona & manca il nodo Slot DefaultSlot nell'AnimGraph \\
\bottomrule
\end{tabular}
\caption{I problemi incontrati durante la lezione, con la causa e il rimedio.}
\label{ue9:tab:errori}
\end{table}

\begin{table}[htbp]
\centering
\small
\begin{tabular}{@{}p{4.8cm}p{9.2cm}@{}}
\toprule
\textbf{Nodo, asset o impostazione} & \textbf{Che cosa fa} \\
\midrule
State Machine (AnimGraph) & organizza le animazioni in stati e transizioni \\
Entry / State / Conduit / Alias & ingresso; stato con una posa; snodo senza posa; alias
di più stati \\
Output Animation Pose & la posa prodotta da uno stato \\
Transition Rule (Result) & la condizione booleana per passare da uno stato all'altro \\
Automatic Rule Based on Sequence Player in State & la transizione scatta a fine
animazione \\
NOT Boolean / AND Boolean & negazione; congiunzione di due condizioni \\
Split Struct Pin & divide un vettore nei suoi componenti X, Y, Z \\
Blend Space 1D / Blend Space (2D) & mescola animazioni su uno o due assi \\
Horizontal / Vertical Axis, Grid Divisions & nome, minimo e massimo degli assi;
divisioni della griglia \\
Smoothing Time & tempo di passaggio fra i campioni \\
Blend Space Player & il nodo che suona un Blend Space nell'AnimGraph \\
Calculate Direction & l'angolo, da $-180$ a 180, fra la velocità e la rotazione \\
Get Actor Rotation / Get Velocity & la rotazione dell'actor; la sua velocità \\
Default Pawn Class (World Settings) & il personaggio del giocatore nel livello \\
Spring Arm: Target Arm Length, Socket Offset, Target Offset & lunghezza del braccio;
spostamento in fondo al braccio; spostamento dell'origine \\
Camera: Field of View & angolo di visuale \\
Use Pawn Control Rotation (Spring Arm) & la camera segue il controller \\
Use Controller Rotation Yaw (Pawn) & il corpo segue la rotazione del controller \\
Orient Rotation to Movement, Rotation Rate (Character Movement) & il corpo si gira verso
il movimento, con quella velocità \\
Max Walk Speed / Set Max Walk Speed & velocità massima a piedi (cm/s) \\
Retarget Animations, Export Animations & trasferisce un'animazione su un altro
scheletro \\
Input Mapping Context (IMC) & associa i tasti alle azioni in un contesto \\
Input Action (IA): Value Type, Triggers & l'azione, il tipo di valore, quando scatta \\
Add Mapping Context & attiva un IMC per il giocatore \\
Evento IA: Triggered, Started, Ongoing, Canceled, Completed & le fasi di un'azione \\
FInterp To & un passo di interpolazione fra Current e Target \\
Get World Delta Seconds & il tempo dell'ultimo frame \\
Animation Montage & animazione ``una tantum'', modulare, con sezioni e slot \\
Notify / Play Sound / Notify Track & eventi sulla linea del tempo; un suono; una traccia
di eventi \\
Montage Section & una parte con nome del Montage \\
Slot (DefaultGroup.DefaultSlot) & dove il Montage si applica (tutto il corpo, o parti) \\
Save Cached Pose / Use Cached Pose & salva e rilegge una posa nell'AnimGraph \\
Play Montage & suona un Montage; On Completed, On Blend Out, On Interrupted, Notify
Begin/End \\
Disable Input / Enable Input & toglie e restituisce l'input all'actor \\
\bottomrule
\end{tabular}
\caption{I nodi, gli asset e le impostazioni della nona lezione.}
\label{ue9:tab:nodi}
\end{table}

% ---------------------------------------------------------------------
\section*{Punti chiave del capitolo}
\addcontentsline{toc}{section}{Punti chiave del capitolo}

\begin{itemize}
  \item Una \textbf{State Machine} nell'AnimGraph organizza le animazioni in
        \textbf{stati} (WalkRun, Jump, Fall), ciascuno con la sua \textbf{Output
        Animation Pose}, e in \textbf{transizioni} con regole booleane: un ``loop di
        animazione secondo delle condizioni''. Va collegata all'Output Pose.
  \item Un \textbf{Blend Space 1D} mescola animazioni lungo un asse: per la locomozione,
        la \textbf{GroundSpeed}. I campioni vanno posti alle velocità reali del
        personaggio (Max Walk Speed, in centimetri al secondo): Idle a 0, Walk a 300,
        corsa a 600. Valori intermedi ``inventati'' producono obbrobri.
  \item Un \textbf{Character} non ha la camera: si aggiungono uno \textbf{Spring Arm} e,
        come suo figlio, una \textbf{Camera}. Distanza e altezza si regolano sullo
        Spring Arm (\textbf{Target Arm Length}, \textbf{Socket Offset}), mai spostando la
        camera; il \textbf{Target Offset} solo sulla Z e poco. Sono il lavoro dei
        \textbf{3C}: Camera, Control, Character.
  \item Tre spunte governano la rotazione: \textbf{Use Pawn Control Rotation} sullo
        Spring Arm (la camera segue il mouse), \textbf{Use Controller Rotation Yaw} sul
        Pawn (il corpo segue la camera), \textbf{Orient Rotation to Movement} nel
        Character Movement (il corpo si gira verso il movimento).
  \item Il \textbf{retargeting} porta un'animazione su un altro scheletro (Retarget
        Animations, Export Animations): comodo, ma il risultato dipende dall'animazione,
        e il root di Mixamo dà problemi.
  \item Fra salto e caduta: \textbf{Automatic Rule Based on Sequence Player in State} per
        passare a fine animazione, \textbf{NOT IsFalling} per tornare a terra, e la
        \textbf{velocità sulla Z} per distinguere un salto (Z $>$ 0) da una caduta
        (Z $<$ 0).
  \item L'\textbf{Enhanced Input} separa i tasti dalle azioni: l'\textbf{Input Action}
        (IA\_Run, Digital bool, con i suoi Triggers) è associata a uno o più tasti
        (Shift, il cerchio del controller) in un \textbf{Input Mapping Context}, che va
        attivato con \textbf{Add Mapping Context}. Lo sprint: Set Max Walk Speed a 600
        sul Triggered, a 300 sul Completed.
  \item \textbf{FInterp To} fa un solo passo per esecuzione: per un passaggio graduale
        va eseguito a ogni frame (Tick, timer, Timeline).
  \item Il \textbf{Blend Space 2D} aggiunge la \textbf{direzione} (\textbf{Calculate
        Direction}, da $-180$ a 180: positivo a destra, negativo a sinistra, $\pm 180$
        indietro): otto animazioni a 600, l'Idle ripetuto a 0, la corsa all'indietro ai
        due estremi, griglia 8$\times$8, Smoothing Time. Serve Use Controller Rotation Yaw
        acceso.
  \item Un \textbf{Animation Montage} è un'animazione ``una tantum'' e modulare, con
        \textbf{Notify} (suoni, VFX in un punto preciso), \textbf{sezioni} e
        \textbf{slot}. Per suonarlo, nell'AnimGraph: State Machine, \textbf{Save Cached
        Pose}, Use Cached Pose, \textbf{Slot DefaultSlot}, Output Pose. Si fa partire con
        \textbf{Play Montage}, bloccando prima movimento, input e rotazioni e
        restituendoli su On Blend Out e On Completed.
  \item Compito facoltativo: un nuovo personaggio da zero con lo \textbf{StackOBot},
        Blueprint e Animation Blueprint completi.
\end{itemize}
